% Generated mechanically from frozen input inputs/p06/ja/duality.tex.
% Do not edit this generated file; edit the bound locale source instead.

% OK, start here.
%

\setcounter{chapter}{47}
\chapter{スキームの双対性}



\phantomsection
\label{duality-section-phantom}


\section{序論}
\label{duality-section-introduction}

\noindent
本章では、スキームの射に対する相対双対性とスキーム上の双対化複体を研究する。
参考文献として \cite{RD} がある。

\medskip\noindent
Noether 環に対する双対化複体は『双対化複体』の第 \ref{dualizing-section-dualizing} 節以下で定義・研究された。本章ではその続きを扱い、
スキーム上の双対化複体を研究する；第 \ref{duality-section-dualizing-schemes} 節を参照する。

\medskip\noindent
本章の大部分は、準連接コホモロジー層をもつ加群層の導来圏という設定で、押し出しの右随伴を
研究することに充てられる。第 \ref{duality-section-twisted-inverse-image} 節、第 \ref{duality-section-restriction-to-opens} 節、第 \ref{duality-section-base-change-map} 節、第 \ref{duality-section-base-change-II} 節、第 \ref{duality-section-trace} 節、第 \ref{duality-section-compare-with-pullback} 節、第 \ref{duality-section-sections-with-exact-support} 節、第 \ref{duality-section-duality-finite} 節、第 \ref{duality-section-perfect-proper} 節、第 \ref{duality-section-dualizing-Cartier} 節および第 \ref{duality-section-examples} 節
を参照する。ここでは論文
\cite{Neeman-Grothendieck}, \cite{LN},
\cite{Lipman-notes}, and \cite{Neeman-improvement}
に従う。

\medskip\noindent
Noether スキーム間の有限型分離射に対する重要で有用な上付きシュリーク函手 $f^!$ を
第 \ref{duality-section-upper-shriek} 節、第 \ref{duality-section-upper-shriek-properties} 節および第 \ref{duality-section-base-change-shriek} 節
で論じ、概観を与える第 \ref{duality-section-duality} 節に至る。

\medskip\noindent
第 \ref{duality-section-glue} 節では、双対化複体を備えた固定された局所 Noether 基底上で
作業するときの、双対性および双対化複体に関する別の理論を説明する
（この節は Hartshorne の著書中の一つの注意に対応する）。

\medskip\noindent
残りの節では、いくつかの応用を与える。

\medskip\noindent
本章の続きは代数空間上の双対性を扱う章にある；『代数空間の双対性』の第 \ref{spaces-duality-section-introduction} 節を参照する。







\section{スキーム上の双対化複体}
\label{duality-section-dualizing-schemes}

\noindent
局所 Noether スキーム上の双対化複体を、アフィン局所的には対応する環上の双対化複体から
得られる複体として定義する。これは完全に標準的な定義ではないが、有限次元 Noether
スキームについて文献に現れるすべての定義と一致する。

\begin{lemma}
\label{duality-lemma-equivalent-definitions}
$X$ を局所 Noether スキームとし、$K$ を $D(\mathcal{O}_X)$ の対象とする。
次は同値である。
\begin{enumerate}
\item 任意のアフィン開集合 $U = \Spec(A) \subset X$ に対して、$A$ の双対化複体
$\omega_A^\bullet$ が存在し、$K|_U$ は関手
$\widetilde{} : D(A) \to D(\mathcal{O}_U)$ による $\omega_A^\bullet$ の像と同型である。
\item アフィン開被覆 $X = \bigcup U_i$, $U_i = \Spec(A_i)$ が存在し、各 $i$ に対して
$A_i$ の双対化複体 $\omega_i^\bullet$ が存在して、$K|_{U_i}$ は関手
$\widetilde{} : D(A_i) \to D(\mathcal{O}_{U_i})$ による $\omega_i^\bullet$ の像と同型である。
\end{enumerate}
\end{lemma}

\begin{proof}
(2) を仮定し、$U = \Spec(A)$ を $X$ のアフィン開集合とする。(2) から
$K$ は $D_\QCoh(\mathcal{O}_X)$ に属するので、対応する準連接層の複体が $K|_U$ と
同型になる対象 $\omega_A^\bullet$ が $D(A)$ に存在する；『スキームの導来圏』の補題 \ref{perfect-lemma-affine-compare-bounded} を参照する。$\omega_A^\bullet$ が
$A$ の双対化複体であることを示せば、証明が完了する。

\medskip\noindent
$X = \bigcup U_i$ は開被覆なので、標準開被覆
$U = D(f_1) \cup \ldots \cup D(f_m)$ で、各 $D(f_j)$ がいずれかのアフィン開集合
$U_i$ の標準開集合となるものを取れる；『スキーム』の補題 \ref{schemes-lemma-standard-open-two-affines} を参照する。
$g_j \in A_{i_j}$ に対して $D(f_j) = D(g_j)$ としよう。このとき
$A_{f_j} \cong (A_{i_j})_{g_j}$ であり、『スキームの導来圏』の補題 \ref{perfect-lemma-affine-compare-bounded} により導来圏で
$$
(\omega_A^\bullet)_{f_j} \cong (\omega_i^\bullet)_{g_j}
$$
である。『双対化複体』の補題 \ref{dualizing-lemma-dualizing-localize} により、
$j = 1, \ldots, m$ に対して複体 $(\omega_A^\bullet)_{f_j}$ は $A_{f_j}$ 上の
双対化複体である。従って『双対化複体』の補題 \ref{dualizing-lemma-dualizing-glue} により $\omega_A^\bullet$ は双対化複体である。
\end{proof}

\begin{definition}
\label{duality-definition-dualizing-scheme}
$X$ を局所 Noether スキームとする。$D(\mathcal{O}_X)$ の対象 $K$ が補題 \ref{duality-lemma-equivalent-definitions} の同値な条件を満たすとき、$K$ を
{\it 双対化複体}と呼ぶ。
\end{definition}

\noindent
本節冒頭の注意も参照されたい。

\begin{lemma}
\label{duality-lemma-affine-duality}
$A$ を Noether 環、$X = \Spec(A)$ とし、$K, L$ を $D(A)$ の対象とする。
$K \in D_{\textit{Coh}}(A)$ であり、$L$ が有限入射次元をもつならば、
$D(\mathcal{O}_X)$ において
$$
R\SheafHom_{\mathcal{O}_X}(\widetilde{K}, \widetilde{L})
=
\widetilde{R\Hom_A(K, L)}
$$
である。
\end{lemma}

\begin{proof}
$L$ は入射 $A$-加群からなる有限複体 $I^\bullet$ により与えられると仮定してよい。
$I^\bullet$ の長さに関する帰納法と、構成が区別三角と両立することにより、
$I$ が入射 $A$-加群である場合の $L = I[0]$ に帰着する。この場合『スキームの導来圏』の補題 \ref{perfect-lemma-quasi-coherence-internal-hom} により、
$R\SheafHom_{\mathcal{O}_X}(\widetilde{K}, \widetilde{L})$ の第 $n$ コホモロジー層は、
前層
$$
D(f) \longmapsto \Ext^n_{A_f}(K \otimes_A A_f, I \otimes_A A_f)
$$
に伴う層である。$A$ は Noether なので、$A_f$-加群 $I \otimes_A A_f$ は入射である
（『双対化複体』の補題 \ref{dualizing-lemma-localization-injective-modules}）。従って
\begin{align*}
\Ext^n_{A_f}(K \otimes_A A_f, I \otimes_A A_f)
& =
\Hom_{A_f}(H^{-n}(K \otimes_A A_f), I \otimes_A A_f) \\
& =
\Hom_{A_f}(H^{-n}(K) \otimes_A A_f, I \otimes_A A_f) \\
& =
\Hom_A(H^{-n}(K), I) \otimes_A A_f
\end{align*}
である。最後の等号は $H^{-n}(K)$ が有限 $A$-加群であることによる；『可換代数』の補題 \ref{algebra-lemma-hom-from-finitely-presented} を参照する。これにより標準写像
$$
\widetilde{R\Hom_A(K, L)}
\longrightarrow
R\SheafHom_{\mathcal{O}_X}(\widetilde{K}, \widetilde{L})
$$
はこの場合に擬同型であり、証明が完了する。
\end{proof}

\begin{lemma}
\label{duality-lemma-internal-hom-evaluate-isom}
$X$ を Noether スキームとし、$K, L, M \in D_\QCoh(\mathcal{O}_X)$ とする。
このとき『層のコホモロジー』の補題 \ref{cohomology-lemma-internal-hom-evaluate} の写像
$$
R\SheafHom(L, M) \otimes_{\mathcal{O}_X}^\mathbf{L} K
\longrightarrow
R\SheafHom(R\SheafHom(K, L), M)
$$
は次の二つの場合に同型である。
\begin{enumerate}
\item $K \in D^-_{\textit{Coh}}(\mathcal{O}_X)$,
$L \in D^+_{\textit{Coh}}(\mathcal{O}_X)$ であり、$M$ がアフィン局所的に有限入射次元をもつ
（証明を参照する）。
\item $K$ と $L$ が $D_{\textit{Coh}}(\mathcal{O}_X)$ に属し、対象
$R\SheafHom(L, M)$ が有限 Tor 次元をもち、$L$ と $M$ がアフィン局所的に有限入射次元をもつ
（特に $L$ と $M$ は有界である）。
\end{enumerate}
\end{lemma}

\begin{proof}
(1) の証明。$X$ がアフィン開集合 $U = \Spec(A)$ による開被覆をもち、$M|_U$ に対応する
$D(A)$ の対象（『スキームの導来圏』の補題 \ref{perfect-lemma-affine-compare-bounded}）が有限入射次元をもつとき、$M$ は
アフィン局所的に有限入射次元をもつという\footnote{この条件は Noether スキーム $X$ の
アフィン開被覆の選択に依存しない。詳細は省略する。}。補題を証明するため $X$ を $U$ で
置き換えてよい。すなわち、ある Noether 環 $A$ に対して $X = \Spec(A)$ と仮定してよい。
『スキームの導来圏』の補題 \ref{perfect-lemma-coherent-internal-hom} により
$R\SheafHom(K, L)$ は $D^+_{\textit{Coh}}(\mathcal{O}_X)$ に属する。また補題 \ref{duality-lemma-affine-duality} と『スキームの導来圏』の補題 \ref{perfect-lemma-quasi-coherence-internal-hom} により、この矢印の左辺と右辺の形成は関手
$D(A) \to D_\QCoh(\mathcal{O}_X)$ と可換である
（正確には、ここで $K$, $L$, $M$ に関する仮定と、今示した
$R\SheafHom(K, L)$ に関する事実を用いる）。最後に『代数の続論』の補題 \ref{more-algebra-lemma-internal-hom-evaluate-isomorphism} の (2) により、この矢印は同型である。

\medskip\noindent
(2) の証明。上と同様に論じる。わずかな相違は、ここではアフィン局所的に
（補題 \ref{duality-lemma-affine-duality} によりこれは許される）『双対化複体』の補題 \ref{dualizing-lemma-finite-ext-into-bounded-injective} を適用できるため、
$R\SheafHom(K, L)$ が $D_{\textit{Coh}}(\mathcal{O}_X)$ に属することである。
最後に『代数の続論』の補題 \ref{more-algebra-lemma-internal-hom-evaluate-isomorphism-technical} により結論する。
\end{proof}

\begin{lemma}
\label{duality-lemma-dualizing-schemes}
$K$ を局所 Noether スキーム $X$ 上の双対化複体とする。このとき $K$ は
$D_{\textit{Coh}}(\mathcal{O}_X)$ の対象であり、
$D = R\SheafHom_{\mathcal{O}_X}(-, K)$ は反同値
$$
D :
D_{\textit{Coh}}(\mathcal{O}_X)
\longrightarrow
D_{\textit{Coh}}(\mathcal{O}_X)
$$
を誘導し、これは標準同型 $\text{id} \to D \circ D$ を備える。$X$ が準コンパクトならば、
$D$ は $D^+_{\textit{Coh}}(\mathcal{O}_X)$ と
$D^-_{\textit{Coh}}(\mathcal{O}_X)$ を交換し、反同値
$D^b_{\textit{Coh}}(\mathcal{O}_X) \to D^b_{\textit{Coh}}(\mathcal{O}_X)$
を誘導する。
\end{lemma}

\begin{proof}
$U \subset X$ をアフィン開集合とする。$U = \Spec(A)$ と書き、補題 \ref{duality-lemma-equivalent-definitions} のように $K|_U$ に対応する $A$ の双対化複体を
$\omega_A^\bullet$ とする。補題 \ref{duality-lemma-affine-duality} により図式
$$
\xymatrix{
D_{\textit{Coh}}(A) \ar[r] \ar[d]_{R\Hom_A(-, \omega_A^\bullet)} &
D_{\textit{Coh}}(\mathcal{O}_U) \ar[d]^{R\SheafHom_{\mathcal{O}_X}(-, K|_U)} \\
D_{\textit{Coh}}(A) \ar[r] &
D(\mathcal{O}_U)
}
$$
は可換である。従って $D$ は $D_{\textit{Coh}}(\mathcal{O}_X)$ を
$D_{\textit{Coh}}(\mathcal{O}_X)$ に送る。さらに標準写像
$$
L
\longrightarrow
R\SheafHom_{\mathcal{O}_X}(K, K) \otimes_{\mathcal{O}_X}^\mathbf{L} L
\longrightarrow
R\SheafHom_{\mathcal{O}_X}(R\SheafHom_{\mathcal{O}_X}(L, K), K)
$$
（第二の矢印には『層のコホモロジー』の補題 \ref{cohomology-lemma-internal-hom-evaluate} を用いる）は、任意の $L$ に対して同型である。
実際、アフィン上では『双対化複体』の補題 \ref{dualizing-lemma-dualizing} によりそうである\footnote{別の方法として、まず
アフィン局所的に計算して
$R\SheafHom_{\mathcal{O}_X}(K, K) = \mathcal{O}_X$ を示し、次に補題 \ref{duality-lemma-internal-hom-evaluate-isom} の (2) を用いて写像が同型であることを示せる。}。
そしてアフィン上では代数の場合の結果が回収されることをすでに見た。
準コンパクトな場合の関手 $D$ の有界性に関する主張も、『双対化複体』の補題 \ref{dualizing-lemma-dualizing} の対応する主張から従う。
\end{proof}

\noindent
$X$ を局所環付き空間とする。$D(\mathcal{O}_X)$ の対象 $L$ が、導来テンソル積により
$D(\mathcal{O}_X)$ に与えられる対称モノイド構造に関して可逆対象であるとき、それは
{\it 可逆}であるという。『層のコホモロジー』の補題 \ref{cohomology-lemma-invertible-derived} で見たように、これは $L$ が完全であり、ある開被覆
$X = \bigcup U_i$ と整数 $n_i$ が存在して
$L|_{U_i} \cong \mathcal{O}_{U_i}[-n_i]$ となることを意味する。この場合、関数
$$
x \mapsto n_x,\quad
\text{where }n_x\text{ is the unique integer such that }
H^{n_x}(L_x) \not = 0
$$
は $X$ 上局所定数である。特に
$L = \bigoplus H^n(L)[-n]$ であり、これは $L$ を表す明確に定まる
$\mathcal{O}_X$-加群の複体（微分は零）を与える。

\begin{lemma}
\label{duality-lemma-dualizing-unique-schemes}
$X$ を局所 Noether スキームとする。$K$ と $K'$ が $X$ 上の双対化複体ならば、
$D(\mathcal{O}_X)$ のある可逆対象 $L$ に対して $K'$ は
$K \otimes_{\mathcal{O}_X}^\mathbf{L} L$ と同型である。
\end{lemma}

\begin{proof}
$$
L = R\SheafHom_{\mathcal{O}_X}(K, K')
$$
とおく。アフィン局所的にはこれが成り立つので、これは $D(\mathcal{O}_X)$ の可逆対象である；
『双対化複体』の補題 \ref{dualizing-lemma-dualizing-unique} とその証明を参照する。
同じ理由により、評価写像 $L \otimes_{\mathcal{O}_X}^\mathbf{L} K \to K'$ は同型である。
\end{proof}

\begin{lemma}
\label{duality-lemma-dimension-function-scheme}
$X$ を局所 Noether スキームとし、$\omega_X^\bullet$ を $X$ 上の双対化複体とする。
このとき $X$ は普遍鎖状であり、
$$
x \longmapsto \delta(x)\text{ such that }
\omega_{X, x}^\bullet[-\delta(x)]
\text{ is a normalized dualizing complex over }
\mathcal{O}_{X, x}
$$
により定まる関数 $X \to \mathbf{Z}$ は次元関数である。
\end{lemma}

\begin{proof}
アフィンの場合である『双対化複体』の補題 \ref{dualizing-lemma-dimension-function} と定義から直ちに従う。
\end{proof}

\begin{lemma}
\label{duality-lemma-sitting-in-degrees}
$X$ を局所 Noether スキームとし、$\omega_X^\bullet$ を対応する次元関数 $\delta$ をもつ
$X$ 上の双対化複体とする。$\mathcal{F}$ を連接 $\mathcal{O}_X$-加群とし、
$\mathcal{E}^i = \SheafExt^{-i}_{\mathcal{O}_X}(\mathcal{F}, \omega_X^\bullet)$
とおく。このとき $\mathcal{E}^i$ は連接 $\mathcal{O}_X$-加群であり、
$x \in X$ に対して次が成り立つ。
\begin{enumerate}
\item $\mathcal{E}^i_x$ が零でないのは
$\delta(x) \leq i \leq \delta(x) + \dim(\text{Supp}(\mathcal{F}_x))$
の場合に限る。
\item $\dim(\text{Supp}(\mathcal{E}^{i + \delta(x)}_x)) \leq i$ である。
\item $\text{depth}(\mathcal{F}_x)$ は
$\mathcal{E}_x^{i + \delta(x)} \not = 0$ となる最小の整数 $i \geq 0$ である。
\item
$x \in \text{Supp}(\bigoplus_{j \leq i} \mathcal{E}^j)
\Leftrightarrow
\text{depth}_{\mathcal{O}_{X, x}}(\mathcal{F}_x) + \delta(x) \leq i$
である。
\end{enumerate}
\end{lemma}

\begin{proof}
補題 \ref{duality-lemma-dualizing-schemes} により $\mathcal{E}^i$ は連接である。$x$ の
アフィン近傍を選び、『スキームの導来圏』の補題 \ref{perfect-lemma-quasi-coherence-internal-hom} と『代数の続論』の補題 \ref{more-algebra-lemma-base-change-RHom} の (3) を用いると
$$
\mathcal{E}^i_x =
\SheafExt^{-i}_{\mathcal{O}_X}(\mathcal{F}, \omega_X^\bullet)_x =
\Ext^{-i}_{\mathcal{O}_{X, x}}(\mathcal{F}_x,
\omega_{X, x}^\bullet) =
\Ext^{\delta(x) - i}_{\mathcal{O}_{X, x}}(\mathcal{F}_x,
\omega_{X, x}^\bullet[-\delta(x)])
$$
を得る。補題 \ref{duality-lemma-dimension-function-scheme} における $\delta$ の構成により、
(1), (2), (3) は『双対化複体』の補題 \ref{dualizing-lemma-sitting-in-degrees} に帰着する。(4) は (3) と (1) の形式的帰結である。
\end{proof}




\section{押し出しの右随伴}
\label{duality-section-twisted-inverse-image}

\noindent
本節および次節の参考文献は
\cite{Neeman-Grothendieck}, \cite{LN},
\cite{Lipman-notes}, および \cite{Neeman-improvement} である。

\medskip\noindent
$f : X \to Y$ をスキームの射とする。本節では関手
$Rf_* : D_\QCoh(\mathcal{O}_X) \to D_\QCoh(\mathcal{O}_Y)$
の右随伴を考える。文献では、この関手が存在するとき $f^{\times}$ と表されることがある。
この記法は広く合意されたものではないので、ここでは用いない。また、このような関手に
$f^!$ という記法も用いない。一般の射 $f$ に対しては \cite{RD} の記法と衝突するからである。

\begin{lemma}
\label{duality-lemma-twisted-inverse-image}
\begin{reference}
これは \cite[Example 4.2]{Neeman-Grothendieck} とほぼ同じである。
\end{reference}
$f : X \to Y$ を準分離かつ準コンパクトなスキームの間の射とする。このとき関手
$Rf_* : D_\QCoh(\mathcal{O}_X) \to D_\QCoh(\mathcal{O}_Y)$
は右随伴をもつ。
\end{lemma}

\begin{proof}
『導来圏』の命題 \ref{derived-proposition-brown} の仮定を検証して、右随伴の存在を示す。
まず圏 $D_\QCoh(\mathcal{O}_X)$ は直和をもつ；『スキームの導来圏』の補題 \ref{perfect-lemma-quasi-coherence-direct-sums} を参照する。また『スキームの導来圏』の定理 \ref{perfect-theorem-bondal-van-den-Bergh} により、圏 $D_\QCoh(\mathcal{O}_X)$ は
コンパクト生成される。$X$ と $Y$ は準コンパクトかつ準分離なので $f$ もそうである；『スキーム』の補題 \ref{schemes-lemma-compose-after-separated} および \ref{schemes-lemma-quasi-compact-permanence} を参照する。従って関手 $Rf_*$ は直和と可換である；
『スキームの導来圏』の補題 \ref{perfect-lemma-quasi-coherence-pushforward-direct-sums} を参照する。以上で証明が完了する。
\end{proof}

\begin{example}
\label{duality-example-affine-twisted-inverse-image}
$A \to B$ を環準同型とする。$Y = \Spec(A)$, $X = \Spec(B)$ とおき、
$f : X \to Y$ を $A \to B$ に対応する射とする。このとき
$Rf_* : D_\QCoh(\mathcal{O}_X) \to D_\QCoh(\mathcal{O}_Y)$ は、同値
$D(B) \to D_\QCoh(\mathcal{O}_X)$ および $D(A) \to D_\QCoh(\mathcal{O}_Y)$
を介して、係数制限 $D(B) \to D(A)$ に対応する。従って右随伴は『双対化複体』の第 \ref{dualizing-section-trivial} 節の関手 $K \longmapsto R\Hom(B, K)$ に対応する。
\end{example}

\begin{example}
\label{duality-example-does-not-preserve-coherent}
$f : X \to Y$ が Noether スキームの分離有限型射であっても、
$Rf_* : D_\QCoh(\mathcal{O}_X) \to D_\QCoh(\mathcal{O}_Y)$ の右随伴は
$D_{\textit{Coh}}(\mathcal{O}_Y)$ を $D_{\textit{Coh}}(\mathcal{O}_X)$
に送るとは限らない。実際、$k$ を体とし、射 $f : \mathbf{A}^1_k \to \Spec(k)$ を考える。
例 \ref{duality-example-affine-twisted-inverse-image} により、これは $A = k$, $B = k[x]$ のときに
$R\Hom(B, -)$ が $D_{\textit{Coh}}(A)$ を $D_{\textit{Coh}}(B)$ に送るかという問題に対応する。
これは成り立たない。なぜなら
$$
R\Hom(k[x], k) = \left(\prod\nolimits_{n \geq 0} k\right)[0]
$$
は有限 $k[x]$-加群ではないからである。従って $a(\mathcal{O}_Y)$ は連接なコホモロジー層をもたない。
\end{example}

\begin{example}
\label{duality-example-does-not-preserve-bounded-above}
$f : X \to Y$ が Noether スキームの固有射、さらには有限射であっても、
$Rf_* : D_\QCoh(\mathcal{O}_X) \to D_\QCoh(\mathcal{O}_Y)$ の右随伴は
$D_\QCoh^-(\mathcal{O}_Y)$ を $D_\QCoh^-(\mathcal{O}_X)$ に送るとは限らない。
実際、$k$ を体、$k[\epsilon]$ を $k$ 上の双対数とし、$X = \Spec(k)$,
$Y = \Spec(k[\epsilon])$ とおく。このとき任意の $i \geq 0$ に対して
$\Ext^i_{k[\epsilon]}(k, k)$ は零でない。従って例 \ref{duality-example-affine-twisted-inverse-image} により、$a(\mathcal{O}_Y)$ は上に有界でない。
\end{example}

\begin{lemma}
\label{duality-lemma-twisted-inverse-image-bounded-below}
$f : X \to Y$ を準コンパクトかつ準分離なスキームの射とする。$a : D_\QCoh(\mathcal{O}_Y) \to D_\QCoh(\mathcal{O}_X)$
を補題 \ref{duality-lemma-twisted-inverse-image} の $Rf_*$ の右随伴とする。このとき $a$ は
$D^+_\QCoh(\mathcal{O}_Y)$ を $D^+_\QCoh(\mathcal{O}_X)$ に送る。実際、ある整数 $N$ が存在して、
$H^i(K) = 0$ が $i \leq c$ に対して成り立てば、$H^i(a(K)) = 0$ が $i \leq c - N$ に対して成り立つ。
\end{lemma}

\begin{proof}
『スキームの導来圏』の補題 \ref{perfect-lemma-quasi-coherence-direct-image}
により、関手 $Rf_*$ は有限コホモロジー次元をもつ。言い換えると、ある整数 $N$ が存在して、
$H^i(L) = 0$ が $i \geq c$ に対して成り立つならば、$H^i(Rf_*L) = 0$ が
$i \geq N + c$ に対して成り立つ。$K \in D^+_\QCoh(\mathcal{O}_Y)$ が
$i \leq c$ に対して $H^i(K) = 0$ を満たすとする。このとき上の議論により
$$
\Hom_{D(\mathcal{O}_X)}(\tau_{\leq c - N}a(K), a(K)) =
\Hom_{D(\mathcal{O}_Y)}(Rf_*\tau_{\leq c - N}a(K), K) = 0
$$
である。これは明らかに、$i \leq c - N$ に対して $H^i(a(K)) = 0$ であることを意味する。
\end{proof}

\noindent
$f : X \to Y$ を準分離かつ準コンパクトなスキームの射とし、$a$ を
$Rf_* : D_\QCoh(\mathcal{O}_X) \to D_\QCoh(\mathcal{O}_Y)$ の右随伴とする。
任意の $K \in D_\QCoh(\mathcal{O}_Y)$ および $L \in D_\QCoh(\mathcal{O}_X)$ に対して標準写像
\begin{equation}
\label{duality-equation-sheafy-trace}
Rf_*R\SheafHom_{\mathcal{O}_X}(L, a(K))
\longrightarrow
R\SheafHom_{\mathcal{O}_Y}(Rf_*L, K)
\end{equation}
を得る。具体的には、この写像は合成
$$
Rf_*R\SheafHom_{\mathcal{O}_X}(L, a(K)) \to
R\SheafHom_{\mathcal{O}_Y}(Rf_*L, Rf_*a(K)) \to
R\SheafHom_{\mathcal{O}_Y}(Rf_*L, K)
$$
として構成される。第一の矢印は『層のコホモロジー』の注意 \ref{cohomology-remark-projection-formula-for-internal-hom} のものであり、第二の矢印は随伴の余単位
$Rf_*a(K) \to K$ である。

\begin{lemma}
\label{duality-lemma-iso-on-RSheafHom}
$f : X \to Y$ を準コンパクトかつ準分離なスキームの射とし、$a$ を
$Rf_* : D_\QCoh(\mathcal{O}_X) \to D_\QCoh(\mathcal{O}_Y)$ の右随伴とする。
$L \in D_\QCoh(\mathcal{O}_X)$, $K \in D_\QCoh(\mathcal{O}_Y)$ とする。このとき写像
(\ref{duality-equation-sheafy-trace})
$$
Rf_*R\SheafHom_{\mathcal{O}_X}(L, a(K))
\longrightarrow
R\SheafHom_{\mathcal{O}_Y}(Rf_*L, K)
$$
は、『スキームの導来圏』の第 \ref{perfect-section-better-coherator} 節で論じた関手
$DQ_Y : D(\mathcal{O}_Y) \to D_\QCoh(\mathcal{O}_Y)$ を適用すると同型になる。
\end{lemma}

\begin{proof}
『スキームの導来圏』の補題 \ref{perfect-lemma-better-coherator} により $DQ_Y$ は存在するので、
この主張には意味がある。$DQ_Y$ は包含関手 $D_\QCoh(\mathcal{O}_Y) \to D(\mathcal{O}_Y)$
の右随伴である。従って補題を示すには、任意の $M \in D_\QCoh(\mathcal{O}_Y)$ に対し写像
(\ref{duality-equation-sheafy-trace}) が全単射
$$
\Hom_Y(M, Rf_*R\SheafHom_{\mathcal{O}_X}(L, a(K)))
\longrightarrow
\Hom_Y(M, R\SheafHom_{\mathcal{O}_Y}(Rf_*L, K))
$$
を誘導することを示せばよい。そのため、次の一連の等式を用いる。
\begin{align*}
\Hom_Y(M, Rf_*R\SheafHom_{\mathcal{O}_X}(L, a(K)))
& =
\Hom_X(Lf^*M, R\SheafHom_{\mathcal{O}_X}(L, a(K))) \\
& =
\Hom_X(Lf^*M \otimes_{\mathcal{O}_X}^\mathbf{L} L, a(K)) \\
& =
\Hom_Y(Rf_*(Lf^*M \otimes_{\mathcal{O}_X}^\mathbf{L} L), K) \\
& =
\Hom_Y(M \otimes_{\mathcal{O}_Y}^\mathbf{L} Rf_*L, K) \\
& =
\Hom_Y(M, R\SheafHom_{\mathcal{O}_Y}(Rf_*L, K))
\end{align*}
第一の等式は『層のコホモロジー』の補題 \ref{cohomology-lemma-adjoint} により、第二の等式は
『層のコホモロジー』の補題 \ref{cohomology-lemma-internal-hom} により成り立つ。第三の等式は $a$ の構成による。
第四の等式は『スキームの導来圏』の補題 \ref{perfect-lemma-cohomology-base-change} による（これが重要な段階である）。第五の等式は
『層のコホモロジー』の補題 \ref{cohomology-lemma-internal-hom} による。
\end{proof}

\begin{example}
\label{duality-example-iso-on-RSheafHom}
補題 \ref{duality-lemma-iso-on-RSheafHom} の主張は「コヒーレータ」$DQ_Y$ を適用しなければ成り立たない。
実際、$Y = \Spec(R)$, $X = \mathbf{A}^1_R$ とし、$L = \mathcal{O}_X$,
$K = \mathcal{O}_Y$ とする。矢印の左辺は $D_\QCoh(\mathcal{O}_Y)$ に属するが、右辺は
準連接でない $\prod_{n \geq 0} \mathcal{O}_Y$ と同型である。
\end{example}

\begin{remark}
\label{duality-remark-iso-on-RSheafHom}
補題 \ref{duality-lemma-iso-on-RSheafHom} の状況では『スキームの導来圏』の補題 \ref{perfect-lemma-pushforward-better-coherator} により
$$
DQ_Y(Rf_*R\SheafHom_{\mathcal{O}_X}(L, a(K))) =
Rf_* DQ_X(R\SheafHom_{\mathcal{O}_X}(L, a(K)))
$$
である。従って $R\SheafHom_{\mathcal{O}_X}(L, a(K)) \in D_\QCoh(\mathcal{O}_X)$ ならば、
矢印の左辺の $DQ_Y$ を「取り除く」ことができる。一方、
$R\SheafHom_{\mathcal{O}_Y}(Rf_*L, K) \in D_\QCoh(\mathcal{O}_Y)$ と分かっているならば、
矢印の右辺から $DQ_Y$ を「取り除く」ことができる。両方が成り立てば
(\ref{duality-equation-sheafy-trace}) は同型である。これと『スキームの導来圏』の補題 \ref{perfect-lemma-quasi-coherence-internal-hom} を組み合わせると、次のいずれかの場合に
$Rf_*R\SheafHom_{\mathcal{O}_X}(L, a(K)) \to
R\SheafHom_{\mathcal{O}_Y}(Rf_*L, K)$ は同型である。
\begin{enumerate}
\item $L$ と $Rf_*L$ が完全である。
\item $K$ が下に有界であり、$L$ と $Rf_*L$ が擬連接である。
\end{enumerate}
(2) では、$K$ が下に有界ならば $a(K)$ も下に有界であることを用いる；補題 \ref{duality-lemma-twisted-inverse-image-bounded-below} を参照する。
\end{remark}

\begin{example}
\label{duality-example-iso-on-RSheafHom-noetherian}
$f : X \to Y$ を Noether スキームの固有射とし、$L \in D^-_{\textit{Coh}}(X)$,
$K \in D^+_{\QCoh}(\mathcal{O}_Y)$ とする。このとき写像
$Rf_*R\SheafHom_{\mathcal{O}_X}(L, a(K)) \to
R\SheafHom_{\mathcal{O}_Y}(Rf_*L, K)$ は同型である。実際、『スキームの導来圏』の補題 \ref{perfect-lemma-identify-pseudo-coherent-noetherian} および \ref{perfect-lemma-direct-image-coherent} により複体 $L$ と $Rf_*L$ は擬連接であり、注意 \ref{duality-remark-iso-on-RSheafHom} の議論が適用できる。
\end{example}

\begin{lemma}
\label{duality-lemma-iso-global-hom}
$f : X \to Y$ を準分離かつ準コンパクトなスキームの射とする。任意の
$L \in D_\QCoh(\mathcal{O}_X)$ および $K \in D_\QCoh(\mathcal{O}_Y)$ に対して、
(\ref{duality-equation-sheafy-trace}) は大域導来 Hom の同型
$R\Hom_X(L, a(K)) \to R\Hom_Y(Rf_*L, K)$ を誘導する。
\end{lemma}

\begin{proof}
『層のコホモロジー』の第 \ref{cohomology-section-global-RHom} 節の構成により
$$
R\Hom_X(L, a(K)) =
R\Gamma(X, R\SheafHom_{\mathcal{O}_X}(L, a(K))) =
R\Gamma(Y, Rf_*R\SheafHom_{\mathcal{O}_X}(L, a(K)))
$$
および
$$
R\Hom_Y(Rf_*L, K) = R\Gamma(Y, R\SheafHom_{\mathcal{O}_Y}(Rf_*L, K))
$$
を得る。従って補題は補題 \ref{duality-lemma-iso-on-RSheafHom} の帰結である。実際、
$D(\mathcal{O}_Y)$ における写像 $E \to E'$ が同型 $DQ_Y(E) \to DQ_Y(E')$ を誘導するならば、
それは擬同型 $R\Gamma(Y, E) \to R\Gamma(Y, E')$ を誘導する。なぜなら
$\mathcal{O}_Y[-i]$ は $D_\QCoh(\mathcal{O}_Y)$ に属し、$DQ_Y$ は包含関手
$D_\QCoh(\mathcal{O}_Y) \to D(\mathcal{O}_Y)$ の右随伴なので、
$H^i(Y, E) = \Ext^i_Y(\mathcal{O}_Y, E) = \Hom(\mathcal{O}_Y[-i], E) =
\Hom(\mathcal{O}_Y[-i], DQ_Y(E))$ だからである。
\end{proof}







\section{押し出しの右随伴と開部分スキームへの制限}
\label{duality-section-restriction-to-opens}

\noindent
本節では、押し出しの右随伴が開部分スキームへの制限とどの程度可換するかを調べる。
これは基底変換の問題なので、まずより一般的に論じる。

\medskip\noindent
押し出しの右随伴が基底変換と可換するかを知りたいことが多い。そこで、準コンパクトかつ
準分離なスキームからなるデカルト平方
\begin{equation}
\label{duality-equation-base-change}
\vcenter{
\xymatrix{
X' \ar[r]_{g'} \ar[d]_{f'} & X \ar[d]^f \\
Y' \ar[r]^g & Y
}
}
\end{equation}
を考える。
$$
a  : D_\QCoh(\mathcal{O}_Y) \to D_\QCoh(\mathcal{O}_X)
\quad\text{and}\quad
a'  : D_\QCoh(\mathcal{O}_{Y'}) \to D_\QCoh(\mathcal{O}_{X'})
$$
を $Rf_*$ および $Rf'_*$ の右随伴とする（補題 \ref{duality-lemma-twisted-inverse-image}）。『層のコホモロジー』の注意 \ref{cohomology-remark-base-change} の基底変換写像を考える。それは準連接コホモロジーをもつ
層の導来圏上の関手の変換
$$
Lg^* \circ Rf_* \longrightarrow Rf'_* \circ L(g')^*
$$
を誘導し、従って反対方向の右随伴間の変換
$$
a \circ Rg_* \longleftarrow Rg'_* \circ a'
$$
を誘導する。

\begin{lemma}
\label{duality-lemma-flat-precompose-pus}
図式 (\ref{duality-equation-base-change}) において $g$ が平坦であるか、より一般に $f$ と $g$ が
Tor 独立であると仮定する。このとき $a \circ Rg_* \leftarrow Rg'_* \circ a'$ は同型である。
\end{lemma}

\begin{proof}
この場合『スキームの導来圏』の補題 \ref{perfect-lemma-compare-base-change} により、任意の $D_\QCoh(\mathcal{O}_X)$ の $K$ に対して
基底変換写像 $Lg^* \circ Rf_* K \longrightarrow Rf'_* \circ L(g')^*K$ は同型である。
従って対応する随伴関手間の変換も同型である。
\end{proof}

\noindent
$f : X \to Y$ を準コンパクトかつ準分離なスキームの射とする。$V \subset Y$ を
準コンパクトな開部分スキームとし、$U = f^{-1}(V)$ とおく。これにより
(\ref{duality-equation-base-change}) のようなデカルト平方
$$
\xymatrix{
U \ar[r]_{j'} \ar[d]_{f|_U} & X \ar[d]^f \\
V \ar[r]^j & Y
}
$$
を得る。補題 \ref{duality-lemma-flat-precompose-pus} により、$a$ と $a'$ をそれぞれ
$Rf_*$ と $R(f|_U)_*$ の右随伴とすると、写像
$\xi : a \circ Rj_* \leftarrow Rj'_* \circ a'$ は同型である。関手
$D_\QCoh(\mathcal{O}_Y) \to D_\QCoh(\mathcal{O}_U)$ の変換
\begin{equation}
\label{duality-equation-sheafy}
(j')^* \circ a \to
(j')^* \circ a \circ Rj_* \circ j^* \xrightarrow{\xi^{-1}}
(j')^* \circ Rj'_* \circ a' \circ j^* \to a' \circ j^*
\end{equation}
を得る。ここで第一の矢印は $\text{id} \to Rj_* \circ j^*$ から、最後の矢印は同型
$(j')^* \circ Rj'_* \to \text{id}$ から得られる。特に、$K$ において評価した
(\ref{duality-equation-sheafy}) が同型であることと、$a(K)|_U \to a(Rj_*(K|_V))|_U$ が同型であることは同値である。

\begin{example}
\label{duality-example-not-supported-on-inverse-image}
Noether スキームの有限射 $f : X \to Y$ であって、ある
$K \in D_{\textit{Coh}}(\mathcal{O}_Y)$ 上で評価した (\ref{duality-equation-sheafy}) が
同型でないものが存在する。実際、$k$ を体、$\epsilon^2 = 0$ とし、
$A = k[x, \epsilon]$, $B = k[x] = A/(\epsilon)$ により
$X = \Spec(B) \to Y = \Spec(A)$ とする。$n \in \mathbf{N}$ に対して
$M_n = A/(\epsilon, x^n)$ とおく。$B$ は $\epsilon$ 倍写像からなる自由周期分解
$\ldots \to A \to A \to A$ をもつので
$$
\Ext^i_A(B, M_n) = M_n,\quad i \geq 0
$$
である。$D_{\textit{Coh}}(A)$ の対象 $K = \bigoplus M_n[n] = \prod M_n[n]$ を考える
（$D(A)$ における等号は『導来圏』の補題 \ref{derived-lemma-direct-sums} および \ref{derived-lemma-products} による）。例 \ref{duality-example-affine-twisted-inverse-image} により $a(K)$ は $R\Hom(B, K)$ に対応し、上の計算から
$$
H^0(R\Hom(B, K)) = \Ext^0_A(B, K) =
\prod\nolimits_{n \geq 1} \Ext^n_A(B. M_n) = 
\prod\nolimits_{n \geq 1} M_n
$$
を得る。しかしこの加群には $x$ のどの冪によっても消されない元がある一方、複体 $K$ の
コホモロジーの各元は $x$ のある冪で消される。言い換えると、$V = D(x)$,
$U = D(x)$ とした写像 (\ref{duality-equation-sheafy}) は複体 $K$ 上で同型にはなりえない。
実際、$(j')^*(a(K))$ は零でなく、$a'(j^*K)$ は零である。
\end{example}

\begin{lemma}
\label{duality-lemma-when-sheafy}
$f : X \to Y$ を準コンパクトかつ準分離なスキームの射とし、$a$ を
$Rf_* : D_\QCoh(\mathcal{O}_X) \to D_\QCoh(\mathcal{O}_Y)$ の右随伴とする。
$V \subset Y$ を準コンパクトな開集合とし、その逆像を $U \subset X$ とする。
\begin{enumerate}
\item $Y \setminus V$ に台をもつ任意の $Q \in D_\QCoh^+(\mathcal{O}_Y)$ に対して
$a(Q)$ が $X \setminus U$ に台をもつことと、$D_\QCoh^+(\mathcal{O}_Y)$ のすべての
$K$ 上で (\ref{duality-equation-sheafy}) が同型であることは同値である。
\item $Y \setminus V$ に台をもつ任意の $Q \in D_\QCoh(\mathcal{O}_Y)$ に対して
$a(Q)$ が $X \setminus U$ に台をもつことと、$D_\QCoh(\mathcal{O}_Y)$ のすべての
$K$ 上で (\ref{duality-equation-sheafy}) が同型であることは同値である。
\item $a$ が直和と可換ならば、(1) の同値な条件は (2) の同値な条件を含意する。
\end{enumerate}
\end{lemma}

\begin{proof}
(1) を示す。$K \in D_\QCoh^+(\mathcal{O}_Y)$ とし、区別三角
$$
K \to Rj_*K|_V \to Q \to K[1]
$$
を選ぶ。$Q$ は $D_\QCoh^+(\mathcal{O}_Y)$ に属し（『スキームの導来圏』の補題 \ref{perfect-lemma-quasi-coherence-direct-image}）、$Y \setminus V$ に台をもつ
（『スキームの導来圏』の定義 \ref{perfect-definition-supported-on}）。
$a$ を適用すると $X$ 上の区別三角
$$
a(K) \to a(Rj_*K|_V) \to a(Q) \to a(K)[1]
$$
を得る。$a(Q)$ が $X \setminus U$ に台をもつならば、$U$ に制限した写像
$a(K)|_U \to a(Rj_*K|_V)|_U$ は同型、すなわち $K$ 上で
(\ref{duality-equation-sheafy}) は同型である。逆は直ちに従う。

\medskip\noindent
(2) の証明は (1) の証明とまったく同じである。

\medskip\noindent
(3) を示す。(1) の同値な条件が成り立つと仮定し、$T = Y \setminus V$ とおく。
$T$ および $f^{-1}(T)$ にコホモロジー層の台をもつ複体の圏をそれぞれ
$D_{\QCoh, T}(\mathcal{O}_Y)$ および $D_{\QCoh, f^{-1}(T)}(\mathcal{O}_X)$ と表す。
$a$ は直和と可換するので、対象
$$
\{Q \in D_{\QCoh, T}(\mathcal{O}_Y) \mid
a(Q) \in D_{\QCoh, f^{-1}(T)}(\mathcal{O}_X)\}
$$
からなる真に充満な飽和三角部分圏 $\mathcal{D}$ は直和で保たれ、従って導来余極限でも保たれる。
一方、圏 $D_{\QCoh, T}(\mathcal{O}_Y)$ は完全対象 $E$ により生成される
（『スキームの導来圏』の補題 \ref{perfect-lemma-generator-with-support} を参照）。
仮定により $E \in \mathcal{D}$ である。『導来圏』の補題 \ref{derived-lemma-write-as-colimit} により、$D_{\QCoh, T}(\mathcal{O}_Y)$ の任意の対象 $Q$ は、
遷移写像の錐が $E$ のシフトの直和である系 $Q_1 \to Q_2 \to Q_3 \to \ldots$ の導来余極限である。
帰納法によりすべての $n$ に対して $Q_n \in \mathcal{D}$ であり、最後に
$Q$ も $\mathcal{D}$ に属する。従って (2) の同値な条件が成り立つ。
\end{proof}

\begin{lemma}
\label{duality-lemma-proper-noetherian}
$Y$ を準コンパクトかつ準分離なスキームとし、$f : X \to Y$ を固有射とする。
次のいずれかを仮定する。\footnote{この証明は、$X$ 上の任意の完全な $P$ に対して
$Rf_*P$ が擬連接となる準コンパクトかつ準分離なスキームの射に対して成り立つ。
Kiehl の定理 \cite{Kiehl} から、$f$ が固有かつ擬連接ならばこの条件が成り立つことが容易に従う。
これが本補題および本章の他のいくつかの結果に対する正しい一般性である。}
\begin{enumerate}
\item $f$ は平坦かつ有限表示である。
\item $Y$ は Noether である。
\end{enumerate}
このとき $Y$ のすべての準コンパクト開集合 $V$ に対して、補題 \ref{duality-lemma-when-sheafy} の (1) の同値な条件が成り立つ。
\end{lemma}

\begin{proof}
$Q \in D^+_\QCoh(\mathcal{O}_Y)$ が $Y \setminus V$ に台をもつとする。背理法により、
$a(Q)$ が $X \setminus U$ に台をもたないと仮定する。このとき $U$ 上の完全複体 $P_U$ と
零でない写像 $P_U \to a(Q)|_U$ が存在する（『スキームの導来圏』の定理 \ref{perfect-theorem-bondal-van-den-Bergh} から従う）。さらに『スキームの導来圏』の補題 \ref{perfect-lemma-lift-map-from-perfect-complex-with-support} を用いると、$X$ 上の完全複体
$P$ と、$U$ への制限が零でない写像 $P \to a(Q)$ が存在すると仮定してよい。
$a$ の定義により、この写像は写像 $Rf_*P \to Q$ に随伴する。

\medskip\noindent
複体 $Rf_*P$ は擬連接である。(1) の場合、これは『スキームの導来圏』の補題 \ref{perfect-lemma-flat-proper-pseudo-coherent-direct-image-general} から従う。(2) の場合、これは
『スキームの導来圏』の補題 \ref{perfect-lemma-direct-image-coherent} および \ref{perfect-lemma-identify-pseudo-coherent-noetherian} から従う。従って『スキームの導来圏』の補題 \ref{perfect-lemma-map-from-pseudo-coherent-to-complex-with-support} を適用すると、
$V$ への制限が同型であり、合成
$I \otimes^\mathbf{L}_{\mathcal{O}_Y} Rf_*P \to Rf_*P \to Q$ が零となる完全複体の写像
$I \to \mathcal{O}_Y$ を得る。『スキームの導来圏』の補題 \ref{perfect-lemma-cohomology-base-change} により
$I \otimes^\mathbf{L}_{\mathcal{O}_Y} Rf_*P =
Rf_*(Lf^*I \otimes^\mathbf{L}_{\mathcal{O}_X} P)$ である。従って合成
$$
Lf^*I \otimes^\mathbf{L}_{\mathcal{O}_X} P \to P \to a(Q)
$$
は零である。しかし $U$ への制限は、零でないと仮定した写像 $P|_U \to a(Q)|_U$ である。
この矛盾により証明が完了する。
\end{proof}












\section{押し出しの右随伴と基底変換 I}
\label{duality-section-base-change-map}

\noindent
写像 (\ref{duality-equation-sheafy}) は基底変換写像の特別な場合である。実際、
(\ref{duality-equation-base-change}) のような、準コンパクトかつ準分離なスキームのデカルト図式
$$
\xymatrix{
X' \ar[r]_{g'} \ar[d]_{f'} & X \ar[d]^f \\
Y' \ar[r]^g & Y
}
$$
があり、$f$ と $g$ が {\bf Tor 独立}であると仮定する。このとき合成で与えられる関手
$D_\QCoh(\mathcal{O}_Y) \to D_\QCoh(\mathcal{O}_{X'})$ の射
\begin{equation}
\label{duality-equation-base-change-map}
L(g')^* \circ a \to
L(g')^* \circ a \circ Rg_* \circ Lg^* \leftarrow
L(g')^* \circ Rg'_* \circ a' \circ Lg^* \to a' \circ Lg^*
\end{equation}
を考えられる。第一の矢印は随伴写像 $\text{id} \to Rg_* Lg^*$ から、最後の矢印は随伴写像
$L(g')^*Rg'_* \to \text{id}$ から得られる。中央の矢印を反転するために Tor 独立性の仮定が必要である；
補題 \ref{duality-lemma-flat-precompose-pus} を参照する。別の見方として、$L(g')^*$ と $R(g')_*$ の
随伴性により (\ref{duality-equation-base-change-map}) を自然変換
$$
a \to a \circ Rg_* \circ Lg^* \leftarrow Rg'_* \circ a' \circ Lg^*
$$
と考えることができ、ここでも第二の矢印は可逆である。$M \in D_\QCoh(\mathcal{O}_X)$ および
$K \in D_\QCoh(\mathcal{O}_Y)$ とすると、Yoneda 関手上でこの写像は
\begin{align*}
\Hom_X(M, a(K))
& =
\Hom_Y(Rf_*M, K) \\
& \to
\Hom_Y(Rf_*M, Rg_* Lg^*K) \\
& =
\Hom_{Y'}(Lg^*Rf_*M, Lg^*K) \\
& \leftarrow
\Hom_{Y'}(Rf'_* L(g')^*M, Lg^*K) \\
& =
\Hom_{X'}(L(g')^*M, a'(Lg^*K)) \\
& =
\Hom_X(M, Rg'_*a'(Lg^*K))
\end{align*}
で与えられる（左向きの矢印は『スキームの導来圏』の補題 \ref{perfect-lemma-compare-base-change} の基底変換定理により可逆である）。これは構成をやや明示的にする。

\medskip\noindent
本節ではまず、通常の基底変換写像に対する『層のコホモロジー』の注意 \ref{cohomology-remark-compose-base-change} および \ref{cohomology-remark-compose-base-change-horizontal} と同様に、平方を重ねることに関する
自然な両立性を基底変換写像が満たすことを示す。初読時には本節の残りを飛ばすことを勧める。

\begin{lemma}
\label{duality-lemma-compose-base-change-maps}
準コンパクトかつ準分離なスキームの可換図式
$$
\xymatrix{
X' \ar[r]_k \ar[d]_{f'} & X \ar[d]^f \\
Y' \ar[r]^l \ar[d]_{g'} & Y \ar[d]^g \\
Z' \ar[r]^m & Z
}
$$
を考える。二つの平方はともにデカルトであり、$f$ と $l$、および $g$ と $m$ は Tor 独立であるとする。
このとき二つの平方に対する写像 (\ref{duality-equation-base-change-map}) の合成は、外側の長方形に対する
基底変換写像である（正確な主張は証明を参照）。
\end{lemma}

\begin{proof}
仮定から $g \circ f$ と $m$ は Tor 独立であることが従う（詳細は省略する）ので、主張には意味がある。
本証明では $Lk^*$ の代わりに $k^*$、$Rf_*$ の代わりに $f_*$ と書く。$f$, $g$, $g \circ f$
に対する補題 \ref{duality-lemma-twisted-inverse-image} の右随伴をそれぞれ $a$, $b$, $c$ とし、
プライム付きのものも同様に定める。上側の平方に対応する矢印は合成
$$
\gamma_{top} :
k^* \circ a \to k^* \circ a \circ l_* \circ l^*
\xleftarrow{\xi_{top}} k^* \circ k_* \circ a' \circ l^* \to a' \circ l^*
$$
である。ここで $\xi_{top} : k_* \circ a' \to a \circ l_*$ は同型（従って反転可能）であり、
基底変換写像 $l^* \circ f_* \to f'_* \circ k^*$ に「双対」な矢印である。外側の矢印は標準写像
$1 \to l_* \circ l^*$ および $k^* \circ k_* \to 1$ から得られる。同様に第二の平方に対して
$$
\gamma_{bot} :
l^* \circ b \to l^* \circ b \circ m_* \circ m^*
\xleftarrow{\xi_{bot}} l^* \circ l_* \circ b' \circ m^* \to b' \circ m^*
$$
を得る。外側の長方形に対しては
$$
\gamma_{rect} :
k^* \circ c \to k^* \circ c \circ m_* \circ m^*
\xleftarrow{\xi_{rect}} k^* \circ k_* \circ c' \circ m^* \to c' \circ m^*
$$
を得る。$(g \circ f)_* = g_* \circ f_*$ なので $c = a \circ b$ であり、同様に
$c' = a' \circ b'$ である。補題の主張は、$\gamma_{rect}$ が合成
$$
k^* \circ c = k^* \circ a \circ b \xrightarrow{\gamma_{top}}
a' \circ l^* \circ b \xrightarrow{\gamma_{bot}}
a' \circ b' \circ m^* = c' \circ m^*
$$
に等しいということである。これを見るため次の図式を考える。
$$
\xymatrix{
& & k^* \circ a \circ b \ar[d] \ar[lldd] \\
& & k^* \circ a \circ l_* \circ l^* \circ b \ar[ld] \\
k^* \circ a \circ b \circ m_* \circ m^* \ar[r] &
k^* \circ a \circ l_* \circ l^* \circ b \circ m_* \circ m^* &
k^* \circ k_* \circ a' \circ l^* \circ b \ar[u]_{\xi_{top}} \ar[d] \ar[ld] \\
& k^*\circ k_* \circ a' \circ l^* \circ b \circ m_* \circ m^*
\ar[u]_{\xi_{top}} \ar[rd] &
a' \circ l^* \circ b \ar[d] \\
k^* \circ k_* \circ a' \circ b' \circ m^* \ar[uu]_{\xi_{rect}} \ar[ddrr] &
k^*\circ k_* \circ a' \circ l^* \circ l_* \circ b' \circ m^*
\ar[u]_{\xi_{bot}} \ar[l] \ar[dr] &
a' \circ l^* \circ b \circ m_* \circ m^* \\
& & a' \circ l^* \circ l_* \circ b' \circ m^* \ar[u]_{\xi_{bot}} \ar[d] \\
& & a' \circ b' \circ m^*
}
$$
右辺を下へたどると上の合成となり、左辺を下へたどると $\gamma_{rect}$ となる。この図式の右側の
すべての四辺形は『圏論』の補題 \ref{categories-lemma-properties-2-cat-cats}、またはより直接には
『圏論』の定義 \ref{categories-definition-horizontal-composition} に先立つ議論により可換である。
従って図式
$$
\xymatrix{
a \circ l_* \circ l^* \circ b \circ m_* &
a \circ b \circ m_* \ar[l] \\
k_* \circ a' \circ l^* \circ b \circ m_* \ar[u]_{\xi_{top}} & \\
k_* \circ a' \circ l^* \circ l_* \circ b' \ar[u]_{\xi_{bot}} \ar[r] &
k_* \circ a' \circ b' \ar[uu]_{\xi_{rect}}
}
$$
が矢印 $\xi_{top}$, $\xi_{bot}$, $\xi_{rect}$ を反転した後に可換となることを示せば十分である
（これは元の図式が可換であることを求めるのとは異なる）。一方、図式
$$
\xymatrix{
& a \circ l_* \circ l^* \circ b \circ m_* \\
a \circ l_* \circ l^* \circ l_* \circ b'
\ar[ru]^{\xi_{bot}} & &
k_* \circ a' \circ l^* \circ b \circ m_* \ar[ul]_{\xi_{top}} \\
& k_* \circ a' \circ l^* \circ l_* \circ b'
\ar[ul]^{\xi_{top}} \ar[ur]_{\xi_{bot}}
}
$$
は『圏論』の補題 \ref{categories-lemma-properties-2-cat-cats} により可換である。また図式
$$
\vcenter{
\xymatrix{
a \circ l_* \circ l^* \circ b \circ m_* & a \circ b \circ m \ar[l] \\
a \circ l_* \circ l^* \circ l_* \circ b' \ar[u] &
a \circ l_* \circ b' \ar[l] \ar[u]
}
}
\quad\text{and}\quad
\vcenter{
\xymatrix{
a \circ l_* \circ l^* \circ l_* \circ b' \ar[r] & a \circ l_* \circ b' \\
k_* \circ a' \circ l^* \circ l_* \circ b' \ar[u] \ar[r] &
k_* \circ a' \circ b' \ar[u]
}
}
$$
は可換であり（引用した文献を参照）、合成 $l_* \to l_* \circ l^* \circ l_* \to l_*$ は恒等写像である。
従って
$$
k \circ a' \circ b' \xrightarrow{\xi_{bot}} a \circ l_* \circ b
\xrightarrow{\xi_{top}} a \circ b \circ m_*
$$
が、同一視 $a \circ b = c$ および $a' \circ b' = c'$ のもとで $\xi_{rect}$ に等しいことを示せばよい。
これは『層のコホモロジー』の注意 \ref{cohomology-remark-compose-base-change} の主張の双対であり、証明が完了する。
\end{proof}

\begin{lemma}
\label{duality-lemma-compose-base-change-maps-horizontal}
準コンパクトかつ準分離なスキームの可換図式
$$
\xymatrix{
X'' \ar[r]_{g'} \ar[d]_{f''} & X' \ar[r]_g \ar[d]_{f'} & X \ar[d]^f \\
Y'' \ar[r]^{h'} & Y' \ar[r]^h & Y
}
$$
を考える。二つの平方はともにデカルトであり、$f$ と $h$、および $f'$ と $h'$ は
Tor 独立であるとする。このとき二つの平方に対する写像
(\ref{duality-equation-base-change-map}) の合成は、外側の長方形に対する基底変換写像である
（正確な主張は証明を参照）。
\end{lemma}

\begin{proof}
仮定から $f$ と $h \circ h'$ は Tor 独立であることが従う（詳細は省略する）ので、主張には意味がある。
本証明では $Lg^*$ の代わりに $g^*$、$Rf_*$ の代わりに $f_*$ と書く。$f$, $f'$, $f''$ に対する
補題 \ref{duality-lemma-twisted-inverse-image} の右随伴をそれぞれ $a$, $a'$, $a''$ とする。
右側の平方に対応する矢印は合成
$$
\gamma_{right} :
g^* \circ a \to g^* \circ a \circ h_* \circ h^*
\xleftarrow{\xi_{right}} g^* \circ g_* \circ a' \circ h^* \to a' \circ h^*
$$
である。ここで $\xi_{right} : g_* \circ a' \to a \circ h_*$ は同型（従って反転可能）であり、
基底変換写像 $h^* \circ f_* \to f'_* \circ g^*$ に「双対」な矢印である。外側の矢印は標準写像
$1 \to h_* \circ h^*$ および $g^* \circ g_* \to 1$ から得られる。同様に左側の平方に対して
$$
\gamma_{left} :
(g')^* \circ a' \to (g')^* \circ a' \circ (h')_* \circ (h')^*
\xleftarrow{\xi_{left}}
(g')^* \circ (g')_* \circ a'' \circ (h')^* \to a'' \circ (h')^*
$$
を得る。外側の長方形に対しては
$$
\gamma_{rect} :
k^* \circ a \to
k^* \circ a \circ m_* \circ m^* \xleftarrow{\xi_{rect}}
k^* \circ k_* \circ a'' \circ m^* \to
a'' \circ m^*
$$
を得る。ここで $k = g \circ g'$ および $m = h \circ h'$ である。
$k^* = (g')^* \circ g^*$ および $m^* = (h')^* \circ h^*$ である。補題の主張は、
$\gamma_{rect}$ が合成
$$
k^* \circ a =
(g')^* \circ g^* \circ a \xrightarrow{\gamma_{right}}
(g')^* \circ a' \circ h^* \xrightarrow{\gamma_{left}}
a'' \circ (h')^* \circ h^* = a'' \circ m^*
$$
に等しいということである。これを見るため次の図式を考える。
$$
\xymatrix{
& (g')^* \circ g^* \circ a \ar[d] \ar[ddl] \\
& (g')^* \circ g^* \circ a \circ h_* \circ h^* \ar[ld] \\
(g')^* \circ g^* \circ a \circ h_* \circ (h')_* \circ (h')^* \circ h^* &
(g')^* \circ g^* \circ g_* \circ a' \circ h^*
\ar[u]_{\xi_{right}} \ar[d] \ar[ld] \\
(g')^* \circ g^* \circ g_* \circ a' \circ (h')_* \circ (h')^* \circ h^*
\ar[u]_{\xi_{right}} \ar[dr] &
(g')^* \circ a' \circ h^* \ar[d] \\
(g')^* \circ g^* \circ g_* \circ (g')_* \circ a'' \circ (h')^* \circ h^*
\ar[u]_{\xi_{left}} \ar[ddr] \ar[dr] &
(g')^* \circ a' \circ (h')_* \circ (h')^* \circ h^* \\
& (g')^*\circ (g')_* \circ a'' \circ (h')^* \circ h^*
\ar[u]_{\xi_{left}} \ar[d] \\
& a'' \circ (h')^* \circ h^*
}
$$
右辺を下へたどると上の合成となり、左辺を下へたどると $\gamma_{rect}$ となる。この図式の右側の
すべての四辺形は『圏論』の補題 \ref{categories-lemma-properties-2-cat-cats}、またはより直接には
『圏論』の定義 \ref{categories-definition-horizontal-composition} に先立つ議論により可換である。
従って
$$
g_* \circ (g')_* \circ a'' \xrightarrow{\xi_{left}}
g_* \circ a' \circ (h')_* \xrightarrow{\xi_{right}}
a \circ h_* \circ (h')_*
$$
が $\xi_{rect}$ に等しいことを示せば十分である。これは『層のコホモロジー』の注意 \ref{cohomology-remark-compose-base-change-horizontal} の主張の双対であり、証明が完了する。
\end{proof}

\begin{remark}
\label{duality-remark-going-around}
準コンパクトかつ準分離なスキームの可換図式
$$
\xymatrix{
X'' \ar[r]_{k'} \ar[d]_{f''} & X' \ar[r]_k \ar[d]_{f'} & X \ar[d]^f \\
Y'' \ar[r]^{l'} \ar[d]_{g''} & Y' \ar[r]^l \ar[d]_{g'} & Y \ar[d]^g \\
Z'' \ar[r]^{m'} & Z' \ar[r]^m & Z
}
$$
を考える。すべての平方はデカルトであり、$(f, l)$, $(g, m)$, $(f', l')$,
$(g', m')$ は Tor 独立な射の組である。$f$, $f'$, $f''$, $g$, $g'$, $g''$ に対する
補題 \ref{duality-lemma-twisted-inverse-image} の右随伴を $a$, $a'$, $a''$, $b$, $b'$, $b''$ とする。
図式の平方を次のように $A$, $B$, $C$, $D$ と名付ける。
$$
\begin{matrix}
A & B \\
C & D
\end{matrix}
$$
各平方に対する写像 (\ref{duality-equation-base-change-map}) は次のとおりである
（ここでは $k^* = Lk^*$ などと書く）。
$$
\begin{matrix}
\gamma_A : (k')^* \circ a' \to a'' \circ (l')^* &
\gamma_B : k^* \circ a \to a' \circ l^* \\
\gamma_C : (l')^* \circ b' \to b'' \circ (m')^* &
\gamma_D : l^* \circ b \to b' \circ m^*
\end{matrix}
$$
$2 \times 1$ および $1 \times 2$ の長方形に対して、さらに四つの基底変換写像がある。
$$
\begin{matrix}
\gamma_{A + B} : (k \circ k')^* \circ a \to a'' \circ (l \circ l')^* \\
\gamma_{C + D} : (l \circ l')^* \circ b \to b'' \circ (m \circ m')^* \\
\gamma_{A + C} : (k')^* \circ (a' \circ b') \to (a'' \circ b'') \circ (m')^* \\
\gamma_{B + D} : k^* \circ (a \circ b) \to (a' \circ b') \circ m^*
\end{matrix}
$$
補題 \ref{duality-lemma-compose-base-change-maps-horizontal} により
$$
\gamma_{A + B} = \gamma_A \circ \gamma_B, \quad
\gamma_{C + D} = \gamma_C \circ \gamma_D
$$
であり、補題 \ref{duality-lemma-compose-base-change-maps} により
$$
\gamma_{A + C} = \gamma_C \circ \gamma_A, \quad
\gamma_{B + D} = \gamma_D \circ \gamma_B
$$
である。より正確には『圏論』の第 \ref{categories-section-formal-cat-cat} 節の記法を用いて
$\gamma_{A + B} = (\gamma_A \star \text{id}_{l^*}) \circ
(\text{id}_{(k')^*} \star \gamma_B)$ と書くべきであり、他も同様である。しかし変換の始域と終域が
分かれば補うべきものが定まるので、補題 \ref{duality-lemma-compose-base-change-maps} および \ref{duality-lemma-compose-base-change-maps-horizontal} の証明と同様に、恒等変換との $\star$ 積を省略する記法を続ける。
以上から、先験的には二つの変換
$$
(k')^* \circ k^* \circ a \circ b
\longrightarrow
a'' \circ b'' \circ (m')^* \circ m^*
$$
すなわち
$$
\gamma_C \circ \gamma_A \circ \gamma_D \circ \gamma_B =
\gamma_{A + C} \circ \gamma_{B + D}
$$
および
$$
\gamma_C \circ \gamma_D \circ \gamma_A \circ \gamma_B =
\gamma_{C + D} \circ \gamma_{A + B}
$$
を得る。本注意の要点は、これらの変換が等しいということである。そのためには
$$
\xymatrix{
(k')^* \circ a' \circ l^* \circ b \ar[r]_{\gamma_D} \ar[d]_{\gamma_A} &
(k')^* \circ a' \circ b' \circ m^* \ar[d]^{\gamma_A} \\
a'' \circ (l')^* \circ l^* \circ b \ar[r]^{\gamma_D} &
a'' \circ (l')^* \circ b' \circ m^*
}
$$
が可換であることを示せば十分である。これは『圏論』の補題 \ref{categories-lemma-properties-2-cat-cats}、またはより直接には『圏論』の定義 \ref{categories-definition-horizontal-composition} に先立つ議論から従う。
\end{remark}







\section{押し出しの右随伴と基底変換 II}
\label{duality-section-base-change-II}

\noindent
本節では、第 \ref{duality-section-base-change-map} 節の基底変換写像がいくつかの場合に同型であることを示す。
まず、押し出しの右随伴の形成が開集合への制限と可換するならば、アフィン開集合上で確認すれば十分であることを述べる。

\begin{remark}
\label{duality-remark-check-over-affines}
準コンパクトかつ準分離なスキームのデカルト図式
$$
\xymatrix{
X' \ar[r]_{g'} \ar[d]_{f'} & X \ar[d]^f \\
Y' \ar[r]^g & Y
}
$$
を考え、$(g, f)$ は Tor 独立であるとする。$V \subset Y$ および $V' \subset Y'$ を
$g(V') \subset V$ を満たすアフィン開集合とする。デカルト図式
$$
\vcenter{
\xymatrix{
U \ar[r] \ar[d] & X \ar[d] \\
V \ar[r] & Y
}
}
\quad\text{and}\quad
\vcenter{
\xymatrix{
U' \ar[r] \ar[d] & X' \ar[d] \\
V' \ar[r] & Y'
}
}
$$
を作る。$K$ と第一の図式に関する (\ref{duality-equation-sheafy})、および $Lg^*K$ と第二の図式に関する
(\ref{duality-equation-sheafy}) が同型であると仮定する。このとき基底変換写像
(\ref{duality-equation-base-change-map})
$$
L(g')^*a(K) \longrightarrow a'(Lg^*K)
$$
の $U'$ への制限は、$K|_V$ とデカルト図式
$$
\xymatrix{
U' \ar[r] \ar[d] & U \ar[d] \\
V' \ar[r] & V
}
$$
に対する基底変換写像 (\ref{duality-equation-base-change-map}) と同型である。これは
(\ref{duality-equation-sheafy}) が基底変換写像 (\ref{duality-equation-base-change-map}) の特別な場合であり、
平方を水平方向に重ねたとき基底変換写像が正しく合成されることから従う；補題 \ref{duality-lemma-compose-base-change-maps-horizontal} を参照する。従って、$U'$ に制限した基底変換写像が
同型であることを確認するには、最後の図式で作業すれば十分である。
\end{remark}

\begin{lemma}
\label{duality-lemma-more-base-change}
図式 (\ref{duality-equation-base-change}) において次を仮定する。
\begin{enumerate}
\item $g : Y' \to Y$ はアフィンスキームの射である。
\item $f : X \to Y$ は固有である。
\item $f$ と $g$ は Tor 独立である。
\end{enumerate}
このとき基底変換写像 (\ref{duality-equation-base-change-map}) は、次の場合に同型
$$
L(g')^*a(K) \longrightarrow a'(Lg^*K)
$$
を誘導する。
\begin{enumerate}
\item $f$ が平坦かつ有限表示ならば、すべての $K \in D_\QCoh(\mathcal{O}_Y)$ に対して。
\item $f$ が完全で $Y$ が Noether ならば、すべての $K \in D_\QCoh(\mathcal{O}_Y)$ に対して。
\item $g$ が有限 Tor 次元をもち $Y$ が Noether ならば、$K \in D_\QCoh^+(\mathcal{O}_Y)$ に対して。
\end{enumerate}
\end{lemma}

\begin{proof}
$Y = \Spec(A)$, $Y' = \Spec(A')$ と書く。アフィン射の基底変換なので $g'$ はアフィンである。
$M$ を $D_\QCoh(\mathcal{O}_X)$ の完全生成対象とする；『スキームの導来圏』の定理 \ref{perfect-theorem-bondal-van-den-Bergh} を参照する。このとき $L(g')^*M$ は
$D_\QCoh(\mathcal{O}_{X'})$ の生成対象である；『スキームの導来圏』の注意 \ref{perfect-remark-pullback-generator} を参照する。従って (\ref{duality-equation-base-change-map}) が大域 Hom 複体の同型
\begin{equation}
\label{duality-equation-iso}
R\Hom_{X'}(L(g')^*M, L(g')^*a(K))
\longrightarrow
R\Hom_{X'}(L(g')^*M, a'(Lg^*K))
\end{equation}
を誘導することを示せば十分である；『層のコホモロジー』の第 \ref{cohomology-section-global-RHom} 節を参照する。実際、これにより
$L(g')^*a(K) \to a'(Lg^*K)$ の錐が零であることが従う。証明の構成は次のとおりである。
まずこれらの Hom 複体が同型であることを示し、証明の最後でその同型が
(\ref{duality-equation-iso}) により誘導されることを示す。

\medskip\noindent
左辺について考える。$M$ は完全なので、標準写像
$$
R\Hom_X(M, a(K)) \otimes^\mathbf{L}_A A'
\longrightarrow
R\Hom_{X'}(L(g')^*M, L(g')^*a(K))
$$
は『スキームの導来圏』の補題 \ref{perfect-lemma-affine-morphism-and-hom-out-of-perfect} により同型である。これを補題 \ref{duality-lemma-iso-global-hom} の同型 $R\Hom_Y(Rf_*M, K) = R\Hom_X(M, a(K))$
と組み合わせると、左辺は $R\Hom_Y(Rf_*M, K) \otimes^\mathbf{L}_A A'$ に等しい。

\medskip\noindent
右辺について考える。まず補題 \ref{duality-lemma-iso-global-hom} の同型
$$
R\Hom_{X'}(L(g')^*M, a'(Lg^*K)) = R\Hom_{Y'}(Rf'_*L(g')^*M, Lg^*K)
$$
を用いる。次に『スキームの導来圏』の補題 \ref{perfect-lemma-compare-base-change} により基底変換写像
$Lg^*Rf_*M \to Rf'_*L(g')^*M$ は同型である。従ってこれは
$R\Hom_{Y'}(Lg^*Rf_*M, Lg^*K)$ と書き直せる。$Y$, $Y'$ はアフィンであり、$K$, $Rf_*M$ は
$D_\QCoh(\mathcal{O}_Y)$ に属する（『スキームの導来圏』の補題 \ref{perfect-lemma-quasi-coherence-direct-image}）ので、$D(A')$ における標準写像
$$
\beta :
R\Hom_Y(Rf_*M, K) \otimes^\mathbf{L}_A A'
\longrightarrow
R\Hom_{Y'}(Lg^*Rf_*M, Lg^*K)
$$
がある。これは More on Algebra, Equation
(\ref{more-algebra-equation-base-change-RHom}) の矢印である。ここでは『スキームの導来圏』の補題 \ref{perfect-lemma-affine-compare-bounded} および \ref{perfect-lemma-quasi-coherence-internal-hom} を用いて代数との間を行き来した。
\begin{enumerate}
\item $f$ が平坦かつ有限表示ならば、『スキームの導来圏』の補題 \ref{perfect-lemma-flat-proper-perfect-direct-image-general} により複体 $Rf_*M$ は $Y$ 上完全であり、
『代数の続論』の補題 \ref{more-algebra-lemma-base-change-RHom} の (1) により $\beta$ は同型である。
\item $f$ が完全で $Y$ が Noether ならば、『射の続論』の補題 \ref{more-morphisms-lemma-perfect-proper-perfect-direct-image} により複体 $Rf_*M$ は $Y$ 上完全であり、
前と同様に $\beta$ は同型である。
\item $g$ が有限 Tor 次元をもち $Y$ が Noether ならば、『スキームの導来圏』の補題 \ref{perfect-lemma-direct-image-coherent} および \ref{perfect-lemma-identify-pseudo-coherent-noetherian} により複体 $Rf_*M$ は $Y$ 上擬連接であり、
『代数の続論』の補題 \ref{more-algebra-lemma-base-change-RHom} の (4) により $\beta$ は同型である。
\end{enumerate}
従って前段落と同じ結果が得られる。

\medskip\noindent
証明の残りでは、第二段落と第三段落で与えた (\ref{duality-equation-iso}) の左辺と右辺の同一視が、
実際に (\ref{duality-equation-iso}) によって与えられることを示す。式を扱いやすくするため
$(-, -)_X = R\Hom_X(-, -)$ とし、$- \otimes_A^\mathbf{L} A'$ の代わりに
$- \otimes A'$ を用い、$g^* = Lg^*$ および $f_* = Rf_*$ と略記する。次の可換図式を考える。
$$
\xymatrix{
((g')^*M, (g')^*a(K))_{X'} \ar[d] &
(M, a(K))_X \otimes A' \ar[l]^-\alpha \ar[d] &
(f_*M, K)_Y \otimes A' \ar@{=}[l] \ar[d] \\
((g')^*M, (g')^*a(g_*g^*K))_{X'} &
(M, a(g_*g^*K))_X \otimes A' \ar[l]^-\alpha &
(f_*M, g_*g^*K)_Y \otimes A' \ar@{=}[l] \ar@/_4pc/[dd]_{\mu'} \\
((g')^*M, (g')^*g'_*a'(g^*K))_{X'} \ar[u] \ar[d] &
(M, g'_*a'(g^*K))_X \otimes A' \ar[u] \ar[l]^-\alpha \ar[ld]^\mu &
(f_*M, K) \otimes A' \ar[d]^\beta \\
((g')^*M, a'(g^*K))_{X'} &
(f'_*(g')^*M, g^*K)_{Y'} \ar@{=}[l] \ar[r] &
(g^*f_*M, g^*K)_{Y'}
}
$$
$\alpha$ と記した矢印は、頂点 $X', X, Y', Y$ をもつ図式に対する『スキームの導来圏』の補題 \ref{perfect-lemma-affine-morphism-and-hom-out-of-perfect} の写像である。水平矢印は各変数に関して
関手的なので、図式の上部は可換である。中央の垂直矢印は補題 \ref{duality-lemma-flat-precompose-pus} の可逆変換 $g'_* \circ a' \to a \circ g_*$ から得られるので、
中央の平方も可換である。左辺を下へたどると (\ref{duality-equation-iso}) となる。上側の水平矢印は証明の
第二段落で用いた同一視を与え、$\beta$ を含む下側の水平矢印は第三段落で用いた同一視を与える。
$E \in D(A)$, $E' \in D(A')$, $c : E \to E'$ in $D(A)$ とする。係数制限と基底変換の随伴性および
$c$ により誘導される写像を $\mu_c : E \otimes A' \to E'$ と表す。$c$ が明らかなときは $c$ を記法から落として
$\mu = \mu_c$ と書く。図式中の写像 $\mu$ は、同一視
$(M, g'_*a(g^*K))_X = ((g')^*M, a'(g^*K))_{X'}$ により与えられる $c$ に対するこの形の写像である。
$\mu$ を含む三角形は『スキームの導来圏』の注意 \ref{perfect-remark-multiplication-map} により可換である。

\medskip\noindent
次の図式に注目する。
$$
\xymatrix{
(M, a(g_*g^*K))_X &
(f_*M, g_* g^*K)_Y \ar@{=}[l] &
(g^*f_*M, g^*K)_{Y'} \ar@{=}[l] \\
(M, g'_* a'(g^*K))_X \ar[u] &
((g')^*M, a'(g^*K))_{X'} \ar@{=}[l] &
(f'_*(g')^*M, g^*K)_{Y'} \ar@{=}[l] \ar[u]
}
$$
これは変換 $g'_* \circ a' \to a \circ g_*$ の定義そのものにより可換である。同一視
$(f_*M, g_*g^*K)_X = (g^*f_*M, g^*K)_{Y'}$ に対応する上記の $\mu'$ を取れば、六角形も可換する。
従って $\beta$ が $(f_*M, K)_Y \otimes A' \to (f_*M, g_*g^*K)_X \otimes A'$ と $\mu'$ の合成に
等しいことを示せば十分である。そのためには、誘導される二つの写像
$(f_*M, K)_Y \to (g^*f_*M, g^*K)_{Y'}$ が同じであることを示せば十分である。言い換えると、
すべての $E, K \in D(A)$ に対して図式
$$
\xymatrix{
R\Hom_A(E, K) \ar[rr]_{\text{induced by }\beta} \ar[rd] & &
R\Hom_{A'}(E \otimes_A^\mathbf{L} A', K \otimes_A^\mathbf{L} A') \\
& R\Hom_A(E, K \otimes_A^\mathbf{L} A') \ar[ru]
}
$$
が可換であることを示せば十分である。これは『代数の続論』の第 \ref{more-algebra-section-base-change-RHom} 節における $\beta$ の構成そのものであるから、証明が完了する。
\end{proof}








\section{押し出しの右随伴とトレース写像}
\label{duality-section-trace}

\noindent
$f : X \to Y$ を準コンパクトかつ準分離なスキームの射とし、$a : D_\QCoh(\mathcal{O}_Y) \to D_\QCoh(\mathcal{O}_X)$
を補題 \ref{duality-lemma-twisted-inverse-image} の右随伴とする。『圏論』の第 \ref{categories-section-adjoint} 節により関手の変換
$$
\text{Tr}_f : Rf_* \circ a \longrightarrow \text{id}
$$
を得る。$K \in D_\QCoh(\mathcal{O}_Y)$ に対する対応する写像
$\text{Tr}_{f, K} : Rf_*a(K) \longrightarrow K$ は、しばしば {\it トレース写像}と呼ばれる。
これは、右随伴を特徴づける $L \in D_\QCoh(\mathcal{O}_X)$ に対する全単射
$$
\Hom_X(L, a(K)) \longrightarrow \Hom_Y(Rf_*L, K)
$$
が
$$
\varphi \longmapsto \text{Tr}_{f, K} \circ Rf_*\varphi
$$
により与えられるという性質をもつ写像である。写像 (\ref{duality-equation-sheafy-trace})
$$
Rf_*R\SheafHom_{\mathcal{O}_X}(L, a(K))
\longrightarrow
R\SheafHom_{\mathcal{O}_Y}(Rf_*L, K)
$$
は $\text{Tr}_{f, K}$ との合成によって生じる。本節で考えるすべてのトレース写像は、このトレース写像の
特別な場合である。いくつかの特別な場合を論じる前に、トレース写像の形成が基底変換と可換することを示す。

\begin{lemma}[トレース写像と基底変換]
\label{duality-lemma-trace-map-and-base-change}
$f$ と $g$ が Tor 独立である図式 (\ref{duality-equation-base-change}) があるとする。このとき写像
$1 \star \text{Tr}_f : Lg^* \circ Rf_* \circ a \to Lg^*$ および
$\text{Tr}_{f'} \star 1 : Rf'_* \circ a' \circ Lg^* \to Lg^*$ は、基底変換写像
$\beta : Lg^* \circ Rf_* \to Rf'_* \circ L(g')^*$
（『層のコホモロジー』の注意 \ref{cohomology-remark-base-change}）および
$\alpha : L(g')^* \circ a \to a' \circ Lg^*$
（\ref{duality-equation-base-change-map}）を介して一致する。より正確には、関手の変換の図式
$$
\xymatrix{
Lg^* \circ Rf_* \circ a
\ar[d]_{\beta \star 1} \ar[r]_-{1 \star \text{Tr}_f} &
Lg^* \\
Rf'_* \circ L(g')^* \circ a \ar[r]^{1 \star \alpha} &
Rf'_* \circ a' \circ Lg^* \ar[u]_{\text{Tr}_{f'} \star 1}
}
$$
は可換である。
\end{lemma}

\begin{proof}
本証明では $Rf_*$ を $f_*$、$Lg^*$ を $g^*$ と書く。また変換の始域と終域から補うべきものが
定まるので、恒等変換との $\star$ 積を省略する。$\beta : g^* \circ f_* \to f'_* \circ (g')^*$ は
同型であり、$\alpha$ は $\beta$ の随伴である同型
$\beta^\vee : g'_* \circ a' \to a \circ g_*$ を用いて定義されることを思い出す；補題 \ref{duality-lemma-flat-precompose-pus} とその証明を参照する。まず、補題の図式の上側の水平矢印は合成
$$
g^* \circ f_* \circ a \to
g^* \circ f_* \circ a \circ g_* \circ g^* \to
g^* \circ g_* \circ g^* \to g^*
$$
に等しい。ここで第一の矢印は $(g^*, g_*)$ の単位、第二の矢印は $\text{Tr}_f$、第三の矢印は
$(g^*, g_*)$ の余単位である。これは、単位と余単位の合成
$g^* \to g^* \circ g_* \circ g^* \to g^*$ が恒等写像であることの単純な帰結である。図式
$$
\xymatrix{
& g^* \circ f_* \circ a \ar[ld]_\beta \ar[d] \ar[r]_{\text{Tr}_f} & g^* \\
f'_* \circ (g')^* \circ a \ar[dr] &
g^* \circ f_* \circ a \circ g_* \circ g^* \ar[d]_\beta \ar[ru] &
g^* \circ f_* \circ g'_* \circ a' \circ g^* \ar[l]_{\beta^\vee} \ar[d]_\beta &
f'_* \circ a' \circ g^* \ar[lu]_{\text{Tr}_{f'}} \\
& f'_* \circ (g')^* \circ a \circ g_* \circ g^* &
f'_* \circ (g')^* \circ g'_* \circ a' \circ g^* \ar[ru] \ar[l]_{\beta^\vee}
}
$$
を考える。この図式の二つの平方は『圏論』の補題 \ref{categories-lemma-properties-2-cat-cats}、またはより直接には『圏論』の定義 \ref{categories-definition-horizontal-composition} に先立つ議論により可換である。三角形は上の議論により可換である。
『圏論』の補題 \ref{categories-lemma-transformation-between-functors-and-adjoints} により平方
$$
\xymatrix{
g^* \circ f_* \circ g'_* \circ a' \ar[d]_{\beta^\vee} \ar[r]_-\beta &
f'_* \circ (g')^* \circ g'_* \circ a' \ar[d] \\
g^* \circ f_* \circ a \circ g_* \ar[r] &
\text{id}
}
$$
は可換であり、従って大きな図式の五角形も可換である。$\beta$ と $\beta^\vee$ は同型であり、
大きな図式の外周をたどる写像は定義により $\text{Tr}_f \circ \alpha \circ \beta$ に等しいので、補題が従う。
\end{proof}

\noindent
$f : X \to Y$ を準コンパクトかつ準分離なスキームの射とし、$a : D_\QCoh(\mathcal{O}_Y) \to D_\QCoh(\mathcal{O}_X)$
を補題 \ref{duality-lemma-twisted-inverse-image} の $Rf_*$ の右随伴とする。『圏論』の第 \ref{categories-section-adjoint} 節により関手の変換
$$
\eta_f : \text{id} \to  a \circ Rf_*
$$
を得る。これは随伴の単位と呼ばれる。

\begin{lemma}
\label{duality-lemma-unit-and-base-change}
$f$ と $g$ が Tor 独立である図式 (\ref{duality-equation-base-change}) があるとする。このとき写像
$1 \star \eta_f : L(g')^* \to L(g')^* \circ a \circ Rf_*$ および
$\eta_{f'} \star 1 : L(g')^* \to a' \circ Rf'_* \circ L(g')^*$ は、基底変換写像
$\beta : Lg^* \circ Rf_* \to Rf'_* \circ L(g')^*$
（『層のコホモロジー』の注意 \ref{cohomology-remark-base-change}）および
$\alpha : L(g')^* \circ a \to a' \circ Lg^*$
（\ref{duality-equation-base-change-map}）を介して一致する。より正確には、関手の変換の図式
$$
\xymatrix{
L(g')^* \ar[r]_-{1 \star \eta_f} \ar[d]_{\eta_{f'} \star 1} &
L(g')^* \circ a \circ Rf_* \ar[d]^\alpha \\
a' \circ Rf'_* \circ L(g')^* &
a' \circ Lg^* \circ Rf_* \ar[l]_-\beta
}
$$
は可換である。
\end{lemma}

\begin{proof}
この証明は補題 \ref{duality-lemma-trace-map-and-base-change} の証明の双対である。本証明では $Rf_*$ を
$f_*$、$Lg^*$ を $g^*$ と書き、恒等変換との $\star$ 積を省略する。
$\beta : g^* \circ f_* \to f'_* \circ (g')^*$ は同型であり、$\alpha$ は $\beta$ の随伴である同型
$\beta^\vee : g'_* \circ a' \to a \circ g_*$ を用いて定義されることを思い出す；補題 \ref{duality-lemma-flat-precompose-pus} とその証明を参照する。まず、補題の図式の左側の垂直矢印は合成
$$
(g')^* \to (g')^* \circ g'_* \circ (g')^* \to
(g')^* \circ g'_* \circ a' \circ f'_* \circ (g')^* \to
a' \circ f'_* \circ (g')^*
$$
に等しい。ここで第一の矢印は $((g')^*, g'_*)$ の単位、第二の矢印は $\eta_{f'}$、第三の矢印は
$((g')^*, g'_*)$ の余単位である。これは単位と余単位の合成
$(g')^* \to (g')^* \circ (g')_* \circ (g')^* \to (g')^*$ が恒等写像であることの単純な帰結である。
図式
$$
\xymatrix{
& (g')^* \circ a \circ f_* \ar[r] &
(g')^* \circ a \circ g_* \circ g^* \circ f_*
\ar[ld]_\beta \\
(g')^* \ar[ru]^{\eta_f} \ar[dd]_{\eta_{f'}} \ar[rd] &
(g')^* \circ a \circ g_* \circ f'_* \circ (g')^* &
(g')^* \circ g'_* \circ a' \circ g^* \circ f_*
\ar[u]_{\beta^\vee} \ar[ld]_\beta \ar[d] \\
& (g')^* \circ g'_* \circ a' \circ f'_* \circ (g')^*
\ar[ld] \ar[u]_{\beta^\vee} &
a' \circ g^* \circ f_* \ar[lld]^\beta \\
a' \circ f'_* \circ (g')^*
}
$$
を考える。この図式の二つの平方は『圏論』の補題 \ref{categories-lemma-properties-2-cat-cats}、またはより直接には『圏論』の定義 \ref{categories-definition-horizontal-composition} に先立つ議論により可換である。三角形は上の議論により可換である。
『圏論』の補題 \ref{categories-lemma-transformation-between-functors-and-adjoints} の双対により平方
$$
\xymatrix{
\text{id} \ar[r] \ar[d] &
g'_* \circ a' \circ g^* \circ f_* \ar[d]^\beta \\
g'_* \circ a' \circ g^* \circ f_* \ar[r]^{\beta^\vee} &
a \circ g_* \circ f'_* \circ (g')^*
}
$$
は可換であり、従って大きな図式の五角形も可換である。$\beta$ と $\beta^\vee$ は同型であり、
大きな図式の外周をたどる写像は定義により $\beta \circ \alpha \circ \eta_f$ に等しいので、補題が従う。
\end{proof}

\begin{example}
\label{duality-example-trace-affine}
$A \to B$ を環準同型とする。$Y = \Spec(A)$, $X = \Spec(B)$ とし、$f : X \to Y$ を
$A \to B$ に対応する射とする。例 \ref{duality-example-affine-twisted-inverse-image} で見たように、
$Rf_* : D_\QCoh(\mathcal{O}_X) \to D_\QCoh(\mathcal{O}_Y)$ の右随伴は
$D(A) = D_\QCoh(\mathcal{O}_Y)$ の対象 $K$ を $D(B) = D_\QCoh(\mathcal{O}_X)$ の
$R\Hom(B, K)$ に送る。トレース写像は、$A$-加群の写像 $A \to B$ により誘導される写像
$$
\text{Tr}_{f, K} : R\Hom(B, K) \longrightarrow R\Hom(A, K) = K
$$
である。
\end{example}




\section{押し出しの右随伴と引き戻し}
\label{duality-section-compare-with-pullback}

\noindent
$f : X \to Y$ を準コンパクトかつ準分離なスキームの射とし、$a$ を補題 \ref{duality-lemma-twisted-inverse-image} の押し出しの右随伴とする。$K, L \in D_\QCoh(\mathcal{O}_Y)$ に対して標準写像
$$
Lf^*K \otimes^\mathbf{L}_{\mathcal{O}_X} a(L)
\longrightarrow
a(K \otimes_{\mathcal{O}_Y}^\mathbf{L} L)
$$
がある。実際、この写像は写像
$$
Rf_*(Lf^*K \otimes^\mathbf{L}_{\mathcal{O}_X} a(L)) =
K \otimes^\mathbf{L}_{\mathcal{O}_Y} Rf_*(a(L))
\longrightarrow
K \otimes^\mathbf{L}_{\mathcal{O}_Y} L
$$
に随伴する（等号は『スキームの導来圏』の補題 \ref{perfect-lemma-cohomology-base-change} による）。ここではトレース写像 $Rf_*a(L) \to L$ を用いる。
$L = \mathcal{O}_Y$ のとき、$K$ に関して関手的で区別三角と両立する写像
\begin{equation}
\label{duality-equation-compare-with-pullback}
Lf^*K \otimes^\mathbf{L}_{\mathcal{O}_X} a(\mathcal{O}_Y) \longrightarrow a(K)
\end{equation}
を得る。

\begin{lemma}
\label{duality-lemma-compare-with-pullback-perfect}
$f : X \to Y$ を準コンパクトかつ準分離なスキームの射とする。上で
$K, L \in D_\QCoh(\mathcal{O}_Y)$ に対して定義した写像
$Lf^*K \otimes^\mathbf{L}_{\mathcal{O}_X} a(L) \to
a(K \otimes_{\mathcal{O}_Y}^\mathbf{L} L)$ は、$K$ が完全ならば同型である。特に、$K$ が完全ならば
(\ref{duality-equation-compare-with-pullback}) は同型である。
\end{lemma}

\begin{proof}
$K^\vee$ を $K$ の「双対」とする；『層のコホモロジー』の補題 \ref{cohomology-lemma-dual-perfect-complex} を参照する。$M \in D_\QCoh(\mathcal{O}_X)$ に対して
\begin{align*}
\Hom_{D(\mathcal{O}_Y)}(Rf_*M, K \otimes^\mathbf{L}_{\mathcal{O}_Y} L)
& =
\Hom_{D(\mathcal{O}_Y)}(
Rf_*M \otimes^\mathbf{L}_{\mathcal{O}_Y} K^\vee, L) \\
& =
\Hom_{D(\mathcal{O}_X)}(
M \otimes^\mathbf{L}_{\mathcal{O}_X} Lf^*K^\vee, a(L)) \\
& =
\Hom_{D(\mathcal{O}_X)}(M,
Lf^*K \otimes^\mathbf{L}_{\mathcal{O}_X} a(L))
\end{align*}
である。第二の等号は $a$ の定義と射影公式（『層のコホモロジー』の補題 \ref{cohomology-lemma-projection-formula-perfect}）、またはより一般に『スキームの導来圏』の補題 \ref{perfect-lemma-cohomology-base-change} による。従って Yoneda の補題から結果が従う。
\end{proof}

\begin{lemma}
\label{duality-lemma-restriction-compare-with-pullback}
$f$ と $g$ が Tor 独立である図式 (\ref{duality-equation-base-change}) があるとし、
$K \in D_\QCoh(\mathcal{O}_Y)$ とする。このとき図式
$$
\xymatrix{
L(g')^*(Lf^*K \otimes^\mathbf{L}_{\mathcal{O}_X} a(\mathcal{O}_Y))
\ar[r] \ar[d] & L(g')^*a(K) \ar[d] \\
L(f')^*Lg^*K \otimes_{\mathcal{O}_{X'}}^\mathbf{L} a'(\mathcal{O}_{Y'})
\ar[r] & a'(Lg^*K)
}
$$
は可換である。ここで水平矢印は $K$ および $Lg^*K$ に対する写像
(\ref{duality-equation-compare-with-pullback}) であり、垂直写像は『層のコホモロジー』の注意 \ref{cohomology-remark-base-change} および (\ref{duality-equation-base-change-map}) を用いて構成される。
\end{lemma}

\begin{proof}
本証明では $Rf_*$ を $f_*$、$Lf^*$ を $f^*$ などと書き、
$\otimes^\mathbf{L}_{\mathcal{O}_X}$ などを $\otimes$ と書く。
(\ref{duality-equation-compare-with-pullback}) を合成
\begin{align*}
f^*K \otimes a(\mathcal{O}_Y)
& \to
a(f_*(f^*K \otimes a(\mathcal{O}_Y))) \\
& \leftarrow
a(K \otimes f_*a(\mathcal{O}_K)) \\
& \to
a(K \otimes \mathcal{O}_Y) \\
& \to
a(K)
\end{align*}
と書く。第一の矢印は単位 $\eta_f$ である。第二の矢印は Cohomology, Equation
(\ref{cohomology-equation-projection-formula-map}) に $a$ を適用したものであり、『スキームの導来圏』の補題 \ref{perfect-lemma-cohomology-base-change} により同型である。第三の矢印は
$\text{id}_K \otimes \text{Tr}_f$ に $a$ を適用したものであり、第四の矢印は同型
$K \otimes \mathcal{O}_Y = K$ に $a$ を適用したものである。補題の証明は、これらの各写像が補題の
主張のような可換平方を生じることを示すことからなる。$\eta_f$ と $\text{Tr}_f$ については補題 \ref{duality-lemma-unit-and-base-change} および \ref{duality-lemma-trace-map-and-base-change} である。Cohomology, Equation
(\ref{cohomology-equation-projection-formula-map}) を用いる矢印については『層のコホモロジー』の注意 \ref{cohomology-remark-compatible-with-diagram} である。乗法写像については明らかである。証明が完了する。
\end{proof}

\begin{lemma}
\label{duality-lemma-compare-on-open}
$f : X \to Y$ を Noether スキームの固有射とし、$V \subset Y$ を
$f^{-1}(V) \to V$ が同型となる開集合とする。このとき $K \in D_\QCoh^+(\mathcal{O}_Y)$ に対する写像
(\ref{duality-equation-compare-with-pullback}) の $f^{-1}(V)$ への制限は同型である。
\end{lemma}

\begin{proof}
補題 \ref{duality-lemma-proper-noetherian} により、写像 (\ref{duality-equation-sheafy}) は
$D_\QCoh^+(\mathcal{O}_Y)$ の対象に対して同型である。従って補題 \ref{duality-lemma-restriction-compare-with-pullback} により、$K$ に対する
(\ref{duality-equation-compare-with-pullback}) の $f^{-1}(V)$ への制限は、$K|_V$ と
$f^{-1}(V) \to V$ に対する写像 (\ref{duality-equation-compare-with-pullback}) である。ゆえに $f$ が恒等射である場合に
写像が同型であることを示せば十分であり、これは明らかである。
\end{proof}

\begin{lemma}
\label{duality-lemma-transitivity-compare-with-pullback}
$f : X \to Y$ および $g : Y \to Z$ を準コンパクトかつ準分離なスキームの合成可能な射とし、
$h = g \circ f$ とおく。$a, b, c$ を $f, g, h$ に対する補題 \ref{duality-lemma-twisted-inverse-image} の随伴とする。任意の $K \in D_\QCoh(\mathcal{O}_Z)$ に対して図式
$$
\xymatrix{
Lf^*(Lg^*K \otimes_{\mathcal{O}_Y}^\mathbf{L}
b(\mathcal{O}_Z)) \otimes_{\mathcal{O}_X}^\mathbf{L} a(\mathcal{O}_Y)
\ar@{=}[d] \ar[r] &
a(Lg^*K \otimes_{\mathcal{O}_Y}^\mathbf{L} b(\mathcal{O}_Z)) \ar[r] &
a(b(K)) \ar@{=}[d] \\
Lh^*K \otimes_{\mathcal{O}_X}^\mathbf{L} Lf^*b(\mathcal{O}_Z)
\otimes_{\mathcal{O}_X}^\mathbf{L} a(\mathcal{O}_Y) \ar[r] &
Lh^*K \otimes_{\mathcal{O}_X}^\mathbf{L} c(\mathcal{O}_Z) \ar[r] &
c(K)
}
$$
は可換である。ここで矢印は (\ref{duality-equation-compare-with-pullback}) であり、
$Lh^* = Lf^* \circ Lg^*$ および $c = a \circ b$ を用いた。
\end{lemma}

\begin{proof}
本証明では $Rf_*$ を $f_*$、$Lf^*$ を $f^*$ などと書き、
$\otimes^\mathbf{L}_{\mathcal{O}_X}$ などを $\otimes$ と書く。上側の矢印の合成は写像
$$
g_*f_*(f^*(g^*K \otimes b(\mathcal{O}_Z)) \otimes a(\mathcal{O}_Y)) \to K
$$
に随伴する。『スキームの導来圏』の補題 \ref{perfect-lemma-cohomology-base-change} により左辺は
$K \otimes g_*f_*(f^*b(\mathcal{O}_Z) \otimes a(\mathcal{O}_Y))$ に等しく、定義を調べると、この写像は
$$
g_*f_*(f^*b(\mathcal{O}_Z) \otimes a(\mathcal{O}_Y))
\xleftarrow{g_*\epsilon}
g_*(b(\mathcal{O}_Z) \otimes f_*a(\mathcal{O}_Y)) \xrightarrow{g_*\alpha}
g_*(b(\mathcal{O}_Z)) \xrightarrow{\beta} \mathcal{O}_Z
$$
と $\text{id}_K$ とのテンソル積から得られる。ここで $\epsilon$ は『スキームの導来圏』の補題 \ref{perfect-lemma-cohomology-base-change} の同型であり、$\beta$ は余単位写像
$g_*b \to \text{id}$ から得られる。同様に、下側の水平矢印の合成は、合成
$$
g_*f_*(f^*b(\mathcal{O}_Z) \otimes a(\mathcal{O}_Y)) \xrightarrow{g_*f_*\delta}
g_*f_*(ab(\mathcal{O}_Z)) \xrightarrow{g_*\gamma}
g_*(b(\mathcal{O}_Z)) \xrightarrow{\beta}
\mathcal{O}_Z
$$
と $\text{id}_K$ とのテンソル積に随伴する。ここで $\gamma$ は余単位写像
$f_*a \to \text{id}$ から得られ、$\delta$ は合成
$$
f_*(f^*b(\mathcal{O}_Z) \otimes a(\mathcal{O}_Y))
\xleftarrow{\epsilon}
b(\mathcal{O}_Z) \otimes f_*a(\mathcal{O}_Y) \xrightarrow{\alpha}
b(\mathcal{O}_Z)
$$
を随伴とする写像である。随伴関手、随伴写像、余単位の一般的性質
（『圏論』の第 \ref{categories-section-adjoint} 節を参照）により、望むように
$\gamma \circ f_*\delta = \alpha \circ \epsilon^{-1}$ である。
\end{proof}





\section{閉埋め込みに対する押し出しの右随伴}
\label{duality-section-sections-with-exact-support}

\noindent
$i : (Z, \mathcal{O}_Z) \to (X, \mathcal{O}_X)$ を環付き空間の射とし、$i$ は閉部分集合への同相写像で、
$i^\sharp : \mathcal{O}_X \to i_*\mathcal{O}_Z$ は全射であるとする（例えばスキームの閉埋め込み）。
$\mathcal{I} = \Ker(i^\sharp)$ とおく。$\mathcal{O}_X$-加群の層 $\mathcal{F}$ に対して層
$$
\SheafHom_{\mathcal{O}_X}(i_*\mathcal{O}_Z, \mathcal{F})
$$
は $\mathcal{I}$ で消される $\mathcal{O}_X$-加群の層である。従って『加群の層』の補題 \ref{modules-lemma-i-star-equivalence} により、$\mathcal{O}_Z$-加群の層
$\SheafHom(\mathcal{O}_Z, \mathcal{F})$ が存在し、$\mathcal{O}_X$-加群として
$$
i_*\SheafHom(\mathcal{O}_Z, \mathcal{F}) =
\SheafHom_{\mathcal{O}_X}(i_*\mathcal{O}_Z, \mathcal{F})
$$
となる。この意味を詳しく述べる。

\begin{lemma}
\label{duality-lemma-compute-sheaf-with-exact-support}
上の記法のもとで、関手 $\SheafHom(\mathcal{O}_Z, -)$ は関手
$i_* : \textit{Mod}(\mathcal{O}_Z) \to \textit{Mod}(\mathcal{O}_X)$ の右随伴である。
$V \subset Z$ が開集合ならば
$$
\Gamma(V, \SheafHom(\mathcal{O}_Z, \mathcal{F})) =
\{s \in \Gamma(U, \mathcal{F}) \mid \mathcal{I}s = 0\}
$$
である。ここで $U \subset X$ は $Z$ との共通部分が $V$ となる開集合である。
\end{lemma}

\begin{proof}
$\mathcal{G}$ を $\mathcal{O}_Z$-加群の層とする。このとき
$$
\Hom_{\mathcal{O}_X}(i_*\mathcal{G}, \mathcal{F}) =
\Hom_{i_*\mathcal{O}_Z}(i_*\mathcal{G},
\SheafHom_{\mathcal{O}_X}(i_*\mathcal{O}_Z, \mathcal{F})) =
\Hom_{\mathcal{O}_Z}(\mathcal{G}, \SheafHom(\mathcal{O}_Z, \mathcal{F}))
$$
である。第一の等号は『加群の層』の補題 \ref{modules-lemma-adjoint-tensor-restrict} により、第二の等号は
$i_*$ の完全忠実性により成り立つ；『加群の層』の補題 \ref{modules-lemma-i-star-equivalence} を参照する。
切断の記述は読者に委ねる。
\end{proof}

\noindent
関手
$$
\textit{Mod}(\mathcal{O}_X)
\longrightarrow
\textit{Mod}(\mathcal{O}_Z),
\quad
\mathcal{F} \longmapsto \SheafHom(\mathcal{O}_Z, \mathcal{F})
$$
は左完全であり、導来拡張
$$
R\SheafHom(\mathcal{O}_Z, -) : D(\mathcal{O}_X) \to D(\mathcal{O}_Z).
$$
をもつ。

\begin{lemma}
\label{duality-lemma-sheaf-with-exact-support-adjoint}
上の記法のもとで、関手 $R\SheafHom(\mathcal{O}_Z, -)$ は関手
$Ri_* : D(\mathcal{O}_Z) \to D(\mathcal{O}_X)$ の右随伴である。
\end{lemma}

\begin{proof}
補題 \ref{duality-lemma-compute-sheaf-with-exact-support} により $i_*$ と
$\SheafHom(\mathcal{O}_Z, -)$ は随伴関手なので、これはその帰結である。『導来圏』の補題 \ref{derived-lemma-derived-adjoint-functors} を参照する。
\end{proof}

\begin{lemma}
\label{duality-lemma-sheaf-with-exact-support-ext}
上の記法のもとで、すべての $D(\mathcal{O}_X)$ の $K$ に対して $D(\mathcal{O}_X)$ において
$$
Ri_*R\SheafHom(\mathcal{O}_Z, K) =
R\SheafHom_{\mathcal{O}_X}(i_*\mathcal{O}_Z, K)
$$
である。
\end{lemma}

\begin{proof}
関手 $R\SheafHom(\mathcal{O}_Z, -)$ の構成から直ちに従う。
\end{proof}

\begin{lemma}
\label{duality-lemma-sheaf-with-exact-support-internal-home}
上の記法のもとで、$M \in D(\mathcal{O}_Z)$ に対し、すべての $D(\mathcal{O}_X)$ の $K$ について
$D(\mathcal{O}_Z)$ において
$$
R\SheafHom_{\mathcal{O}_X}(Ri_*M, K) =
Ri_*R\SheafHom_{\mathcal{O}_Z}(M, R\SheafHom(\mathcal{O}_Z, K))
$$
である。
\end{lemma}

\begin{proof}
これは関手 $R\SheafHom(\mathcal{O}_Z, -)$ の構成、および $\mathcal{K}^\bullet$ が
$\mathcal{O}_X$-加群の K-入射複体ならば $\SheafHom(\mathcal{O}_Z, \mathcal{K}^\bullet)$ が
$\mathcal{O}_Z$-加群の K-入射複体であることから直ちに従う；『導来圏』の補題 \ref{derived-lemma-adjoint-preserve-K-injectives} を参照する。
\end{proof}

\begin{lemma}
\label{duality-lemma-sheaf-with-exact-support-quasi-coherent}
$i : Z \to X$ をスキームの擬連接閉埋め込みとする
（$X$ が局所 Noether ならば任意の閉埋め込みでよい）。このとき
\begin{enumerate}
\item $R\SheafHom(\mathcal{O}_Z, -)$ は $D^+_\QCoh(\mathcal{O}_X)$ を
$D^+_\QCoh(\mathcal{O}_Z)$ に送り、
\item $X = \Spec(A)$ および $Z = \Spec(B)$ ならば図式
$$
\xymatrix{
D^+(B) \ar[r] & D_\QCoh^+(\mathcal{O}_Z) \\
D^+(A) \ar[r] \ar[u]^{R\Hom(B, -)} &
D_\QCoh^+(\mathcal{O}_X) \ar[u]_{R\SheafHom(\mathcal{O}_Z, -)}
}
$$
は可換である。
\end{enumerate}
\end{lemma}

\begin{proof}
括弧内の注意を説明する。$X$ が局所 Noether ならば『射の続論』の補題 \ref{more-morphisms-lemma-Noetherian-pseudo-coherent} により $i$ は擬連接である。

\medskip\noindent
$K$ を $D^+_\QCoh(\mathcal{O}_X)$ の対象とする。(1) を示すには、『スキームの射』の補題 \ref{morphisms-lemma-i-star-equivalence} により、$H^n(R\SheafHom(\mathcal{O}_Z, K))$ に
$i_*$ を適用すると $X$ 上の準連接加群が得られることを示せば十分である。補題 \ref{duality-lemma-sheaf-with-exact-support-ext} により、これは
$R\SheafHom_{\mathcal{O}_X}(i_*\mathcal{O}_Z, K)$ が $D_\QCoh(\mathcal{O}_X)$ に属することを
示すということである。$i$ は擬連接なので、層 $\mathcal{O}_Z$ は擬連接 $\mathcal{O}_X$-加群である。
従って結果は『スキームの導来圏』の補題 \ref{perfect-lemma-quasi-coherence-internal-hom} から従う。

\medskip\noindent
(2) のように $X = \Spec(A)$ および $Z = \Spec(B)$ と仮定する。$D^+(A)$ の対象 $K$ を表す、
入射 $A$-加群からなる下に有界な複体を $I^\bullet$ とする。このとき
$R\Hom(B, K) = \Hom_A(B, I^\bullet)$ を $B$-加群の複体とみなせる。擬同型
$$
\widetilde{I^\bullet} \longrightarrow \mathcal{I}^\bullet
$$
を選ぶ。ここで $\mathcal{I}^\bullet$ は入射 $\mathcal{O}_X$-加群からなる下に有界な複体である。
補題 \ref{duality-lemma-compute-sheaf-with-exact-support} における関手
$\SheafHom(\mathcal{O}_Z, -)$ の記述から、写像
$$
\Hom_A(B, I^\bullet)
\longrightarrow
\Gamma(Z, \SheafHom(\mathcal{O}_Z, \mathcal{I}^\bullet))
$$
がある。$\SheafHom(\mathcal{O}_Z, \mathcal{I}^\bullet)$ は
$R\SheafHom(\mathcal{O}_Z, \widetilde{K})$ を表す。$\widetilde{\ }$ 関手の普遍性を適用すると、
$D(\mathcal{O}_Z)$ における写像
$$
\widetilde{\Hom_A(B, I^\bullet)}
\longrightarrow
R\SheafHom(\mathcal{O}_Z, \widetilde{K})
$$
を得る。この写像が $D(\mathcal{O}_Z)$ における同型であることは、$i_*$ を適用した後に確認してよい。
しかし $i_*$ を適用すると、『スキームの導来圏』の補題 \ref{perfect-lemma-quasi-coherence-internal-hom} の同型を補題 \ref{duality-lemma-sheaf-with-exact-support-ext} の同一視を介して得る。
\end{proof}

\begin{lemma}
\label{duality-lemma-sheaf-with-exact-support-coherent}
$i : Z \to X$ をスキームの閉埋め込みとし、$X$ は局所 Noether であると仮定する。このとき
$R\SheafHom(\mathcal{O}_Z, -)$ は $D^+_{\textit{Coh}}(\mathcal{O}_X)$ を
$D^+_{\textit{Coh}}(\mathcal{O}_Z)$ に送る。
\end{lemma}

\begin{proof}
問題は $X$ 上局所的なので、$X$ はアフィンであると仮定してよい。$A$ を Noether とし、
$A \to B$ を全射として $X = \Spec(A)$, $Z = \Spec(B)$ と書く。この場合補題 \ref{duality-lemma-sheaf-with-exact-support-quasi-coherent} を適用して問題を代数に翻訳できる。
対応する代数の結果は『双対化複体』の補題 \ref{dualizing-lemma-exact-support-coherent} の帰結である。
\end{proof}

\begin{lemma}
\label{duality-lemma-twisted-inverse-image-closed}
$X$ を準コンパクトかつ準分離なスキームとし、$i : Z \to X$ を擬連接閉埋め込みとする
（$X$ が Noether ならば任意の閉埋め込みは擬連接である）。$a : D_\QCoh(\mathcal{O}_X) \to D_\QCoh(\mathcal{O}_Z)$
を $Ri_*$ の右随伴とする。このとき $K \in D_\QCoh^+(\mathcal{O}_X)$ に対して関手的同型
$$
a(K) = R\SheafHom(\mathcal{O}_Z, K)
$$
がある。
\end{lemma}

\begin{proof}
括弧内の主張は『射の続論』の補題 \ref{more-morphisms-lemma-Noetherian-pseudo-coherent} から従う。補題 \ref{duality-lemma-sheaf-with-exact-support-adjoint} により、関手 $R\SheafHom(\mathcal{O}_Z, -)$ は
$Ri_* : D(\mathcal{O}_Z) \to D(\mathcal{O}_X)$ の右随伴である。さらに補題 \ref{duality-lemma-sheaf-with-exact-support-quasi-coherent} および補題 \ref{duality-lemma-twisted-inverse-image-bounded-below} により、$R\SheafHom(\mathcal{O}_Z, -)$ と $a$ はともに
$D_\QCoh^+(\mathcal{O}_X)$ を $D_\QCoh^+(\mathcal{O}_Z)$ に送る。従って随伴関手の一意性から同型を得る。
\end{proof}

\begin{example}
\label{duality-example-trace-closed-immersion}
$i : Z \to X$ が Noether スキームの閉埋め込みならば、$K \in D_\QCoh^+(\mathcal{O}_X)$ に対して図式
$$
\xymatrix{
i_*a(K) \ar[rr]_-{\text{Tr}_{i, K}} \ar@{=}[d] & &
K \ar@{=}[d] \\
i_*R\SheafHom(\mathcal{O}_Z, K) \ar@{=}[r] &
R\SheafHom_{\mathcal{O}_X}(i_*\mathcal{O}_Z, K)
\ar[r] & K
}
$$
は可換である。ここで水平の等号は補題 \ref{duality-lemma-sheaf-with-exact-support-ext} であり、下側の水平矢印は写像
$\mathcal{O}_X \to i_*\mathcal{O}_Z$ により誘導される。図式の可換性は補題 \ref{duality-lemma-twisted-inverse-image-closed} の帰結である。
\end{example}




\section{閉埋め込みに対する押し出しの右随伴と基底変換}
\label{duality-section-sections-with-exact-support-base-change}

\noindent
スキームのデカルト図式
$$
\xymatrix{
Z' \ar[r]_{i'} \ar[d]_g & X' \ar[d]^f \\
Z \ar[r]^i & X
}
$$
を考える。ここで $i$ は閉埋め込みである。$Z$ と $X'$ が $X$ 上 Tor 独立ならば、
$D(\mathcal{O}_X)$ の $K$ に関して関手的な $D(\mathcal{O}_{Z'})$ における標準基底変換写像
\begin{equation}
\label{duality-equation-base-change-exact-support}
Lg^*R\SheafHom(\mathcal{O}_Z, K)
\longrightarrow
R\SheafHom(\mathcal{O}_{Z'}, Lf^*K)
\end{equation}
がある。実際、補題 \ref{duality-lemma-sheaf-with-exact-support-adjoint} の随伴性により、このような矢印は
$D(\mathcal{O}_{X'})$ における写像
$$
Ri'_*Lg^*R\SheafHom(\mathcal{O}_Z, K)
\longrightarrow
Lf^*K
$$
と同じことである。Tor 独立性により $Ri'_* \circ Lg^* = Lf^* \circ Ri_*$ である
（『スキームの導来圏』の補題 \ref{perfect-lemma-compare-base-change-closed-immersion} を参照）。従ってこれは写像
$$
Lf^*Ri_*R\SheafHom(\mathcal{O}_Z, K)
\longrightarrow
Lf^*K
$$
と同じことである。この写像として $Lf^*(can)$ を用いればよい。ここで
$can : Ri_* R\SheafHom(\mathcal{O}_Z, K) \to K$ は随伴の余単位である。

\begin{lemma}
\label{duality-lemma-check-base-change-is-iso}
上の状況で、写像 (\ref{duality-equation-base-change-exact-support}) が同型であることと、『層のコホモロジー』の注意 \ref{cohomology-remark-prepare-fancy-base-change} の基底変換写像
$$
Lf^*R\SheafHom_{\mathcal{O}_X}(\mathcal{O}_Z, K)
\longrightarrow
R\SheafHom_{\mathcal{O}_{X'}}(\mathcal{O}_{Z'}, Lf^*K)
$$
が同型であることは同値である。
\end{lemma}

\begin{proof}
仮定した Tor 独立性により $\mathcal{O}_{Z'} = Lf^*\mathcal{O}_Z$ なので、主張には意味がある。
$i'_*$ は完全かつ忠実なので、写像 (\ref{duality-equation-base-change-exact-support}) に $Ri'_*$ を適用したものが
同型であることを示せば十分である。仮定した Tor 独立性と『スキームの導来圏』の補題 \ref{perfect-lemma-compare-base-change-closed-immersion} により
$Ri'_* \circ Lg^* = Lf^* \circ Ri_*$ なので、写像
$$
Lf^*Ri_*R\SheafHom(\mathcal{O}_Z, K)
\longrightarrow
Ri'_*R\SheafHom(\mathcal{O}_{Z'}, Lf^*K)
$$
を得る。補題 \ref{duality-lemma-sheaf-with-exact-support-ext} により、その始域と終域は補題の主張に現れるものと
同じである。この写像が『層のコホモロジー』の注意 \ref{cohomology-remark-prepare-fancy-base-change} で構成されたものと同じであることの確認は省略する。
\end{proof}

\begin{lemma}
\label{duality-lemma-flat-bc-sheaf-with-exact-support}
上の状況で、$f$ は平坦で $i$ は擬連接であると仮定する。このとき
(\ref{duality-equation-base-change-exact-support}) は $D^+_\QCoh(\mathcal{O}_X)$ の $K$ に対して同型である。
\end{lemma}

\begin{proof}
第一の証明。この写像が同型であることを示すには局所的に作業してよい。従って
$X$, $X'$, $Z$, $Z'$ はアフィンで、それぞれ環 $A$, $A'$, $B$, $B'$ に対応すると仮定してよい。
このとき $B$ と $A'$ は $A$ 上 Tor 独立である。補題 \ref{duality-lemma-check-base-change-is-iso} により、
すべての $K \in D^+(A)$ に対して $D(A')$ において
$$
R\Hom_A(B, K) \otimes_A^\mathbf{L} A' =
R\Hom_{A'}(B', K \otimes_A^\mathbf{L} A')
$$
であることを確認すれば十分である。ここでは『スキームの導来圏』の補題 \ref{perfect-lemma-quasi-coherence-internal-hom}、および $B$, resp.\ $B'$ が
$A$-加群, resp.\ $A'$-加群として擬連接であることを用いて、環とスキームの水準の導来 Hom を比較する。
表示した等号は『代数の続論』の補題 \ref{more-algebra-lemma-internal-hom-evaluate-tensor-isomorphism} の (3) から従う。『双対化複体』の第 \ref{dualizing-section-base-change-trivial-duality} 節の議論も参照する。

\medskip\noindent
第二の証明\footnote{この証明は、$K$ が $D^+(\mathcal{O}_X)$ に属すると仮定すれば十分であることを示す。}。
$z' \in Z'$ とし、その像を $z \in Z$ とする。まず、$z'$ における茎上の
(\ref{duality-equation-base-change-exact-support}) が Dualizing Complexes, Equation
(\ref{dualizing-equation-base-change}) の写像
$$
R\Hom(\mathcal{O}_{Z, z}, K_z)
\otimes_{\mathcal{O}_{Z, x}}^\mathbf{L} \mathcal{O}_{Z', z'}
\longrightarrow
R\Hom(\mathcal{O}_{Z', z'},
K_z \otimes_{\mathcal{O}_{X, z}}^\mathbf{L} \mathcal{O}_{X', z'})
$$
を誘導することを示す。実際、これらの写像の構成は同一である。次に『双対化複体』の補題 \ref{dualizing-lemma-flat-bc-surjection} を適用する。
\end{proof}

\begin{lemma}
\label{duality-lemma-sheaf-with-exact-support-tensor}
$i : Z \to X$ をスキームの擬連接閉埋め込みとする。$M \in D_\QCoh(\mathcal{O}_X)$ は局所的に
$[a, \infty)$ に Tor 振幅をもち、$K \in D_\QCoh^+(\mathcal{O}_X)$ とする。このとき $D(\mathcal{O}_Z)$ における標準同型
$$
R\SheafHom(\mathcal{O}_Z, K) \otimes_{\mathcal{O}_Z}^\mathbf{L} Li^*M =
R\SheafHom(\mathcal{O}_Z, K \otimes_{\mathcal{O}_X}^\mathbf{L} M)
$$
がある。
\end{lemma}

\begin{proof}
左辺から右辺への写像は、補題 \ref{duality-lemma-sheaf-with-exact-support-adjoint} および \ref{duality-lemma-sheaf-with-exact-support-ext} により写像
$$
Ri_*R\SheafHom(\mathcal{O}_Z, K) \otimes_{\mathcal{O}_X}^\mathbf{L} M
\longrightarrow
K \otimes_{\mathcal{O}_X}^\mathbf{L} M
$$
と同じことである。この写像として余単位 $Ri_*R\SheafHom(\mathcal{O}_Z, K) \to K$ と
$\text{id}_M$ とのテンソル積を取る。与えた仮定のもとでこの写像が同型であることを見るには、補題 \ref{duality-lemma-sheaf-with-exact-support-quasi-coherent} を用いて代数に翻訳し、例えば『代数の続論』の補題 \ref{more-algebra-lemma-internal-hom-evaluate-tensor-isomorphism} の (3) を用いる。補題 \ref{duality-lemma-sheaf-with-exact-support-quasi-coherent} を用いる代わりに、補題 \ref{duality-lemma-flat-bc-sheaf-with-exact-support} の第二の証明のように茎を見てもよい。
\end{proof}



\section{有限射に対する押し出しの右随伴}
\label{duality-section-duality-finite}

\noindent
$i : Z \to X$ がスキームの閉埋め込みならば、関手
$i_* : \textit{Mod}(\mathcal{O}_Z) \to \textit{Mod}(\mathcal{O}_X)$ の右随伴
$\SheafHom(\mathcal{O}_Z, -)$ があり、その導来拡張 $R\SheafHom(\mathcal{O}_Z, -)$ は
$Ri_* : D(\mathcal{O}_Z) \to D(\mathcal{O}_X)$ の右随伴である。第 \ref{duality-section-sections-with-exact-support} 節を参照する。有限射 $f : Y \to X$ の場合、この方針は使えない。
一般に関手 $f_* : \textit{Mod}(\mathcal{O}_Y) \to \textit{Mod}(\mathcal{O}_X)$ は完全でなく、従って
右随伴をもたないからである。代わりに完全関手
$\textit{Mod}(f_*\mathcal{O}_Y) \to \textit{Mod}(\mathcal{O}_X)$ と、それに対応する右随伴および導来拡張を考える。

\medskip\noindent
$f : Y \to X$ をスキームのアフィン射とする。$\mathcal{O}_X$-加群の層 $\mathcal{F}$ に対して層
$$
\SheafHom_{\mathcal{O}_X}(f_*\mathcal{O}_Y, \mathcal{F})
$$
は $f_*\mathcal{O}_Y$-加群の層である。関手
$\textit{Mod}(\mathcal{O}_X) \to \textit{Mod}(f_*\mathcal{O}_Y)$ を得て、これを
$\SheafHom(f_*\mathcal{O}_Y, -)$ と表す。

\begin{lemma}
\label{duality-lemma-compute-sheafhom-affine}
上の記法のもとで、関手 $\SheafHom(f_*\mathcal{O}_Y, -)$ は係数制限関手
$\textit{Mod}(f_*\mathcal{O}_Y) \to \textit{Mod}(\mathcal{O}_X)$ の右随伴である。
アフィン開集合 $U \subset X$ に対して
$$
\Gamma(U, \SheafHom(f_*\mathcal{O}_Y, \mathcal{F})) =
\Hom_A(B, \mathcal{F}(U))
$$
である。ここで $A = \mathcal{O}_X(U)$ および $B = \mathcal{O}_Y(f^{-1}(U))$ である。
\end{lemma}

\begin{proof}
随伴性は『加群の層』の補題 \ref{modules-lemma-adjoint-tensor-restrict} から従う。$f$ はアフィンなので、
$f_*\mathcal{O}_Y$ は $A$-加群とみなした $B$ に対応する準連接層である。従って $U$ 上の切断の記述は
『スキーム』の補題 \ref{schemes-lemma-compare-constructions} から従う。
\end{proof}

\noindent
関手 $\SheafHom(f_*\mathcal{O}_Y, -)$ は左完全である。その導来拡張を
$$
R\SheafHom(f_*\mathcal{O}_Y, -) :
D(\mathcal{O}_X)
\longrightarrow
D(f_*\mathcal{O}_Y)
$$
とする。

\begin{lemma}
\label{duality-lemma-sheafhom-affine-adjoint}
上の記法のもとで、関手 $R\SheafHom(f_*\mathcal{O}_Y, -)$ は関手
$D(f_*\mathcal{O}_Y) \to D(\mathcal{O}_X)$ の右随伴である。
\end{lemma}

\begin{proof}
補題 \ref{duality-lemma-compute-sheafhom-affine} および『導来圏』の補題 \ref{derived-lemma-derived-adjoint-functors} から従う。
\end{proof}

\begin{lemma}
\label{duality-lemma-sheafhom-affine-ext}
上の記法のもとで、合成
$$
D(\mathcal{O}_X) \xrightarrow{R\SheafHom(f_*\mathcal{O}_Y, -)}
D(f_*\mathcal{O}_Y) \to D(\mathcal{O}_X)
$$
は関手 $K \mapsto R\SheafHom_{\mathcal{O}_X}(f_*\mathcal{O}_Y, K)$ である。
\end{lemma}

\begin{proof}
構成から直ちに従う。
\end{proof}

\begin{lemma}
\label{duality-lemma-finite-twisted}
$f : Y \to X$ をスキームの有限擬連接射とする
（Noether スキームの有限射は擬連接である）。関手 $R\SheafHom(f_*\mathcal{O}_Y, -)$ は
$D_\QCoh^+(\mathcal{O}_X)$ を $D_\QCoh^+(f_*\mathcal{O}_Y)$ に送る。$X$ が準コンパクトかつ準分離ならば図式
$$
\xymatrix{
D_\QCoh^+(\mathcal{O}_X) \ar[rr]_a \ar[rd]_{R\SheafHom(f_*\mathcal{O}_Y, -)}
& & D_\QCoh^+(\mathcal{O}_Y) \ar[ld]^\Phi \\
& D_\QCoh^+(f_*\mathcal{O}_Y)
}
$$
は可換である。ここで $a$ は $f$ に対する補題 \ref{duality-lemma-twisted-inverse-image} の右随伴であり、
$\Phi$ は『スキームの導来圏』の補題 \ref{perfect-lemma-affine-morphism-equivalence} の同値である。
\end{lemma}

\begin{proof}
括弧内の注意は『射の続論』の補題 \ref{more-morphisms-lemma-Noetherian-pseudo-coherent} から従う。$f$ は擬連接なので、
$\mathcal{O}_X$-加群 $f_*\mathcal{O}_Y$ は擬連接である；『射の続論』の補題 \ref{more-morphisms-lemma-finite-pseudo-coherent} を参照する。従って
$R\SheafHom(f_*\mathcal{O}_Y, -)$ は $D_\QCoh^+(\mathcal{O}_X)$ を
$D_\QCoh^+(f_*\mathcal{O}_Y)$ に送る；『スキームの導来圏』の補題 \ref{perfect-lemma-quasi-coherence-internal-hom} を参照する。さらに $\Phi \circ a$ と
$R\SheafHom(f_*\mathcal{O}_Y, -)$ は $D_\QCoh^+(\mathcal{O}_X)$ 上で一致する。これらはともに係数制限関手
$D_\QCoh^+(f_*\mathcal{O}_Y) \to D_\QCoh^+(\mathcal{O}_X)$ の右随伴だからである。これを見るには補題 \ref{duality-lemma-twisted-inverse-image-bounded-below} および \ref{duality-lemma-sheafhom-affine-adjoint} を用いる。
\end{proof}

\begin{remark}
\label{duality-remark-trace-map-finite}
$f : Y \to X$ が Noether スキームの有限射ならば、$K \in D_\QCoh^+(\mathcal{O}_X)$ に対して図式
$$
\xymatrix{
Rf_*a(K) \ar[r]_-{\text{Tr}_{f, K}} \ar@{=}[d] & K \ar@{=}[d] \\
R\SheafHom_{\mathcal{O}_X}(f_*\mathcal{O}_Y, K) \ar[r] & K
}
$$
は可換である。これは補題 \ref{duality-lemma-finite-twisted} から従う。下側の水平矢印は写像
$\mathcal{O}_X \to f_*\mathcal{O}_Y$ により誘導され、上側の水平矢印は第 \ref{duality-section-trace} 節で論じたトレース写像である。
\end{remark}






\section{固有平坦射に対する押し出しの右随伴}
\label{duality-section-proper-flat}

\noindent
準コンパクトかつ準分離なスキームの固有・平坦・有限表示な射に対して、
押し出しの右随伴はいくつかの顕著な性質をもつ。

\begin{lemma}
\label{duality-lemma-proper-flat}
$Y$ を準コンパクトかつ準分離なスキームとする。
$f : X \to Y$ を固有・平坦・有限表示なスキームの射とする。
補題 \ref{duality-lemma-twisted-inverse-image} における
$Rf_* : D_\QCoh(\mathcal{O}_X) \to D_\QCoh(\mathcal{O}_Y)$ の右随伴を
$a$ とする。このとき $a$ は直和と可換である。
\end{lemma}

\begin{proof}
$P$ を $D(\mathcal{O}_X)$ の完全対象とする。『スキームの導来圏』の補題 \ref{perfect-lemma-flat-proper-perfect-direct-image-general} により、
複体 $Rf_*P$ は $Y$ 上完全である。
$K_i$ を $D_\QCoh(\mathcal{O}_Y)$ の対象の族とする。このとき
\begin{align*}
\Hom_{D(\mathcal{O}_X)}(P, a(\bigoplus K_i))
& =
\Hom_{D(\mathcal{O}_Y)}(Rf_*P, \bigoplus K_i) \\
& =
\bigoplus \Hom_{D(\mathcal{O}_Y)}(Rf_*P, K_i) \\
& =
\bigoplus \Hom_{D(\mathcal{O}_X)}(P, a(K_i))
\end{align*}
である。これは完全対象がコンパクトだからである（『スキームの導来圏』の命題 \ref{perfect-proposition-compact-is-perfect}）。
$D_\QCoh(\mathcal{O}_X)$ は完全生成対象をもつので（『スキームの導来圏』の定理 \ref{perfect-theorem-bondal-van-den-Bergh}）、写像
$\bigoplus a(K_i) \to a(\bigoplus K_i)$ は同型であると結論される。すなわち、$a$ は直和と可換である。
\end{proof}

\begin{lemma}
\label{duality-lemma-proper-flat-relative}
$Y$ を準コンパクトかつ準分離なスキームとする。
$f : X \to Y$ を固有・平坦・有限表示なスキームの射とする。
補題 \ref{duality-lemma-twisted-inverse-image} における
$Rf_* : D_\QCoh(\mathcal{O}_X) \to D_\QCoh(\mathcal{O}_Y)$ の右随伴を
$a$ とする。このとき
\begin{enumerate}
\item 任意の閉部分集合 $T \subset Y$ に対し、$Q \in D_\QCoh(Y)$ が $T$ 上に台をもつならば、
$a(Q)$ は $f^{-1}(T)$ 上に台をもつ。
\item 任意の準コンパクトな開部分集合 $V \subset Y$ と任意の
$K \in D_\QCoh(\mathcal{O}_Y)$ に対して、写像 (\ref{duality-equation-sheafy}) は同型である。
\end{enumerate}
\end{lemma}

\begin{proof}
これは補題 \ref{duality-lemma-when-sheafy} および \ref{duality-lemma-proper-noetherian}, および
\ref{duality-lemma-proper-flat} から従う。厳密に言えば、(1) は補集合が準コンパクトな開集合である
$T$ に対してのみ直接従う。しかし任意の閉部分集合はそのような閉部分集合の共通部分なので、
一般の場合も従う。
\end{proof}

\begin{lemma}
\label{duality-lemma-compare-with-pullback-flat-proper}
$Y$ を準コンパクトかつ準分離なスキームとする。
$f : X \to Y$ を固有・平坦・有限表示なスキームの射とする。
$D_\QCoh(\mathcal{O}_Y)$ の任意の対象 $K$ に対して、写像
(\ref{duality-equation-compare-with-pullback}) は同型である。
\end{lemma}

\begin{proof}
補題 \ref{duality-lemma-proper-flat} により、$a$ は直和と可換である。
従って、(\ref{duality-equation-compare-with-pullback}) が同型となる
$D_\QCoh(\mathcal{O}_Y)$ の対象全体は、$D_\QCoh(\mathcal{O}_Y)$ の真に充満な飽和三角部分圏をなし、さらに直和を取る操作で
保たれる。$D_\QCoh(\mathcal{O}_Y)$ は単一の完全対象によって生成される加群圏なので
（『スキームの導来圏』の定理 \ref{perfect-theorem-DQCoh-is-Ddga} および『スキームの導来圏』の定理 \ref{perfect-theorem-bondal-van-den-Bergh}）、『代数の続論』の注意 \ref{more-algebra-remark-P-resolution} と同様に論じることにより、
(\ref{duality-equation-compare-with-pullback}) が一つの完全対象に対して同型であることを示せば十分である。
しかし完全対象に対する主張は補題 \ref{duality-lemma-compare-with-pullback-perfect} により成り立つ。
\end{proof}

\noindent
次の補題は、有限表示な固有平坦射に対して基底変換写像
(\ref{duality-equation-base-change-map}) が同型であることを示す。
例 \ref{duality-example-base-change-wrong} では、完全固有射についてこの主張がもはや成り立たず、
その場合には Tor 独立性の条件を課す必要があることを見る。

\begin{lemma}
\label{duality-lemma-proper-flat-base-change}
$g : Y' \to Y$ を準コンパクトかつ準分離なスキームの射とする。
$f : X \to Y$ を有限表示な固有平坦射とする。
このとき、すべての $K \in D_\QCoh(\mathcal{O}_Y)$ に対して基底変換写像
(\ref{duality-equation-base-change-map}) は同型である。
\end{lemma}

\begin{proof}
補題 \ref{duality-lemma-proper-flat-relative} により、関手 $a$ および $a'$ の構成は
$Y$ および $Y'$ の開集合への制限と可換である。従って注意 \ref{duality-remark-check-over-affines} により、$Y' \to Y$ はアフィンスキームの射であると仮定してよい。
この場合、主張は補題 \ref{duality-lemma-more-base-change} から従う。
\end{proof}

\begin{remark}
\label{duality-remark-relative-dualizing-complex}
$Y$ を準コンパクトかつ準分離なスキームとする。
$f : X \to Y$ を有限表示な固有平坦射とする。
$a$ を $f$ に対する補題 \ref{duality-lemma-twisted-inverse-image} の随伴とする。
この状況で、$\omega_{X/Y}^\bullet = a(\mathcal{O}_Y)$ はしばしば
{\it 相対双対化複体}と呼ばれる。補題 \ref{duality-lemma-compare-with-pullback-flat-proper} により、
$K \in D_\QCoh(\mathcal{O}_Y)$ に対して関手的同型
$a(K) = Lf^*K \otimes_{\mathcal{O}_X}^\mathbf{L} \omega_{X/Y}^\bullet$
がある。さらに、第 \ref{duality-section-trace} 節のトレース写像
$$
\text{Tr}_{f, \mathcal{O}_Y} : Rf_*\omega_{X/Y}^\bullet \to \mathcal{O}_Y
$$
は $D_\QCoh(\mathcal{O}_Y)$ のすべての $K$ に対するトレース写像を誘導する。
より正確には、次の図式が可換である。
$$
\xymatrix{
Rf_*a(K) \ar[rrr]_{\text{Tr}_{f, K}} \ar@{=}[d] & & &
K \ar@{=}[d] \\
Rf_*(Lf^*K \otimes_{\mathcal{O}_X}^\mathbf{L} \omega_{X/Y}^\bullet)
\ar@{=}[r] &
K \otimes_{\mathcal{O}_Y}^\mathbf{L} Rf_*\omega_{X/Y}^\bullet
\ar[rr]^-{\text{id}_K \otimes \text{Tr}_{f, \mathcal{O}_Y}} & & K
}
$$
ここで右下の等号は『スキームの導来圏』の補題 \ref{perfect-lemma-cohomology-base-change} である。
$g : Y' \to Y$ を準コンパクトかつ準分離なスキームの射とし、
$X' = Y' \times_Y X$ とする。このとき補題 \ref{duality-lemma-proper-flat-base-change} により、
射影を $g' : X' \to X$ と書けば
$\omega_{X'/Y'}^\bullet = L(g')^*\omega_{X/Y}^\bullet$ である。また補題 \ref{duality-lemma-trace-map-and-base-change} により、$f' : X' \to Y'$ に対するトレース写像
$$
\text{Tr}_{f', \mathcal{O}_{Y'}} :
Rf'_*\omega_{X'/Y'}^\bullet \to \mathcal{O}_{Y'}
$$
は、基底変換同型を介した $\text{Tr}_{f, \mathcal{O}_Y}$ の基底変換である。
\end{remark}

\begin{remark}
\label{duality-remark-relative-dualizing-complex-relative-cup-product}
$f : X \to Y$, $\omega^\bullet_{X/Y}$, および $\text{Tr}_{f, \mathcal{O}_Y}$ を
注意 \ref{duality-remark-relative-dualizing-complex} のとおりとする。
$K$ と $M$ を $D_\QCoh(\mathcal{O}_X)$ の対象とし、$M$ は擬連接であるとする
（例えば完全であるとする）。Cohomology, Equation
(\ref{cohomology-equation-internal-hom}) を介して同型
$K \to R\SheafHom_{\mathcal{O}_X}(M, \omega^\bullet_{X/Y})$ に対応する写像
$K \otimes_{\mathcal{O}_X}^\mathbf{L} M \to \omega^\bullet_{X/Y}$ が与えられているとする。
このとき相対カップ積（『層のコホモロジー』の注意 \ref{cohomology-remark-cup-product}）
$$
Rf_*K \otimes_{\mathcal{O}_Y}^\mathbf{L} Rf_*M
\to
Rf_*(K \otimes_{\mathcal{O}_X}^\mathbf{L} M)
\to
Rf_*\omega^\bullet_{X/Y}
\xrightarrow{\text{Tr}_{f, \mathcal{O}_Y}}
\mathcal{O}_Y
$$
は同型
$Rf_*K \to R\SheafHom_{\mathcal{O}_Y}(Rf_*M, \mathcal{O}_Y)$ を定める。
実際、$\omega^\bullet_{X/Y} = a(\mathcal{O}_Y)$ なので、標準写像
(\ref{duality-equation-sheafy-trace})
$$
Rf_*R\SheafHom_{\mathcal{O}_X}(M, \omega^\bullet_{X/Y}) \to
R\SheafHom_{\mathcal{O}_Y}(Rf_*M, \mathcal{O}_Y)
$$
は補題 \ref{duality-lemma-iso-on-RSheafHom}, 注意 \ref{duality-remark-iso-on-RSheafHom},
および $M$ と $Rf_*M$ が擬連接であることにより同型である。後者については
『スキームの導来圏』の補題 \ref{perfect-lemma-flat-proper-pseudo-coherent-direct-image-general} を参照する。
この同型が相対カップ積によって誘導されることは、『層のコホモロジー』の注意 \ref{cohomology-remark-relative-cup-and-composition} の図式の可換性から従う。
\end{remark}

\begin{lemma}
\label{duality-lemma-properties-relative-dualizing}
$Y$ を準コンパクトかつ準分離なスキームとする。
$f : X \to Y$ を固有・平坦・有限表示なスキームの射とし、その相対双対化複体を
$\omega_{X/Y}^\bullet$ とする（注意 \ref{duality-remark-relative-dualizing-complex}）。
このとき
\begin{enumerate}
\item $\omega_{X/Y}^\bullet$ は $D(\mathcal{O}_X)$ の $Y$-完全対象である。
\item $Rf_*\omega_{X/Y}^\bullet$ の正次数のコホモロジー層は消滅する。
\item $\mathcal{O}_X \to
R\SheafHom_{\mathcal{O}_X}(\omega_{X/Y}^\bullet, \omega_{X/Y}^\bullet)$
は同型である。
\end{enumerate}
\end{lemma}

\begin{proof}
$\omega_{X/Y}^\bullet$ の構成は基底変換と可換なので（注意 \ref{duality-remark-relative-dualizing-complex} を参照）、$Y$ はアフィンであると仮定してよいし、実際そう仮定する。
$E$ を $D(\mathcal{O}_X)$ の完全対象とすると
\begin{align*}
Rf_*(E \otimes_{\mathcal{O}_X}^\mathbf{L} \omega_{X/Y}^\bullet)
& =
Rf_*R\SheafHom_{\mathcal{O}_X}(E^\vee, \omega_{X/Y}^\bullet) \\
& =
R\SheafHom_{\mathcal{O}_Y}(Rf_*E^\vee, \mathcal{O}_Y) \\
& =
(Rf_*E^\vee)^\vee
\end{align*}
である。第一の等号については『層のコホモロジー』の補題 \ref{cohomology-lemma-dual-perfect-complex} を参照する。第二の等号については補題 \ref{duality-lemma-iso-on-RSheafHom}, 注意 \ref{duality-remark-iso-on-RSheafHom}, および
『スキームの導来圏』の補題 \ref{perfect-lemma-flat-proper-perfect-direct-image-general} を参照する。
第三の等号は双対の定義である。特に、これらの参照結果は得られた対象が
$D(\mathcal{O}_Y)$ の完全対象であることも示す。従って『射の続論』の補題 \ref{more-morphisms-lemma-characterize-relatively-perfect} により、
$\omega_{X/Y}^\bullet$ は $Y$-完全である。これで (1) が示された。

\medskip\noindent
$M$ を $D_\QCoh(\mathcal{O}_Y)$ の対象とする。このとき
\begin{align*}
\Hom_Y(M, Rf_*\omega_{X/Y}^\bullet) & =
\Hom_X(Lf^*M, \omega_{X/Y}^\bullet) \\
& =
\Hom_Y(Rf_*Lf^*M, \mathcal{O}_Y) \\
& =
\Hom_Y(M \otimes_{\mathcal{O}_Y}^\mathbf{L} Rf_*\mathcal{O}_X, \mathcal{O}_Y)
\end{align*}
である。第一の等号は『層のコホモロジー』の補題 \ref{cohomology-lemma-adjoint} による。
第二の等号は $a$ の構成による。第三の等号は『スキームの導来圏』の補題 \ref{perfect-lemma-cohomology-base-change} による。
$Rf_*\mathcal{O}_X$ はある $N$ に対して $[0, N]$ に Tor 振幅をもつ完全対象であることを想起する。
『スキームの導来圏』の補題 \ref{perfect-lemma-flat-proper-perfect-direct-image-general} を参照する。
従って $Rf_*\mathcal{O}_X$ は次数 $[0, N]$ に位置する有限射影加群の複体で表せる
（『代数の続論』の補題 \ref{more-algebra-lemma-perfect} および $Y$ がアフィンであることを用いる）。
ゆえに、ある $i > 0$ に対して $M = \mathcal{O}_Y[-i]$ ならば最後の群は零である。
$Y$ はアフィンなので、$i > 0$ に対して
$H^i(Rf_*\omega_{X/Y}^\bullet) = 0$ と結論される。これで (2) が示された。

\medskip\noindent
$E$ を $D_\QCoh(\mathcal{O}_X)$ の完全対象とする。このとき
\begin{align*}
\Hom_X(E, R\SheafHom_{\mathcal{O}_X}(\omega_{X/Y}^\bullet, \omega_{X/Y}^\bullet)
& =
\Hom_X(E \otimes_{\mathcal{O}_X}^\mathbf{L} \omega_{X/Y}^\bullet,
\omega_{X/Y}^\bullet) \\
& =
\Hom_Y(Rf_*(E \otimes_{\mathcal{O}_X}^\mathbf{L} \omega_{X/Y}^\bullet),
\mathcal{O}_Y) \\
& =
\Hom_Y(Rf_*(R\SheafHom_{\mathcal{O}_X}(E^\vee, \omega_{X/Y}^\bullet)),
\mathcal{O}_Y) \\
& =
\Hom_Y(R\SheafHom_{\mathcal{O}_Y}(Rf_*E^\vee, \mathcal{O}_Y),
\mathcal{O}_Y) \\
& =
R\Gamma(Y, Rf_*E^\vee) \\
& =
\Hom_X(E, \mathcal{O}_X)
\end{align*}
を得る。第一の等号は『層のコホモロジー』の補題 \ref{cohomology-lemma-internal-hom} による。第二の等号は
$\omega_{X/Y}^\bullet$ の定義である。第三の等号は双対完全複体 $E^\vee$ の構成から得られる。
『層のコホモロジー』の補題 \ref{cohomology-lemma-dual-perfect-complex} を参照する。
第四の等号は、証明の第一段落で示した等式
$Rf_*R\SheafHom_{\mathcal{O}_X}(E^\vee, \omega_{X/Y}^\bullet) =
R\SheafHom_{\mathcal{O}_Y}(Rf_*E^\vee, \mathcal{O}_Y)$ から従う。
第五の等号は完全複体に対する二重双対性（『層のコホモロジー』の補題 \ref{cohomology-lemma-dual-perfect-complex}）と、『スキームの導来圏』の補題 \ref{perfect-lemma-flat-proper-perfect-direct-image-general} により $Rf_*E$ が完全であることから従う。
最後の等号は $f$ に対する Leray である。
この一連の等式は、本質的に米田の補題によって (3) が成り立つことを示す。実際、対象
$R\SheafHom(\omega_{X/Y}^\bullet, \omega_{X/Y}^\bullet)$ は『スキームの導来圏』の補題 \ref{perfect-lemma-quasi-coherence-internal-hom} により $D_\QCoh(\mathcal{O}_X)$ に属する。
上で $E = \mathcal{O}_X$ とおくと、
$\text{id}_{\mathcal{O}_X} \in \Hom_X(\mathcal{O}_X, \mathcal{O}_X)$ に対応する写像
$\alpha : \mathcal{O}_X \to
R\SheafHom_{\mathcal{O}_X}(\omega_{X/Y}^\bullet, \omega_{X/Y}^\bullet)$
を得る。上のすべての同型は $E$ に関して関手的なので、$\alpha$ の錐は
$D_\QCoh(\mathcal{O}_X)$ の対象 $C$ であり、すべての完全 $E$ に対して
$\Hom(E, C) = 0$ を満たす。完全対象は生成するので（『スキームの導来圏』の定理 \ref{perfect-theorem-bondal-van-den-Bergh}）、$\alpha$ は同型である。
\end{proof}

\begin{lemma}[剛性]
\label{duality-lemma-van-den-bergh}
$Y$ を準コンパクトかつ準分離なスキームとする。
$f : X \to Y$ を有限表示な固有平坦射とし、その相対双対化複体を
$\omega_{X/Y}^\bullet$ とする（注意 \ref{duality-remark-relative-dualizing-complex}）。
標準同型
\begin{equation}
\label{duality-equation-pre-rigid}
\mathcal{O}_X =
c(L\text{pr}_1^*\omega_{X/Y}^\bullet) =
c(L\text{pr}_2^*\omega_{X/Y}^\bullet)
\end{equation}
および標準同型
\begin{equation}
\label{duality-equation-rigid}
\omega_{X/Y}^\bullet =
c\left(L\text{pr}_1^*\omega_{X/Y}^\bullet
\otimes_{\mathcal{O}_{X \times_Y X}}^\mathbf{L}
L\text{pr}_2^*\omega_{X/Y}^\bullet\right)
\end{equation}
がある。ここで $c$ は対角射 $\Delta : X \to X \times_Y X$ に対する
補題 \ref{duality-lemma-twisted-inverse-image} の右随伴である。
\end{lemma}

\begin{proof}
$a$ を補題 \ref{duality-lemma-twisted-inverse-image} における $Rf_*$ の右随伴とする。
デカルト平方
$$
\xymatrix{
X \times_Y X \ar[r]_q \ar[d]_p & X \ar[d]_f \\
X \ar[r]^f & Y
}
$$
を考える。$b$ を補題 \ref{duality-lemma-twisted-inverse-image} における $p$ の右随伴とする。このとき
\begin{align*}
\omega_{X/Y}^\bullet
& =
c(b(\omega_{X/Y}^\bullet)) \\
& =
c(Lp^*\omega_{X/Y}^\bullet
\otimes_{\mathcal{O}_{X \times_Y X}}^\mathbf{L} b(\mathcal{O}_X)) \\
& =
c(Lp^*\omega_{X/Y}^\bullet
\otimes_{\mathcal{O}_{X \times_Y X}}^\mathbf{L}
Lq^*a(\mathcal{O}_Y)) \\
& =
c(Lp^*\omega_{X/Y}^\bullet
\otimes_{\mathcal{O}_{X \times_Y X}}^\mathbf{L}
Lq^*\omega_{X/Y}^\bullet)
\end{align*}
となり、これは (\ref{duality-equation-rigid}) である。各等号の説明は次のとおりである。
\begin{enumerate}
\item $\text{id}_X = p \circ \Delta$ なので $\text{id} = c \circ b$ であり、第一の等号が成り立つ。
\item 第二の等号は補題 \ref{duality-lemma-compare-with-pullback-flat-proper} による。
\item 第三の等号は補題 \ref{duality-lemma-proper-flat-base-change} および
$\mathcal{O}_X = Lf^*\mathcal{O}_Y$ による。
\item 第四の等号は $\omega_{X/Y}^\bullet = a(\mathcal{O}_Y)$ による。
\end{enumerate}
Equation (\ref{duality-equation-pre-rigid}) もまったく同じ方法で示される。
\end{proof}

\begin{remark}
\label{duality-remark-van-den-bergh}
補題 \ref{duality-lemma-van-den-bergh} は、ここでの相対双対化複体が
\cite{vdB-rigid} で導入された概念に類似する意味で {\it rigid} であることを意味する。
実際、(\ref{duality-equation-rigid}) の右辺の関手は $\omega_{X/Y}^\bullet$ について「二次的」であり、
(\ref{duality-equation-rigid}) の左辺の関手は「一次的」なので、これは複体
$\omega_{X/Y}^\bullet$ をある程度「一意に定める」。rigid な（相対）双対化複体を用いる
双対理論へのアプローチがあり、例えば \cite{Neeman-rigid}, \cite{Yekutieli-rigid},
および \cite{Yekutieli-Zhang} を参照する。第 \ref{duality-section-relative-dualizing-complexes} 節でこれに戻る。
\end{remark}





\section{完全固有射に対する押し出しの右随伴}
\label{duality-section-perfect-proper}

\noindent
この節にとって適切な一般性は、準コンパクトかつ準分離なスキームの完全固有射を
考えることであろう。\cite{LN} を参照する。

\begin{lemma}
\label{duality-lemma-proper-flat-noetherian}
$f : X \to Y$ を Noether スキームの完全固有射とする。
補題 \ref{duality-lemma-twisted-inverse-image} における
$Rf_* : D_\QCoh(\mathcal{O}_X) \to D_\QCoh(\mathcal{O}_Y)$ の右随伴を
$a$ とする。このとき $a$ は直和と可換である。
\end{lemma}

\begin{proof}
$P$ を $D(\mathcal{O}_X)$ の完全対象とする。『射の続論』の補題 \ref{more-morphisms-lemma-perfect-proper-perfect-direct-image} により、
複体 $Rf_*P$ は $Y$ 上完全である。
$K_i$ を $D_\QCoh(\mathcal{O}_Y)$ の対象の族とする。このとき
\begin{align*}
\Hom_{D(\mathcal{O}_X)}(P, a(\bigoplus K_i))
& =
\Hom_{D(\mathcal{O}_Y)}(Rf_*P, \bigoplus K_i) \\
& =
\bigoplus \Hom_{D(\mathcal{O}_Y)}(Rf_*P, K_i) \\
& =
\bigoplus \Hom_{D(\mathcal{O}_X)}(P, a(K_i))
\end{align*}
である。これは完全対象がコンパクトだからである（『スキームの導来圏』の命題 \ref{perfect-proposition-compact-is-perfect}）。
$D_\QCoh(\mathcal{O}_X)$ は完全生成対象をもつので（『スキームの導来圏』の定理 \ref{perfect-theorem-bondal-van-den-Bergh}）、写像
$\bigoplus a(K_i) \to a(\bigoplus K_i)$ は同型であると結論される。すなわち、$a$ は直和と可換である。
\end{proof}

\begin{lemma}
\label{duality-lemma-proper-flat-noetherian-relative}
$f : X \to Y$ を Noether スキームの完全固有射とする。
補題 \ref{duality-lemma-twisted-inverse-image} における
$Rf_* : D_\QCoh(\mathcal{O}_X) \to D_\QCoh(\mathcal{O}_Y)$ の右随伴を
$a$ とする。このとき
\begin{enumerate}
\item 任意の閉部分集合 $T \subset Y$ に対し、$Q \in D_\QCoh(Y)$ が $T$ 上に台をもつならば、
$a(Q)$ は $f^{-1}(T)$ 上に台をもつ。
\item 任意の開部分集合 $V \subset Y$ と任意の $K \in D_\QCoh(\mathcal{O}_Y)$ に対して、
写像 (\ref{duality-equation-sheafy}) は同型である。
\end{enumerate}
\end{lemma}

\begin{proof}
これは補題 \ref{duality-lemma-when-sheafy} および \ref{duality-lemma-proper-noetherian}, および
\ref{duality-lemma-proper-flat-noetherian} から従う。
\end{proof}

\begin{lemma}
\label{duality-lemma-compare-with-pullback-flat-proper-noetherian}
$f : X \to Y$ を Noether スキームの完全固有射とする。
$D_\QCoh(\mathcal{O}_Y)$ の任意の対象 $K$ に対して、写像
(\ref{duality-equation-compare-with-pullback}) は同型である。
\end{lemma}

\begin{proof}
補題 \ref{duality-lemma-proper-flat-noetherian} により、$a$ は直和と可換である。
従って、(\ref{duality-equation-compare-with-pullback}) が同型となる
$D_\QCoh(\mathcal{O}_Y)$ の対象全体は、$D_\QCoh(\mathcal{O}_Y)$ の真に充満な飽和三角部分圏をなし、
さらに直和を取る操作で保たれる。$D_\QCoh(\mathcal{O}_Y)$ は単一の完全対象によって
生成される加群圏なので（『スキームの導来圏』の定理 \ref{perfect-theorem-DQCoh-is-Ddga} および『スキームの導来圏』の定理 \ref{perfect-theorem-bondal-van-den-Bergh}）、『代数の続論』の注意 \ref{more-algebra-remark-P-resolution} と同様に論じることにより、
(\ref{duality-equation-compare-with-pullback}) が一つの完全対象に対して同型であることを示せば十分である。
しかし完全対象に対する主張は補題 \ref{duality-lemma-compare-with-pullback-perfect} により成り立つ。
\end{proof}

\begin{lemma}
\label{duality-lemma-proper-perfect-base-change}
$f : X \to Y$ を Noether スキームの完全固有射とする。
$g : Y' \to Y$ を射とし、$Y'$ は Noether であるとする。
$X$ と $Y'$ が $Y$ 上 Tor 独立ならば、すべての
$K \in D_\QCoh(\mathcal{O}_Y)$ に対して基底変換写像
(\ref{duality-equation-base-change-map}) は同型である。
\end{lemma}

\begin{proof}
補題 \ref{duality-lemma-proper-flat-noetherian-relative} により、関手 $a$ および $a'$ の構成は
$Y$ および $Y'$ の開集合への制限と可換である。従って注意 \ref{duality-remark-check-over-affines} により、$Y' \to Y$ はアフィンスキームの射であると仮定してよい。
この場合、主張は補題 \ref{duality-lemma-more-base-change} から従う。
\end{proof}



\section{有効 Cartier 因子に対する押し出しの右随伴}
\label{duality-section-dualizing-Cartier}

\noindent
$X$ をスキームとし、$i : D \to X$ を有効 Cartier 因子の包含とする。
$\mathcal{N} = i^*\mathcal{O}_X(D)$ を $i$ の法層と書く。『スキームの射』の第 \ref{morphisms-section-conormal-sheaf} 節および『因子』の第 \ref{divisors-section-effective-Cartier-divisors} 節を参照する。
$R\SheafHom(\mathcal{O}_D, -)$ は
$i_* : D(\mathcal{O}_D) \to D(\mathcal{O}_X)$ の右随伴を表し、
$i_*R\SheafHom(\mathcal{O}_D, -) =
R\SheafHom_{\mathcal{O}_X}(i_*\mathcal{O}_D, -)$ を満たすことを想起する。
第 \ref{duality-section-sections-with-exact-support} 節を参照する。

\begin{lemma}
\label{duality-lemma-compute-for-effective-Cartier}
上と同様に、$X$ をスキームとし、$D \subset X$ を有効 Cartier 因子とする。
$D(\mathcal{O}_D)$ に標準同型
$R\SheafHom(\mathcal{O}_D, \mathcal{O}_X) = \mathcal{N}[-1]$ がある。
\end{lemma}

\begin{proof}
同値な言い方をすれば、$R\SheafHom(\mathcal{O}_D, \mathcal{O}_X)$ は次数 $1$ にのみ
非零なコホモロジー層をもち、その層は $\mathcal{N}$ と同型であると主張している。
$i_*$ は完全かつ完全忠実なので、
$i_*R\SheafHom(\mathcal{O}_D, \mathcal{O}_X)$ が $i_*\mathcal{N}[-1]$ と同型であることを
示せば十分である。補題 \ref{duality-lemma-sheaf-with-exact-support-ext} により
$i_*R\SheafHom(\mathcal{O}_D, \mathcal{O}_X) =
R\SheafHom_{\mathcal{O}_X}(i_*\mathcal{O}_D, \mathcal{O}_X)$ である。
$\mathcal{I}$ を $D$ のイデアル層とすると、計算に用いることのできる分解
$$
0 \to \mathcal{I} \to \mathcal{O}_X \to i_*\mathcal{O}_D \to 0
$$
がある。局所計算により
$R\SheafHom_{\mathcal{O}_X}(\mathcal{O}_X, \mathcal{O}_X) = \mathcal{O}_X$ かつ
$R\SheafHom_{\mathcal{O}_X}(\mathcal{I}, \mathcal{O}_X) = \mathcal{O}_X(D)$ なので、
$$
R\SheafHom_{\mathcal{O}_X}(i_*\mathcal{O}_D, \mathcal{O}_X) =
(\mathcal{O}_X \to \mathcal{O}_X(D))
$$
を得る。右辺では $\mathcal{O}_X$ が次数 $0$ に、$\mathcal{O}_X(D)$ が次数 $1$ にある。
主張は短完全列
$$
0 \to \mathcal{O}_X \to \mathcal{O}_X(D) \to i_*\mathcal{N} \to 0
$$
から従う。この列は $D$ が $\mathcal{O}_X(D)$ の標準切断の零点スキームであること、および
$\mathcal{N} = i^*\mathcal{O}_X(D)$ から得られる。
\end{proof}

\noindent
$D(\mathcal{O}_X)$ の任意の対象 $K$ に対し、$K$ に関して関手的で識別三角形と両立する
$D(\mathcal{O}_D)$ における標準写像
\begin{equation}
\label{duality-equation-map-effective-Cartier}
Li^*K
\otimes_{\mathcal{O}_D}^\mathbf{L}
R\SheafHom(\mathcal{O}_D, \mathcal{O}_X)
\longrightarrow
R\SheafHom(\mathcal{O}_D, K)
\end{equation}
がある。実際、この写像は
$$
i_*(Li^*K \otimes^\mathbf{L}_{\mathcal{O}_D}
R\SheafHom(\mathcal{O}_D, \mathcal{O}_X)) =
K \otimes^\mathbf{L}_{\mathcal{O}_X}
R\SheafHom_{\mathcal{O}_X}(i_*\mathcal{O}_D, \mathcal{O}_X)
\longrightarrow K
$$
の随伴である。ここで等号は『層のコホモロジー』の補題 \ref{cohomology-lemma-projection-formula-closed-immersion} であり、矢印は
$\mathcal{O}_X \to i_*\mathcal{O}_D$ から誘導される標準写像
$R\SheafHom_{\mathcal{O}_X}(i_*\mathcal{O}_D, \mathcal{O}_X) \to \mathcal{O}_X$ から得られる。

\medskip\noindent
$K \in D_\QCoh(\mathcal{O}_X)$ ならば、
(\ref{duality-equation-map-effective-Cartier}) は (\ref{duality-equation-compare-with-pullback}) に等しい。
ここでは補題 \ref{duality-lemma-twisted-inverse-image-closed} の同一視
$a(K) = R\SheafHom(\mathcal{O}_D, K)$ を用いる。
$K \in D_\QCoh(\mathcal{O}_X)$ かつ $X$ が Noether ならば、次の補題は補題 \ref{duality-lemma-compare-with-pullback-flat-proper-noetherian} の特別な場合である。

\begin{lemma}
\label{duality-lemma-sheaf-with-exact-support-effective-Cartier}
上と同様に、$X$ をスキームとし、$D \subset X$ を有効 Cartier 因子とする。
(\ref{duality-equation-map-effective-Cartier}) と補題 \ref{duality-lemma-compute-for-effective-Cartier} を組み合わせると、$D(\mathcal{O}_X)$ の $K$ に関して
関手的な同型
$$
Li^*K \otimes_{\mathcal{O}_D}^\mathbf{L} \mathcal{N}[-1]
\longrightarrow
R\SheafHom(\mathcal{O}_D, K)
$$
が定まる。
\end{lemma}

\begin{proof}
$i_*$ は加群上で完全かつ完全忠実なので、写像が同型であることを示すには、$i_*$ を適用した後に
同型となることを示せば十分である。補題 \ref{duality-lemma-compute-for-effective-Cartier} の証明で用いた
短完全列 $0 \to \mathcal{I} \to \mathcal{O}_X \to i_*\mathcal{O}_D \to 0$ および
$0 \to \mathcal{O}_X \to \mathcal{O}_X(D) \to i_*\mathcal{N} \to 0$ を、以下では断らずに用いる。
『層のコホモロジー』の補題 \ref{cohomology-lemma-projection-formula-closed-immersion} により、
これは写像 (\ref{duality-equation-map-effective-Cartier}) の定義に用いた結果であるが、左辺は
$$
K \otimes_{\mathcal{O}_X}^\mathbf{L} i_*\mathcal{N}[-1] =
K \otimes_{\mathcal{O}_X}^\mathbf{L} (\mathcal{O}_X \to \mathcal{O}_X(D))
$$
となる。右辺は
\begin{align*}
R\SheafHom_{\mathcal{O}_X}(i_*\mathcal{O}_D, K) & =
R\SheafHom_{\mathcal{O}_X}((\mathcal{I} \to \mathcal{O}_X), K) \\
& =
R\SheafHom_{\mathcal{O}_X}((\mathcal{I} \to \mathcal{O}_X), \mathcal{O}_X)
\otimes_{\mathcal{O}_X}^\mathbf{L} K
\end{align*}
となり、最後の等号は『層のコホモロジー』の補題 \ref{cohomology-lemma-dual-perfect-complex} による。写像は同型
$$
R\SheafHom_{\mathcal{O}_X}((\mathcal{I} \to \mathcal{O}_X), \mathcal{O}_X)
= (\mathcal{O}_X \to \mathcal{O}_X(D))
$$
から得られるので、補題は明らかである。
\end{proof}





\section{例における押し出しの右随伴}
\label{duality-section-examples}

\noindent
この節では、いくつかの例について押し出しの右随伴を計算する。同型は標準的ではあるが、
可能な限り弱い意味においてのみ標準的である。すなわち、これらの同型が基底変換や射の合成などの
さまざまな操作と両立することは証明も主張もしない。この種の問題には膨大な文献がある。
読者は \cite{RD}, \cite{Conrad-GD} から始めることができる（これらの文献は双対性について
異なる出発点を採るが、相対双対化複体の標準的代表を構成する問題を扱っている）。その後、
Joseph Lipman と共同研究者による諸研究へ進むことができる。

\begin{lemma}
\label{duality-lemma-upper-shriek-P1}
$Y$ を準コンパクトかつ準分離なスキームとする。
$\mathcal{E}$ を階数 $n + 1$ の有限局所自由 $\mathcal{O}_Y$-加群とし、その行列式を
$\mathcal{L} = \wedge^{n + 1}(\mathcal{E})$ とする。
$f : X = \mathbf{P}(\mathcal{E}) \to Y$ を射影とする。
補題 \ref{duality-lemma-twisted-inverse-image} における
$Rf_* : D_\QCoh(\mathcal{O}_X) \to D_\QCoh(\mathcal{O}_Y)$ の右随伴を $a$ とする。
このとき同型
$$
c : f^*\mathcal{L}(-n - 1)[n] \longrightarrow a(\mathcal{O}_Y)
$$
がある。特に、$\mathcal{E} = \mathcal{O}_Y^{\oplus n + 1}$ ならば
$X = \mathbf{P}^n_Y$ であり、$a(\mathcal{O}_Y) = \mathcal{O}_X(-n - 1)[n]$ を得る。
\end{lemma}

\begin{proof}
『スキームのコホモロジー』の補題 \ref{coherent-lemma-cohomology-projective-bundle} の証明において、標準同型
$$
R^nf_*(f^*\mathcal{L}(-n - 1)) \longrightarrow \mathcal{O}_Y
$$
を構成した。さらに
$Rf_*(f^*\mathcal{L}(-n - 1))[n] = R^nf_*(f^*\mathcal{L}(-n - 1))$、
すなわち他の高次順像は零である。従って同型
$$
Rf_*(f^*\mathcal{L}(-n - 1)[n]) \longrightarrow \mathcal{O}_Y
$$
を得る。$a$ は $Rf_*$ の右随伴なので、この同型は補題の主張における $c$ を定める。
補題 \ref{duality-lemma-proper-noetherian} により $a$ の構成は底上局所的である。
特に、$c$ が同型であることを確かめるには $Y$ 上局所的に作業してよい。
言い換えれば、$Y$ はアフィンであり $\mathcal{E} = \mathcal{O}_Y^{\oplus n + 1}$ であると仮定してよい。
この場合、層 $\mathcal{O}_X, \mathcal{O}_X(-1), \ldots,
\mathcal{O}_X(-n)$ は $D_\QCoh(X)$ を生成する。『スキームの導来圏』の補題 \ref{perfect-lemma-generator-P1} を参照する。従って、
$c : \mathcal{O}_X(-n - 1)[n] \to a(\mathcal{O}_Y)$ が
$$
F_{i, p}(-) = \Hom_{D(\mathcal{O}_X)}(\mathcal{O}_X(i), (-)[p])
$$
という関手のもとで同型に移されることを、$i \in \{-n, \ldots, 0\}$ および
$p \in \mathbf{Z}$ について示せば十分である。
$F_{0, p}$ については、これは矢印 $c$ の構成から成り立つ。
$i \in \{-n, \ldots, -1\}$ に対しては
$$
\Hom_{D(\mathcal{O}_X)}(\mathcal{O}_X(i), \mathcal{O}_X(-n - 1)[n + p]) =
H^p(X, \mathcal{O}_X(-n - 1 - i)) = 0
$$
である。これは射影空間のコホモロジーの計算による（『スキームのコホモロジー』の補題 \ref{coherent-lemma-cohomology-projective-space-over-ring}）。また
$$
\Hom_{D(\mathcal{O}_X)}(\mathcal{O}_X(i), a(\mathcal{O}_Y)[p]) =
\Hom_{D(\mathcal{O}_Y)}(Rf_*\mathcal{O}_X(i), \mathcal{O}_Y[p]) = 0
$$
である。これは同じ補題により $Rf_*\mathcal{O}_X(i) = 0$ だからである。
従って $F_{i, p}(c)$ の始域と終域はいずれも消滅し、$F_{i, p}(c)$ は必然的に同型である。
これで証明が完了する。
\end{proof}

\begin{example}
\label{duality-example-base-change-wrong}
$f$ が完全固有で $g$ が完全であるとき、基底変換写像
(\ref{duality-equation-base-change-map}) は同型とは限らない。
$k$ を体とする。$Y = \mathbf{A}^2_k$ とし、$f : X \to Y$ を $Y$ の原点における
ブローアップとする。$E \subset X$ を例外因子と書く。このとき $f$ は
$$
X \xrightarrow{i} \mathbf{P}^1_Y \xrightarrow{p} Y
$$
と分解できる。これにより分解 $a = c \circ b$ を得る。ここで $a$, $b$, $c$ はそれぞれ
$Rf_*$, $Rp_*$, $Ri_*$ に対する補題 \ref{duality-lemma-twisted-inverse-image} の右随伴である。
$\mathcal{O}(n)$ を $\mathbf{P}^1_Y$ 上の構造層の Serre 捩りと書き、
$\mathcal{O}_X(n)$ をその $X$ への制限と書く。
$X \subset \mathbf{P}^1_Y$ は一次方程式で切り出されるので、
$\mathcal{O}(X) = \mathcal{O}(1)$ である。補題 \ref{duality-lemma-upper-shriek-P1} により
$b(\mathcal{O}_Y) = \mathcal{O}(-2)[1]$ である。補題 \ref{duality-lemma-twisted-inverse-image-closed} により
$$
a(\mathcal{O}_Y) = c(b(\mathcal{O}_Y)) =
c(\mathcal{O}(-2)[1]) =
R\SheafHom(\mathcal{O}_X, \mathcal{O}(-2)[1]) =
\mathcal{O}_X(-1)
$$
である。最後の等号は補題 \ref{duality-lemma-sheaf-with-exact-support-effective-Cartier} による。
$Y' = \Spec(k)$ を $Y$ の原点とする。$a(\mathcal{O}_Y)$ の
$X' = E = \mathbf{P}^1_k$ への制限は、コホモロジー次数 $0$ に置かれた次数 $-1$ の可逆層である。
一方、$a'(\mathcal{O}_{\Spec(k)}) = \mathcal{O}_E(-2)[1]$ はコホモロジー次数 $-1$ に
置かれた次数 $-2$ の可逆層であり、両者は異なる。この例では補題 \ref{duality-lemma-more-base-change} の Tor 独立性の仮定が破れている。
\end{example}

\begin{lemma}
\label{duality-lemma-ext}
$Y$ を環付き空間とし、$\mathcal{I} \subset \mathcal{O}_Y$ をイデアル層とする。
$\mathcal{O}_X = \mathcal{O}_Y/\mathcal{I}$ および
$\mathcal{N} =
\SheafHom_{\mathcal{O}_Y}(\mathcal{I}/\mathcal{I}^2, \mathcal{O}_X)$ とおく。
標準同型
$c : \mathcal{N} \to
\SheafExt^1_{\mathcal{O}_Y}(\mathcal{O}_X, \mathcal{O}_X)
$ がある。
\end{lemma}

\begin{proof}
標準短完全列
\begin{equation}
\label{duality-equation-second-order-thickening}
0 \to \mathcal{I}/\mathcal{I}^2 \to \mathcal{O}_Y/\mathcal{I}^2 \to
\mathcal{O}_X \to 0
\end{equation}
を考える。$U \subset X$ を開集合とし、$s \in \mathcal{N}(U)$ とする。
(\ref{duality-equation-second-order-thickening}) を $s$ に沿って押し出すと、
$\mathcal{O}_X|_U$ の $\mathcal{O}_X|_U$ による拡大 $E_s$ を得る。
これはさらに $U$ 上の
$\SheafExt^1_{\mathcal{O}_Y}(\mathcal{O}_X, \mathcal{O}_X)$ の切断 $c(s)$ を定める。
『層のコホモロジー』の補題 \ref{cohomology-lemma-section-RHom-over-U} および
『導来圏』の補題 \ref{derived-lemma-ext-1} を参照する。
逆に、$\mathcal{O}_U$-加群の拡大
$$
0 \to \mathcal{O}_X|_U \to \mathcal{E} \to \mathcal{O}_X|_U \to 0
$$
が与えられたとする。開被覆 $U = \bigcup U_i$ と、
$1 \in \mathcal{O}_X(U_i)$ に写る切断 $e_i \in \mathcal{E}(U_i)$ を取ることができる。
このとき $e_i$ は、核が $\mathcal{I}^2$ を含む写像
$\mathcal{O}_Y|_{U_i} \to \mathcal{E}|_{U_i}$ を定める。
従って $\mathcal{E}|_{U_i}$ は上と同じ押し出しから得られる。
これにより $c$ は全射である。単射性の証明は省略する。
\end{proof}

\begin{lemma}
\label{duality-lemma-regular-ideal-ext}
$Y$ を環付き空間とし、$\mathcal{I} \subset \mathcal{O}_Y$ をイデアル層とする。
$\mathcal{O}_X = \mathcal{O}_Y/\mathcal{I}$ とおく。
$\mathcal{I}$ が Koszul 正則ならば（『因子』の定義 \ref{divisors-definition-regular-ideal-sheaf}）、
$R\SheafHom_{\mathcal{O}_Y}(\mathcal{O}_X, \mathcal{O}_X)$ 上の合成は、すべての $i$ に対して同型
$$
\wedge^i(\SheafExt^1_{\mathcal{O}_Y}(\mathcal{O}_X, \mathcal{O}_X))
\longrightarrow
\SheafExt^i_{\mathcal{O}_Y}(\mathcal{O}_X, \mathcal{O}_X)
$$
を定める。
\end{lemma}

\begin{proof}
合成とは、『層のコホモロジー』の補題 \ref{cohomology-lemma-internal-hom-composition} の写像
$$
R\SheafHom_{\mathcal{O}_Y}(\mathcal{O}_X, \mathcal{O}_X)
\otimes_{\mathcal{O}_Y}^\mathbf{L}
R\SheafHom_{\mathcal{O}_Y}(\mathcal{O}_X, \mathcal{O}_X)
\longrightarrow
R\SheafHom_{\mathcal{O}_Y}(\mathcal{O}_X, \mathcal{O}_X)
$$
を意味する。これは乗法写像
$$
\SheafExt^a_{\mathcal{O}_Y}(\mathcal{O}_X, \mathcal{O}_X)
\otimes_{\mathcal{O}_Y}
\SheafExt^b_{\mathcal{O}_Y}(\mathcal{O}_X, \mathcal{O}_X)
\longrightarrow
\SheafExt^{a + b}_{\mathcal{O}_Y}(\mathcal{O}_X, \mathcal{O}_X)
$$
を誘導する。More on Algebra, Equation
(\ref{more-algebra-equation-simple-tor-product}) と比較されたい。
補題の主張は、誘導写像
$$
\SheafExt^1_{\mathcal{O}_Y}(\mathcal{O}_X, \mathcal{O}_X)
\otimes \ldots \otimes
\SheafExt^1_{\mathcal{O}_Y}(\mathcal{O}_X, \mathcal{O}_X)
\longrightarrow
\SheafExt^i_{\mathcal{O}_Y}(\mathcal{O}_X, \mathcal{O}_X)
$$
が外積を経由し、その後同型を誘導するという意味である。
これを示すには $Y$ 上局所的に作業してよい。従って、$\mathcal{I}$ を生成し Koszul 正則列をなす
$\mathcal{O}_Y$ の大域切断 $f_1, \ldots, f_r$ があると仮定してよい。
$$
\mathcal{A} = \mathcal{O}_Y\langle \xi_1, \ldots, \xi_r\rangle
$$
を、狭義可換微分次数付き $\mathcal{O}_Y$-代数の層であって、$\mathcal{O}_Y$ 上次数 $-1$ の
$\xi_1, \ldots, \xi_r$ に関する（分切冪）多項式代数であり、微分 $\text{d}$ が規則
$\text{d}\xi_i = f_i$ で与えられるものとする。
$\mathcal{A}^\bullet$ をその基礎にある $\mathcal{O}_Y$-加群の複体と書く。これは上記の Koszul 複体である。
従って標準写像 $\mathcal{A}^\bullet \to \mathcal{O}_X$ は擬同型である。
『層のコホモロジー』の補題 \ref{cohomology-lemma-Rhom-strictly-perfect} により擬同型
$$
R\SheafHom_{\mathcal{O}_Y}(\mathcal{O}_X, \mathcal{O}_X) \to
\SheafHom^\bullet(\mathcal{A}^\bullet, \mathcal{A}^\bullet) \to
\SheafHom^\bullet(\mathcal{A}^\bullet, \mathcal{O}_X)
$$
を得る。最後の複体の微分は零なので
$$
\SheafExt^i_{\mathcal{O}_Y}(\mathcal{O}_X, \mathcal{O}_X)
\cong \SheafHom_{\mathcal{O}_Y}(\mathcal{A}^{-i}, \mathcal{O}_X)
$$
である。$j \in \{1, \ldots, r\}$ に対し、$\delta_j : \mathcal{A} \to \mathcal{A}$ を
$\delta_j(\xi_i) = \delta_{ij}$（Kronecker のデルタ）を満たす次数 $1$ の導分とする。
計算により $\delta_j \circ \text{d} = - \text{d} \circ \delta_j$ であるから、複体の射
$$
\delta_j : \mathcal{A}^\bullet \to \mathcal{A}^\bullet[1].
$$
を得る。従って $\delta_j$ は対応する $\SheafExt$-層の切断を定める。
別の計算により、$\delta_1, \ldots, \delta_r$ は $\mathcal{O}_X$ 上の
$\SheafHom_{\mathcal{O}_Y}(\mathcal{A}^{-1}, \mathcal{O}_X)$ の基底に写る。
$\mathcal{A}$ の自己準同型として、従って $\SheafExt$-層において、明らかに
$\delta_j \circ \delta_j = 0$ かつ
$\delta_j \circ \delta_{j'} = - \delta_{j'} \circ \delta_j$ である。
従って上の写像は外冪を経由する。所望の同型を得ることを見るには、読者は
$$
\delta_{j_1} \circ \ldots \circ \delta_{j_i}
$$
が $j_1 < \ldots < j_i$ に対して、$\mathcal{O}_X$ 上の層
$\SheafHom_{\mathcal{O}_Y}(\mathcal{A}^{-i}, \mathcal{O}_X)$ の基底に写ることを確かめればよい。
\end{proof}

\begin{lemma}
\label{duality-lemma-regular-immersion-ext}
$Y$ を環付き空間とし、$\mathcal{I} \subset \mathcal{O}_Y$ をイデアル層とする。
$\mathcal{O}_X = \mathcal{O}_Y/\mathcal{I}$ および
$\mathcal{N} =
\SheafHom_{\mathcal{O}_Y}(\mathcal{I}/\mathcal{I}^2, \mathcal{O}_X)$ とおく。
$\mathcal{I}$ が Koszul 正則ならば（『因子』の定義 \ref{divisors-definition-regular-ideal-sheaf}）、
$$
R\SheafHom_{\mathcal{O}_Y}(\mathcal{O}_X, \mathcal{O}_Y) =
\wedge^r \mathcal{N}[-r]
$$
である。ここで $r : Y \to \{1, 2, 3, \ldots \}$ は、$y$ を $y$ の近傍で
$\mathcal{I}$ を生成するのに必要な最小生成元数へ送る。
\end{lemma}

\begin{proof}
補題 \ref{duality-lemma-ext} および \ref{duality-lemma-regular-ideal-ext} により、$i \geq 0$ に対する同型
$\wedge^i\mathcal{N} \to
\SheafExt^i_{\mathcal{O}_Y}(\mathcal{O}_X, \mathcal{O}_X)$ がある。
従って、写像 $\mathcal{O}_Y \to \mathcal{O}_X$ が同型
$$
\SheafExt^r_{\mathcal{O}_Y}(\mathcal{O}_X, \mathcal{O}_Y)
\longrightarrow
\SheafExt^r_{\mathcal{O}_Y}(\mathcal{O}_X, \mathcal{O}_X)
$$
を誘導すること、および $i \not = r$ に対して
$\SheafExt^i_{\mathcal{O}_Y}(\mathcal{O}_X, \mathcal{O}_Y)$ が零であることを示せば十分である。
これらの主張は $Y$ 上局所的である。従って、$\mathcal{I}$ を生成し Koszul 正則列をなす
$\mathcal{O}_Y$ の大域切断 $f_1, \ldots, f_r$ があると仮定してよい。
$\mathcal{A}^\bullet$ を、補題 \ref{duality-lemma-regular-ideal-ext} の証明で導入した
$f_1, \ldots, f_r$ 上の Koszul 複体とする。このとき『層のコホモロジー』の補題 \ref{cohomology-lemma-Rhom-strictly-perfect} により
$$
R\SheafHom_{\mathcal{O}_Y}(\mathcal{O}_X, \mathcal{O}_Y) =
\SheafHom^\bullet(\mathcal{A}^\bullet, \mathcal{O}_Y)
$$
である。$1 : \mathcal{A}^\bullet \to \mathcal{O}_Y$ を、次数 $0$ における恒等写像
$\mathcal{A}^0 = \mathcal{O}_Y \to \mathcal{O}_Y$ で与えられる微分次数付き
$\mathcal{O}_Y$-代数の写像とする。補題 \ref{duality-lemma-regular-ideal-ext} の証明における
$\delta_j$ を用いると、次数付き $\mathcal{O}_Y$-加群の同型
$$
\mathcal{O}_Y\langle \delta_1, \ldots, \delta_r\rangle
\longrightarrow
\SheafHom^\bullet(\mathcal{A}^\bullet, \mathcal{O}_Y)
$$
を得る。これは $\delta_{j_1} \ldots \delta_{j_i}$ を次数 $i$ の
$1 \circ \delta_{j_1} \circ \ldots \circ \delta_{j_i}$ に写す。
この同型を介して、右辺の微分は左辺の微分 $\text{d}$ を誘導する。
符号規約により $\text{d}(1) = - \sum f_j \delta_j$ である。
$\delta_j : \mathcal{A}^\bullet \to \mathcal{A}^\bullet[1]$ は複体の射なので
$$
\text{d}(\delta_{j_1} \ldots \delta_{j_i}) =
(- \sum f_j \delta_j )\delta_{j_1} \ldots \delta_{j_i}
$$
となる。微分次数付き代数 $\mathcal{A}$ 上では $\text{d} = \sum f_j \delta_j$ であることに注意する。
従って規則
$$
1 \circ \delta_{j_1} \ldots \delta_{j_i} \longmapsto
(\delta_{j_1} \circ \ldots \circ \delta_{j_i})(\xi_1 \ldots \xi_r)
$$
で定められる写像は、$r$ が奇数なら複体の同型
$$
\SheafHom^\bullet(\mathcal{A}^\bullet, \mathcal{O}_Y)
\longrightarrow \mathcal{A}^\bullet[-r]
$$
を定め、$r$ が偶数なら符号を除いて微分と可換する。いずれの場合にも両複体のコホモロジーは
同型であり、所望の消滅が従う。写像
$$
\SheafHom^\bullet(\mathcal{A}^\bullet, \mathcal{O}_Y)
\longrightarrow
\SheafHom^\bullet(\mathcal{A}^\bullet, \mathcal{O}_X)
$$
のもとでの次数 $r$ のコホモロジー上の同型も、これから直ちに従う。
（Koszul 複体に関する規約の選択は、同型
$R\SheafHom_{\mathcal{O}_X}(\mathcal{O}_X, \mathcal{O}_Y) =
\wedge^r \mathcal{N}[-r]$ の定義に実際に関与することに注意する。）
\end{proof}

\begin{lemma}
\label{duality-lemma-regular-immersion}
$Y$ を準コンパクトかつ準分離なスキームとする。
$i : X \to Y$ を Koszul 正則閉埋め込みとする。
補題 \ref{duality-lemma-twisted-inverse-image} における
$Ri_* : D_\QCoh(\mathcal{O}_X) \to D_\QCoh(\mathcal{O}_Y)$ の右随伴を $a$ とする。
このとき同型
$$
\wedge^r\mathcal{N}[-r] \longrightarrow a(\mathcal{O}_Y)
$$
がある。ここで
$\mathcal{N} = \SheafHom_{\mathcal{O}_X}(\mathcal{C}_{X/Y}, \mathcal{O}_X)$
は $i$ の法層であり（『スキームの射』の第 \ref{morphisms-section-conormal-sheaf} 節）、$r$ は $X$ 上の局所定数関数とみなしたその階数である。
\end{lemma}

\begin{proof}
補題 \ref{duality-lemma-twisted-inverse-image-closed} および
\ref{duality-lemma-sheaf-with-exact-support-ext} により、$a(\mathcal{O}_Y)$ は
$D_\QCoh(\mathcal{O}_X)$ の対象であり、その $Y$ への押し出しは
$R\SheafHom_{\mathcal{O}_Y}(i_*\mathcal{O}_X, \mathcal{O}_Y)$ であることを想起する。
従って主張は補題 \ref{duality-lemma-regular-immersion-ext} から従う。
\end{proof}

\begin{lemma}
\label{duality-lemma-smooth-proper}
$S$ を Noether スキームとする。
$f : X \to S$ を相対次元 $d$ の滑らかな固有射とする。
補題 \ref{duality-lemma-twisted-inverse-image} における
$Rf_* : D_\QCoh(\mathcal{O}_X) \to D_\QCoh(\mathcal{O}_S)$ の右随伴を $a$ とする。
このとき $D(\mathcal{O}_X)$ に同型
$$
\wedge^d \Omega_{X/S}[d] \longrightarrow a(\mathcal{O}_S)
$$
がある。
\end{lemma}

\begin{proof}
注意 \ref{duality-remark-relative-dualizing-complex} のとおり
$\omega_{X/S}^\bullet = a(\mathcal{O}_S)$ とおく。
$c$ を $\Delta : X \to X \times_S X$ に対する補題 \ref{duality-lemma-twisted-inverse-image} の右随伴とする。
$\Delta$ は滑らかな射の対角射なので Koszul 正則埋め込みである。『因子』の補題 \ref{divisors-lemma-immersion-smooth-into-smooth-regular-immersion} を参照する。
特に $\Delta$ は完全固有射であり（『射の続論』の補題 \ref{more-morphisms-lemma-regular-immersion-perfect}）、次を得る。
\begin{align*}
\mathcal{O}_X
& =
c(L\text{pr}_1^*\omega_{X/S}^\bullet) \\
& =
L\Delta^*(L\text{pr}_1^*\omega_{X/S}^\bullet)
\otimes_{\mathcal{O}_X}^\mathbf{L}
c(\mathcal{O}_{X \times_S X}) \\
& =
\omega_{X/S}^\bullet \otimes_{\mathcal{O}_X}^\mathbf{L}
c(\mathcal{O}_{X \times_S X}) \\
& =
\omega_{X/S}^\bullet
\otimes_{\mathcal{O}_X}^\mathbf{L}
\wedge^d(\mathcal{N}_\Delta)[-d]
\end{align*}
第一の等号は $\omega_{X/S}^\bullet = a(\mathcal{O}_S)$ であることと
(\ref{duality-equation-pre-rigid}) による。第二の等号は補題 \ref{duality-lemma-compare-with-pullback-flat-proper-noetherian} による。
第三の等号は $\text{pr}_1 \circ \Delta = \text{id}_X$ による。
第四の等号は補題 \ref{duality-lemma-regular-immersion} による。
$\wedge^d(\mathcal{N}_\Delta)$ は可逆 $\mathcal{O}_X$-加群であることに注意する。
従って $\wedge^d(\mathcal{N}_\Delta)[-d]$ は $D(\mathcal{O}_X)$ の可逆対象であり、
$a(\mathcal{O}_S) = \omega_{X/S}^\bullet = \wedge^d(\mathcal{C}_\Delta)[d]$ と結論される。
$\Delta$ の余法層 $\mathcal{C}_\Delta$ は『スキームの射』の補題 \ref{morphisms-lemma-differentials-diagonal} により $\Omega_{X/S}$ なので、証明が完了する。
\end{proof}







\section{上付きシュリーク函手}
\label{duality-section-upper-shriek}

\noindent
この節では、固定した Noether 底上で有限型かつ分離なスキーム間の射に対する函手 $f^!$ を、
コンパクト化を用いて構成する。連接双対性では通常のことだが、可換であると示さなければならない
図式が多数ある。読者には、構成を読んだ後は補題の検証を飛ばし、上付きシュリーク函手の性質を
論じる次節へ進むことを勧める。

\begin{situation}
\label{duality-situation-shriek}
ここで $S$ は Noether スキームであり、$\textit{FTS}_S$ は次の圏である。
\begin{enumerate}
\item 対象は $S$ 上のスキーム $X$ であって、構造射 $X \to S$ が分離かつ有限型であるもの。
\item 対象間の射 $f : X \to Y$ は $S$ 上のスキームの射である。
\end{enumerate}
\end{situation}

\noindent
状況 \ref{duality-situation-shriek} において、$\textit{FTS}_S$ の射 $f : X \to Y$ が与えられたとする。
三角圏の完全函手
$$
f^! : D_\QCoh^+(\mathcal{O}_Y) \to D_\QCoh^+(\mathcal{O}_X)
$$
を定義する。まず $Y$ 上のコンパクト化 $X \to \overline{X}$ を選ぶ。
これは『平坦性の続論』の定理 \ref{flat-theorem-nagata} および補題 \ref{flat-lemma-compactifyable} により可能である。構造射を
$\overline{f} : \overline{X} \to Y$ と書く。$\overline{a} : D_\QCoh(\mathcal{O}_Y) \to D_\QCoh(\mathcal{O}_{\overline{X}})$
を、補題 \ref{duality-lemma-twisted-inverse-image} で構成した $R\overline{f}_*$ の右随伴とする。
$K \in D_\QCoh^+(\mathcal{O}_Y)$ に対して
$$
f^!K  = \overline{a}(K)|_X
$$
とおく。補題 \ref{duality-lemma-twisted-inverse-image-bounded-below} により、得られるものは
$D_\QCoh^+(\mathcal{O}_X)$ の対象である。

\begin{lemma}
\label{duality-lemma-shriek-well-defined}
状況 \ref{duality-situation-shriek} において $f : X \to Y$ を $\textit{FTS}_S$ の射とする。
函手 $f^!$ は、標準同型を除いてコンパクト化の選択に依存しない。
\end{lemma}

\begin{proof}
$Y$ 上の $X$ のコンパクト化の圏は『平坦性の続論』の第 \ref{flat-section-compactify} 節で定義される。『平坦性の続論』の定理 \ref{flat-theorem-nagata} および補題 \ref{flat-lemma-compactifyable} により、この圏は空でない。
コンパクト化
$$
j : X \to \overline{X},\quad \overline{f} : \overline{X} \to Y
$$
を選ぶごとに、上の構成は函手 $j^* \circ \overline{a} :
D_\QCoh^+(\mathcal{O}_Y) \to D_\QCoh^+(\mathcal{O}_X)$ を対応させる。
ここで $\overline{a}$ は補題 \ref{duality-lemma-twisted-inverse-image} で構成した
$R\overline{f}_*$ の右随伴である。

\medskip\noindent
$Y$ 上のコンパクト化 $j_i : X \to \overline{X}_i$ の間の射
$g : \overline{X}_1 \to \overline{X}_2$ で、
$g^{-1}(j_2(X)) = j_1(X)$ を満たすものが与えられたとする\footnote{これは、ここでの
コンパクト化の定義では成り立たないことがある。『平坦性の続論』の第 \ref{flat-section-compactify} 節を参照する。}。
$\overline{c}$ を $g$ に対する補題 \ref{duality-lemma-twisted-inverse-image} の右随伴とする。
これらの函手は
$R\overline{f}_{2, *} \circ Rg_* = R(\overline{f}_2 \circ g)_*$ の随伴なので、
$\overline{c} \circ \overline{a}_2 = \overline{a}_1$ である。
(\ref{duality-equation-sheafy}) により、函手
$D^+_\QCoh(\mathcal{O}_{\overline{X}_2}) \to D^+_\QCoh(\mathcal{O}_X)$ の標準変換
$$
j_1^* \circ \overline{c} \longrightarrow j_2^*
$$
がある。これは補題 \ref{duality-lemma-proper-noetherian} により同型である。合成
$$
j_1^* \circ \overline{a}_1 \longrightarrow
j_1^* \circ \overline{c} \circ \overline{a}_2 \longrightarrow
j_2^* \circ \overline{a}_2
$$
は函手の同型であり、これを $\alpha_g$ と書く。

\medskip\noindent
$Y$ 上の $X$ の二つのコンパクト化 $j_i : X \to \overline{X}_i$, $i = 1, 2$ を考える。
『平坦性の続論』の補題 \ref{flat-lemma-compactifications-cofiltered} part (b) により、
稠密像をもつコンパクト化 $j : X \to \overline{X}$ とコンパクト化の射
$g_i : \overline{X} \to \overline{X}_i$ を取ることができる。
『平坦性の続論』の補題 \ref{flat-lemma-compactifications-cofiltered} part (c) により
$g_i^{-1}(j_i(X)) = j(X)$ である。従って前段落により同型
$$
\alpha_{g_i} :
j^* \circ \overline{a}
\longrightarrow
j_i^* \circ \overline{a}_i
$$
を得る。よって同型
$$
\alpha_{g_2} \circ \alpha_{g_1}^{-1} :
j_1^* \circ \overline{a}_1 \to j_2^* \circ \overline{a}_2
$$
を得る。証明を終えるには、これらの同型が良定義であることを示さなければならない。
前段落で構成した同型の合成が再びその種の同型であることを示せば十分だと主張する
（正確な主張は次段落を参照する）。読者はこれが正しいことを紙片上で確かめてもよいが、
この段落の残りで完全に書き下す。
稠密像をもつ別のコンパクト化 $j' : X \to \overline{X}'$ と、コンパクト化の射
$g'_i : \overline{X}' \to \overline{X}_i$ を考える。『平坦性の続論』の補題 \ref{flat-lemma-compactifications-cofiltered} により、稠密像をもつコンパクト化
$j'' : X \to \overline{X}''$ とコンパクト化の射
$h : \overline{X}'' \to \overline{X}$ および
$h' : \overline{X}'' \to \overline{X}'$ を取ることができる。さらに
$g_1 \circ h = g'_1 \circ h'$ かつ $g_2 \circ h = g'_2 \circ h'$ と仮定してよい。
次段落の結果から
$$
\alpha_{g_i} \circ \alpha_h = \alpha_{g_i \circ h} =
\alpha_{g'_i \circ h'} = \alpha_{g'_i} \circ \alpha_{h'}
$$
を $i = 1, 2$ に対して得る。これらはすべて函手の同型なので、所望のとおり
$\alpha_{g_2} \circ \alpha_{g_1}^{-1} =
\alpha_{g'_2} \circ \alpha_{g'_1}^{-1}$ と結論される。

\medskip\noindent
$i = 1, 2, 3$ に対するコンパクト化 $j_i : X \to \overline{X}_i$ が与えられ、
コンパクト化の射 $g : \overline{X}_1 \to \overline{X}_2$ および
$h : \overline{X}_2 \to \overline{X}_3$ が、
$g^{-1}(j_2(X)) = j_1(X)$ かつ $h^{-1}(j_2(X)) = j_3(X)$ を満たすとする。
$\overline{a}_i$ を上のとおりとする。上の主張は
$$
\alpha_g \circ \alpha_h = \alpha_{g \circ h} :
j_1^* \circ \overline{a}_1 \to j_3^* \circ \overline{a}_3
$$
を意味する。$\overline{c}$ と\ $\overline{d}$ をそれぞれ $g$ と\ $h$ に対する
補題 \ref{duality-lemma-twisted-inverse-image} の右随伴とする。このとき
$\overline{c} \circ \overline{a}_2 = \overline{a}_1$ および
$\overline{d} \circ \overline{a}_3 = \overline{a}_2$ であり、上と同じ理由により函手
$D^+_\QCoh(\mathcal{O}_{\overline{X}_2}) \to D^+_\QCoh(\mathcal{O}_X)$ および
$D^+_\QCoh(\mathcal{O}_{\overline{X}_3}) \to D^+_\QCoh(\mathcal{O}_X)$ の標準変換
$$
j_1^* \circ \overline{c} \longrightarrow j_2^*
\quad\text{and}\quad
j_2^* \circ \overline{d} \longrightarrow j_3^*
$$
がある。$\overline{e}$ を $h \circ g$ に対する補題 \ref{duality-lemma-twisted-inverse-image} の右随伴と書く。(\ref{duality-equation-sheafy}) により、函手
$D^+_\QCoh(\mathcal{O}_{\overline{X}_3}) \to D^+_\QCoh(\mathcal{O}_X)$ の標準変換
$$
j_1^* \circ \overline{e} \longrightarrow j_3^*
$$
がある。書き下すと、合成
$$
\alpha_h \circ \alpha_g :
j_1^* \circ \overline{a}_1 \to
j_1^* \circ \overline{c} \circ \overline{a}_2 \to
j_2^* \circ \overline{a}_2 \to
j_2^* \circ \overline{d} \circ \overline{a}_3 \to
j_3^* \circ \overline{a}_3
$$
が合成
$$
\alpha_{h \circ g} :
j_1^* \circ \overline{a}_1 \to
j_1^* \circ \overline{e} \circ \overline{a}_3 \to
j_3^* \circ \overline{a}_3
$$
に等しいことを示せばよい。これを二つに分ける。第一は、下の水平矢印が同一視
$\overline{e} = \overline{c} \circ \overline{d}$ から得られる次の図式が可換であることを示すことである。
$$
\xymatrix{
\overline{a}_1 \ar[r] \ar[d] & \overline{c} \circ \overline{a}_2 \ar[d] \\
\overline{e} \circ \overline{a}_3 \ar[r] &
\overline{c} \circ \overline{d} \circ \overline{a}_3
}
$$
これは対応する全導来順像函手の図式
$$
\xymatrix{
R\overline{f}_{1, *} \ar[r] \ar[d] & Rg_* \circ R\overline{f}_{2, *} \ar[d] \\
R(h \circ g)_* \circ R\overline{f}_{3, *} \ar[r] &
Rg_* \circ Rh_* \circ R\overline{f}_{3, *}
}
$$
が可換だからである（将来の参照をここに挿入）。第二は、合成
$$
j_1^* \circ \overline{c} \circ \overline{d} \to
j_2^* \circ \overline{d} \to j_3^*
$$
が、同一視 $\overline{e} = \overline{c} \circ \overline{d}$ のもとで写像
$$
j_1^* \circ \overline{e} \to j_3^*
$$
に等しいことを示すことである。これは補題 \ref{duality-lemma-compose-base-change-maps} で示された
（現在の場合、その補題の射 $f', g'$ は $\text{id}_X$ に等しいことに注意する）。
\end{proof}

\begin{lemma}
\label{duality-lemma-upper-shriek-composition}
状況 \ref{duality-situation-shriek} において、$f : X \to Y$ と $g : Y \to Z$ を
$\textit{FTS}_S$ の合成可能な射とする。このとき標準同型
$(g \circ f)^! \to f^! \circ g^!$ がある。
\end{lemma}

\begin{proof}
$Z$ 上の $Y$ のコンパクト化 $i : Y \to \overline{Y}$ を選ぶ。
$\overline{Y}$ 上の $X$ のコンパクト化 $X \to \overline{X}$ を選ぶ。
ここでは『平坦性の続論』の定理 \ref{flat-theorem-nagata} と補題 \ref{flat-lemma-compactifyable} を二度用いる。
$\overline{a}$ を $\overline{X} \to \overline{Y}$ に対する補題 \ref{duality-lemma-twisted-inverse-image} の右随伴とし、$\overline{b}$ を
$\overline{Y} \to Z$ に対する補題 \ref{duality-lemma-twisted-inverse-image} の右随伴とする。
このとき $\overline{a} \circ \overline{b}$ は、合成 $\overline{X} \to Z$ に対する補題 \ref{duality-lemma-twisted-inverse-image} の右随伴である。従って
$g^! = i^* \circ \overline{b}$ かつ
$(g \circ f)^! = (X \to \overline{X})^* \circ \overline{a} \circ \overline{b}$ である。
$U$ を $\overline{X}$ における $Y$ の逆像とすると、可換図式
$$
\xymatrix{
X \ar[r]_j \ar[d] & U \ar[dl] \ar[r]_{j'} & \overline{X} \ar[dl] \\
Y \ar[r]_i \ar[d] & \overline{Y} \ar[dl] \\
Z
}
$$
を得る。$\overline{a}'$ を $U \to Y$ に対する補題 \ref{duality-lemma-twisted-inverse-image} の右随伴とする。このとき
$f^! = j^* \circ \overline{a}'$ である。(\ref{duality-equation-sheafy}) により
$$
\gamma : (j')^* \circ \overline{a} \to \overline{a}' \circ i^*
$$
を得るので、これを用いて
$$
(g \circ f)^! =
(j' \circ j)^* \circ \overline{a} \circ \overline{b} =
j^* \circ (j')^* \circ \overline{a} \circ \overline{b}
\to
j^* \circ \overline{a}' \circ i^* \circ \overline{b} =
f^! \circ g^!
$$
を定義できる。これは補題 \ref{duality-lemma-proper-noetherian} により
$D_\QCoh^+(\mathcal{O}_Z)$ の対象上で同型である。証明を終えるため、この同型が選択に依存しないことを示す。

\medskip\noindent
二つの図式
$$
\vcenter{
\xymatrix{
X \ar[r]_{j_1} \ar[d] & U_1 \ar[dl] \ar[r]_{j'_1} & \overline{X}_1 \ar[dl] \\
Y \ar[r]_{i_1} \ar[d] & \overline{Y}_1 \ar[dl] \\
Z
}
}
\quad\text{and}\quad
\vcenter{
\xymatrix{
X \ar[r]_{j_2} \ar[d] & U_2 \ar[dl] \ar[r]_{j'_2} & \overline{X}_2 \ar[dl] \\
Y \ar[r]_{i_2} \ar[d] & \overline{Y}_2 \ar[dl] \\
Z
}
}
$$
があるとする。まず、$Z$ 上で $Y$ の像が稠密であり、$\overline{Y}_1$ と
$\overline{Y}_2$ の両方を支配するコンパクト化 $i : Y \to \overline{Y}$ を選べる。
『平坦性の続論』の補題 \ref{flat-lemma-compactifications-cofiltered} を参照する。
『平坦性の続論』の補題 \ref{flat-lemma-right-multiplicative-system} および『圏論』の補題 \ref{categories-lemma-morphisms-right-localization} および \ref{categories-lemma-equality-morphisms-right-localization} により、$\overline{Y}$ 上で $X$ の像が稠密な
コンパクト化 $X \to \overline{X}$ と射 $\overline{X} \to \overline{X}_1$ および
$\overline{X} \to \overline{X}_2$ を選べる。さらに、合成
$\overline{X} \to \overline{Y} \to \overline{Y}_1$ は合成
$\overline{X} \to \overline{X}_1 \to \overline{Y}_1$ に等しく、合成
$\overline{X} \to \overline{Y} \to \overline{Y}_2$ は合成
$\overline{X} \to \overline{X}_2 \to \overline{Y}_2$ に等しい。
従って、次の可換図式がある場合に、各図式によって定まる写像を比較すれば十分である。
$$
\xymatrix{
X \ar[rr]_{j_1} \ar@{=}[d] & &
U_1 \ar[d] \ar[ddll] \ar[rr]_{j'_1} & &
\overline{X}_1 \ar[d] \ar[ddll] \\
X \ar'[r][rr]^-{j_2} \ar[d] & &
U_2 \ar'[dl][ddll] \ar'[r][rr]^-{j'_2} & &
\overline{X}_2 \ar[ddll] \\
Y \ar[rr]^{i_1} \ar@{=}[d] & & \overline{Y}_1 \ar[d] \\
Y \ar[rr]^{i_2} \ar[d] & & \overline{Y}_2 \ar[dll] \\
Z
}
$$
さらに、コンパクト化 $X \to \overline{X}_1$ および $Y \to \overline{Y}_2$ の像は稠密である。
$\overline{a}_i$, $\overline{a}'_i$, $\overline{c}$, $\overline{c}'$ をそれぞれ
$\overline{X}_i \to \overline{Y}_i$, $U_i \to Y$,
$\overline{X}_1 \to \overline{X}_2$, $U_1 \to U_2$ に対する補題 \ref{duality-lemma-twisted-inverse-image} の右随伴とする。各平方
$$
\xymatrix{
X \ar[r] \ar[d] \ar@{}[dr]|A & U_1 \ar[d] \\
X \ar[r] & U_2
}
\quad
\xymatrix{
U_2 \ar[r] \ar[d] \ar@{}[dr]|B & \overline{X}_2 \ar[d] \\
Y \ar[r] & \overline{Y}_2
}
\quad
\xymatrix{
U_1 \ar[r] \ar[d] \ar@{}[dr]|C & \overline{X}_1 \ar[d] \\
Y \ar[r] & \overline{Y}_1
}
\quad
\xymatrix{
Y \ar[r] \ar[d] \ar@{}[dr]|D & \overline{Y}_1 \ar[d] \\
Y \ar[r] & \overline{Y}_2
}
\quad
\xymatrix{
X \ar[r] \ar[d] \ar@{}[dr]|E & \overline{X}_1 \ar[d] \\
X \ar[r] & \overline{X}_2
}
$$
はデカルトであり（A, D, E については『平坦性の続論』の補題 \ref{flat-lemma-compactifications-cofiltered} part (c) を参照し、B, C については $U_i$ が
$\overline{X}_i \to \overline{Y}_i$ による $Y$ の逆像であることを想起する）、従って次の基底変換写像
(\ref{duality-equation-sheafy}) を与える。
$$
\begin{matrix}
\gamma_A : j_1^* \circ \overline{c}' \to j_2^* &
\gamma_B : (j_2')^* \circ \overline{a}_2 \to \overline{a}'_2 \circ i_2^* &
\gamma_C : (j_1')^* \circ \overline{a}_1 \to \overline{a}'_1 \circ i_1^* \\
\gamma_D : i_1^* \circ \overline{d} \to i_2^* &
\gamma_E : (j'_1 \circ j_1)^* \circ \overline{c} \to (j'_2 \circ j_2)^*
\end{matrix}
$$
$f_1^! = j_1^* \circ \overline{a}'_1$,
$f_2^! = j_2^* \circ \overline{a}'_2$,
$g_1^! = i_1^* \circ \overline{b}_1$,
$g_2^! = i_2^* \circ \overline{b}_2$,
$(g \circ f)_1^! =
(j_1' \circ j_1)^* \circ \overline{a}_1 \circ \overline{b}_1$, および
$(g \circ f)^!_2 =
(j_2' \circ j_2)^* \circ \overline{a}_2 \circ \overline{b}_2$ と書く。
証明の第一段落および補題 \ref{duality-lemma-shriek-well-defined} の構成では
\begin{enumerate}
\item $(g \circ f)^!_1 \to f_1^! \circ g_1^!$ に $\gamma_C$ を用いる。
\item $(g \circ f)^!_2 \to f_2^! \circ g_2^!$ に $\gamma_B$ を用いる。
\item $f_1^! \to f_2^!$ に $\gamma_A$ を用いる。
\item $g_1^! \to g_2^!$ に $\gamma_D$ を用いる。
\item $(g \circ f)^!_1 \to (g \circ f)^!_2$ に $\gamma_E$ を用いる。
\end{enumerate}
図式
$$
\xymatrix{
(g \circ f)^!_1 \ar[r]_{\gamma_E} \ar[d]_{\gamma_C} &
(g \circ f)^!_2 \ar[d]_{\gamma_B} \\
f_1^! \circ g_1^! \ar[r]^{\gamma_A \circ \gamma_D} & f_2^! \circ g_2^!
}
$$
が可換であることを示さなければならない。補題 \ref{duality-lemma-compose-base-change-maps} および \ref{duality-lemma-compose-base-change-maps-horizontal} と、
注意 \ref{duality-remark-going-around} と同じ記法を濫用して用いる（特に恒等変換との
$\star$ 積を記法から省く）。$\gamma_E = \gamma_A \circ \gamma_F$ と書ける。ここで
$$
\xymatrix{
U_1 \ar[r] \ar[d] \ar@{}[rd]|F & \overline{X}_1 \ar[d] \\
U_2 \ar[r] & \overline{X}_2
}
$$
である。従って
$$
\gamma_B \circ \gamma_E = \gamma_B \circ \gamma_A  \circ \gamma_F
= \gamma_A \circ \gamma_B \circ \gamma_F
$$
である。最後の等号は、二つの平方 $A$ と $B$ が一点でしか交わらないためである
（注意 \ref{duality-remark-going-around} の最後の議論と同様）。従って
$\gamma_D \circ \gamma_C = \gamma_B \circ \gamma_F$ を示せば十分である。
両辺は平方
$$
\xymatrix{
U_1 \ar[r] \ar[d] & \overline{X}_1 \ar[d] \\
Y \ar[r] & \overline{Y}_2
}
$$
に対する写像 (\ref{duality-equation-sheafy}) に等しいので、結論を得る。
\end{proof}

\begin{lemma}
\label{duality-lemma-pseudo-functor}
状況 \ref{duality-situation-shriek} において、補題 \ref{duality-lemma-shriek-well-defined} および \ref{duality-lemma-upper-shriek-composition} の構成は、
圏 $\textit{FTS}_S$ から圏の $2$-圏への擬関手を定める（『圏論』の定義 \ref{categories-definition-functor-into-2-category} を参照する）。
\end{lemma}

\begin{proof}
これを示すには、射 $f : X \to Y$, $g : Y \to Z$, $h : Z \to T$ が与えられたとき、図式
$$
\xymatrix{
(h \circ g \circ f)^! \ar[r]_{\gamma_{A + B}} \ar[d]_{\gamma_{B + C}} &
f^! \circ (h \circ g)^! \ar[d]^{\gamma_C} \\
(g \circ f)^! \circ h^! \ar[r]^{\gamma_A} & f^! \circ g^! \circ h^!
}
$$
が可換であることを示さなければならない（$\gamma$ の意味は以下を参照する）。
そのために、$T$ 上の $Z$ のコンパクト化 $\overline{Z}$、次に $\overline{Z}$ 上の
$Y$ のコンパクト化 $\overline{Y}$、さらに $\overline{Y}$ 上の $X$ のコンパクト化
$\overline{X}$ を選ぶ。ここでは『平坦性の続論』の定理 \ref{flat-theorem-nagata} および補題 \ref{flat-lemma-compactifyable} を用いる。
$W \subset \overline{Y}$ を $\overline{Y} \to \overline{Z}$ による $Z$ の逆像とし、
$U \subset V \subset \overline{X}$ を $\overline{X} \to \overline{Y}$ による
$Y \subset W$ の逆像とする。これにより図式
$$
\xymatrix{
X \ar[d]_f \ar[r] & U \ar[r] \ar[d] \ar@{}[dr]|A &
V \ar[d] \ar[r] \ar@{}[rd]|B & \overline{X} \ar[d] \\
Y \ar[d]_g \ar[r] & Y \ar[r] \ar[d] & W \ar[r] \ar[d] \ar@{}[rd]|C &
\overline{Y} \ar[d] \\
Z \ar[d]_h \ar[r] & Z \ar[d] \ar[r] & Z \ar[d] \ar[r] & \overline{Z} \ar[d] \\
T \ar[r] & T \ar[r] & T \ar[r] & T
}
$$
を得る。大量の記法を導入せず補題 \ref{duality-lemma-upper-shriek-composition} の証明と
まったく同様に論じると、最初の図式の写像は、示した長方形 $A + B$, $B + C$, $A$, $C$ に対する
(\ref{duality-equation-sheafy}) の写像を用いる。補題 \ref{duality-lemma-compose-base-change-maps} および \ref{duality-lemma-compose-base-change-maps-horizontal} により
$\gamma_{A + B} = \gamma_A \circ \gamma_B$ かつ
$\gamma_{B + C} = \gamma_C \circ \gamma_B$ なので、所望の等式を得るには
$\gamma_A \circ \gamma_C = \gamma_C \circ \gamma_A$ であれば十分である。
これは二つの平方 $A$ と $C$ が一点でしか交わらないため成り立つ
（注意 \ref{duality-remark-going-around} の最後の議論と同様）。
\end{proof}

\begin{lemma}
\label{duality-lemma-map-pullback-to-shriek-well-defined}
状況 \ref{duality-situation-shriek} において、$f : X \to Y$ を $\textit{FTS}_S$ の射とする。
$D^+_\QCoh(\mathcal{O}_Y)$ の $K$ に関して関手的な標準写像
$$
\mu_{f, K} :
Lf^*K \otimes_{\mathcal{O}_X}^\mathbf{L} f^!\mathcal{O}_Y
\longrightarrow
f^!K
$$
がある。$g : Y \to Z$ を $\textit{FTS}_S$ の別の射とすると、すべての
$K \in D^+_\QCoh(\mathcal{O}_Z)$ に対して図式
$$
\xymatrix{
Lf^*(Lg^*K \otimes_{\mathcal{O}_Y}^\mathbf{L} g^!\mathcal{O}_Z)
\otimes_{\mathcal{O}_X}^\mathbf{L} f^!\mathcal{O}_Y
\ar@{=}[d] \ar[r]_-{\mu_f} &
f^!(Lg^*K \otimes_{\mathcal{O}_Y}^\mathbf{L} g^!\mathcal{O}_Z)
\ar[r]_-{f^!\mu_g} &
f^!g^!K \ar@{=}[d] \\
Lf^*Lg^*K \otimes_{\mathcal{O}_X}^\mathbf{L} Lf^* g^!\mathcal{O}_Z
\otimes_{\mathcal{O}_X}^\mathbf{L} f^!\mathcal{O}_Y \ar[r]^-{\mu_f} &
Lf^*Lg^*K \otimes_{\mathcal{O}_X}^\mathbf{L} f^!g^!\mathcal{O}_Z
\ar[r]^-{\mu_{g \circ f}} & f^!g^!K
}
$$
は可換である。
\end{lemma}

\begin{proof}
$f$ が固有ならば $f^! = a$ であり、(\ref{duality-equation-compare-with-pullback}) を用いることができる。
$g$ も固有ならば補題 \ref{duality-lemma-transitivity-compare-with-pullback} が、より一般的な形で
図式の可換性を示す。

\medskip\noindent
写像 $\mu_{f, K}$ を定義しよう。$Y$ 上の $X$ のコンパクト化
$j : X \to \overline{X}$ を選ぶ。$f^!$ は $j^* \circ \overline{a}$ として定義されるので、
$\mu_{f, K}$ は写像 (\ref{duality-equation-compare-with-pullback})
$$
L\overline{f}^*K \otimes_{\mathcal{O}_{\overline{X}}}^\mathbf{L}
\overline{a}(\mathcal{O}_Y)
\longrightarrow
\overline{a}(K)
$$
の $X$ への制限として得られる。これがコンパクト化の選択に依存しないことを見るには、
補題 \ref{duality-lemma-shriek-well-defined} の証明と同様に論じる。読者には先にその補題の証明を読むことを勧める。

\medskip\noindent
$Y$ 上のコンパクト化 $j_i : X \to \overline{X}_i$ の間の射
$g : \overline{X}_1 \to \overline{X}_2$ で、$g^{-1}(j_2(X)) = j_1(X)$ を満たすものが
与えられたとする。$\overline{c}$ を、射 $g$ に対する補題 \ref{duality-lemma-twisted-inverse-image} の押し出しの右随伴とする。写像
$$
L\overline{f}_1^*K \otimes_{\mathcal{O}_{\overline{X}}}^\mathbf{L}
\overline{a}_1(\mathcal{O}_Y)
\longrightarrow
\overline{a}_1(K)
\quad\text{and}\quad
L\overline{f}_2^*K \otimes_{\mathcal{O}_{\overline{X}}}^\mathbf{L}
\overline{a}_2(\mathcal{O}_Y)
\longrightarrow
\overline{a}_2(K)
$$
は可換図式
$$
\xymatrix{
Lg^*(L\overline{f}_2^*K \otimes^\mathbf{L}
\overline{a}_2(\mathcal{O}_Y))
\otimes^\mathbf{L} \overline{c}(\mathcal{O}_{\overline{X}_2})
\ar@{=}[d] \ar[r]_-\sigma &
\overline{c}(L\overline{f}_2^*K \otimes^\mathbf{L}
\overline{a}_2(\mathcal{O}_Y)) \ar[r] &
\overline{c}(\overline{a}_2(K)) \ar@{=}[d] \\
L\overline{f}_1^*K \otimes^\mathbf{L} Lg^*\overline{a}_2(\mathcal{O}_Y)
\otimes^\mathbf{L} \overline{c}(\mathcal{O}_{\overline{X}_2})
\ar[r]^-{1 \otimes \tau} &
L\overline{f}_1^*K \otimes^\mathbf{L} \overline{a}_1(\mathcal{O}_Y) \ar[r] &
\overline{a}_1(K)
}
$$
に入る。これは補題 \ref{duality-lemma-transitivity-compare-with-pullback} による。
補題 \ref{duality-lemma-compare-on-open} により、写像 $\sigma$ と $\tau$ は $X$ 上で同型に制限される。
実際、さらに詳しく述べられる。補題 \ref{duality-lemma-shriek-well-defined} の証明では、写像
(\ref{duality-equation-sheafy}) $\gamma : j_1^* \circ \overline{c} \to j_2^*$ を用いて同型
$\alpha_g : j_1^* \circ \overline{a}_1 \to j_2^* \circ \overline{a}_2$ を構成したことを想起する。
写像 $\sigma$ を $j_1$ で引き戻すと、$j_1^* \circ \overline{c} \to j_2^*$ を介して
$j_1^*\overline{c}(\mathcal{O}_{\overline{X}_2})$ を $\mathcal{O}_X$ と同一視するならば、
$j_2^*\left(L\overline{f}_2^*K \otimes^\mathbf{L}
\overline{a}_2(\mathcal{O}_Y)\right)$ 上の恒等写像を得る。補題 \ref{duality-lemma-restriction-compare-with-pullback} を参照する。同様に、写像
$\tau : Lg^*\overline{a}_2(\mathcal{O}_Y)
\otimes^\mathbf{L} \overline{c}(\mathcal{O}_{\overline{X}_2}) \to
\overline{a}_1(\mathcal{O}_Y) = \overline{c}(\overline{a}_2(\mathcal{O}_Y))$
は $j_2^*\overline{a}_2(\mathcal{O}_Y)$ 上の恒等写像に引き戻される。
従って $j_1$ で引き戻し、可能な箇所で $\gamma$ を適用すると、可換図式
$$
\xymatrix{
j_2^*\left(L\overline{f}_2^*K \otimes^\mathbf{L}
\overline{a}_2(\mathcal{O}_Y)\right) \ar[r] \ar[d] &
j_2^*\overline{a}_2(K) \\
j_1^*L\overline{f}_1^*K \otimes^\mathbf{L} j_2^*\overline{a}_2(\mathcal{O}_Y) &
j_1^*(L\overline{f}_1^*K \otimes^\mathbf{L} \overline{a}_1(\mathcal{O}_Y))
\ar[r] \ar[l]_{1 \otimes \alpha_g} &
j_1^* \overline{a}_1(K) \ar[lu]_{\alpha_g}
}
$$
を得る。この図式の可換性は、コンパクト化 $\overline{X}_1$ を用いて構成した写像
$\mu_{f, K}$ が、補題 \ref{duality-lemma-shriek-well-defined} の証明で用いた同一視 $\alpha_g$ を介して、
コンパクト化 $\overline{X}_2$ を用いて構成した写像 $\mu_{f, K}$ と等しいことをまさに示す。
補題 \ref{duality-lemma-shriek-well-defined} の証明とまったく同じ圏論的議論により、
$\mu_{f, K}$ は良定義である（小さな細部は省略する）。

\medskip\noindent
以上を踏まえると、補題の主張における図式の可換性は、同型
$(g \circ f)^! \to f^! \circ g^!$ の構成（補題 \ref{duality-lemma-upper-shriek-composition} の証明の第一部で
$\overline{X} \to \overline{Y} \to Z$ を用いたもの）と、
$\overline{X} \to \overline{Y} \to Z$ に対する補題 \ref{duality-lemma-transitivity-compare-with-pullback} の結果から従う。
\end{proof}



\section{上付きシュリーク函手の性質}
\label{duality-section-upper-shriek-properties}

\noindent
上付きシュリーク函手のいくつかの性質を述べる。

\begin{lemma}
\label{duality-lemma-shriek-open-immersion}
状況 \ref{duality-situation-shriek} において、$Y$ を $\textit{FTS}_S$ の対象とし、
$j : X \to Y$ を開埋め込みとする。このとき函手の標準同型 $j^! = j^*$ がある。
\end{lemma}

\noindent
$\textit{FTS}_S$ の \'etale 射 $f : X \to Y$ に対しても $f^* \cong f^!$ である。
補題 \ref{duality-lemma-shriek-etale} を参照する。

\begin{proof}
この場合、コンパクト化として $\overline{X} = Y$ を選べる。
$\text{id} : Y \to Y$ に対する補題 \ref{duality-lemma-twisted-inverse-image} の右随伴は恒等函手なので、
定義により $j^! = j^*$ である。
\end{proof}

\begin{lemma}
\label{duality-lemma-restrict-before-or-after}
状況 \ref{duality-situation-shriek} において、
$$
\xymatrix{
U \ar[r]_j \ar[d]_g & X \ar[d]^f \\
V \ar[r]^{j'} & Y
}
$$
を $\textit{FTS}_S$ の可換図式とし、$j$ と $j'$ は開埋め込みであるとする。
このとき函手 $D^+_\QCoh(\mathcal{O}_Y) \to D^+(\mathcal{O}_U)$ として
$j^* \circ f^! = g^! \circ (j')^*$ である。
\end{lemma}

\begin{proof}
$h = f \circ j = j' \circ g$ とおく。補題 \ref{duality-lemma-upper-shriek-composition} により
$h^! = j^! \circ f^! = g^! \circ (j')^!$ である。補題 \ref{duality-lemma-shriek-open-immersion} により $j^! = j^*$ かつ $(j')^! = (j')^*$ である。
\end{proof}

\begin{lemma}
\label{duality-lemma-shriek-affine-line}
状況 \ref{duality-situation-shriek} において、$Y$ を $\textit{FTS}_S$ の対象とし、
$f : X = \mathbf{A}^1_Y \to Y$ を射影とする。このとき函手の（非標準的）同型
$f^!(-) \cong Lf^*(-) [1]$ がある。
\end{lemma}

\begin{proof}
$X = \mathbf{A}^1_Y \subset \mathbf{P}^1_Y$ かつ
$\mathcal{O}_{\mathbf{P}^1_Y}(-2)|_X \cong \mathcal{O}_X$ なので、主張は補題 \ref{duality-lemma-upper-shriek-P1} および \ref{duality-lemma-compare-with-pullback-flat-proper-noetherian} から従う。
\end{proof}

\begin{lemma}
\label{duality-lemma-shriek-closed-immersion}
状況 \ref{duality-situation-shriek} において、$Y$ を $\textit{FTS}_S$ の対象とし、
$i : X \to Y$ を閉埋め込みとする。このとき函手の標準同型
$i^!(-) = R\SheafHom(\mathcal{O}_X, -)$ がある。
\end{lemma}

\begin{proof}
これは補題 \ref{duality-lemma-twisted-inverse-image-closed} の言い換えである。
\end{proof}

\begin{remark}[上付きシュリークの局所記述]
\label{duality-remark-local-calculation-shriek}
状況 \ref{duality-situation-shriek} において、$f : X \to Y$ を $\textit{FTS}_S$ の射とする。
上の補題を用いると、$f^!$ を次のように局所的に計算できる。アフィン開集合の図式
$$
\xymatrix{
U \ar[r]_j \ar[d]_g & X \ar[d]^f \\
V \ar[r]^i & Y
}
$$
が与えられたとする。補題 \ref{duality-lemma-upper-shriek-composition} により
$j^! \circ f^! = g^! \circ i^!$ であり、補題 \ref{duality-lemma-shriek-open-immersion} により
$j^!$ と $i^!$ は制限で与えられるので、任意の $E \in D^+_\QCoh(\mathcal{O}_X)$ に対して
$$
(f^!E)|_U = g^!(E|_V)
$$
である。$U = \Spec(A)$ および $V = \Spec(R)$ と書き、
$\varphi : R \to A$ を $g$ に対応する有限型環写像とする。
$P = R[x_1, \ldots, x_n]$ が $R$ 上 $n$ 変数の多項式代数である表示 $A = P/I$ を選ぶ。
$E|_V$ に対応する対象 $K \in D^+(R)$ を選ぶ（『スキームの導来圏』の補題 \ref{perfect-lemma-affine-compare-bounded}）。このとき $f^!E|_U$ は
$$
\varphi^!(K) = R\Hom(A, K \otimes_R^\mathbf{L} P)[n]
$$
に対応すると主張する。ここで $R\Hom(A, -) : D(P) \to D(A)$ は『双対化複体』の第 \ref{dualizing-section-trivial} 節の函手であり、$\varphi^! : D(R) \to D(A)$ は
『双対化複体』の第 \ref{dualizing-section-relative-dualizing-complex-algebraic} 節の函手である。
実際、表示の選択は分解
$$
U \rightarrow \mathbf{A}^n_V \to \mathbf{A}^{n - 1}_V \to \ldots \to
\mathbf{A}^1_V \to V
$$
を与える。補題 \ref{duality-lemma-shriek-affine-line} をちょうど $n$ 回適用すると、
$(\mathbf{A}^n_V \to V)^!(E|_V)$ は $K \otimes_R^\mathbf{L} P[n]$ に対応する。
補題 \ref{duality-lemma-sheaf-with-exact-support-quasi-coherent} および \ref{duality-lemma-shriek-closed-immersion} により、最後の段階は $R\Hom(A, -)$ の適用に対応する。
\end{remark}

\begin{lemma}
\label{duality-lemma-shriek-coherent}
状況 \ref{duality-situation-shriek} において、$f : X \to Y$ を $\textit{FTS}_S$ の射とする。
このとき $f^!$ は $D_{\textit{Coh}}^+(\mathcal{O}_Y)$ を
$D_{\textit{Coh}}^+(\mathcal{O}_X)$ に写す。
\end{lemma}

\begin{proof}
問題は $X$ 上局所的なので、$X$ と $Y$ はアフィンスキームであると仮定してよい。
この場合 $f : X \to Y$ は
$$
X \xrightarrow{i} \mathbf{A}^n_Y \to \mathbf{A}^{n - 1}_Y \to \ldots \to
\mathbf{A}^1_Y \to Y
$$
と分解できる。ここで $i$ は閉埋め込みである。補題は補題 \ref{duality-lemma-shriek-affine-line} および \ref{duality-lemma-sheaf-with-exact-support-coherent},
『双対化複体』の補題 \ref{dualizing-lemma-dualizing-polynomial-ring} と帰納法から従う。
\end{proof}

\begin{lemma}
\label{duality-lemma-shriek-dualizing}
状況 \ref{duality-situation-shriek} において、$f : X \to Y$ を $\textit{FTS}_S$ の射とする。
$K$ が $Y$ の双対化複体ならば、$f^!K$ は $X$ の双対化複体である。
\end{lemma}

\begin{proof}
問題は $X$ 上局所的なので、$X$ と $Y$ はアフィンスキームであると仮定してよい。
この場合 $f : X \to Y$ は
$$
X \xrightarrow{i} \mathbf{A}^n_Y \to \mathbf{A}^{n - 1}_Y \to \ldots \to
\mathbf{A}^1_Y \to Y
$$
と分解できる。ここで $i$ は閉埋め込みである。補題 \ref{duality-lemma-shriek-affine-line},
『双対化複体』の補題 \ref{dualizing-lemma-dualizing-polynomial-ring} と帰納法により、
$p : \mathbf{A}^n_Y \to Y$ を射影とすると $p^!K$ は $\mathbf{A}^n_Y$ 上の双対化複体である。
同様に、『双対化複体』の補題 \ref{dualizing-lemma-dualizing-quotient} および補題 \ref{duality-lemma-sheaf-with-exact-support-quasi-coherent} および \ref{duality-lemma-shriek-closed-immersion} により、$i^!$ は双対化複体を双対化複体に写す。
\end{proof}

\begin{lemma}
\label{duality-lemma-shriek-via-duality}
状況 \ref{duality-situation-shriek} において、$f : X \to Y$ を $\textit{FTS}_S$ の射とし、
$K$ を $Y$ 上の双対化複体とする。$M \in D_{\textit{Coh}}(\mathcal{O}_Y)$ に対して
$D_Y(M) = R\SheafHom_{\mathcal{O}_Y}(M, K)$ とおき、
$E \in D_{\textit{Coh}}(\mathcal{O}_X)$ に対して
$D_X(E) = R\SheafHom_{\mathcal{O}_X}(E, f^!K)$ とおく。
このとき $M \in D_{\textit{Coh}}^+(\mathcal{O}_Y)$ に対して標準同型
$$
f^!M \longrightarrow D_X(Lf^*D_Y(M))
$$
がある。
\end{lemma}

\begin{proof}
$Y$ 上の $X$ のコンパクト化 $j : X \subset \overline{X}$ を選ぶ
（『平坦性の続論』の定理 \ref{flat-theorem-nagata} および補題 \ref{flat-lemma-compactifyable}）。$a$ を $\overline{X} \to Y$ に対する補題 \ref{duality-lemma-twisted-inverse-image} の右随伴とする。
$E \in D_{\textit{Coh}}(\mathcal{O}_{\overline{X}})$ に対して
$D_{\overline{X}}(E) = R\SheafHom_{\mathcal{O}_{\overline{X}}}(E, a(K))$ とおく。
$R\SheafHom$ の構成は開集合への制限と可換し、$f^! = j^* \circ a$ なので、
$M \in D_{\textit{Coh}}(\mathcal{O}_Y)$ に対して標準同型
$$
a(M) \longrightarrow D_{\overline{X}}(L\overline{f}^*D_Y(M))
$$
があることを示せば十分である。$F \in D_\QCoh(\mathcal{O}_X)$ に対して
\begin{align*}
\Hom_{\overline{X}}(
F, D_{\overline{X}}(L\overline{f}^*D_Y(M)))
& =
\Hom_{\overline{X}}(
F \otimes_{\mathcal{O}_X}^\mathbf{L} L\overline{f}^*D_Y(M), a(K)) \\
& =
\Hom_Y(
R\overline{f}_*(F \otimes_{\mathcal{O}_X}^\mathbf{L} L\overline{f}^*D_Y(M)),
K) \\
& =
\Hom_Y(
R\overline{f}_*(F) \otimes_{\mathcal{O}_Y}^\mathbf{L} D_Y(M),
K) \\
& =
\Hom_Y(
R\overline{f}_*(F), D_Y(D_Y(M))) \\
& =
\Hom_Y(R\overline{f}_*(F), M) \\
& = \Hom_{\overline{X}}(F, a(M))
\end{align*}
である。第一の等号は『層のコホモロジー』の補題 \ref{cohomology-lemma-internal-hom} による。
第二の等号は $a$ の定義による。第三の等号は『スキームの導来圏』の補題 \ref{perfect-lemma-cohomology-base-change} による。第四の等号は『層のコホモロジー』の補題 \ref{cohomology-lemma-internal-hom} および $D_Y$ の定義による。
第五の等号は補題 \ref{duality-lemma-dualizing-schemes} による。最後の等号は $a$ の定義による。
従って米田の補題により
$a(M) = D_{\overline{X}}(L\overline{f}^*D_Y(M))$ である。
\end{proof}

\begin{lemma}
\label{duality-lemma-perfect-comparison-shriek}
状況 \ref{duality-situation-shriek} において、$f : X \to Y$ を
$\textit{FTS}_S$ の射とする。$f$ は完全である（例えば平坦である）と仮定する。このとき
\begin{enumerate}
\item[(a)] $f^!\mathcal{O}_Y$ は $D_{\textit{Coh}}^b(\mathcal{O}_X)$ に属する。
\item[(b)] $f^!\mathcal{O}_Y$ は $D(f^{-1}\mathcal{O}_Y)$ において有限 Tor 次元をもつ。
\item[(c)] $\mathcal{O}_X \to
R\SheafHom_{\mathcal{O}_X}(f^!\mathcal{O}_Y, f^!\mathcal{O}_Y)$
は同型である。
\item[(d)] $f^!$ は $D_{\textit{Coh}}^b(\mathcal{O}_Y)$ を
$D_{\textit{Coh}}^b(\mathcal{O}_X)$ に写す。
\item[(e)] 補題 \ref{duality-lemma-map-pullback-to-shriek-well-defined} の写像
$\mu_{f,  K} :
Lf^*K \otimes_{\mathcal{O}_X}^\mathbf{L} f^!\mathcal{O}_Y
\to
f^!K$
は、すべての $K \in D_\QCoh^+(\mathcal{O}_Y)$ に対して同型である。
\end{enumerate}
\end{lemma}

\begin{proof}
（有限表示な平坦射は完全である。『射の続論』の補題 \ref{more-morphisms-lemma-flat-finite-presentation-perfect} を参照する。）
主張 (a), (b), (c) は $X$ 上局所的である。従って $X$ と $Y$ はアフィンであると仮定してよい。
すると注意 \ref{duality-remark-local-calculation-shriek} により、(a), (b), (c) はそれぞれ
『双対化複体』の補題 \ref{dualizing-lemma-relative-dualizing-algebraic} の
(1)(a), (1)(b), (1)(c) に帰着する。
（主張を対応させるには『スキームの導来圏』の補題 \ref{perfect-lemma-tor-dimension-rel-affine}、\ref{perfect-lemma-pseudo-coherent} および \ref{perfect-lemma-identify-pseudo-coherent-noetherian}
\ref{perfect-lemma-quasi-coherence-internal-hom} を用いる。）

\medskip\noindent
(d) と (e) を証明するため、まずいくつかの予備的注意を述べる。
\begin{enumerate}
\item $f^!$ が $D_{\textit{Coh}}^+(\mathcal{O}_Y)$ を
$D_{\textit{Coh}}^+(\mathcal{O}_X)$ に写すことは既知である。補題 \ref{duality-lemma-shriek-coherent} を参照する。
\item $f$ が開埋め込みならば $f^! = f^*$ であり、$f^!$ と $\mu_f$ の構成において
$\overline{X} = Y$ と取れるため、(d) と (e) は成り立つ。補題 \ref{duality-lemma-shriek-open-immersion} を参照する。
\item $f$ が完全固有射ならば、補題 \ref{duality-lemma-compare-with-pullback-flat-proper-noetherian} により (e) は成り立つ。
\item 開被覆 $X = \bigcup U_i$ が存在し、$U_i \to Y$ に対して (d) が成り立つならば、
$X \to Y$ に対しても (d) が成り立つ。(e) についても同様である。
これは $f^!$ と $\mu_f$ の構成が開部分スキームへの移行と可換だからである。
\item $g : Y \to Z$ が $\textit{FTS}_S$ の第二の完全射であり、$f$ と $g$ に対して
(e) が成り立つならば、
$f^!g^!\mathcal{O}_Z =
Lf^*g^!\mathcal{O}_Z \otimes_{\mathcal{O}_X}^\mathbf{L} f^!\mathcal{O}_Y$
であり、補題 \ref{duality-lemma-map-pullback-to-shriek-well-defined} の可換図式により
$g \circ f$ に対しても (e) が成り立つ。
\item $f$ と $g$ の双方に対して (d) と (e) が成り立つならば、
$g \circ f$ に対しても (d) と (e) が成り立つ。実際、このとき $f^!g^!\mathcal{O}_Z$ は
（前項により）上に有界であり、$L(g \circ f)^*$ は有限コホモロジー次元をもつので、
上で見た (e) から (d) が従う。
\end{enumerate}
これらから、$X$ がアフィンの場合に (d) と (e) を示せば十分であることが分かる。
埋め込み $X \to \mathbf{A}^n_Y$ を選ぶ
（『スキームの射』の補題 \ref{morphisms-lemma-quasi-affine-finite-type-over-S}）。これを
$X \to U \to \mathbf{A}^n_Y \to Y$ と分解する。ここで $X \to U$ は閉埋め込みであり、
$U \subset \mathbf{A}^n_Y$ は開である。『射の続論』の補題 \ref{more-morphisms-lemma-perfect-permanence} により、$X \to U$ は完全閉埋め込みであることに注意する。
従って、完全閉埋め込みと射影 $\mathbf{A}^n_Y \to Y$ に対して補題を示せば十分である。

\medskip\noindent
$f : X \to Y$ を完全閉埋め込みとする。(d) が成り立つことはすでに分かっている。
$K \in D^b_{\textit{Coh}}(\mathcal{O}_Y)$ とする。このとき
$f^!K = R\SheafHom(\mathcal{O}_X, K)$
（補題 \ref{duality-lemma-shriek-closed-immersion}）および
$f_*f^!K = R\SheafHom_{\mathcal{O}_Y}(f_*\mathcal{O}_X, K)$ である。
$f$ は完全なので、複体 $f_*\mathcal{O}_X$ は完全であり、従って
$R\SheafHom_{\mathcal{O}_Y}(f_*\mathcal{O}_X, K)$ は上に有界である。
これで (d) が成り立つことが示された。細部は省略する。

\medskip\noindent
$f : \mathbf{A}^n_Y \to Y$ を射影とする。このとき補題 \ref{duality-lemma-shriek-affine-line} を繰り返し適用すれば (d) が成り立つ。
最後に、(e) は $\mathbf{P}^n_Y \to Y$ に対して成り立ち（平坦かつ固有）、かつ
$\mathbf{A}^n_Y \subset \mathbf{P}^n_Y$ が開であるから、(e) も成り立つ。
\end{proof}

\begin{lemma}
\label{duality-lemma-flat-shriek-relatively-perfect}
状況 \ref{duality-situation-shriek} において、$f : X \to Y$ を
$\textit{FTS}_S$ の射とする。$f$ が平坦ならば、
$f^!\mathcal{O}_Y$ は $D(\mathcal{O}_X)$ の $Y$-完全対象であり、
$\mathcal{O}_X \to
R\SheafHom_{\mathcal{O}_X}(f^!\mathcal{O}_Y, f^!\mathcal{O}_Y)$
は同型である。
\end{lemma}

\begin{proof}
いずれの主張も $X$ 上局所的である。従って $X$ と $Y$ はアフィンであると仮定してよい。
すると注意 \ref{duality-remark-local-calculation-shriek} により、この補題は代数の補題、すなわち
『双対化複体』の補題 \ref{dualizing-lemma-relative-dualizing-algebraic} に帰着する。
（用語を対応させるには『スキームの導来圏』の補題 \ref{perfect-lemma-affine-locally-rel-perfect} を用いる。）
\end{proof}

\begin{lemma}
\label{duality-lemma-lci-shriek}
状況 \ref{duality-situation-shriek} において、$f : X \to Y$ を
$\textit{FTS}_S$ の射とする。$f : X \to Y$ は局所完全交叉射であると仮定する。このとき
\begin{enumerate}
\item $f^!\mathcal{O}_Y$ は $D(\mathcal{O}_X)$ の可逆対象である。
\item $f^!$ は完全複体を完全複体に写す。
\end{enumerate}
\end{lemma}

\begin{proof}
局所完全交叉射が完全であることを想起する。『射の続論』の補題 \ref{more-morphisms-lemma-lci-properties} を参照する。
補題 \ref{duality-lemma-perfect-comparison-shriek} により、
$f^!\mathcal{O}_Y$ が $D(\mathcal{O}_X)$ の可逆対象であることを示せば十分である。
この問題は $X$ と $Y$ 上局所的である。従って、$X \to Y$ は
$X \to \mathbf{A}^n_Y \to Y$ と分解し、最初の矢印は Koszul 正則埋め込みであると仮定してよい。
『射の続論』の第 \ref{more-morphisms-section-lci} 節を参照する。
補題 \ref{duality-lemma-shriek-affine-line} により、結果は $\mathbf{A}^n_Y \to Y$ に対して成り立つ。
従って $f$ が Koszul 正則埋め込みである場合に補題を示せば十分である。
もう一度局所的に考えることにより、$X = \Spec(A)$, $Y = \Spec(B)$ であり、
ある正則列 $f_1, \ldots, f_r \in B$ に対して $A = B/(f_1, \ldots, f_r)$ である場合に帰着する
（Noether 局所環に対して Koszul 正則と正則の概念は同じであることを用いる。『代数の続論』の補題 \ref{more-algebra-lemma-noetherian-finite-all-equivalent} を参照する）。
従って $X \to Y$ は合成
$$
X = X_r \to X_{r - 1} \to \ldots \to X_1 \to X_0 = Y
$$
であり、各矢印は有効 Cartier 因子の包含である。このようにして、有効 Cartier 因子の包含
$i : D \to X$ の場合に帰着する。この場合、補題 \ref{duality-lemma-compute-for-effective-Cartier} により
$i^!\mathcal{O}_X = \mathcal{N}[1]$ であり、証明は完了する。
\end{proof}







\section{上付きシュリークに対する基底変換}
\label{duality-section-base-change-shriek}

\noindent
状況 \ref{duality-situation-shriek} において、
$$
\xymatrix{
X' \ar[r]_{g'} \ar[d]_{f'} & X \ar[d]^f \\
Y' \ar[r]^g & Y
}
$$
を $\textit{FTS}_S$ における Cartesian 図式とし、$X$ と $Y'$ は $Y$ 上 Tor 独立であるとする。
現在の設定は、この一般性で基底変換写像
$L(g')^* \circ f^! \to (f')^! \circ Lg^*$ を構成するには不十分である。
その理由は、一般には $Y$ 上のコンパクト化 $j : X \to \overline{X}$ を、
$\overline{X}$ と $Y'$ が $Y$ 上 Tor 独立となるように選ぶことができず、従って
第 \ref{duality-section-base-change-map} 節における基底変換写像の構成を適用できないからである\footnote{
導来代数幾何に通じた読者ならば、これは「本質的な」問題ではないことが分かるであろう。
実際、$\overline{X}'$ を $\overline{X}$ と $Y'$ の $Y$ 上の導来ファイバー積とすれば、
補題 \ref{duality-lemma-base-change-shriek-flat} の証明とまったく同様に論じてこの写像を定義できる。
結局、$X$ と $Y'$ の Tor 独立性により、$X'$ は導来スキーム
$\overline{X}'$ の開部分スキームとなる。}。

\medskip\noindent
部分的な解決は第 \ref{duality-section-relative-dualizing-complexes} 節で得られる。
すなわち、射 $f$ が平坦ならば相対双対化複体の適切な概念があり、補題 \ref{duality-lemma-compactifyable-relative-dualizing}
\ref{duality-lemma-base-change-relative-dualizing},
and \ref{duality-lemma-perfect-comparison-shriek}
を用いて標準的な基底変換同型を構成できる。これを使用する必要が生じれば、
本章の後段で精密な主張と証明を加える。

\begin{lemma}
\label{duality-lemma-base-change-shriek-flat}
状況 \ref{duality-situation-shriek} において、
$$
\xymatrix{
X' \ar[r]_{g'} \ar[d]_{f'} & X \ar[d]^f \\
Y' \ar[r]^g & Y
}
$$
を $g$ が平坦であるような $\textit{FTS}_S$ の Cartesian 図式とする。
このとき $D_\QCoh^+(\mathcal{O}_Y)$ 上に同型
$L(g')^* \circ f^! \to (f')^! \circ Lg^*$ が存在する。
\end{lemma}

\begin{proof}
$g$ は平坦なので、$Y$ 上の $X$ の任意のコンパクト化
$j : X \to \overline{X}$ に対して、スキーム $\overline{X}$ は $Y'$ と Tor 独立である。
$j$ の基底変換を $j' : X' \to \overline{X}'$ と書き、射影を
$\overline{g}' : \overline{X}' \to \overline{X}$ と書く。基底変換写像を合成
$$
L(g')^* \circ f^! = L(g')^* \circ j^* \circ a =
(j')^* \circ L(\overline{g}')^* \circ a \longrightarrow
(j')^* \circ a' \circ Lg^* = (f')^! \circ Lg^*
$$
として定義する。ここで中央の矢印は基底変換写像
(\ref{duality-equation-base-change-map}) であり、$a$ と $a'$ はそれぞれ
$\overline{X} \to Y$ と $\overline{X}' \to Y'$ に対する補題 \ref{duality-lemma-twisted-inverse-image} の押し出しの右随伴である。
この構成はコンパクト化の選択に依存しない
（この結果を使用する必要が生じれば、精密な補題を定式化して証明する）。

\medskip\noindent
証明を終えるには、基底変換写像 $L(g')^* \circ a \to a' \circ Lg^*$ が
$D_\QCoh^+(\mathcal{O}_Y)$ 上の同型であることを示せば十分である。
補題 \ref{duality-lemma-proper-noetherian} により、$a$ と $a'$ の構成は
$Y$ と $Y'$ のアフィン開集合への制限と可換である。従って注意 \ref{duality-remark-check-over-affines} により、$Y$ と $Y'$ はアフィンであると仮定してよい。
このとき結果は補題 \ref{duality-lemma-more-base-change} から従う。
\end{proof}

\begin{lemma}
\label{duality-lemma-shriek-etale}
状況 \ref{duality-situation-shriek} において、$f : X \to Y$ を
$\textit{FTS}_S$ の \'etale 射とする。このとき $D^+_\QCoh(\mathcal{O}_Y)$ 上の関手として
$f^! \cong f^*$ である。
\end{lemma}

\begin{proof}
\'etale 射が平坦かつ syntomic であり、局所完全交叉射でもあることを用いる
（『スキームの射』の補題 \ref{morphisms-lemma-etale-syntomic} および \ref{morphisms-lemma-etale-flat} and
『射の続論』の補題 \ref{more-morphisms-lemma-flat-lci}）。
補題 \ref{duality-lemma-perfect-comparison-shriek} により、
$f^!\mathcal{O}_Y = \mathcal{O}_X$ を示せば十分である。
補題 \ref{duality-lemma-lci-shriek} により、$f^!\mathcal{O}_Y$ は可逆加群である。
可換図式
$$
\xymatrix{
X \times_Y X \ar[r]_{p_2} \ar[d]_{p_1} & X \ar[d]^f \\
X \ar[r]^f & Y
}
$$
と対角射 $\Delta : X \to X \times_Y X$ を考える。$\Delta$ は開埋め込みなので
（『スキームの射』の補題 \ref{morphisms-lemma-diagonal-unramified-morphism} および \ref{morphisms-lemma-etale-smooth-unramified}）、補題 \ref{duality-lemma-shriek-open-immersion} により $\Delta^! = \Delta^*$ である。
補題 \ref{duality-lemma-upper-shriek-composition} により
$\Delta^! \circ p_1^! \circ f^! = f^!$ である。
この図式に補題 \ref{duality-lemma-base-change-shriek-flat} を適用すると
$p_1^!\mathcal{O}_X = p_2^*f^!\mathcal{O}_Y$ である。従って
$$
f^!\mathcal{O}_Y = \Delta^!p_1^!f^!\mathcal{O}_Y =
\Delta^*(p_1^*f^!\mathcal{O}_Y \otimes p_1^!\mathcal{O}_X) =
\Delta^*(p_2^*f^!\mathcal{O}_Y \otimes p_1^*f^!\mathcal{O}_Y) =
(f^!\mathcal{O}_Y)^{\otimes 2}
$$
を得る。第二段階では補題 \ref{duality-lemma-perfect-comparison-shriek} をもう一度用いた。
ゆえに所望のとおり $f^!\mathcal{O}_Y = \mathcal{O}_X$ である。
\end{proof}


\noindent
本節の残りでは、基底変換写像の適切な理論から帰結する、証明容易な結果をいくつか定式化する。

\begin{lemma}[暫定的基底変換]
\label{duality-lemma-base-change-locally}
状況 \ref{duality-situation-shriek} において、
$$
\xymatrix{
X' \ar[r]_{g'} \ar[d]_{f'} & X \ar[d]^f \\
Y' \ar[r]^g & Y
}
$$
を $\textit{FTS}_S$ の Cartesian 図式とする。
$E \in D^+_\QCoh(\mathcal{O}_Y)$ を、$Lg^*E$ が $D^+(\mathcal{O}_Y)$ に属するような対象とする。
$f$ が平坦ならば、$Y$, $Y'$, $X$ のアフィン開集合へ写る任意のアフィン開集合
$U' \subset X'$ に対し、$L(g')^*f^!E$ と $(f')^!Lg^*E$ の制限は
$D(\mathcal{O}_{U'})$ の同型な対象である。
\end{lemma}

\begin{proof}
仮定により、直ちに $X$, $Y$, $Y'$, $X'$ がアフィンである場合に帰着する。
$Y = \Spec(R)$, $Y' = \Spec(R')$, $X = \Spec(A)$, $X' = \Spec(A')$ と書く。
このとき $A' = A \otimes_R R'$ である。$E$ は $K \in D^+(R)$ に対応するとする。
与えられた写像を $\varphi : R \to A$, $\varphi' : R' \to A'$ と書くと、注意 \ref{duality-remark-local-calculation-shriek} により、$L(g')^*f^!E$ と $(f')^!Lg^*E$ はそれぞれ
$\varphi^!(K) \otimes_A^\mathbf{L} A'$ と
$(\varphi')^!(K \otimes_R^\mathbf{L} R')$ に対応する。ここで $\varphi^!$ と
$(\varphi')^!$ は『双対化複体』の第 \ref{dualizing-section-relative-dualizing-complex-algebraic} 節の関手である。
結果は『双対化複体』の補題 \ref{dualizing-lemma-bc-flat} から従う。
\end{proof}

\begin{lemma}
\label{duality-lemma-relative-dualizing-fibres}
状況 \ref{duality-situation-shriek} において、$f : X \to Y$ を
$\textit{FTS}_S$ の射とし、$f$ は平坦であると仮定する。
$D^b_{\textit{Coh}}(X)$ において $\omega_{X/Y}^\bullet = f^!\mathcal{O}_Y$ とおく。
$y \in Y$ とし、$h : X_y \to X$ を射影とする。このとき
$Lh^*\omega_{X/Y}^\bullet$ は $X_y$ 上の双対化複体である。
\end{lemma}

\begin{proof}
補題 \ref{duality-lemma-perfect-comparison-shriek} により、複体
$\omega_{X/Y}^\bullet$ は $D^b_{\textit{Coh}}$ に属する。
双対化複体であることは局所的性質である。従って補題 \ref{duality-lemma-base-change-locally} により、$(X_y \to y)^!\mathcal{O}_y$ が
$X_y$ 上の双対化複体であることを示せば十分である。
これは補題 \ref{duality-lemma-shriek-dualizing} から従う。
\end{proof}








\section{双対理論}
\label{duality-section-duality}

\noindent
本節では、上で得た非常に一般的な結果が、固定した Noether 基底スキーム上の有限型分離スキームに
どのような双対理論を与えるかを明示する。

\medskip\noindent
Noether スキーム $X$ 上の双対化複体とは、アフィン局所的に対応する環の双対化複体を与える
$D(\mathcal{O}_X)$ の対象であったことを想起する。定義 \ref{duality-definition-dualizing-scheme} を参照する。

\medskip\noindent
Noether スキーム $S$ に対し、$S$ 上有限型かつ分離なスキームの圏を
$\textit{FTS}_S$ と書く。このとき次が成り立つ。
\begin{enumerate}
\item 関手 $f^!$ により、$D_\QCoh^+$ は $\textit{FTS}_S$ 上の擬関手となる。
\item $f : X \to Y$ が $\textit{FTS}_S$ の固有射ならば、$f^!$ は
$Rf_* : D_\QCoh(\mathcal{O}_X) \to D_\QCoh(\mathcal{O}_Y)$ の右随伴を
$D_\QCoh^+(\mathcal{O}_Y)$ に制限したものであり、すべての
$K \in D_{\textit{Coh}}^-(\mathcal{O}_X)$ と
$M \in D_\QCoh^+(\mathcal{O}_Y)$ に対して標準同型
$$
Rf_*R\SheafHom_{\mathcal{O}_X}(K, f^!M)
\to
R\SheafHom_{\mathcal{O}_Y}(Rf_*K, M)
$$
がある。
\item $\textit{FTS}_S$ の対象 $X$ が双対化複体 $\omega_X^\bullet$ をもつならば、関手
$D_X = R\SheafHom_{\mathcal{O}_X}(-, \omega_X^\bullet)$ は
$D_{\textit{Coh}}(\mathcal{O}_X)$ の対合を定め、
$D_{\textit{Coh}}^+(\mathcal{O}_X)$ と $D_{\textit{Coh}}^-(\mathcal{O}_X)$ を交換し、
$D_{\textit{Coh}}^b(\mathcal{O}_X)$ を保つ。
\item $f : X \to Y$ が $\textit{FTS}_S$ の射であり、$\omega_Y^\bullet$ が
$Y$ 上の双対化複体ならば、
\begin{enumerate}
\item $\omega_X^\bullet = f^!\omega_Y^\bullet$ は $X$ の双対化複体である。
\item $M \in D_{\textit{Coh}}^+(\mathcal{O}_Y)$ に対して標準的に
$f^!M = D_X(Lf^*D_Y(M))$ である。
\item さらに $f$ が固有ならば、$D^-_{\textit{Coh}}(\mathcal{O}_X)$ の $K$ に対して
$$
Rf_*R\SheafHom_{\mathcal{O}_X}(K, \omega_X^\bullet) =
R\SheafHom_{\mathcal{O}_Y}(Rf_*K, \omega_Y^\bullet)
$$
である。
\end{enumerate}
\item $f : X \to Y$ が $\textit{FTS}_S$ の閉埋め込みならば、
$f^!(-) = R\SheafHom(\mathcal{O}_X, -)$ である。
\item $f : Y \to X$ が $\textit{FTS}_S$ の有限射ならば、
$f_*f^!(-) = R\SheafHom_{\mathcal{O}_X}(f_*\mathcal{O}_Y, -)$ である。
\item $f : X \to Y$ が $\textit{FTS}_S$ の対象への有効 Cartier 因子の包含ならば、
$f^!(-) = Lf^*(-) \otimes_{\mathcal{O}_X} f^*\mathcal{O}_Y(X)[-1]$ である。
\item $f : X \to Y$ が $\textit{FTS}_S$ の対象への余次元 $c$ の Koszul 正則埋め込みならば、
$f^!(-) \cong Lf^*(-) \otimes_{\mathcal{O}_X} \wedge^c\mathcal{N}[-c]$ である。
\item $f : X \to Y$ が $\textit{FTS}_S$ における相対次元 $d$ の滑らかな固有射ならば、
$f^!(-) \cong Lf^*(-)  \otimes_{\mathcal{O}_X} \Omega^d_{X/Y}[d]$ である。
\end{enumerate}
これは補題 \ref{duality-lemma-dualizing-schemes}、\ref{duality-lemma-iso-on-RSheafHom}、\ref{duality-lemma-twisted-inverse-image-closed}、\ref{duality-lemma-finite-twisted}、\ref{duality-lemma-sheaf-with-exact-support-effective-Cartier}、\ref{duality-lemma-regular-immersion}、\ref{duality-lemma-smooth-proper}、\ref{duality-lemma-upper-shriek-composition}、\ref{duality-lemma-pseudo-functor}、\ref{duality-lemma-shriek-closed-immersion}、\ref{duality-lemma-shriek-dualizing}、\ref{duality-lemma-shriek-via-duality} および \ref{duality-lemma-perfect-comparison-shriek} and
例 \ref{duality-example-iso-on-RSheafHom-noetherian}
から従う。我々の関手は非常に抽象的な手続きによって得られており、最終的には存在定理
（『導来圏』の命題 \ref{derived-proposition-brown}）の適用に依存している。
従って一般には、関手 $f^!$ の明示的記述をもたない。これはときに問題となるが、実際には
双対化複体の存在と双対性同型を知るだけで $f^!$ を決定できる場合が多い。






\section{双対化複体の貼り合わせ}
\label{duality-section-glue}

\noindent
ここでは双対化複体の貼り合わせを用いて、状況 \ref{duality-situation-dualizing} のような組
$(S, \omega_S^\bullet)$ が与えられたときに、$S$ 上のすべての有限型スキームに対して機能する理論を得る。
これは \cite[Remark on page 310]{RD} と類似している。

\begin{situation}
\label{duality-situation-dualizing}
ここで $S$ は Noether スキームであり、$\omega_S^\bullet$ は双対化複体である。
\end{situation}

\noindent
状況 \ref{duality-situation-dualizing} において、$X$ を $S$ 上有限型なスキームとする。
$\mathcal{U} : X = \bigcup_{i = 1, \ldots, n} U_i$ を、$\textit{FTS}_S$ の対象による
$X$ の有限開被覆とする。状況 \ref{duality-situation-shriek} を参照する。
これは単に射 $U_i \to S$ が分離であることを意味する（有限型であることはすでに仮定されている）。
『スキーム』の補題 \ref{schemes-lemma-compose-after-separated} により、$S$ 上有限型な任意の
アフィンスキームは $\textit{FTS}_S$ の対象なので、そのような開被覆は確かに存在する。
このとき各 $i, j, k \in \{1, \ldots, n\}$ に対して、射
$p_i : U_i \to S$, $p_{ij} : U_i \cap U_j \to S$,
$p_{ijk} : U_i \cap U_j \cap U_k \to S$ は分離であり、これらの各スキームは
$\textit{FTS}_S$ の対象である。このような開被覆から次を得る。
\begin{enumerate}
\item $\omega_i^\bullet = p_i^!\omega_S^\bullet$ は $U_i$ 上の双対化複体である。
第 \ref{duality-section-duality} 節を参照する。
\item 各 $i, j$ に対して標準同型
$\varphi_{ij} :
\omega_i^\bullet|_{U_i \cap U_j} \to \omega_j^\bullet|_{U_i \cap U_j}$ がある。
\item
\label{duality-item-cocycle-glueing}
各 $i, j, k$ に対して
$$
\varphi_{ik}|_{U_i \cap U_j \cap U_k} =
\varphi_{jk}|_{U_i \cap U_j \cap U_k} \circ
\varphi_{ij}|_{U_i \cap U_j \cap U_k}
$$
が $D(\mathcal{O}_{U_i \cap U_j \cap U_k})$ において成り立つ。
\end{enumerate}
ここで (2) では、$(U_i \cap U_j \to U_i)^!$ が制限によって与えられること
（補題 \ref{duality-lemma-shriek-open-immersion}）と、補題 \ref{duality-lemma-upper-shriek-composition} による標準同型
$$
(U_i \cap U_j \to U_i)^! \circ p_i^! = p_{ij}^! =
(U_i \cap U_j \to U_j)^! \circ p_j^!
$$
を用いる。(3) には、補題 \ref{duality-lemma-pseudo-functor} により上付きシュリーク関手が
擬関手をなすことを用いる。

\medskip\noindent
上で述べた状況において、{\it $\omega_S^\bullet$ および $\mathcal{U}$ に関して正規化された双対化複体}
とは、$K \in D(\mathcal{O}_X)$ と同型
$\alpha_i : K|_{U_i} \to \omega_i^\bullet$ からなる組 $(K, \alpha_i)$ であって、
$\varphi_{ij}$ が
$\alpha_j|_{U_i \cap U_j} \circ \alpha_i^{-1}|_{U_i \cap U_j}$ によって与えられるものをいう。
スキーム上の双対化複体であることは局所的性質なので、$\omega_S^\bullet$ および
$\mathcal{U}$ に関して正規化された双対化複体は実際に双対化複体である。

\begin{lemma}
\label{duality-lemma-good-dualizing-unique}
状況 \ref{duality-situation-dualizing} において、$X$ を $S$ 上有限型なスキームとし、
$\mathcal{U}$ を $S$ 上分離なスキームによる $X$ の有限開被覆とする。
$\omega_S^\bullet$ および $\mathcal{U}$ に関して正規化された双対化複体が存在するならば、
それは一意な同型を除いて一意である。
\end{lemma}

\begin{proof}
$(K, \alpha_i)$ と $(K', \alpha_i')$ が二つあるとし、
$L = R\SheafHom_{\mathcal{O}_X}(K, K')$ を考える。
補題 \ref{duality-lemma-dualizing-unique-schemes} とその証明により、これは
$D(\mathcal{O}_X)$ の可逆対象である。$\alpha_i$ と $\alpha'_i$ を用いると同型
$$
\alpha_i^t \otimes \alpha'_i :
L|_{U_i} \longrightarrow
R\SheafHom_{\mathcal{O}_X}(\omega_i^\bullet, \omega_i^\bullet) =
\mathcal{O}_{U_i}[0]
$$
を得る。これはすでに $D(\mathcal{O}_X)$ において $L = H^0(L)[0]$ であることを意味する。
さらに $H^0(L)$ は、$X$ の開集合 $U_i$ 上に指定された自明化をもつ可逆層である。
最後に、
$\alpha_j|_{U_i \cap U_j} \circ \alpha_i^{-1}|_{U_i \cap U_j}$ と
$\alpha'_j|_{U_i \cap U_j} \circ (\alpha'_i)^{-1}|_{U_i \cap U_j}$
の双方が $\varphi_{ij}$ を与えるという条件から、遷移写像は $1$ であり、同型
$H^0(L) = \mathcal{O}_X$ を得る。
\end{proof}

\begin{lemma}
\label{duality-lemma-good-dualizing-independence-covering}
状況 \ref{duality-situation-dualizing} において、$X$ を $S$ 上有限型なスキームとし、
$\mathcal{U}$, $\mathcal{V}$ を $S$ 上分離なスキームによる $X$ の二つの有限開被覆とする。
$\omega_S^\bullet$ および $\mathcal{U}$ に関して正規化された双対化複体が存在するならば、
$\omega_S^\bullet$ および $\mathcal{V}$ に関して正規化された双対化複体も存在し、
これらの複体は標準的に同型である。
\end{lemma}

\begin{proof}
$\mathcal{U}$ が開集合 $U_1, \ldots, U_n$ で、$\mathcal{V}$ が開集合
$U_1, \ldots, U_{n + m}$ で与えられる場合を示せば十分である。実際、$m = 1$ と仮定してよいし、そうする。
$\omega_S^\bullet$ と $\mathcal{V}$ に関して正規化された双対化複体 $(K, \alpha_i)$ から、
$\omega_S^\bullet$ と $\mathcal{U}$ に関して正規化された双対化複体へ移るには、
$i = n + 1$ に対する $\alpha_i$ を忘れればよい。逆に、$(K, \alpha_i)$ を
$\omega_S^\bullet$ と $\mathcal{U}$ に関して正規化された双対化複体とする。
証明を終えるには、所要の条件を満たす写像
$\alpha_{n + 1} : K|_{U_{n + 1}} \to \omega_{n + 1}^\bullet$ を構成しなければならない。
このため、$U_{n + 1} = \bigcup U_i \cap U_{n + 1}$ が開被覆であることに注意する。
$(K|_{U_{n + 1}}, \alpha_i|_{U_i \cap U_{n + 1}})$ が $\omega_S^\bullet$ と被覆
$U_{n + 1} = \bigcup U_i \cap U_{n + 1}$ に関して正規化された双対化複体であることは明らかである。
他方、条件 (\ref{duality-item-cocycle-glueing}) により、組
$(\omega_{n + 1}^\bullet|_{U_{n + 1}}, \varphi_{n + 1i})$ も $\omega_S^\bullet$ と被覆
$U_{n + 1} = \bigcup U_i \cap U_{n + 1}$ に関して正規化された双対化複体である。
補題 \ref{duality-lemma-good-dualizing-unique} により、与えられた局所同型と両立する一意な同型
$$
\alpha_{n + 1} : K|_{U_{n + 1}} \longrightarrow \omega_{n + 1}^\bullet
$$
を得る。これが所要の性質を満たすことを示すのは快い演習である。
\end{proof}

\begin{lemma}
\label{duality-lemma-existence-good-dualizing}
状況 \ref{duality-situation-dualizing} において、$X$ を $S$ 上有限型なスキームとし、
$\mathcal{U}$ を $S$ 上分離なスキームによる $X$ の有限開被覆とする。
このとき $\omega_S^\bullet$ および $\mathcal{U}$ に関して正規化された双対化複体が存在する。
\end{lemma}

\begin{proof}
$\mathcal{U} : X = \bigcup_{i = 1, \ldots, n} U_i$ と書く。
$n$ に関する帰納法で補題を示す。$n = 1$ の場合は直ちに従う。
$n > 1$ と仮定し、$X' = U_1 \cup \ldots \cup U_{n - 1}$ とおく。
$(K', \{\alpha'_i\}_{i = 1, \ldots, n - 1})$ を、$\omega_S^\bullet$ と
$\mathcal{U}' : X' = \bigcup_{i = 1, \ldots, n - 1} U_i$ に関して正規化された双対化複体とする。
$(K'|_{X' \cap U_n}, \alpha'_i|_{U_i \cap U_n})$ が $\omega_S^\bullet$ と被覆
$X' \cap U_n = \bigcup_{i = 1, \ldots, n - 1} U_i \cap U_n$ に関して正規化された双対化複体であることは明らかである。
他方、条件 (\ref{duality-item-cocycle-glueing}) により、組
$(\omega_n^\bullet|_{X' \cap U_n}, \varphi_{ni})$ も $\omega_S^\bullet$ と被覆
$X' \cap U_n = \bigcup_{i = 1, \ldots, n - 1} U_i \cap U_n$ に関して正規化された双対化複体である。
補題 \ref{duality-lemma-good-dualizing-unique} により、与えられた局所同型と両立する一意な同型
$$
\epsilon : K'|_{X' \cap U_n} \longrightarrow \omega_i^\bullet|_{X' \cap U_n}
$$
を得る。『層のコホモロジー』の補題 \ref{cohomology-lemma-glue} により、
$K \in D(\mathcal{O}_X)$ と同型 $\beta : K|_{X'} \to K'$, 
$\gamma : K|_{U_n} \to \omega_n^\bullet$ であって
$\epsilon = \gamma|_{X'\cap U_n} \circ \beta|_{X' \cap U_n}^{-1}$ を満たすものを得る。
そこで
$$
\alpha_i = \alpha'_i \circ \beta|_{U_i}, i = 1, \ldots, n - 1,
\text{ and }
\alpha_n = \gamma
$$
と定める。なお、$\varphi_{ij}$ が
$\alpha_j|_{U_i \cap U_j} \circ \alpha_i^{-1}|_{U_i \cap U_j}$ によって与えられることを確認する必要がある。
$i, j \leq n - 1$ ならば、これは $\alpha_i'$ に対する対応する条件から従う。
$i = j = n$ の場合も明らかである。$i < j = n$ ならば
$$
\alpha_n|_{U_i \cap U_n} \circ \alpha_i^{-1}|_{U_i \cap U_n} =
\gamma|_{U_i \cap U_n} \circ \beta^{-1}|_{U_i \cap U_n}
\circ (\alpha'_i)^{-1}|_{U_i \cap U_n} =
\epsilon|_{U_i \cap U_n} \circ (\alpha'_i)^{-1}|_{U_i \cap U_n}
$$
を得る。これは、$\epsilon$ が写像 $\alpha_i'$ および $\alpha_{ni}$ と両立する一意な写像であることから、
ちょうど $\alpha_{in}$ に等しい。
\end{proof}

\noindent
$(S, \omega_S^\bullet)$ を状況 \ref{duality-situation-dualizing} のとおりとする。
上の補題群から得られる結論は次のとおりである。$S$ 上有限型な任意のスキーム $X$ に対し、
一意な同型を除いて与えられる組 $(K, \alpha_U)$ がある。これは対象
$K \in D(\mathcal{O}_X)$ と、$S$ 上分離な各開部分スキーム $U \subset X$ に対する同型
$\alpha_U : K|_U \to \omega_U^\bullet$ からなる。ここで
$\omega_U^\bullet = (U \to S)^!\omega_S^\bullet$ は $U$ 上の双対化複体である。
第 \ref{duality-section-duality} 節を参照する。さらに、
$\mathcal{U} : X = \bigcup U_i$ が $S$ 上分離な開集合による有限開被覆ならば、
$(K, \alpha_{U_i})$ は $\omega_S^\bullet$ と $\mathcal{U}$ に関して正規化された双対化複体である。
実際、一意な同型を除く一意性は補題 \ref{duality-lemma-good-dualizing-unique}、
一つの開被覆に対する存在は補題 \ref{duality-lemma-existence-good-dualizing}、そして
$K$ がすべての開被覆に対して機能することは補題 \ref{duality-lemma-good-dualizing-independence-covering} による。

\begin{definition}
\label{duality-definition-good-dualizing}
$S$ を Noether スキームとし、$\omega_S^\bullet$ を $S$ 上の双対化複体とする。
$X$ を $S$ 上有限型なスキームとする。上で構成した複体 $K$ を
{\it $\omega_S^\bullet$ に関して正規化された双対化複体}といい、$\omega_X^\bullet$ と書く。
\end{definition}

\noindent
用語が示すように、$\omega_S^\bullet$ に関して正規化された双対化複体は、単なる
$X$ の導来圏の対象ではなく、上で述べた局所同型を備えている。
$X$ が $S$ 上分離であり $p : X \to S$ が構造射ならば、これを
$\omega_X^\bullet = p^!\omega_S^\bullet$ とおくことと矛盾しない。より一般に次の整合性確認がある。

\begin{lemma}
\label{duality-lemma-good-over-both}
$(S, \omega_S^\bullet)$ を状況 \ref{duality-situation-dualizing} のとおりとする。
$f : X \to Y$ を $S$ 上有限型なスキームの射とする。
$\omega_X^\bullet$ と $\omega_Y^\bullet$ を $\omega_S^\bullet$ に関して正規化された双対化複体とする。
このとき $\omega_X^\bullet$ は $\omega_Y^\bullet$ に関して正規化された双対化複体である。
\end{lemma}

\begin{proof}
これは記号を整理するだけである。有限アフィン開被覆
$\mathcal{V} : Y = \bigcup V_j$ を選ぶ。各 $j$ に対して有限アフィン開被覆
$f^{-1}(V_j) = U_{ji}$ を選び、$\mathcal{U} : X = \bigcup U_{ji}$ とおく。
スキーム $V_j$ と $U_{ji}$ は $S$ 上分離なので、
$q_j : V_j \to S$, $p_{ji} : U_{ji} \to S$,
$f_{ji} : U_{ji} \to V_j$, $f_{ji}' : U_{ji} \to Y$ に対する上付きシュリーク関手がある。
$(L, \beta_j)$ を $\omega_S^\bullet$ と $\mathcal{V}$ に関して正規化された双対化複体とし、
$(K, \gamma_{ji})$ を $\omega_S^\bullet$ と $\mathcal{U}$ に関して正規化された双対化複体とする。
（言い換えれば、$L = \omega_Y^\bullet$, $K = \omega_X^\bullet$ である。）次のように定義できる。
$$
\alpha_{ji} :
K|_{U_{ji}} \xrightarrow{\gamma_{ji}}
p_{ji}^!\omega_S^\bullet = f_{ji}^!q_j^!\omega_S^\bullet
\xrightarrow{f_{ji}^!\beta_j^{-1}} f_{ji}^!(L|_{V_j}) =
(f_{ji}')^!(L)
$$
証明を終えるには、
$\alpha_{ji}|_{U_{ji} \cap U_{j'i'}}
\circ \alpha_{j'i'}^{-1}|_{U_{ji} \cap U_{j'i'}}$
が標準同型
$(f_{ji}')^!(L)|_{U_{ji} \cap U_{j'i'}} \to
(f_{j'i'}')^!(L)|_{U_{ji} \cap U_{j'i'}}$
であることを示さなければならない。これは形式的であり、細部は省略する。
\end{proof}

\begin{lemma}
\label{duality-lemma-open-immersion-good-dualizing-complex}
$(S, \omega_S^\bullet)$ を状況 \ref{duality-situation-dualizing} のとおりとする。
$j : X \to Y$ を $S$ 上有限型なスキームの開埋め込みとする。
$\omega_X^\bullet$ と $\omega_Y^\bullet$ を $\omega_S^\bullet$ に関して正規化された双対化複体とする。
このとき標準同型 $\omega_X^\bullet = \omega_Y^\bullet|_X$ がある。
\end{lemma}

\begin{proof}
定義 \ref{duality-definition-good-dualizing} の直前で与えた正規化双対化複体の構成から直ちに従う。
\end{proof}

\begin{lemma}
\label{duality-lemma-proper-map-good-dualizing-complex}
$(S, \omega_S^\bullet)$ を状況 \ref{duality-situation-dualizing} のとおりとする。
$f : X \to Y$ を $S$ 上有限型なスキームの固有射とする。
$\omega_X^\bullet$ と $\omega_Y^\bullet$ を $\omega_S^\bullet$ に関して正規化された双対化複体とする。
$a$ を $f$ に対する補題 \ref{duality-lemma-twisted-inverse-image} の右随伴とする。
このとき標準同型 $a(\omega_Y^\bullet) = \omega_X^\bullet$ がある。
\end{lemma}

\begin{proof}
$p : X \to S$ と $q : Y \to S$ を構造射とする。
$X$ と $Y$ が $S$ 上分離ならば、$\omega_X^\bullet = p^!\omega_S^\bullet$,
$\omega_Y^\bullet = q^!\omega_S^\bullet$, $f^! = a$ であり、
$f^! \circ q^! = p^!$ であること（補題 \ref{duality-lemma-upper-shriek-composition}）から従う。
一般の場合、まず補題 \ref{duality-lemma-good-over-both} により $Y = S$ の場合に帰着する。
この場合 $X$ と $Y$ は $S$ 上分離であり、結果は上で見たとおりである。
\end{proof}

\noindent
$(S, \omega_S^\bullet)$ を状況 \ref{duality-situation-dualizing} のとおりとする。
$S$ 上有限型なスキーム $X$ に対し、$\omega_X^\bullet$ を
$\omega_S^\bullet$ に関して正規化された $X$ の双対化複体と書く。
補題 \ref{duality-lemma-dualizing-schemes} と同様に
$D_X(-) = R\SheafHom_{\mathcal{O}_X}(-, \omega_X^\bullet)$ と定める。
$f : X \to Y$ を $S$ 上有限型なスキームの射とし、
$$
f_{new}^! = D_X \circ Lf^* \circ D_Y :
D_{\textit{Coh}}^+(\mathcal{O}_Y)
\to
D_{\textit{Coh}}^+(\mathcal{O}_X)
$$
と定める。$f : X \to Y$ と $g : Y \to Z$ が $S$ 上有限型なスキーム間の合成可能な射ならば、
\begin{align*}
(g \circ f)^!_{new} & = D_X \circ L(g \circ f)^* \circ D_Z \\
& = D_X \circ Lf^* \circ Lg^* \circ D_Z \\
& \to D_X \circ Lf^* \circ D_Y \circ D_Y \circ Lg^* \circ D_Z \\
& = f^!_{new} \circ g^!_{new}
\end{align*}
と定める。ここで矢印は補題 \ref{duality-lemma-dualizing-schemes} で定義されたものである。
結果を次の補題にまとめる。

\begin{lemma}
\label{duality-lemma-duality-bootstrap}
$(S, \omega_S^\bullet)$ を状況 \ref{duality-situation-dualizing} のとおりとする。
$S$ 上有限型なすべてのスキーム（の射）に対し $f^!_{new}$ と $\omega_X^\bullet$ を上のように定めると、
\begin{enumerate}
\item 関手 $f^!_{new}$ と矢印
$(g \circ f)^!_{new} \to f^!_{new} \circ g^!_{new}$ により、
$D_{\textit{Coh}}^+$ は $S$ 上有限型なスキームの圏から圏の $2$-圏への擬関手となる。
\item $\omega_X^\bullet = (X \to S)^!_{new} \omega_S^\bullet$ である。
\item 関手 $D_X$ は $D_{\textit{Coh}}(\mathcal{O}_X)$ の対合を定め、
$D_{\textit{Coh}}^+(\mathcal{O}_X)$ と $D_{\textit{Coh}}^-(\mathcal{O}_X)$ を交換し、
$D_{\textit{Coh}}^b(\mathcal{O}_X)$ を保つ。
\item $f : X \to Y$ が $S$ 上有限型なスキームの射ならば、
$\omega_X^\bullet = f^!_{new}\omega_Y^\bullet$ である。
\item $M \in D_{\textit{Coh}}^+(\mathcal{O}_Y)$ に対して
$f^!_{new}M = D_X(Lf^*D_Y(M))$ である。
\item さらに $f$ が固有ならば、$f^!_{new}$ は
$Rf_* : D_\QCoh(\mathcal{O}_X) \to D_\QCoh(\mathcal{O}_Y)$ の右随伴を
$D_{\textit{Coh}}^+(\mathcal{O}_Y)$ に制限した関手と同型であり、
$K \in D^-_{\textit{Coh}}(\mathcal{O}_X)$ と
$M \in D_{\textit{Coh}}^+(\mathcal{O}_Y)$ に対して標準同型
$$
Rf_*R\SheafHom_{\mathcal{O}_X}(K, f^!_{new}M)
\to
R\SheafHom_{\mathcal{O}_Y}(Rf_*K, M)
$$
があり、$K \in D^-_{\textit{Coh}}(\mathcal{O}_X)$ に対して
$$
Rf_*R\SheafHom_{\mathcal{O}_X}(K, \omega_X^\bullet) =
R\SheafHom_{\mathcal{O}_Y}(Rf_*K, \omega_Y^\bullet)
$$
であり、
\end{enumerate}
$X$ が $S$ 上分離ならば、$\omega_X^\bullet$ は標準的に
$(X \to S)^!\omega_S^\bullet$ と同型である。また $f$ が $S$ 上分離なスキーム間の射ならば、
標準同型\footnote{これらが同型
$(g \circ f)^! \to f^! \circ g^!$ および
$(g \circ f)^!_{new} \to f^!_{new} \circ g^!_{new}$ と両立することは確認していない。
後で必要になればここで確認する。}
$f_{new}^!K = f^!K$ が $D_{\textit{Coh}}^+$ の $K$ に対して存在する。
\end{lemma}

\begin{proof}
$f : X \to Y$, $g : Y \to Z$, $h : Z \to T$ を $S$ 上有限型なスキームの射とする。
示すべきことは
$$
\xymatrix{
(h \circ g \circ f)^!_{new} \ar[r] \ar[d] &
f^!_{new} \circ (h \circ g)^!_{new} \ar[d] \\
(g \circ f)^!_{new} \circ h^!_{new} \ar[r] &
f^!_{new} \circ g^!_{new} \circ h^!_{new}
}
$$
が可換であることである。$\eta_Y : \text{id} \to D_Y^2$ と
$\eta_Z : \text{id} \to D_Z^2$ を補題 \ref{duality-lemma-dualizing-schemes} の標準同型とする。
『圏論』の補題 \ref{categories-lemma-properties-2-cat-cats} を用いた計算（省略）により、二つの矢印
$(h \circ g \circ f)^!_{new} \to f^!_{new} \circ g^!_{new} \circ h^!_{new}$
はいずれも
$$
1 \star \eta_Y \star 1 \star \eta_Z \star 1 :
D_X \circ Lf^* \circ Lg^* \circ Lh^* \circ D_T
\longrightarrow
D_X \circ Lf^* \circ D_Y^2 \circ Lg^* \circ D_Z^2 \circ Lh^* \circ D_T
$$
で与えられる。これで (1) が示された。(2) は $(X \to S)^!_{new}$ の定義と
$D_S(\omega_S^\bullet) = \mathcal{O}_S$ から直ちに従う。
(3) は補題 \ref{duality-lemma-dualizing-schemes} である。(4) は (2) と同じ議論から従う。
(5) は $f^!_{new}$ の定義である。

\medskip\noindent
(6) を証明する。$a$ を、$S$ 上有限型なスキームの固有射
$f : X \to Y$ に対する補題 \ref{duality-lemma-twisted-inverse-image} の右随伴とする。
問題は、$X$ または $Y$ が $S$ 上分離であることが分からず（一般には実際そうではない）、
$S$ 上の $f$ に補題 \ref{duality-lemma-shriek-via-duality} を直ちに適用できないことである。
これを回避するため、補題 \ref{duality-lemma-proper-map-good-dualizing-complex} の標準同一視
$\omega_X^\bullet = a(\omega_Y^\bullet)$ を用いる。従って、基底スキーム $Y$ 上の
$f : X \to Y$ に補題 \ref{duality-lemma-shriek-via-duality} を適用することにより、
$f^!_{new}$ は $a$ の $D_{\textit{Coh}}^+(\mathcal{O}_Y)$ への制限である。
表示した等式は例 \ref{duality-example-iso-on-RSheafHom-noetherian} により成り立つ。

\medskip\noindent
最後の主張群は、正規化双対化複体の構成と、すでに用いた補題 \ref{duality-lemma-shriek-via-duality} から従う。
\end{proof}

\begin{remark}
\label{duality-remark-independent-omega-S}
$S$ を双対化複体をもつ Noether スキームとする。
$f : X \to Y$ を $S$ 上有限型なスキームの射とする。このとき関手
$$
f_{new}^! : D^+_{Coh}(\mathcal{O}_Y) \to D^+_{Coh}(\mathcal{O}_X)
$$
は、標準同型を除いて双対化複体 $\omega_S^\bullet$ の選択に依存しない。
証明を概説する。任意の第二の双対化複体は
$\omega_S^\bullet \otimes_{\mathcal{O}_S}^\mathbf{L} \mathcal{L}$ の形である。
ここで $\mathcal{L}$ は $D(\mathcal{O}_S)$ の可逆対象である。補題 \ref{duality-lemma-dualizing-unique-schemes} を参照する。
有限型分離射 $p : U \to S$ に対して、補題 \ref{duality-lemma-compare-with-pullback-perfect} により
$p^!(\omega_S^\bullet \otimes^\mathbf{L}_{\mathcal{O}_S} \mathcal{L}) =
p^!(\omega_S^\bullet) \otimes^\mathbf{L}_{\mathcal{O}_U} Lp^*\mathcal{L}$
である。従って $\omega_X^\bullet$ と $\omega_Y^\bullet$ が
$\omega_S^\bullet$ に関して正規化された双対化複体ならば、
$\omega_X^\bullet \otimes_{\mathcal{O}_X}^\mathbf{L} La^*\mathcal{L}$ と
$\omega_Y^\bullet \otimes_{\mathcal{O}_Y}^\mathbf{L} Lb^*\mathcal{L}$ は
$\omega_S^\bullet \otimes_{\mathcal{O}_S}^\mathbf{L} \mathcal{L}$ に関して正規化された双対化複体である
（ここで $a : X \to S$ と $b : Y \to S$ は構造射である）。そこで結果は、
$K \in D^+_{Coh}(\mathcal{O}_Y)$ に対する次の計算から従う。
\begin{align*}
& R\SheafHom_{\mathcal{O}_X}(Lf^*R\SheafHom_{\mathcal{O}_Y}(K,
\omega_Y^\bullet \otimes_{\mathcal{O}_Y}^\mathbf{L} Lb^*\mathcal{L}),
\omega_X^\bullet \otimes_{\mathcal{O}_X}^\mathbf{L} La^*\mathcal{L}) \\
& = R\SheafHom_{\mathcal{O}_X}(Lf^*R(\SheafHom_{\mathcal{O}_Y}(K,
\omega_Y^\bullet) \otimes_{\mathcal{O}_Y}^\mathbf{L} Lb^*\mathcal{L}),
\omega_X^\bullet \otimes_{\mathcal{O}_X}^\mathbf{L} La^*\mathcal{L}) \\
& = R\SheafHom_{\mathcal{O}_X}(Lf^*R\SheafHom_{\mathcal{O}_Y}(K,
\omega_Y^\bullet) \otimes_{\mathcal{O}_X}^\mathbf{L} La^*\mathcal{L},
\omega_X^\bullet \otimes_{\mathcal{O}_X}^\mathbf{L} La^*\mathcal{L}) \\
& = R\SheafHom_{\mathcal{O}_X}(Lf^*R\SheafHom_{\mathcal{O}_Y}(K,
\omega_Y^\bullet), \omega_X^\bullet)
\end{align*}
最後の等号は $La^*\mathcal{L}$ が $D(\mathcal{O}_X)$ の可逆対象だからである。
\end{remark}


\begin{example}
\label{duality-example-trace-proper}
$S$ を Noether スキームとし、$\omega_S^\bullet$ を双対化複体とする。
$f : X \to Y$ を $S$ 上有限型なスキームの固有射とする。
$\omega_X^\bullet$ と $\omega_Y^\bullet$ を $\omega_S^\bullet$ に関して正規化された双対化複体とする。
この状況では $a(\omega_Y^\bullet) = \omega_X^\bullet$ であり
（補題 \ref{duality-lemma-proper-map-good-dualizing-complex}）、従ってトレース写像
（第 \ref{duality-section-trace} 節）は標準的な矢印
$$
\text{Tr}_f : Rf_*\omega_X^\bullet \longrightarrow \omega_Y^\bullet
$$
であって、$D_\QCoh(\mathcal{O}_X)$ の $L$ に対して同型
（補題 \ref{duality-lemma-duality-bootstrap}）
$$
\Hom_X(L, \omega_X^\bullet) = \Hom_Y(Rf_*L, \omega_Y^\bullet)
$$
および
$$
Rf_*R\SheafHom_{\mathcal{O}_X}(L, \omega_X^\bullet) =
R\SheafHom_{\mathcal{O}_Y}(Rf_*L, \omega_Y^\bullet)
$$
を与える。
\end{example}

\begin{remark}
\label{duality-remark-dualizing-finite}
$S$ を Noether スキームとし、$\omega_S^\bullet$ を双対化複体とする。
$f : X \to Y$ を $S$ 上有限型なスキーム間の有限射とする。
$\omega_X^\bullet$ と $\omega_Y^\bullet$ を $\omega_S^\bullet$ に関して正規化された双対化複体とする。
このとき補題 \ref{duality-lemma-finite-twisted} および \ref{duality-lemma-proper-map-good-dualizing-complex} により、$D_\QCoh^+(f_*\mathcal{O}_X)$ において
$$
f_*\omega_X^\bullet = R\SheafHom(f_*\mathcal{O}_X, \omega_Y^\bullet)
$$
であり、例 \ref{duality-example-trace-proper} のトレース写像は
$$
\text{Tr}_f : Rf_*\omega_X^\bullet = f_*\omega_X^\bullet =
R\SheafHom(f_*\mathcal{O}_X, \omega_Y^\bullet) \longrightarrow
\omega_Y^\bullet
$$
という写像である。これはしばしば「$1$ での評価」と呼ばれる。
\end{remark}

\begin{remark}
\label{duality-remark-relative-dualizing-complex-shriek}
$f : X \to Y$ を状況 \ref{duality-situation-dualizing} のような組
$(S, \omega_S^\bullet)$ 上の有限型スキームの平坦固有射とする。
相対双対化複体（注意 \ref{duality-remark-relative-dualizing-complex}）は
$\omega_{X/Y}^\bullet = a(\mathcal{O}_Y)$ である。補題 \ref{duality-lemma-proper-map-good-dualizing-complex} により、$D(\mathcal{O}_X)$ における
$$
\omega_X^\bullet = a(\omega_Y^\bullet) =
Lf^*\omega_Y^\bullet \otimes_{\mathcal{O}_X}^\mathbf{L} \omega_{X/Y}^\bullet
$$
の第一の標準同型を得る。第二の標準同型は注意 \ref{duality-remark-relative-dualizing-complex} の議論から従う。
\end{remark}




\section{次元関数}
\label{duality-section-dimension-functions}

\noindent
あるスキーム上有限型なスキームへ移るときに、次元関数がどのように変化するかについて
もう少し情報が必要である。

\begin{lemma}
\label{duality-lemma-good-dualizing-normalized}
$S$ を Noether スキームとし、$\omega_S^\bullet$ を双対化複体とする。
$X$ を $S$ 上有限型なスキームとし、$\omega_X^\bullet$ を
$\omega_S^\bullet$ に関して正規化された双対化複体とする。
$x \in X$ が $S$ の閉点 $s$ の上にある閉点ならば、$\omega_{S, s}^\bullet$ が
$\mathcal{O}_{S, s}$ 上の正規化双対化複体である限り、$\omega_{X, x}^\bullet$ は
$\mathcal{O}_{X, x}$ 上の正規化双対化複体である。
\end{lemma}

\begin{proof}
$X$ を $x$ のアフィン近傍で置き換えてよいので、$f : X \to S$ は分離であると仮定する。
すると $\omega_X^\bullet = f^!\omega_S^\bullet$ である。示すべきことは
$R\Hom_{\mathcal{O}_{X, x}}(\kappa(x), \omega_{X, x}^\bullet)$ が次数 $0$ に位置することである。
$i_x : x \to X$ を包含射とする。$x$ は閉点なので、これは閉埋め込みである。
従って補題 \ref{duality-lemma-shriek-closed-immersion} により、
$R\Hom_{\mathcal{O}_{X, x}}(\kappa(x), \omega_{X, x}^\bullet)$ は
$i_x^!\omega_X^\bullet$ を表す。可換図式
$$
\xymatrix{
x \ar[r]_{i_x} \ar[d]_\pi & X \ar[d]^f \\
s \ar[r]^{i_s} & S
}
$$
を考える。『スキームの射』の補題 \ref{morphisms-lemma-closed-point-fibre-locally-finite-type} により、拡大
$\kappa(x)/\kappa(s)$ は有限であり、従って $\pi$ は有限射である。ゆえに
$$
i_x^!\omega_X^\bullet = i_x^! f^! \omega_S^\bullet =
\pi^! i_s^! \omega_S^\bullet
$$
である。従って $\omega_{S, s}^\bullet$ が $\mathcal{O}_{S, s}$ 上の正規化双対化複体ならば、
上と同じ議論により $i_s^!\omega_S^\bullet = \kappa(s)[0]$ である。また
$$
R\pi_*(\pi^!(\kappa(s)[0])) =
R\SheafHom_{\mathcal{O}_s}(R\pi_*(\kappa(x)[0]), \kappa(s)[0]) =
\widetilde{\Hom_{\kappa(s)}(\kappa(x), \kappa(s))}
$$
である。第一の等号には $L = \kappa(x)[0]$ として例 \ref{duality-example-iso-on-RSheafHom-noetherian} を適用した。第二の等号は $\pi_*$ が完全だからである。
従って $\pi^!(\kappa(s)[0])$ は次数 $0$ に台をもち、主張が従う。
\end{proof}

\begin{lemma}
\label{duality-lemma-good-dualizing-dimension-function}
$S$ を Noether スキームとし、$\omega_S^\bullet$ を双対化複体とする。
$f : X \to S$ を有限型とし、$\omega_X^\bullet$ を $\omega_S^\bullet$ に関して正規化された双対化複体とする。
すべての $x \in X$ に対して
$$
\delta_X(x) - \delta_S(f(x)) = \text{trdeg}_{\kappa(f(x))}(\kappa(x))
$$
である。ここで $\delta_S$, resp.\ $\delta_X$ は
$\omega_S^\bullet$, resp.\ $\omega_X^\bullet$ の次元関数である。補題 \ref{duality-lemma-dimension-function-scheme} を参照する。
\end{lemma}

\begin{proof}
$X$ を $x$ のアフィン近傍で置き換えてよい。従って $S$ 上のコンパクト化
$X \subset \overline{X}$ があると仮定する。さらに $X$ を $\overline{X}$ で置き換え、
$X$ は $S$ 上固有であると仮定してよい。また $X$ を $x$ を含む $X$ の連結成分で置き換えて、
$X$ は連結であると仮定してよい。
次に、$\delta_X$ と関数
$x \mapsto \delta_S(f(x)) + \text{trdeg}_{\kappa(f(x))}(\kappa(x))$
はいずれも $X$ 上の次元関数であることを想起する。『スキームの射』の補題 \ref{morphisms-lemma-dimension-function-propagates} を参照する
（また補題 \ref{duality-lemma-dimension-function-scheme} により $S$ は普遍的鎖状である）。
『位相』の補題 \ref{topology-lemma-dimension-function-unique} により、その差は局所定数であり、
$X$ は連結なので定数である。従って $X$ の任意の一点で等式を示せば十分である。
『スキームの性質』の補題 \ref{properties-lemma-locally-Noetherian-closed-point} により、
$X$ は閉点 $x$ をもつ。$X \to S$ は固有なので、$x$ の像 $s$ は $S$ で閉である。
従って補題 \ref{duality-lemma-good-dualizing-normalized} を適用して結論を得る。
\end{proof}

\begin{lemma}
\label{duality-lemma-shriek}
状況 \ref{duality-situation-shriek} において、$f : X \to Y$ を
$\textit{FTS}_S$ の射とする。$x \in X$ の像を $y \in Y$ とする。このとき
$$
H^i(f^!\mathcal{O}_Y)_x \not = 0
\Rightarrow - \dim_x(X_y) \leq i.
$$
\end{lemma}

\begin{proof}
主張は $X$ 上局所的なので、$X$ と $Y$ はアフィンスキームであると仮定してよい。
$X = \Spec(A)$, $Y = \Spec(R)$ と書く。$f^!\mathcal{O}_Y$ は
『双対化複体』の第 \ref{dualizing-section-relative-dualizing-complexes-Noetherian} 節の相対双対化複体
$\omega_{A/R}^\bullet$ に対応する。これは注意 \ref{duality-remark-local-calculation-shriek} による。従って補題は『双対化複体』の補題 \ref{dualizing-lemma-relative-dualizing-trivial-vanishing} から従う。
\end{proof}

\begin{lemma}
\label{duality-lemma-flat-shriek}
状況 \ref{duality-situation-shriek} において、$f : X \to Y$ を
$\textit{FTS}_S$ の射とする。$x \in X$ の像を $y \in Y$ とする。
$f$ が平坦ならば
$$
H^i(f^!\mathcal{O}_Y)_x \not = 0
\Rightarrow - \dim_x(X_y) \leq i \leq 0.
$$
実際、$f$ のすべてのファイバーの次元が $\leq d$ ならば、$f^!\mathcal{O}_Y$ は
$D(X, f^{-1}\mathcal{O}_Y)$ の対象として $[-d, 0]$ に Tor 振幅をもつ。
\end{lemma}

\begin{proof}
補題 \ref{duality-lemma-shriek} の証明とまったく同様に論じれば、これは『双対化複体』の補題 \ref{dualizing-lemma-relative-dualizing-flat-vanishing} から従う。
\end{proof}

\begin{lemma}
\label{duality-lemma-shriek-over-CM}
状況 \ref{duality-situation-shriek} において、$f : X \to Y$ を
$\textit{FTS}_S$ の射とする。$x \in X$ の像を $y \in Y$ とする。次を仮定する。
\begin{enumerate}
\item $\mathcal{O}_{Y, y}$ は Cohen-Macaulay である。
\item $x$ を含む $X$ の既約成分の任意の生成点 $\xi$ に対して
$\text{trdeg}_{\kappa(f(\xi))}(\kappa(\xi)) \leq r$ である。
\end{enumerate}
このとき
$$
H^i(f^!\mathcal{O}_Y)_x \not = 0
\Rightarrow - r \leq i
$$
であり、茎 $H^{-r}(f^!\mathcal{O}_Y)_x$ は $\mathcal{O}_{X, x}$-加群として $(S_2)$ である。
\end{lemma}

\begin{proof}
$X$ を $x$ の開近傍で置き換えることにより、$X$ のすべての既約成分が $x$ を通ると仮定してよい。
すると補題 \ref{duality-lemma-shriek} の証明とまったく同様に論じれば、これは『双対化複体』の補題 \ref{dualizing-lemma-relative-dualizing-CM-vanishing} から従う。
\end{proof}

\begin{lemma}
\label{duality-lemma-flat-quasi-finite-shriek}
状況 \ref{duality-situation-shriek} において、$f : X \to Y$ を
$\textit{FTS}_S$ の射とする。$f$ が平坦かつ準有限ならば
$$
f^!\mathcal{O}_Y = \omega_{X/Y}[0]
$$
である。ここで $\omega_{X/Y}$ は $Y$ 上平坦なある連接 $\mathcal{O}_X$-加群である。
\end{lemma}

\begin{proof}
補題 \ref{duality-lemma-flat-shriek} と、補題 \ref{duality-lemma-shriek-coherent} によって
$f^!\mathcal{O}_Y$ のコホモロジー層が連接であることから従う。
\end{proof}

\begin{lemma}
\label{duality-lemma-CM-shriek}
状況 \ref{duality-situation-shriek} において、$f : X \to Y$ を
$\textit{FTS}_S$ の射とする。$f$ が Cohen-Macaulay ならば
（『射の続論』の定義 \ref{more-morphisms-definition-CM}）、
$$
f^!\mathcal{O}_Y = \omega_{X/Y}[d]
$$
である。ここで $\omega_{X/Y}$ は $Y$ 上平坦なある連接 $\mathcal{O}_X$-加群であり、
$d$ は $Y$ 上の $X$ の相対次元を与える $X$ 上局所定数な関数である。
\end{lemma}

\begin{proof}
『スキームの射』の補題 \ref{morphisms-lemma-flat-finite-presentation-CM-fibres-relative-dimension} により、
相対次元 $d$ はよく定義され局所定数である。補題 \ref{duality-lemma-shriek-coherent} により、
$f^!\mathcal{O}_Y$ のコホモロジー層は連接である。
$f^!\mathcal{O}_Y$ の他のコホモロジー層が零であることを示せれば、補題 \ref{duality-lemma-flat-shriek} から $\omega_{X/Y}$ の平坦性を得る。

\medskip\noindent
問題は $X$ 上局所的なので、$X$ と $Y$ はアフィンであり、射は相対次元 $d$ をもつと仮定してよい。
$d = 0$ ならば、結果は補題 \ref{duality-lemma-flat-quasi-finite-shriek} から直ちに従う。
$d > 0$ ならば、分解
$$
X \xrightarrow{g} \mathbf{A}^d_Y \xrightarrow{p} Y
$$
で $g$ が準有限かつ平坦なものがあると仮定してよい。『射の続論』の補題 \ref{more-morphisms-lemma-flat-finite-presentation-characterize-CM} を参照する。
このとき $f^! = g^! \circ p^!$ である。補題 \ref{duality-lemma-shriek-affine-line} により
$p^!\mathcal{O}_Y \cong \mathcal{O}_{\mathbf{A}^d_Y}[-d]$ である。
$d = 0$ の場合から結論を得る。
\end{proof}

\begin{remark}
\label{duality-remark-the-same-is-true}
$S$ を双対化複体 $\omega_S^\bullet$ を備えた Noether スキームとする。この場合補題 \ref{duality-lemma-shriek}、\ref{duality-lemma-flat-shriek}、\ref{duality-lemma-flat-quasi-finite-shriek} および \ref{duality-lemma-CM-shriek} は、
$S$ 上有限型なスキームの任意の射 $f : X \to Y$ に対し、$f^!$ を $f_{new}^!$ に置き換えれば成り立つ。
実際、いずれの場合も証明は直ちにアフィンの場合へ帰着し、その場合補題 \ref{duality-lemma-duality-bootstrap} により $f^! = f_{new}^!$ である。
\end{remark}




\section{双対化加群}
\label{duality-section-dualizing-module}


\noindent
本節は『双対化複体』の第 \ref{dualizing-section-dualizing-module} 節の続きである。

\medskip\noindent
$X$ を Noether スキームとし、$\omega_X^\bullet$ を双対化複体とする。
$H^n(\omega_X^\bullet)$ が非零となる最小の整数を $n \in \mathbf{Z}$ とする。
言い換えれば、$-n$ は $\omega_X^\bullet$ に伴う次元関数の最大値である
（補題 \ref{duality-lemma-dimension-function-scheme}）。
$H^n(\omega_X^\bullet)$ は $X$ の {\it 双対化加群}または {\it 双対化層}と呼ばれることがあり、
その場合しばしば $\omega_X$ と書かれる。「$\omega_X$ を双対化加群とする」と述べるときは、
以上の意味で用いる。

\medskip\noindent
Noether スキーム $X$ 上で双対化加群 $\omega_X$ を用いる際には注意が必要である。
\begin{enumerate}
\item $X$ から $X$ の開集合 $U$ へ移ると整数 $n$ が変わり得るため、
$\omega_X|_U = \omega_U$ が成り立つとは限らない。
\item 双対化複体は一意ではなく、双対化加群は可逆加群とのテンソル積を除いてのみ一意である。
\end{enumerate}
第二の問題はしばしば無関係となる。実際、我々は固定された双対化複体
$\omega_S^\bullet$ を備えた基底スキーム $S$ 上有限型な $X$ を扱い、
$\omega_X^\bullet$ は $\omega_S^\bullet$ に関して正規化された双対化複体となる。
第一の問題は $X$ が等次元ならば、より精密には $\omega_X^\bullet$ に伴う次元関数
（補題 \ref{duality-lemma-dimension-function-scheme}）が $X$ のすべての生成点を同じ整数に写すならば生じない。

\begin{example}
\label{duality-example-proper-over-local}
$S = \Spec(A)$ とし、$(A, \mathfrak m, \kappa)$ は Noether 局所環であり、
$\omega_S^\bullet$ は正規化双対化複体 $\omega_A^\bullet$ に対応するとする。
$f : X \to S$ が $S$ 上固有であり、$\omega_X^\bullet = f^!\omega_S^\bullet$ ならば、連接層
$$
\omega_X = H^{-\dim(X)}(\omega_X^\bullet)
$$
は双対化加群であり、しばしば $X$ の双対化加群と呼ばれる
（$S$ と $\omega_S^\bullet$ は了解されているものとする）。これが良い性質をもつことを後で見る。
\end{example}

\begin{example}
\label{duality-example-equidimensional-over-field}
$X$ を体 $k$ 上有限型な等次元スキームとする。このとき通常、$\omega_X^\bullet$ を
$k[0]$ に関して正規化された双対化複体に取り、
$$
\omega_X = H^{-\dim(X)}(\omega_X^\bullet)
$$
を $X$ の双対化加群と呼ぶ。$X$ が $k$ 上分離ならば、補題 \ref{duality-lemma-duality-bootstrap} により、構造射 $f : X \to \Spec(k)$ に対して
$\omega_X^\bullet = f^!\mathcal{O}_{\Spec(k)}$ である。
$X$ が $k$ 上固有ならば、これは例 \ref{duality-example-proper-over-local} の特別な場合である。
\end{example}

\begin{lemma}
\label{duality-lemma-dualizing-module}
$X$ を連結 Noether スキームとし、$\omega_X$ を $X$ 上の双対化加群とする。
$\omega_X$ の台は、任意の次元関数に関して最大次元をもつ既約成分の合併であり、
$\omega_X$ は性質 $(S_2)$ をもつ連接 $\mathcal{O}_X$-加群である。
\end{lemma}

\begin{proof}
上で述べた規約により、$\omega_X$ が最も左の非零コホモロジー層となる双対化複体
$\omega_X^\bullet$ が存在する。$X$ は連結なので、任意の二つの次元関数の差は定数である
（『位相』の補題 \ref{topology-lemma-dimension-function-unique}）。
従って $\omega_X^\bullet$ に伴う次元関数を用いてよい
（補題 \ref{duality-lemma-dimension-function-scheme}）。以上を踏まえると、補題は
『双対化複体』の補題 \ref{dualizing-lemma-depth-dualizing-module} と定義群
（特に『スキームのコホモロジー』の定義 \ref{coherent-definition-depth}）から従う。
\end{proof}

\begin{lemma}
\label{duality-lemma-vanishing-good-dualizing}
$X/A$, $\omega_X^\bullet$, $\omega_X$ を例 \ref{duality-example-proper-over-local} のとおりとする。このとき
\begin{enumerate}
\item $H^i(\omega_X^\bullet) \not = 0 \Rightarrow
i \in \{-\dim(X), \ldots, 0\}$ である。
\item $H^i(\omega_X^\bullet)$ の台の次元は高々 $-i$ である。
\item $\text{Supp}(\omega_X)$ は次元 $\dim(X)$ の成分の合併である。
\item $\omega_X$ は性質 $(S_2)$ をもつ。
\end{enumerate}
\end{lemma}

\begin{proof}
$\delta_X$ と $\delta_S$ を、補題 \ref{duality-lemma-good-dualizing-dimension-function} のように
$\omega_X^\bullet$ と $\omega_S^\bullet$ に伴う次元関数とする。
$X$ は $A$ 上固有なので、$X$ の任意の閉部分スキームは、閉点 $s \in S$ に写る閉点 $x$ を含み、
$\delta_X(x) = \delta_S(s) = 0$ である。従って任意の点 $\xi \in X$ に対して
$\delta_X(\xi) = \dim(\overline{\{\xi\}})$ である。
ゆえに補題の各主張は $\Spec(\mathcal{O}_{X, x})$ 上で起こることを調べれば確認でき、
その場合、結果は『双対化複体』の補題 \ref{dualizing-lemma-sitting-in-degrees} および \ref{dualizing-lemma-depth-dualizing-module} から従う。細部は省略する。
最後の二つの主張は補題 \ref{duality-lemma-dualizing-module} からも導ける。
\end{proof}

\begin{lemma}
\label{duality-lemma-dualizing-module-proper-over-A}
$X/A$ と双対化加群 $\omega_X$ を例 \ref{duality-example-proper-over-local} のとおりとする。
$d = \dim(X_s)$ を閉ファイバーの次元とする。$\dim(X) = d + \dim(A)$ ならば、
双対化加群 $\omega_X$ は連接 $\mathcal{O}_X$-加群の圏上の関手
$$
\mathcal{F} \longmapsto \Hom_A(H^d(X, \mathcal{F}), \omega_A)
$$
を表現する。
\end{lemma}

\begin{proof}
次を得る。
\begin{align*}
\Hom_X(\mathcal{F}, \omega_X)
& =
\Ext^{-\dim(X)}_X(\mathcal{F}, \omega_X^\bullet) \\
& =
\Hom_X(\mathcal{F}[\dim(X)], \omega_X^\bullet) \\
& =
\Hom_X(\mathcal{F}[\dim(X)], f^!(\omega_A^\bullet)) \\
& =
\Hom_S(Rf_*\mathcal{F}[\dim(X)], \omega_A^\bullet) \\
& =
\Hom_A(H^d(X, \mathcal{F}), \omega_A)
\end{align*}
第一の等号は、補題 \ref{duality-lemma-vanishing-good-dualizing} および『導来圏』の補題 \ref{derived-lemma-negative-exts} により、$i < -\dim(X)$ に対して
$H^i(\omega_X^\bullet) = 0$ だからである。第二の等号は Ext 群の定義から従う。
第三の等号は $\omega_X^\bullet$ の選択による。第四の等号は、$f^!$ が $f$ に対する
補題 \ref{duality-lemma-twisted-inverse-image} の右随伴だからである。第 \ref{duality-section-duality} 節を参照する。最後の等号は、$i > d$ に対して
$R^if_*\mathcal{F}$ が零であり（『スキームのコホモロジー』の補題 \ref{coherent-lemma-higher-direct-images-zero-above-dimension-fibre}）、
$j < -\dim(A)$ に対して $H^j(\omega_A^\bullet)$ が零だからである。
\end{proof}








\section{Cohen-Macaulay スキーム}
\label{duality-section-CM}

\noindent
本節は『双対化複体』の第 \ref{dualizing-section-CM} 節の続きである。
Cohen-Macaulay スキームに対して双対性は特に簡単な形をとる。

\begin{lemma}
\label{duality-lemma-dualizing-module-CM-scheme}
$X$ を双対化複体 $\omega_X^\bullet$ をもつ局所 Noether スキームとする。
\begin{enumerate}
\item $X$ が Cohen-Macaulay であること $\Leftrightarrow$ $\omega_X^\bullet$ が局所的に一意な非零コホモロジー層をもつこと。
\item $\mathcal{O}_{X, x}$ が Cohen-Macaulay であること $\Leftrightarrow$
$\omega_{X, x}^\bullet$ が一意な非零コホモロジーをもつこと。
\item $U = \{x \in X \mid \mathcal{O}_{X, x}\text{ is Cohen-Macaulay}\}$
は開であり Cohen-Macaulay である。
\end{enumerate}
$X$ が連結かつ Cohen-Macaulay ならば、ある整数 $n$ と連接 Cohen-Macaulay
$\mathcal{O}_X$-加群 $\omega_X$ が存在して $\omega_X^\bullet = \omega_X[-n]$ となる。
\end{lemma}

\begin{proof}
定義および『双対化複体』の補題 \ref{dualizing-lemma-dualizing-localize} により、
各 $x \in X$ に対して複体 $\omega_{X, x}^\bullet$ は $\mathcal{O}_{X, x}$ 上の双対化複体である。
従って『双対化複体』の補題 \ref{dualizing-lemma-apply-CM} により (2) が成り立つ。

\medskip\noindent
(3) を示すため、$\mathcal{O}_{X, x}$ は Cohen-Macaulay であると仮定する。
$H^{n_{x}}(\omega_{X, x}^\bullet)$ が非零となる一意な整数を $n_x$ とする。
$x$ のアフィン近傍 $V \subset X$ に対して、
$\omega_X^\bullet|_V$ は $D^b_{\textit{Coh}}(\mathcal{O}_V)$ に属するので、非零な連接加群
$H^i(\omega_X^\bullet)|_V$ は有限個しかない。従って $V$ を縮小することにより、
$H^{n_x}$ のみが非零であると仮定してよい。『加群の層』の補題 \ref{modules-lemma-finite-type-stalk-zero} を参照する。
これにより、すべての $v \in V$ に対して $\mathcal{O}_{X, v}$ は Cohen-Macaulay である。
従って $U$ は開であり、Cohen-Macaulay スキームでもある。

\medskip\noindent
(1) を証明する。含意 $\Leftarrow$ は (2) から従う。
含意 $\Rightarrow$ は前段落の議論から従う。そこでは、$\mathcal{O}_{X, x}$ が Cohen-Macaulay ならば、
$x$ の近傍で複体 $\omega_X^\bullet$ はただ一つの非零コホモロジー層をもつことを示した。

\medskip\noindent
$X$ は連結かつ Cohen-Macaulay であると仮定する。上の議論により、写像
$x \mapsto n_x$ は局所定数である。$X$ は連結なので定数であり、その値を $n$ とする。
$\omega_X = H^n(\omega_X^\bullet)$ とおけば補題が成り立つ。実際、『双対化複体』の補題 \ref{dualizing-lemma-apply-CM} により $\omega_X$ は Cohen-Macaulay である
（『スキームのコホモロジー』の定義 \ref{coherent-definition-Cohen-Macaulay} も参照する）。
\end{proof}

\begin{lemma}
\label{duality-lemma-has-dualizing-module-CM-scheme}
$X$ を局所 Noether スキームとする。$\omega_X[0]$ が $X$ 上の双対化複体となる連接層
$\omega_X$ が存在するならば、$X$ は Cohen-Macaulay スキームである。
\end{lemma}

\begin{proof}
『双対化複体』の補題 \ref{dualizing-lemma-has-dualizing-module-CM} と定義から直ちに従う。
\end{proof}

\begin{lemma}
\label{duality-lemma-affine-flat-Noetherian-CM}
状況 \ref{duality-situation-shriek} において、$f : X \to Y$ を
$\textit{FTS}_S$ の射とし、$x \in X$ とする。$f$ が平坦ならば、次は同値である。
\begin{enumerate}
\item $f$ は $x$ において Cohen-Macaulay である。
\item $f^!\mathcal{O}_Y$ は $x$ の近傍で一意な非零コホモロジー層をもつ。
\end{enumerate}
\end{lemma}

\begin{proof}
一方の向きは補題 \ref{duality-lemma-CM-shriek} から従う。
逆を示すため、$f^!\mathcal{O}_Y$ は一意な非零コホモロジー層をもつと仮定する。
$y = f(x)$ とおく。$\xi_1, \ldots, \xi_n \in X_y$ を、$x$ に特殊化するファイバー
$X_y$ の生成点とする。対応する $X_y$ の既約成分の次元を $d_1, \ldots, d_n$ とする。
『射の続論』の補題 \ref{more-morphisms-lemma-flat-finite-presentation-CM-open} により、射 $f : X \to Y$ は
$\eta_i$ において Cohen-Macaulay である。従って補題 \ref{duality-lemma-CM-shriek} により
$d_1 = \ldots = d_n$ である。共通値を $d$ とすれば $d = \dim_x(X_y)$ である。
$X$ を縮小することにより、すべてのファイバーの次元は高々 $d$ と仮定してよい
（『スキームの射』の補題 \ref{morphisms-lemma-openness-bounded-dimension-fibres}）。
すると唯一の非零コホモロジー層 $\omega = H^{-d}(f^!\mathcal{O}_Y)$ は
補題 \ref{duality-lemma-flat-shriek} により $Y$ 上平坦である。
従って $h : X_y \to X$ を標準射と書けば、『スキームの導来圏』の補題 \ref{perfect-lemma-tor-independence-and-tor-amplitude} により
$Lh^*(f^!\mathcal{O}_Y) = Lh^*(\omega[d]) = (h^*\omega)[d]$ である。
補題 \ref{duality-lemma-relative-dualizing-fibres} により $h^*\omega[d]$ は $X_y$ の双対化複体である。
従って補題 \ref{duality-lemma-dualizing-module-CM-scheme} により $X_y$ は Cohen-Macaulay である。
所望のとおり、これで $f$ は $x$ において Cohen-Macaulay であることが示された。
\end{proof}

\begin{remark}
\label{duality-remark-CM-morphism-compare-dualizing}
状況 \ref{duality-situation-shriek} において、$f : X \to Y$ を
$\textit{FTS}_S$ の射とする。$f$ は相対次元 $d$ の Cohen-Macaulay 射であると仮定する。
$\omega_{X/Y} = H^{-d}(f^!\mathcal{O}_Y)$ を $f^!\mathcal{O}_Y$ の一意な非零コホモロジー層とする。
補題 \ref{duality-lemma-CM-shriek} を参照する。このとき、すべての
$K \in D^+_\QCoh(\mathcal{O}_Y)$ に対して標準同型
$$
f^!K = Lf^*K \otimes_{\mathcal{O}_X}^\mathbf{L} \omega_{X/Y}[d]
$$
がある。補題 \ref{duality-lemma-perfect-comparison-shriek} を参照する。特に、$S$ が双対化複体
$\omega_S^\bullet$ をもち、
$\omega_Y^\bullet = (Y \to S)^!\omega_S^\bullet$, かつ
$\omega_X^\bullet = (X \to S)^!\omega_S^\bullet$
ならば
$$
\omega_X^\bullet =
Lf^*\omega_Y^\bullet \otimes_{\mathcal{O}_X}^\mathbf{L} \omega_{X/Y}[d]
$$
である。さらに $X$ と $Y$ が連結かつ Cohen-Macaulay であり、$\omega_Y$ と $\omega_X$ がそれぞれ
$\omega_Y^\bullet$ と $\omega_X^\bullet$ の一意な非零コホモロジー層を表すならば
$$
\omega_X = f^*\omega_Y \otimes_{\mathcal{O}_X} \omega_{X/Y}.
$$
同様の結果は、$S$ 上有限型な任意の $X$ と $Y$（すなわち $S$ 上分離とは限らない）に対しても、
第 \ref{duality-section-glue} 節のように $\omega_S^\bullet$ に関して正規化された双対化複体を用いれば成り立つ。
\end{remark}





\section{Gorenstein スキーム}
\label{duality-section-gorenstein}

\noindent
本節は『双対化複体』の第 \ref{dualizing-section-gorenstein} 節の続きである。

\begin{definition}
\label{duality-definition-gorenstein}
$X$ をスキームとする。$X$ が局所 Noether であり、すべての $x \in X$ に対して
$\mathcal{O}_{X, x}$ が Gorenstein であるとき、$X$ は {\it Gorenstein} であるという。
\end{definition}

\noindent
この定義は意味をもつ。実際、Noether 環が Gorenstein であるとは、そのすべての局所環が
Gorenstein であることをいう。『双対化複体』の定義 \ref{dualizing-definition-gorenstein} を参照する。

\begin{lemma}
\label{duality-lemma-gorenstein-CM}
Gorenstein スキームは Cohen-Macaulay である。
\end{lemma}

\begin{proof}
アフィン局所的に見れば、代数における対応する結果、すなわち『双対化複体』の補題 \ref{dualizing-lemma-gorenstein-CM} から従う。
\end{proof}

\begin{lemma}
\label{duality-lemma-regular-gorenstein}
正則スキームは Gorenstein である。
\end{lemma}

\begin{proof}
アフィン局所的に見れば、代数における対応する結果、すなわち『双対化複体』の補題 \ref{dualizing-lemma-regular-gorenstein} から従う。
\end{proof}

\begin{lemma}
\label{duality-lemma-gorenstein}
$X$ を局所 Noether スキームとする。
\begin{enumerate}
\item $X$ が双対化複体 $\omega_X^\bullet$ をもつならば、
\begin{enumerate}
\item $X$ が Gorenstein であること $\Leftrightarrow$ $\omega_X^\bullet$ が
$D(\mathcal{O}_X)$ の可逆対象であること。
\item $\mathcal{O}_{X, x}$ が Gorenstein であること $\Leftrightarrow$
$\omega_{X, x}^\bullet$ が $D(\mathcal{O}_{X, x})$ の可逆対象であること。
\item $U = \{x \in X \mid \mathcal{O}_{X, x}\text{ is Gorenstein}\}$ は開 Gorenstein 部分スキームである。
\end{enumerate}
\item $X$ が Gorenstein ならば、$X$ が双対化複体をもつことと
$\mathcal{O}_X[0]$ が双対化複体であることは同値である。
\end{enumerate}
\end{lemma}

\begin{proof}
アフィン局所的に見れば、代数における対応する結果、すなわち『双対化複体』の補題 \ref{dualizing-lemma-gorenstein} から従う。
\end{proof}

\begin{lemma}
\label{duality-lemma-gorenstein-lci}
$f : Y \to X$ が局所完全交叉射であり、$X$ が Gorenstein スキームならば、$Y$ は Gorenstein である。
\end{lemma}

\begin{proof}
『射の続論』の補題 \ref{more-morphisms-lemma-affine-lci} により、
環準同型に関する対応する主張を示せば十分である。これは『双対化複体』の補題 \ref{dualizing-lemma-gorenstein-lci} である。
\end{proof}

\begin{lemma}
\label{duality-lemma-gorenstein-local-syntomic}
性質 $\mathcal{P}(S) =$``$S$ is Gorenstein'' は syntomic 位相に関して局所的である。
\end{lemma}

\begin{proof}
$\{S_i \to S\}$ を syntomic 被覆とする。『降下』の補題 \ref{descent-lemma-Noetherian-local-fppf} により、スキーム $S$ が局所 Noether であることと、
各 $S_i$ が Noether であることは同値である。従って $S$ と $S_i$ は局所 Noether であると仮定してよい。
$S$ が Gorenstein ならば、補題 \ref{duality-lemma-gorenstein-lci} により各 $S_i$ は Gorenstein である。
逆に各 $S_i$ が Gorenstein ならば、各点 $s \in S$ に対し、$s$ に写るある $i$ と
$t \in S_i$ を選べる。このとき
$\mathcal{O}_{S, s} \to \mathcal{O}_{S_i, t}$ は平坦な局所環準同型であり、
$\mathcal{O}_{S_i, t}$ は Gorenstein である。従って『双対化複体』の補題 \ref{dualizing-lemma-flat-under-gorenstein} により $\mathcal{O}_{S, s}$ は Gorenstein である。
\end{proof}






\section{Gorenstein 射}
\label{duality-section-gorenstein-morphisms}

\noindent
本節は一連の節の一つである。正規射、正則射、および Cohen-Macaulay 射に対応する節は、
『射の続論』の第 \ref{more-morphisms-section-normal} 節、第 \ref{more-morphisms-section-regular} 節および第 \ref{more-morphisms-section-CM} 節
にある。

\medskip\noindent
次の補題によれば、幾何学的 Gorenstein スキームを定義しても意味がない。
そのようなスキームは Gorenstein スキームと同じものになるからである。

\begin{lemma}
\label{duality-lemma-gorenstein-base-change}
$X$ を体 $k$ 上の局所 Noether スキームとする。
$k'/k$ を有限生成体拡大とする。
$x \in X$ を点とし、$x$ 上にある点 $x' \in X_{k'}$ をとる。このとき
$$
\mathcal{O}_{X, x}\text{ is Gorenstein}
\Leftrightarrow
\mathcal{O}_{X_{k'}, x'}\text{ is Gorenstein}
$$
が成り立つ。$X$ が $k$ 上局所有限型ならば、任意の体拡大 $k'/k$ に対して同じことが成り立つ。
\end{lemma}

\begin{proof}
いずれの場合も、環準同型
$\mathcal{O}_{X, x} \to \mathcal{O}_{X_{k'}, x'}$
は Noether 局所環の忠実平坦な局所準同型である。従って
$\mathcal{O}_{X_{k'}, x'}$ が Gorenstein ならば、『双対化複体』の補題 \ref{dualizing-lemma-flat-under-gorenstein} により
$\mathcal{O}_{X, x}$ も Gorenstein である。逆向きにも『双対化複体』の補題 \ref{dualizing-lemma-flat-under-gorenstein} を用いる。従って示すべきことは
$$
\mathcal{O}_{X_{k'}, x'}/\mathfrak m_x \mathcal{O}_{X_{k'}, x'} =
\kappa(x) \otimes_k k'
$$
が Gorenstein であることである。第1の場合には $k \to k'$ が有限生成であり、
第2の場合には $k \to \kappa(x)$ が有限生成であることに注意する。
従って Gorenstein に対して性質 (A) が成り立つことから従う。『双対化複体』の補題 \ref{dualizing-lemma-formal-fibres-gorenstein} を参照する。
\end{proof}

\noindent
上の補題により、次が Gorenstein 射の正しい定義であることが保証される。

\begin{definition}
\label{duality-definition-gorenstein-morphism}
$f : X \to Y$ をスキームの射とする。
すべてのファイバー $X_y$ が局所 Noether スキームであると仮定する。
\begin{enumerate}
\item $x \in X$ とし、$y = f(x)$ とおく。$f$ が $x$ において平坦であり、
スキーム $X_y$ の $x$ における局所環が Gorenstein であるとき、$f$ は
{\it $x$ において Gorenstein} であるという。
\item $f$ が $X$ のすべての点で Gorenstein であるとき、$f$ を
{\it Gorenstein 射}という。
\end{enumerate}
\end{definition}

\noindent
これは次のように言い換えられる。

\begin{lemma}
\label{duality-lemma-gorenstein-morphism}
$f : X \to Y$ をスキームの射とし、$f$ のすべてのファイバーが局所 Noether であると仮定する。
次は同値である。
\begin{enumerate}
\item $f$ は Gorenstein である。
\item $f$ は平坦であり、そのファイバーは Gorenstein スキームである。
\end{enumerate}
\end{lemma}

\begin{proof}
定義から直ちに従う。
\end{proof}

\begin{lemma}
\label{duality-lemma-gorenstein-CM-morphism}
Gorenstein 射は Cohen-Macaulay である。
\end{lemma}

\begin{proof}
補題 \ref{duality-lemma-gorenstein-CM} と定義から従う。
\end{proof}

\begin{lemma}
\label{duality-lemma-lci-gorenstein}
syntomic 射は Gorenstein である。同値な言い方をすれば、平坦な局所完全交叉射は Gorenstein である。
\end{lemma}

\begin{proof}
syntomic 射は平坦であり、そのファイバーは体上の局所完全交叉であることを思い出す。
『スキームの射』の補題 \ref{morphisms-lemma-syntomic-flat-fibres} を参照する。
体上の局所完全交叉は補題 \ref{duality-lemma-gorenstein-lci} により Gorenstein スキームなので結論を得る。
性質 ``syntomic'' と ``flat and local
complete intersection morphism'' は『射の続論』の補題 \ref{more-morphisms-lemma-flat-lci} により同値である。
\end{proof}

\begin{lemma}
\label{duality-lemma-composition-gorenstein}
$f : X \to Y$ および $g : Y \to Z$ を射とする。$f$, $g$, $g \circ f$ のファイバー
$X_y$, $Y_z$, $X_z$ は局所 Noether であると仮定する。
\begin{enumerate}
\item $f$ が $x$ において Gorenstein であり、$g$ が $f(x)$ において Gorenstein ならば、
$g \circ f$ は $x$ において Gorenstein である。
\item $f$ と $g$ が Gorenstein ならば、$g \circ f$ は Gorenstein である。
\item $g \circ f$ が $x$ において Gorenstein であり、$f$ が $x$ において平坦ならば、
$f$ は $x$ において Gorenstein であり、$g$ は $f(x)$ において Gorenstein である。
\item $g \circ f$ が Gorenstein であり、$f$ が平坦ならば、$f$ は Gorenstein であり、
$g$ は $f$ の像のすべての点において Gorenstein である。
\end{enumerate}
\end{lemma}

\begin{proof}
代数に言い換えれば、『双対化複体』の補題 \ref{dualizing-lemma-flat-under-gorenstein} から従う。
\end{proof}

\begin{lemma}
\label{duality-lemma-flat-morphism-from-gorenstein-scheme}
\begin{slogan}
平坦ファイブレーションの全空間が Gorenstein ならば、底空間とファイバーも同様である
\end{slogan}
$f : X \to Y$ を局所 Noether スキームの平坦射とする。
$X$ が Gorenstein ならば、$f$ は Gorenstein であり、すべての $x \in X$ に対して
$\mathcal{O}_{Y, f(x)}$ は Gorenstein である。
\end{lemma}

\begin{proof}
代数に言い換えれば、『双対化複体』の補題 \ref{dualizing-lemma-flat-under-gorenstein} から従う。
\end{proof}

\begin{lemma}
\label{duality-lemma-base-change-gorenstein}
$f : X \to Y$ をスキームの射とする。
すべてのファイバー $X_y$ が局所 Noether スキームであると仮定する。
$Y' \to Y$ は局所有限型であるとし、$f' : X' \to Y'$ を $f$ の基底変換とする。
$x' \in X'$ を、その像が $x \in X$ である点とする。
\begin{enumerate}
\item $f$ が $x$ において Gorenstein ならば、
$f' : X' \to Y'$ は $x'$ において Gorenstein である。
\item $f$ が $x$ において平坦であり、$f'$ が $x'$ において Gorenstein ならば、
$f$ は $x$ において Gorenstein である。
\item $Y' \to Y$ が $f'(x')$ において平坦であり、$f'$ が $x'$ において Gorenstein ならば、
$f$ は $x$ において Gorenstein である。
\end{enumerate}
\end{lemma}

\begin{proof}
$Y' \to Y$ に関する仮定により、$y \in Y$ に写る $y' \in Y'$ に対して
体拡大 $\kappa(y')/\kappa(y)$ は有限生成であることに注意する。従ってすべてのファイバー
$X'_{y'} = (X_y)_{\kappa(y')}$ も局所 Noether である。『多様体』の補題 \ref{varieties-lemma-locally-Noetherian-base-change} を参照する。
従って補題の主張には意味がある。$y' = f'(x')$ および $y = f(x)$ とおく。
このとき局所環の次の可換図式を得る。
$$
\xymatrix{
\mathcal{O}_{X', x'} & \mathcal{O}_{X, x} \ar[l] \\
\mathcal{O}_{Y', y'} \ar[u] & \mathcal{O}_{Y, y} \ar[l] \ar[u]
}
$$
ここで左上の環は、右下の環上で右上の環と左下の環をテンソル積したものの局所化である。

\medskip\noindent
$f$ が $x$ において Gorenstein であると仮定する。
$\mathcal{O}_{Y, y} \to \mathcal{O}_{X, x}$ の平坦性から
$\mathcal{O}_{Y', y'} \to \mathcal{O}_{X', x'}$ の平坦性が従う。『可換代数』の補題 \ref{algebra-lemma-base-change-flat-up-down} を参照する。
$\mathcal{O}_{X, x}/\mathfrak m_y\mathcal{O}_{X, x}$ が Gorenstein であることから、
補題 \ref{duality-lemma-gorenstein-base-change} により
$\mathcal{O}_{X', x'}/\mathfrak m_{y'}\mathcal{O}_{X', x'}$
も Gorenstein である。従って $f'$ は $x'$ において Gorenstein である。

\medskip\noindent
$f$ が $x$ において平坦であり、$f'$ が $x'$ において Gorenstein であると仮定する。
$\mathcal{O}_{X', x'}/\mathfrak m_{y'}\mathcal{O}_{X', x'}$
が Gorenstein であることから、補題 \ref{duality-lemma-gorenstein-base-change} により
$\mathcal{O}_{X, x}/\mathfrak m_y\mathcal{O}_{X, x}$
も Gorenstein である。従って $f$ は $x$ において Gorenstein である。

\medskip\noindent
$Y' \to Y$ が $y'$ において平坦であり、$f'$ が $x'$ において Gorenstein であると仮定する。
$\mathcal{O}_{Y', y'} \to \mathcal{O}_{X', x'}$ と
$\mathcal{O}_{Y, y} \to \mathcal{O}_{Y', y'}$ の平坦性から
$\mathcal{O}_{Y, y} \to \mathcal{O}_{X, x}$ の平坦性が従う。『可換代数』の補題 \ref{algebra-lemma-base-change-flat-up-down} を参照する。
$\mathcal{O}_{X', x'}/\mathfrak m_{y'}\mathcal{O}_{X', x'}$
が Gorenstein であることから、補題 \ref{duality-lemma-gorenstein-base-change} により
$\mathcal{O}_{X, x}/\mathfrak m_y\mathcal{O}_{X, x}$
も Gorenstein である。従って $f$ は $x$ において Gorenstein である。
\end{proof}

\begin{lemma}
\label{duality-lemma-flat-lft-base-change-gorenstein}
$f : X \to Y$ を平坦かつ局所有限型なスキームの射とする。このとき集合
$\{x \in X \mid f\text{ is Gorenstein at }x\}$
の形成は任意の基底変換と可換である。
\end{lemma}

\begin{proof}
仮定により、$f$ の任意のファイバーは体上局所有限型であり、従って局所 Noether である。
任意の基底変換についても同様である。従って主張には意味がある。
ファイバーを考えることで、次の問題に帰着する。$X$ を体 $k$ 上局所有限型なスキーム、
$K/k$ を体拡大、$x_K \in X_K$ をその像が $x \in X$ である点とする。
問題は、$\mathcal{O}_{X_K, x_K}$ が Gorenstein であることと
$\mathcal{O}_{X, x}$ が Gorenstein であることが同値であると示すことである。

\medskip\noindent
この問題は、少し代数を用いて次のように解ける。
$x$ を含むアフィン開集合 $\Spec(A) \subset X$ を選ぶ。
$x$ は $\mathfrak p \subset A$ に対応するとする。
$A_K = A \otimes_k K$ とおけば、$\Spec(A_K) \subset X_K$ は $x_K$ を含む。
$x_K$ は $\mathfrak p_K \subset A_K$ に対応するとする。
$\omega_A^\bullet$ を $A$ の双対化複体とする。
『双対化複体』の補題 \ref{dualizing-lemma-base-change-dualizing-over-field} により
$\omega_A^\bullet \otimes_A A_K$ は $A_K$ の双対化複体である。
ここで $A_\mathfrak p \to (A_K)_{\mathfrak p_K}$ は Noether 環の平坦な局所準同型なので、
$(\omega_A^\bullet)_\mathfrak p$ が $D(A_\mathfrak p)$ の可逆対象であることと
$(\omega_A^\bullet)_\mathfrak p \otimes_{A_\mathfrak p} (A_K)_{\mathfrak p_K}$
が $D((A_K)_{\mathfrak p_K})$ の可逆対象であることは同値であり、結論を得る。
細部は省略する。ヒントはコホモロジー加群を調べることである。
\end{proof}

\begin{lemma}
\label{duality-lemma-affine-flat-Noetherian-gorenstein}
状況 \ref{duality-situation-shriek} において、$f : X \to Y$ を
$\textit{FTS}_S$ の射とし、$x \in X$ とする。$f$ が平坦ならば、次は同値である。
\begin{enumerate}
\item $f$ は $x$ において Gorenstein である。
\item $f^!\mathcal{O}_Y$ は $x$ の近傍で可逆対象に同型である。
\end{enumerate}
特に、$f$ が Gorenstein である点の集合は $X$ の開集合である。
\end{lemma}

\begin{proof}
$\omega^\bullet = f^!\mathcal{O}_Y$ とおく。補題 \ref{duality-lemma-relative-dualizing-fibres} により、これは連接コホモロジー層をもつ有界複体であり、
ファイバー $X_y$ への導来制限 $Lh^*\omega^\bullet$ は $X_y$ 上の双対化複体である。
点の包含を $i : x \to X_y$ と書く。このとき次は同値である。
\begin{enumerate}
\item $f$ は $x$ において Gorenstein である。
\item $\mathcal{O}_{X_y, x}$ は Gorenstein である。
\item $Lh^*\omega^\bullet$ は $x$ の近傍で可逆である。
\item $Li^* Lh^* \omega^\bullet$ は $\kappa(x)$ 上次元 $1$ の非零コホモロジーをちょうど一つもつ。
\item $L(h \circ i)^* \omega^\bullet$ は $\kappa(x)$ 上次元 $1$ の非零コホモロジーをちょうど一つもつ。
\item $\omega^\bullet$ は $x$ の近傍で可逆である。
\end{enumerate}
(1) と (2) の同値性は定義による（$f$ は平坦である）。
(2) と (3) の同値性は補題 \ref{duality-lemma-gorenstein} から従う。
(3) と (4) の同値性は『代数の続論』の補題 \ref{more-algebra-lemma-lift-bounded-pseudo-coherent-to-perfect} から従う。
(4) と (5) の同値性は $Li^* Lh^* = L(h \circ i)^*$ から従う。
(5) と (6) の同値性は『代数の続論』の補題 \ref{more-algebra-lemma-lift-bounded-pseudo-coherent-to-perfect} から従う。
以上で補題は明らかである。
\end{proof}

\begin{lemma}
\label{duality-lemma-flat-finite-presentation-characterize-gorenstein}
$f : X \to S$ を平坦かつ局所有限表示なスキームの射とする。
$x \in X$ の像を $s \in S$ とし、$d = \dim_x(X_s)$ とおく。次は同値である。
\begin{enumerate}
\item $f$ は $x$ において Gorenstein である。
\item $x$ の開近傍 $U \subset X$ と、$S$ 上の局所準有限射
$U \to \mathbf{A}^d_S$ で $x$ において Gorenstein なものが存在する。
\item $x$ の開近傍 $U \subset X$ と、$S$ 上の局所準有限 Gorenstein 射
$U \to \mathbf{A}^d_S$ が存在する。
\item $x$ の任意の開近傍 $U \subset X$ からの任意の $S$-射
$g : U \to \mathbf{A}^d_S$ に対して、
$g$ が $x$ において準有限であること $\Rightarrow$ $g$ が $x$ において Gorenstein であること。
\end{enumerate}
特に、$f$ が Gorenstein である点の集合は $X$ の開集合である。
\end{lemma}

\begin{proof}
$x \in U$ を満たすアフィン開集合 $U = \Spec(A) \subset X$ と、
$f(U) \subset V$ を満たす $V = \Spec(R) \subset S$ を選ぶ。このとき
$R \to A$ は有限表示の平坦環準同型である。$\mathfrak p \subset A$ を $x$ に対応する素イデアルとする。
$A$ をある単項局所化で置き換えれば、準有限準同型
$R[x_1, \ldots, x_d]  \to A$ が存在すると仮定してよい。『可換代数』の補題 \ref{algebra-lemma-quasi-finite-over-polynomial-algebra} を参照する。
従って、$x$ の開近傍 $U \subset X$ と局所\footnote{$S$
が準分離ならば、$g$ は準有限になる。}準有限射
$g : U \to \mathbf{A}^d_S$ からなる組 $(U, g)$ が少なくとも一つ存在する。

\medskip\noindent
以上を踏まえると、補題は次の代数的問題に言い換えられる
（言い換えの詳細は省略する）。$R \to A$ は平坦かつ有限表示であり、
$\mathfrak p \subset A$ は素イデアルであり、
$\varphi : R[x_1, \ldots, x_d] \to A$ は $\mathfrak p$ において準有限であるとする。
このとき次は同値である。
\begin{enumerate}
\item[(a)] $\Spec(\varphi)$ は $\mathfrak p$ において Gorenstein である。
\item[(b)] $\Spec(A) \to \Spec(R)$ は $\mathfrak p$ において Gorenstein である。
\item[(c)] $\Spec(A) \to \Spec(R)$ は $\mathfrak p$ のある開近傍で Gorenstein である。
\end{enumerate}
いずれの場合にも $R[x_1, \ldots, x_n] \to A$ は $\mathfrak p$ において平坦である。
従って平坦性の開性
（『可換代数』の定理 \ref{algebra-theorem-openness-flatness}）により、
$A$ を適当な単項局所化で置き換えて
$R[x_1, \ldots, x_n] \to A$ が平坦であると仮定してよい。
『可換代数』の補題 \ref{algebra-lemma-flat-finite-presentation-limit-flat} により、
$R_0$ が $\mathbf{Z}$ 上有限型で、$R_0 \to A_0$ が有限型であり、
$R_0[x_1, \ldots, x_n] \to A_0$ が平坦となる
$R_0 \subset R$ と $R_0[x_1, \ldots, x_n] \to A_0$ が存在する。
平坦有限型射が Gorenstein である点の集合は、補題 \ref{duality-lemma-base-change-gorenstein} により基底変換と可換であることに注意する。
このようにして $R$ が Noether である場合に帰着する。

\medskip\noindent
従って $X$ と $S$ はアフィンであり、$f$ は
$$
X \xrightarrow{g} \mathbf{A}^n_S \xrightarrow{p} S
$$
と分解され、$g$ は平坦かつ準有限で、$S$ は Noether であると仮定してよい。
このとき $X$ と $\mathbf{A}^n_S$ は $S$ 上分離的であり、上付きシュリーク関手の既知の性質
（補題 \ref{duality-lemma-upper-shriek-composition} および \ref{duality-lemma-shriek-affine-line}）により
$$
f^!\mathcal{O}_S = g^!p^!\mathcal{O}_S = g^!\mathcal{O}_{\mathbf{A}^n_S}[n]
$$
を得る。従って補題 \ref{duality-lemma-affine-flat-Noetherian-gorenstein} により
(a), (b), (c) は同値である。
\end{proof}

\begin{lemma}
\label{duality-lemma-gorenstein-local-source-and-target}
性質
$\mathcal{P}(f)=$「$f$ のファイバーは局所 Noether であり、$f$ は Gorenstein である」
は標的上の fppf 位相に関して局所的であり、始域上の syntomic 位相に関して局所的である。
\end{lemma}

\begin{proof}
$\mathcal{P}(f) =
\mathcal{P}_1(f) \wedge \mathcal{P}_2(f)$
と書ける。ここで
$\mathcal{P}_1(f)=$「$f$ は平坦である」、
$\mathcal{P}_2(f)=$「$f$ のファイバーは局所 Noether かつ Gorenstein である」
とする。$\mathcal{P}_1$ は始域および標的上の fppf 位相に関して局所的であることが分かっている。
『降下』の補題 \ref{descent-lemma-descending-property-flat} および \ref{descent-lemma-flat-fpqc-local-source} を参照する。
従って $\mathcal{P}_2$ を扱えばよい。

\medskip\noindent
$f : X \to Y$ をスキームの射とする。
$\{\varphi_i : Y_i \to Y\}_{i \in I}$ を $Y$ の fppf 被覆とする。
$f_i : X_i \to Y_i$ を $\varphi_i$ による $f$ の基底変換とする。
$i \in I$ と点 $y_i \in Y_i$ をとり、$y = \varphi_i(y_i)$ とおく。このとき
$$
X_{i, y_i} = \Spec(\kappa(y_i)) \times_{\Spec(\kappa(y))} X_y.
$$
であり、$\kappa(y_i)/\kappa(y)$ は有限生成体拡大である。
従って $X_y$ が局所 Noether ならば、『多様体』の補題 \ref{varieties-lemma-locally-Noetherian-base-change} により
$X_{i, y_i}$ も局所 Noether である。さらに $X_y$ が Gorenstein ならば、
補題 \ref{duality-lemma-gorenstein-base-change} により
$X_{i, y_i}$ も Gorenstein である。
従って $\mathcal{P}_2$ は標的上 fppf 局所的である。

\medskip\noindent
$\{X_i \to X\}$ を $X$ の syntomic 被覆とし、$y \in Y$ とする。
この場合 $\{X_{i, y} \to X_y\}$ はファイバーの syntomic 被覆である。
従って始域上の syntomic 位相に関する $\mathcal{P}_2$ の局所性は
補題 \ref{duality-lemma-gorenstein-local-syntomic} から従う。
\end{proof}




\section{双対化複体についての補足}
\label{duality-section-more-dualizing}

\noindent
ほかのどこにも収まりにくい補題をいくつか述べる。

\begin{lemma}
\label{duality-lemma-descent-ascent}
$f : X \to Y$ を局所 Noether スキームの射とする。次のいずれかを仮定する。
\begin{enumerate}
\item $f$ は syntomic かつ全射である。
\item $f$ は全射な平坦局所完全交叉射である。
\item $f$ は有限型の全射 Gorenstein 射である。
\end{enumerate}
このとき $K \in D_\QCoh(\mathcal{O}_Y)$ が $Y$ 上の双対化複体であることと、
$Lf^*K$ が $X$ 上の双対化複体であることは同値である。
\end{lemma}

\begin{proof}
アフィン開集合をとり、『スキームの導来圏』の補題 \ref{perfect-lemma-affine-compare-bounded} を用いれば、これは
『双対化複体』の補題 \ref{dualizing-lemma-descent-ascent}
に言い換えられる。
\end{proof}







\section{体上固有なスキームに対する双対性}
\label{duality-section-duality-proper-over-field}

\noindent
本節では、双対化複体に関する上の非常に一般的な議論から、体上固有なスキームの双対性について
得られる帰結を詳しく述べる。

\begin{lemma}
\label{duality-lemma-duality-proper-over-field}
$X$ を体 $k$ 上固有なスキームとする。次の性質をもつ双対化複体
$\omega_X^\bullet$ が存在する。
\begin{enumerate}
\item $H^i(\omega_X^\bullet)$ が非零となり得るのは $i \in [-\dim(X), 0]$ のときに限る。
\item $\omega_X = H^{-\dim(X)}(\omega_X^\bullet)$ は連接
$(S_2)$-加群であり、その台は次元 $\dim(X)$ の既約成分からなる。
\item $H^i(\omega_X^\bullet)$ の台の次元は高々 $-i$ である。
\item 閉点 $x \in X$ に対し、加群
$H^i(\omega_{X, x}^\bullet) \oplus \ldots \oplus H^0(\omega_{X, x}^\bullet)$
が非零であることと $\text{depth}(\mathcal{O}_{X, x}) \leq -i$ は同値である。
\item $K \in D_\QCoh(\mathcal{O}_X)$ に対して、シフトおよび完全三角形と両立する関手的同型
\footnote{Yoneda の補題により、この性質は $D_\QCoh(\mathcal{O}_X)$ において
$\omega_X^\bullet$ を一意同型を除いて特徴づける。$\omega_X^\bullet$ は
$D^b_{\textit{Coh}}(\mathcal{O}_X)$ に属するので、実際には
$K \in D^b_{\textit{Coh}}(\mathcal{O}_X)$ だけを考えれば十分である。}
$$
\Ext^i_X(K, \omega_X^\bullet) = \Hom_k(H^{-i}(X, K), k)
$$
が存在する。
\item $X$ 上準連接な $\mathcal{F}$ に対して、関手的同型
$\Hom(\mathcal{F}, \omega_X) = \Hom_k(H^{\dim(X)}(X, \mathcal{F}), k)$
が存在する。
\item $X \to \Spec(k)$ が相対次元 $d$ の滑らかな射ならば、
$\omega_X^\bullet \cong \wedge^d\Omega_{X/k}[d]$ かつ
$\omega_X \cong \wedge^d\Omega_{X/k}$ である。
\end{enumerate}
\end{lemma}

\begin{proof}
構造射を $f : X \to \Spec(k)$ と書く。この射の押し出しの右随伴を $a$ とする。
補題 \ref{duality-lemma-twisted-inverse-image} を参照する。相対双対化複体
$$
\omega_X^\bullet = a(\mathcal{O}_{\Spec(k)})
$$
を考える。注意 \ref{duality-remark-relative-dualizing-complex} と比較せよ。
$f$ は固有なので、定義により
$f^!(\mathcal{O}_{\Spec(k)}) = a(\mathcal{O}_{\Spec(k)})$
である。第 \ref{duality-section-upper-shriek} 節を参照する。
補題 \ref{duality-lemma-shriek-dualizing} を適用すると、
$\omega_X^\bullet$ は双対化複体である。また、
$\omega_X^\bullet$ と $\omega_X$ は例 \ref{duality-example-proper-over-local} および例 \ref{duality-example-equidimensional-over-field} におけるものと同じである。

\medskip\noindent
(1), (2), (3) は補題 \ref{duality-lemma-vanishing-good-dualizing} から従う。

\medskip\noindent
閉点 $x \in X$ に対し、$\omega_{X, x}^\bullet$ は
$\mathcal{O}_{X, x}$ 上の正規化双対化複体である。補題 \ref{duality-lemma-good-dualizing-normalized} を参照する。
従って (4) は『双対化複体』の補題 \ref{dualizing-lemma-depth-in-terms-dualizing-complex} から従う。

\medskip\noindent
(5) は構成から従う。実際、$a$ は
$Rf_* : D_\QCoh(\mathcal{O}_X) \to D(\mathcal{O}_{\Spec(k)}) = D(k)$
の右随伴であり、この関手は $K \mapsto R\Gamma(X, K)$ と同一視できる。
また、$k$-加群の導来圏 $D(k)$ は次数付き $k$-ベクトル空間の圏と同じであることを用いる。

\medskip\noindent
連接な $\mathcal{F}$ に対する (6) は補題 \ref{duality-lemma-dualizing-module-proper-over-A} から従う。
一般の場合は、$K = \mathcal{F}[0]$ および $i = -\dim(X)$ として (5) を展開すれば従う。

\medskip\noindent
(7) は補題 \ref{duality-lemma-smooth-proper} から従う。
\end{proof}

\begin{remark}
\label{duality-remark-duality-proper-over-field}
$k$, $X$, $\omega_X^\bullet$ は補題 \ref{duality-lemma-duality-proper-over-field} と同じものとする。
複体 $\omega_X^\bullet$ の恒等射は、(5) の関手的同型を通じて写像
$$
t : H^0(X, \omega_X^\bullet) \longrightarrow k
$$
に対応する。$D_\QCoh(\mathcal{O}_X)$ の任意の $K$ に対し、(5) における
$\Hom(K, \omega_X^\bullet)$ と $H^0(X, K)^\vee$ の同一視は、ペアリング
$$
\Hom_X(K, \omega_X^\bullet) \times H^0(X, K) \longrightarrow k,\quad
(\alpha, \beta) \longmapsto t(\alpha(\beta))
$$
に対応する。これは (5) の同型の関手性から従う。同様に、任意の
$i \in \mathbf{Z}$ に対してペアリング
$$
\Ext^i_X(K, \omega_X^\bullet) \times H^{-i}(X, K) \longrightarrow k,\quad
(\alpha, \beta) \longmapsto t(\alpha(\beta))
$$
を得る。ここで $\alpha(\beta)$ を定義するため、$\alpha$ を射
$K[-i] \to \omega_X^\bullet$ とみなし、$\beta$ を
$H^0(X, K[-i])$ の元とみなす。
$K$ が一般の場合、このペアリングが一方の側で非退化であることしか分からないことに注意する。
すなわち、ペアリングは $\Hom_X(K, \omega_X^\bullet)$,
resp.\ $\Ext^i_X(K, \omega_X^\bullet)$ と $H^0(X, K)$,
resp.\ $H^{-i}(X, K)$ の $k$-線形双対との同型を誘導するが、一般には逆向きの同型を誘導しない。
$K$ が $D^b_{\textit{Coh}}(\mathcal{O}_X)$ に属するならば、
$\Hom_X(K, \omega_X^\bullet)$, $\Ext_X(K, \omega_X^\bullet)$,
$H^0(X, K)$, $H^i(X, K)$ は有限次元 $k$-ベクトル空間であり
（『スキームの導来圏』の補題 \ref{perfect-lemma-coherent-internal-hom} および \ref{perfect-lemma-direct-image-coherent-bdd-below} による）、
ペアリングは通常の意味で完全である。
\end{remark}

\begin{remark}
\label{duality-remark-coherent-duality-proper-over-field}
注意 \ref{duality-remark-duality-proper-over-field} の議論を続け、同じ記号
$k$, $X$, $\omega_X^\bullet$, $t$ を用いる。
$\mathcal{F}$ が連接 $\mathcal{O}_X$-加群ならば、有限次元 $k$-ベクトル空間の完全ペアリング
$$
\langle -, - \rangle :
\Ext^i_X(\mathcal{F}, \omega_X^\bullet) \times H^{-i}(X,\mathcal{F})
\longrightarrow k,\quad
(\alpha, \beta) \longmapsto t(\alpha(\beta))
$$
を得る。これらのペアリングは次の（明らかな）関手性を満たす。
$\varphi : \mathcal{F} \to \mathcal{G}$ が連接 $\mathcal{O}_X$-加群の準同型ならば、
$$
\langle \alpha \circ \varphi, \beta \rangle =
\langle \alpha, \varphi(\beta) \rangle
$$
が $\alpha \in \Ext^i_X(\mathcal{G}, \omega_X^\bullet)$ および
$\beta \in H^{-i}(X, \mathcal{F})$ に対して成り立つ。言い換えれば、$\varphi$ が誘導する
$k$-線形写像
$\Ext^i_X(\mathcal{G}, \omega_X^\bullet) \to
\Ext^i_X(\mathcal{F}, \omega_X^\bullet)$
は、これらのペアリングを通じて、$\varphi$ が誘導する $k$-線形写像
$H^{-i}(X, \mathcal{F}) \to H^{-i}(X, \mathcal{G})$ の $k$-線形双対である。
この形で述べれば、$\varphi$ が準連接 $\mathcal{O}_X$-加群の準同型である場合にも同じことが成り立つ。
\end{remark}

\begin{lemma}
\label{duality-lemma-duality-proper-over-field-perfect}
$k$, $X$, $\omega_X^\bullet$ は補題 \ref{duality-lemma-duality-proper-over-field} と同じものとする。
$t : H^0(X, \omega_X^\bullet) \to k$ は注意 \ref{duality-remark-duality-proper-over-field} と同じものとする。
$E \in D(\mathcal{O}_X)$ を完全とする。このとき、すべての $i$ に対してペアリング
$$
H^i(X, \omega_X^\bullet \otimes_{\mathcal{O}_X}^\mathbf{L} E^\vee)
\times
H^{-i}(X, E) \longrightarrow k, \quad
(\xi, \eta) \longmapsto
t((1_{\omega_X^\bullet} \otimes \epsilon)(\xi \cup \eta))
$$
は完全である。ここで $\cup$ は『層のコホモロジー』の第 \ref{cohomology-section-cup-product} 節のカップ積を表し、
$\epsilon : E^\vee \otimes_{\mathcal{O}_X}^\mathbf{L} E \to \mathcal{O}_X$
は『層のコホモロジー』の例 \ref{cohomology-example-dual-derived} のものとする。
\end{lemma}

\begin{proof}
$E$ を $E[-i]$ で置き換えることにより、$i = 0$ の場合に帰着する。
『層のコホモロジー』の補題 \ref{cohomology-lemma-ext-composition-is-cup} により、
このペアリングは注意 \ref{duality-remark-duality-proper-over-field} で論じたものと同じである。
従って同 Remark の議論から結果が従う。
\end{proof}

\begin{lemma}
\label{duality-lemma-duality-proper-over-field-CM}
$X$ を体 $k$ 上固有で、Cohen-Macaulay かつ次元 $d$ の等次元スキームとする。
補題 \ref{duality-lemma-duality-proper-over-field} の加群 $\omega_X$ は次の性質をもつ。
\begin{enumerate}
\item $\omega_X$ は $X$ 上の双対化加群である
（第 \ref{duality-section-dualizing-module} 節）。
\item $\omega_X$ は台が $X$ である連接 Cohen-Macaulay 加群である。
\item $K \in D_\QCoh(X)$ に対して、シフトおよび完全三角形と両立する関手的同型
$\Ext^i_X(K, \omega_X[d]) = \Hom_k(H^{-i}(X, K), k)$
が存在する。
\item $X$ 上準連接な $\mathcal{F}$ に対して、関手的同型
$\Ext^{d - i}(\mathcal{F}, \omega_X) = \Hom_k(H^i(X, \mathcal{F}), k)$
が存在する。
\end{enumerate}
\end{lemma}

\begin{proof}
補題 \ref{duality-lemma-duality-proper-over-field} から、$\omega_X$ が双対化加群であることは明らかである
（双対化複体の最も左にある非零コホモロジー層だからである）。
$\omega_X^\bullet = \omega_X[d]$ である。また、$X$ は Cohen-Macaulay なので、
補題 \ref{duality-lemma-dualizing-module-CM-scheme} により $\omega_X$ は Cohen-Macaulay である。
ほかの主張は、これと補題 \ref{duality-lemma-duality-proper-over-field} の対応する主張を組み合わせれば従う。
\end{proof}

\begin{remark}
\label{duality-remark-rework-duality-locally-free-CM}
$X$ を体 $k$ 上固有な Cohen-Macaulay スキームで、次元 $d$ の等次元スキームとする。
$\omega_X^\bullet$ と $\omega_X$ は補題 \ref{duality-lemma-duality-proper-over-field} と同じものとする。
補題 \ref{duality-lemma-duality-proper-over-field-CM} により
$\omega_X^\bullet = \omega_X[d]$ である。
$t : H^d(X, \omega_X) \to k$ を注意 \ref{duality-remark-duality-proper-over-field} の写像とする。
$\mathcal{E}$ を有限局所自由 $\mathcal{O}_X$-加群とし、その双対を
$\mathcal{E}^\vee$ とする。このとき完全ペアリング
$$
H^i(X, \omega_X \otimes_{\mathcal{O}_X} \mathcal{E}^\vee)
\times
H^{d - i}(X, \mathcal{E})
\longrightarrow
k,\quad
(\xi, \eta) \longmapsto t(1 \otimes \epsilon)(\xi \cup \eta))
$$
を得る。ここで $\cup$ はカップ積であり、
$\epsilon : \mathcal{E}^\vee \otimes_{\mathcal{O}_X} \mathcal{E}
\to \mathcal{O}_X$ は評価写像である。
これは補題 \ref{duality-lemma-duality-proper-over-field-perfect} の特別な場合である。
\end{remark}

\noindent
双対化複体に対する整合性確認を一つ行う。

\begin{lemma}
\label{duality-lemma-sanity-check-duality}
$X$ を体 $k$ 上固有なスキームとする。$\omega_X^\bullet$ と
$\omega_X$ は補題 \ref{duality-lemma-duality-proper-over-field} と同じものとする。
\begin{enumerate}
\item ある体 $k'$ に対して $X \to \Spec(k)$ が
$X \to \Spec(k') \to \Spec(k)$ と分解されるならば、
$\omega_X^\bullet$ と $\omega_X$ は、射 $X \to \Spec(k')$ に対して
補題 \ref{duality-lemma-duality-proper-over-field} で与えられるものと同じである。
\item $K/k$ が体拡大ならば、$\omega_X^\bullet$ と $\omega_X$ の基底変換
$X_K$ への引き戻しは、射 $X_K \to \Spec(K)$ に対して
補題 \ref{duality-lemma-duality-proper-over-field} で与えられるものと同じである。
\end{enumerate}
\end{lemma}

\begin{proof}
構造射を $f : X \to \Spec(k)$ と書き、与えられた分解を
$f' : X \to \Spec(k')$ と書く。
補題 \ref{duality-lemma-duality-proper-over-field} の証明では
$\omega_X^\bullet = a(\mathcal{O}_{\Spec(k)})$ ととった。ここで $a$ は
$f$ に対する補題 \ref{duality-lemma-twisted-inverse-image} の右随伴である。
従って、$f'$ に対する補題 \ref{duality-lemma-twisted-inverse-image} の右随伴を $a'$ とするとき、
$a(\mathcal{O}_{\Spec(k)}) \cong a'(\mathcal{O}_{\Spec(k)})$
を示せばよい。
$k' \subset H^0(X, \mathcal{O}_X)$ なので、$k'/k$ は有限拡大である
（『スキームのコホモロジー』の補題 \ref{coherent-lemma-proper-over-affine-cohomology-finite}）。
随伴の一意性により $a = a' \circ b$ である。ここで $b$ は
$g : \Spec(k') \to \Spec(k)$ に対する補題 \ref{duality-lemma-twisted-inverse-image} の右随伴である。
別の言い方をすれば、$f^! = (f')^! \circ g^!$ である。
従って、例 \ref{duality-example-affine-twisted-inverse-image} により、
$k'$-加群として $\Hom_k(k', k) \cong k'$ を示せば十分である。
これは両辺が同じ次元、すなわち次元 $1$ の $k'$-ベクトル空間だから成り立つ。

\medskip\noindent
(2) の証明。補題 \ref{duality-lemma-more-base-change} により $a$ に対する基底変換があるので成り立つ。
注意 \ref{duality-remark-relative-dualizing-complex} の議論を参照する。
\end{proof}








\section{相対双対化複体}
\label{duality-section-relative-dualizing-complexes}

\noindent
固有・平坦・有限表示な射に対しては、剛体的な相対双対化複体がある。注意 \ref{duality-remark-relative-dualizing-complex} および補題 \ref{duality-lemma-van-den-bergh} を参照する。Noether スキームの分離的有限型射
$f : X \to Y$ に対しては、$f^!\mathcal{O}_Y$ を考えられる。
本節では、（必ずしも準分離でも準コンパクトでもない）平坦かつ局所有限表示な射に対する
相対双対化複体を、（必ずしも局所 Noether でない）スキームの間で定義する。
そのような複体が存在し、一意同型を除いて一意であり、上で述べた場合と一致することを示す。
本節を読む前に、『双対化複体』の第 \ref{dualizing-section-relative-dualizing-complexes} 節を読まれたい。

\begin{definition}
\label{duality-definition-relative-dualizing-complex}
$X \to S$ を平坦かつ局所有限表示なスキームの射とする。
$W \subset X \times_S X$ を、対角射 $\Delta_{X/S} : X \to X \times_S X$
が閉埋め込み $\Delta : X \to W$ を経由するような任意の開集合とする。
{\it 相対双対化複体}とは、対象 $K \in D(\mathcal{O}_X)$ と
$D(\mathcal{O}_W)$ における写像
$$
\xi : \Delta_*\mathcal{O}_X \longrightarrow L\text{pr}_1^*K|_W
$$
からなる組 $(K, \xi)$ で、次を満たすものをいう。
\begin{enumerate}
\item $K$ は $S$-完全である（『スキームの導来圏』の定義 \ref{perfect-definition-relatively-perfect}）。
\item $\xi$ は $\Delta_*\mathcal{O}_X$ と
$R\SheafHom_{\mathcal{O}_W}(
\Delta_*\mathcal{O}_X, L\text{pr}_1^*K|_W)$
との同型を定める。
\end{enumerate}
\end{definition}

\noindent
補題 \ref{duality-lemma-sheaf-with-exact-support-ext} により、条件 (2) は、
$D(\mathcal{O}_X)$ における同型
$$
\mathcal{O}_X \longrightarrow
R\SheafHom(\mathcal{O}_X, L\text{pr}_1^*K|_W)
$$
で、その $\Delta$ による押し出しが $\xi$ に等しいものの存在と同値である。
$R\SheafHom(\mathcal{O}_X, L\text{pr}_1^*K|_W)$ は開集合 $W$ の選択に依存しないので、
組 $(K, \xi)$ の圏もその選択に依存しない。
$X \to S$ が分離的ならば、$W = X \times_S X$ と選べる。
次の補題を用いて、多くの議論を環の場合に帰着する。

\begin{lemma}
\label{duality-lemma-relative-dualizing-complex-algebra}
$X \to S$ を平坦かつ局所有限表示なスキームの射とし、
$(K, \xi)$ を相対双対化複体とする。横の矢印が開埋め込みである任意の可換図式
$$
\xymatrix{
\Spec(A) \ar[d] \ar[r] & X \ar[d] \\
\Spec(R) \ar[r] & S
}
$$
に対し、$K$ の $\Spec(A)$ への制限は、『スキームの導来圏』の補題 \ref{perfect-lemma-affine-compare-bounded} を通じて、
『双対化複体』の定義 \ref{dualizing-definition-relative-dualizing-complex}
の意味での $R \to A$ に対する相対双対化複体に対応する。
\end{lemma}

\begin{proof}
$R\SheafHom$ の形成は開集合への制限と可換なので、
$X = \Spec(A)$ および $S = \Spec(R)$ と仮定してよい。
$D(\mathcal{O}_X)$ の相対完全対象は擬連接であり、従って
$D_\QCoh(\mathcal{O}_X)$ に属することに注意する
（『スキームの導来圏』の補題 \ref{perfect-lemma-pseudo-coherent}）。
従って主張には意味がある。
$\Delta_*$, $L\text{pr}_1^*$, $R\SheafHom$ をとる操作は、『スキームの導来圏』の補題 \ref{perfect-lemma-quasi-coherence-pushforward}、\ref{perfect-lemma-quasi-coherence-pullback} および \ref{perfect-lemma-quasi-coherence-internal-hom}
により代数側の操作と両立する。最後の補題については、
$L\text{pr}_1^*K$ が $S$-完全（従って下に有界）であり、
$\Delta_*\mathcal{O}_X$ が $D(\mathcal{O}_W)$ の擬連接対象であることに注意する。
代数に言い換えると、後者は $A$ が $A \otimes_R A$-加群として擬連接であることを意味する。
これは『代数の続論』の補題 \ref{more-algebra-lemma-more-relative-pseudo-coherent-is-moot}
を $R \to A \otimes_R A \to A$ に適用すれば従う。
従って『双対化複体』の定義 \ref{dualizing-definition-relative-dualizing-complex}
の条件をちょうど復元する。
\end{proof}

\begin{lemma}
\label{duality-lemma-relative-dualizing-RHom}
$X \to S$ を平坦かつ局所有限表示なスキームの射とし、
$(K, \xi)$ を相対双対化複体とする。このとき
$\mathcal{O}_X \to R\SheafHom_{\mathcal{O}_X}(K, K)$
は同型である。
\end{lemma}

\begin{proof}
アフィン局所的に見れば、補題 \ref{duality-lemma-relative-dualizing-complex-algebra} により代数の場合に帰着し、
これは『双対化複体』の補題 \ref{dualizing-lemma-relative-dualizing-RHom} である。
\end{proof}

\begin{lemma}
\label{duality-lemma-uniqueness-relative-dualizing}
$X \to S$ を平坦かつ局所有限表示なスキームの射とする。
$(K, \xi)$ と $(L, \eta)$ が $X/S$ 上の二つの相対双対化複体ならば、
$\xi$ を $\eta$ に写す一意な同型 $K \to L$ が存在する。
\end{lemma}

\begin{proof}
$U \subset X$ を、$S$ のあるアフィン開集合に写るアフィン開集合とする。
補題 \ref{duality-lemma-relative-dualizing-complex-algebra} および『双対化複体』の補題 \ref{dualizing-lemma-uniqueness-relative-dualizing}
により同型 $K|_U \to L|_U$ が存在する。
読者は同補題の議論をスキームの場合に再利用して、本場合の証明を得ることもできる。
ここでは代わりに貼り合わせの議論を用いる。

\medskip\noindent
同型 $\alpha : K \to L$ があるとする。このとき、ある可逆切断
$u \in H^0(W, \Delta_*\mathcal{O}_X) = H^0(X, \mathcal{O}_X)$
に対して $\alpha(\xi) = u \eta$ である
（仮定により、$\eta$ と $\alpha(\xi)$ はともに可逆な
$\Delta_*\mathcal{O}_X$-加群の生成元だからである）。
従って $\alpha$ を $u^{-1}\alpha$ で置き換えれば
$\alpha(\xi) = \eta$ となる。
補題 \ref{duality-lemma-relative-dualizing-RHom} により $K$ の自己同型群は
$H^0(X, \mathcal{O}_X^*)$ なので、そのような $\alpha$ は高々一つである。

\medskip\noindent
$\mathcal{B}$ を、$S$ のあるアフィン開集合に写る $X$ のアフィン開集合全体とする。
前二段落の議論により、各 $U \in \mathcal{B}$ に対して
$\xi$ を $\eta$ に写す一意な同型 $\alpha_U : K|_U \to L|_U$ がある。
補題 \ref{duality-lemma-relative-dualizing-RHom} により、$X$ の任意の開集合 $U$ と
$i < 0$ に対して $\text{Ext}^i(K|_U, K|_U) = 0$ であることに注意する。
『層のコホモロジー』の補題 \ref{cohomology-lemma-vanishing-and-glueing}
を $\text{id} : X \to X$ に適用すると、すべての $U \in \mathcal{B}$ に対して
$\alpha_U$ と一致する一意な射 $\alpha : K \to L$ を得る。
これは局所的に成り立つので、$\alpha$ は $\xi$ を $\eta$ に写す。
\end{proof}

\begin{lemma}
\label{duality-lemma-existence-relative-dualizing}
$X \to S$ を平坦かつ局所有限表示なスキームの射とする。
相対双対化複体 $(K, \xi)$ が存在する。
\end{lemma}

\begin{proof}
$\mathcal{B}$ を、$S$ のあるアフィン開集合に写る $X$ のアフィン開集合全体とする。
各 $U$ に対して、$S$ 上の $U$ の相対双対化複体 $(K_U, \xi_U)$ がある。
実際、アフィン開集合 $V \subset S$ を選び、$U \to X \to S$ が $V$ を経由するとする。
$U = \Spec(A)$ および $V = \Spec(R)$ と書く。
『双対化複体』の補題 \ref{dualizing-lemma-base-change-relative-dualizing} により、
$R \to A$ に対する相対双対化複体 $K_A \in D(A)$ が存在する。
補題 \ref{duality-lemma-relative-dualizing-complex-algebra} の証明を逆にたどると、
これは $V$-完全対象 $K_U \in D(\mathcal{O}_U)$ と、
$D(\mathcal{O}_{U \times_V U})$ における写像
$$
\xi : \Delta_*\mathcal{O}_U \to L\text{pr}_1^*K_U
$$
を定める。$V$-完全であることは $S$-完全であることと同じであり、
$U \times_V U = U \times_S U$ なので、$(K_U, \xi_U)$ は所望のものである。

\medskip\noindent
$U' \subset U \subset X$ かつ $U', U \in \mathcal{B}$ ならば、
補題 \ref{duality-lemma-uniqueness-relative-dualizing} により、
$D(\mathcal{O}_{U'})$ において
$\xi_U|_{U' \times_S U'}$ を $\xi_{U'}$ に写す一意な同型
$\rho_{U'}^U : K_U|_{U'} \to K_{U'}$ がある
（相対双対化複体の開集合への制限が相対双対化複体であることは明らかである）。
一意性により、$\mathcal{B}$ の $U'' \subset U' \subset U$ に対して
$\rho^U_{U''} = \rho^V_{U''} \circ \rho ^U_{U'}|_{U''}$
が保証される。
補題 \ref{duality-lemma-relative-dualizing-RHom} を $U/S$ と $K_U$ に適用すると、
$U \in \mathcal{B}$ および $i < 0$ に対して
$\text{Ext}^i(K_U, K_U) = 0$ である。
従って BBD 貼り合わせ補題
（『層のコホモロジー』の定理 \ref{cohomology-theorem-glueing-bbd-general}）により
一意な解、すなわち対象 $K \in D(\mathcal{O}_X)$ と同型
$\rho_U : K|_U \to K_U$ で、すべての $U' \subset U$,
$U, U' \in \mathcal{B}$ に対して
$\rho^U_{U'} \circ \rho_U|_{U'} = \rho_{U'}$ を満たすものが存在する。

\medskip\noindent
証明を完結するには、$D(\mathcal{O}_W)$ において
$\Delta_*\mathcal{O}_X$ から
$R\SheafHom_{\mathcal{O}_W}(\Delta_*\mathcal{O}_X, L\text{pr}_1^*K|_W)$
への同型を誘導する写像
$$
\xi : \Delta_*\mathcal{O}_X \longrightarrow L\text{pr}_1^*K|_W
$$
を構成しなければならない。$W$ は取り替えてよいので、
$W = \bigcup_{U \in \mathcal{B}} U \times_S U$ と選ぶ。
$\rho_U$ を用いて同型
$$
R\SheafHom_{\mathcal{O}_W}(
\Delta_*\mathcal{O}_X, L\text{pr}_1^*K|_W)|_{U \times_S U}
\xrightarrow{\rho_U}
R\SheafHom_{\mathcal{O}_{U \times_S U}}(
\Delta_*\mathcal{O}_U, L\text{pr}_1^*K_U)
$$
を得る。$W$ は開集合 $U \times_S U$ で被覆されるので、
$R\SheafHom_{\mathcal{O}_W}(\Delta_*\mathcal{O}_X, L\text{pr}_1^*K|_W)$
のコホモロジー層は次数 $0$ を除いて零である。さらに同型
$$
H^0\left(U \times_S U, R\SheafHom_{\mathcal{O}_W}(\Delta_*\mathcal{O}_X,
L\text{pr}_1^*K|_W)\right)
\xrightarrow{\rho_U}
H^0\left((R\SheafHom_{\mathcal{O}_{U \times_S U}}(
\Delta_*\mathcal{O}_U, L\text{pr}_1^*K_U)\right)
$$
を得る。左辺の元で、この写像により $\xi_U$ に写るものを $\tau_U$ とする。
開集合 $U' \subset U \subset X$、$U', U \in \mathcal{B}$ に対する
$\rho^U_{U'}$, $\xi_U$, $\xi_{U'}$, $\rho_U$, $\rho_{U'}$ の両立性により、
$\tau_U|_{U' \times_S U'} = \tau_{U'}$ である。
従って第 $0$ コホモロジー層
$H^0(R\SheafHom_{\mathcal{O}_W}(\Delta_*\mathcal{O}_X, L\text{pr}_1^*K|_W))$
の大域切断 $\tau$ を得る。
$R\SheafHom_{\mathcal{O}_W}(\Delta_*\mathcal{O}_X, L\text{pr}_1^*K|_W)$
のほかのコホモロジー層は零なので、この大域切断 $\tau$ は所望の射 $\xi$ を定める。
$\xi$ の $U \times_S U$ への制限は $\xi_U$ を与えるので、
定義 \ref{duality-definition-relative-dualizing-complex} の最後の条件を満たす。
\end{proof}

\begin{lemma}
\label{duality-lemma-base-change-relative-dualizing}
スキームの Cartesian 正方形
$$
\xymatrix{
X' \ar[d]_{f'} \ar[r]_{g'} & X \ar[d]^f \\
S' \ar[r]^g & S
}
$$
を考える。$X \to S$ は平坦かつ局所有限表示であると仮定する。
$(K, \xi)$ を $f$ に対する相対双対化複体とし、
$K' = L(g')^*K$ とおく。$\xi'$ を $\xi$ の導来基底変換とする
（証明を参照する）。このとき $(K', \xi')$ は $f'$ に対する相対双対化複体である。
\end{lemma}

\begin{proof}
Cartesian 正方形
$$
\xymatrix{
X' \ar[d]_{\Delta_{X'/S'}} \ar[r] & X \ar[d]^{\Delta_{X/S}} \\
X' \times_{S'} X' \ar[r]^{g' \times g'} & X \times_S X
}
$$
を考える。$\Delta_{X/S}$ が閉埋め込み $\Delta : X \to W$ を経由するような開集合
$W \subset X \times_S X$ を選ぶ。
$\Delta_{X'/S'}$ が閉埋め込み $\Delta' : X \to W'$ を経由し、かつ
$(g' \times g')(W') \subset W$ となるような開集合
$W' \subset X' \times_{S'} X'$ を選ぶ。誘導される射も引き続き
$g' \times g' : W' \to W$ と書く。このとき
$$
L(g' \times g')^*\Delta_*\mathcal{O}_X =
\Delta'_*\mathcal{O}_{X'}
\quad\text{and}\quad
L(g' \times g')^*L\text{pr}_1^*K|_W =
L\text{pr}_1^*K'|_{W'}
$$
を得る。第一の等号は、$X$ と $X' \times_{S'} X'$ が
$X \times_S X$ 上 Tor 独立だから成り立つ
（例えば『射の続論』の補題 \ref{more-morphisms-lemma-case-of-tor-independence} を参照する）。
第二の等号は導来引き戻しの推移性による
（『層のコホモロジー』の補題 \ref{cohomology-lemma-derived-pullback-composition}）。
従って $\xi' = L(g' \times g')^*\xi$ は写像
$$
\xi' : \Delta'_*\mathcal{O}_{X'} \longrightarrow L\text{pr}_1^*K'|_{W'}
$$
とみなせる。以上を踏まえれば補題の証明は直接的である。
まず『スキームの導来圏』の補題 \ref{perfect-lemma-base-change-relatively-perfect} により
$K'$ は $S'$-完全である。
$\xi'$ が $\Delta'_*\mathcal{O}_{X'}$ から
$R\SheafHom_{\mathcal{O}_{W'}}(
\Delta'_*\mathcal{O}_{X'}, L\text{pr}_1^*K'|_{W'})$
への同型を誘導することは、アフィン局所的に確認できる。
補題 \ref{duality-lemma-relative-dualizing-complex-algebra} により対応する代数的主張に帰着し、
これは『双対化複体』の補題 \ref{dualizing-lemma-base-change-relative-dualizing} で証明されている。
\end{proof}

\begin{lemma}
\label{duality-lemma-flat-proper-relative-dualizing}
$S$ を準コンパクトかつ準分離なスキームとする。
$f : X \to S$ を固有・平坦・有限表示な射とする。
注意 \ref{duality-remark-relative-dualizing-complex} の相対双対化複体
$\omega_{X/S}^\bullet$ は、(\ref{duality-equation-pre-rigid}) とともに
定義 \ref{duality-definition-relative-dualizing-complex}
の意味での相対双対化複体である。
\end{lemma}

\begin{proof}
補題 \ref{duality-lemma-properties-relative-dualizing} において、
$\omega_{X/S}^\bullet$ が $S$-完全であることを示した。
対角射 $\Delta : X \to X \times_S X$ に対する補題 \ref{duality-lemma-twisted-inverse-image} の右随伴を $c$ とする。
そこで (\ref{duality-equation-pre-rigid}) に $\Delta_*$ を適用すると同型
$$
\Delta_*\mathcal{O}_X \to
\Delta_*(c(L\text{pr}_1^*\omega_{X/S}^\bullet)) =
R\SheafHom_{\mathcal{O}_{X \times_S X}}(
\Delta_*\mathcal{O}_X, L\text{pr}_1^*\omega_{X/S}^\bullet)
$$
を得る。等号は補題 \ref{duality-lemma-twisted-inverse-image-closed} および \ref{duality-lemma-sheaf-with-exact-support-ext} による。これで証明が完了する。
\end{proof}

\begin{remark}
\label{duality-remark-relative-dualizing-complex-bis}
$X \to S$ を平坦・固有・有限表示なスキームの射とする。
補題 \ref{duality-lemma-existence-relative-dualizing} により、定義 \ref{duality-definition-relative-dualizing-complex} の意味での相対双対化複体
$(\omega_{X/S}^\bullet, \xi)$ が存在する。
$S'$ が準コンパクトかつ準分離である任意の射 $g : S' \to S$
（例えば $S$ のアフィン開集合）を考える。
補題 \ref{duality-lemma-base-change-relative-dualizing} により、
$(L(g')^*\omega_{X/S}^\bullet, L(g')^*\xi)$ は基底変換
$f' : X' \to S'$ に対する定義 \ref{duality-definition-relative-dualizing-complex} の意味での相対双対化複体である。
$\omega_{X'/S'}^\bullet$ を注意 \ref{duality-remark-relative-dualizing-complex} の意味での
$X' \to S'$ に対する相対双対化複体とする。
補題 \ref{duality-lemma-flat-proper-relative-dualizing} および \ref{duality-lemma-uniqueness-relative-dualizing} を組み合わせると、
(\ref{duality-equation-pre-rigid}) および $L(g')^*\xi$ と両立する一意な同型
$$
\omega_{X'/S'}^\bullet \longrightarrow L(g')^*\omega_{X/S}^\bullet
$$
が存在する。これらの同型は、$S$ 上の準コンパクトかつ準分離なスキームの間の射、
および補題 \ref{duality-lemma-proper-flat-base-change} の基底変換同型と両立する
（この両立性が必要になれば、ここで注意深く定式化して証明する）。
\end{remark}

\begin{lemma}
\label{duality-lemma-compactifyable-relative-dualizing}
状況 \ref{duality-situation-shriek} において、$f : X \to Y$ を
$\textit{FTS}_S$ の射とする。$f$ が平坦ならば、
$f^!\mathcal{O}_Y$ は定義 \ref{duality-definition-relative-dualizing-complex} の意味での
$Y$ 上の $X$ の相対双対化複体（の第一成分）である。
\end{lemma}

\begin{proof}
補題 \ref{duality-lemma-flat-shriek-relatively-perfect} により
$f^!\mathcal{O}_Y$ は $Y$-完全である。
$f$ は分離的なので、対角射 $\Delta : X \to X \times_Y X$ は閉埋め込みであり、
補題 \ref{duality-lemma-twisted-inverse-image-closed} および \ref{duality-lemma-sheaf-with-exact-support-ext} により
$\Delta_*\Delta^!(-) =
R\SheafHom_{\mathcal{O}_{X \times_Y X}}(\mathcal{O}_X, -)$
である。従って証明を完結するには
$\Delta^!(L\text{pr}_1^*f^!(\mathcal{O}_Y)) \cong \mathcal{O}_X$
を示せば十分である。ここで $\text{pr}_1 : X \times_Y X \to X$
は第一射影である。実際、
$$
\mathcal{O}_X = \Delta^! \text{pr}_1^!\mathcal{O}_X =
\Delta^! \text{pr}_1^! L\text{pr}_2^*\mathcal{O}_Y =
\Delta^!(L\text{pr}_1^* f^!\mathcal{O}_Y)
$$
である。ここで $\text{pr}_2 : X \times_Y X \to X$ は第二射影であり、
補題 \ref{duality-lemma-base-change-shriek-flat} の基底変換同型
$\text{pr}_1^! \circ L\text{pr}_2^* = L\text{pr}_1^* \circ f^!$
を用いた。
\end{proof}

\begin{lemma}
\label{duality-lemma-relative-dualizing-composition}
$f : Y \to X$ および $X \to S$ を平坦かつ有限表示なスキームの射とする。
$(K, \xi)$ と $(M, \eta)$ を、それぞれ $X \to S$ と $Y \to X$
に対する相対双対化複体とする。
$E = M \otimes_{\mathcal{O}_Y}^\mathbf{L} Lf^*K$ とおく。
このとき適当な $\zeta$ に対して、$(E, \zeta)$ は $Y \to S$
に対する相対双対化複体である。
\end{lemma}

\begin{proof}
補題 \ref{duality-lemma-relative-dualizing-complex-algebra} と本補題の代数版
（『双対化複体』の補題 \ref{dualizing-lemma-relative-dualizing-composition}）を用いると、
$E$ はアフィン局所的に相対双対化複体の第一成分であることが分かる。
特に、これはアフィン局所的に確認できるので $E$ は $S$-完全である。
『スキームの導来圏』の補題 \ref{perfect-lemma-affine-locally-rel-perfect} を参照する。

\medskip\noindent
まず、射 $X \to S$ と $Y \to X$ が分離的であり、
$\Delta_{X/S}$, $\Delta_{Y/X}$, $\Delta_{Y/S}$ が閉埋め込みである場合に
$\zeta$ の存在を示す。次の図式を考える。
$$
\xymatrix{
& & Y \ar@{=}[r] & Y \ar[d]^f \\
Y \ar[r]_{\Delta_{Y/X}} &
Y \times_X Y \ar[d]_m \ar[r]_\delta \ar[ru]_q &
Y \times_S Y \ar[d]^{f \times f} \ar[ru]_p & X\\
& X \ar[r]^{\Delta_{X/S}} & X \times_S X \ar[ru]_r
}
$$
ここで $p$, $q$, $r$ は第一射影である。
補題 \ref{duality-lemma-sheaf-with-exact-support-internal-home} により
$$
R\SheafHom_{\mathcal{O}_{Y \times_S Y}}(
\Delta_{Y/S, *}\mathcal{O}_Y, Lp^*E) =
R\delta_*\left(R\SheafHom_{\mathcal{O}_{Y \times_X Y}}(
\Delta_{Y/X, *}\mathcal{O}_Y,
R\SheafHom(\mathcal{O}_{Y \times_X Y}, Lp^*E))\right)
$$
である。補題 \ref{duality-lemma-sheaf-with-exact-support-tensor} により
$$
R\SheafHom(\mathcal{O}_{Y \times_X Y}, Lp^*E) =
R\SheafHom(\mathcal{O}_{Y \times_X Y}, L(f \times f)^*Lr^*K)
\otimes_{\mathcal{O}_{Y \times_S Y}}^\mathbf{L} Lq^*M
$$
である。補題 \ref{duality-lemma-flat-bc-sheaf-with-exact-support} により
$$
R\SheafHom(\mathcal{O}_{Y \times_X Y}, L(f \times f)^*Lr^*K) =
Lm^*R\SheafHom(\mathcal{O}_X, Lr^*K)
$$
である。最後の式は $\xi$ を通じて
$Lm^*\mathcal{O}_X = \mathcal{O}_{Y \times_X Y}$
に同型である。従ってその直前の式は $Lq^*M$ に同型である。ゆえに
$$
R\SheafHom_{\mathcal{O}_{Y \times_S Y}}(
\Delta_{Y/S, *}\mathcal{O}_Y, Lp^*E) =
R\delta_*\left(R\SheafHom_{\mathcal{O}_{Y \times_X Y}}(
\Delta_{Y/X, *}\mathcal{O}_Y, Lq^*M)\right)
$$
である。括弧内の対象は $\eta$ を通じて
$\Delta_{Y/X, *}*\mathcal{O}_X$ に同型である。
これで分離的な場合の証明が完了する。

\medskip\noindent
一般の場合、$\Delta_{X/S}$ が閉埋め込み $\Delta : X \to W$
を経由するような開集合 $W \subset X \times_S X$ と、
$\Delta_{Y/X}$ が閉埋め込み $\Delta' : Y \to V$
を経由するような開集合 $V \subset Y \times_X Y$ を選ぶ。
最後に、$Y \times_X Y$ との交わりが $V$ であり、$W$ に写る開集合
$W' \subset Y \times_S Y$ を選ぶ。このとき図式
$$
\xymatrix{
& & Y \ar@{=}[r] & Y \ar[d]^f \\
Y \ar[r]_{\Delta'} &
V \ar[d]_m \ar[r]_\delta \ar[ru]_q &
W' \ar[d]^{f \times f} \ar[ru]_p & X\\
& X \ar[r]^\Delta & W \ar[ru]_r
}
$$
を考え、先ほどとまったく同じ議論を用いる。
\end{proof}






\section{lci 射の基本類}
\label{duality-section-fundamental-class}

\noindent
本節では第 \ref{duality-section-examples} 節で行った計算を用いる。
従って、本節の結果は同節と同種の非一意性をもつ。

\begin{lemma}
\label{duality-lemma-determinant}
$X$ を局所環付き空間とし、
$$
\mathcal{E}_1 \xrightarrow{\alpha} \mathcal{E}_0 \to \mathcal{F} \to 0
$$
を $\mathcal{O}_X$-加群の短完全列とする。
$\mathcal{E}_1$ と $\mathcal{E}_0$ はそれぞれ階数 $r_1, r_0$
の局所自由加群であると仮定する。このとき標準写像
$$
\wedge^{r_0 - r_1}\mathcal{F}
\longrightarrow
\wedge^{r_1}(\mathcal{E}_1^\vee) \otimes \wedge^{r_0}\mathcal{E}_0
$$
が存在し、これが $x \in X$ における茎上の同型であることと、
$\mathcal{F}$ が $x$ のある開近傍で階数 $r_0 - r_1$ の局所自由加群であることは同値である。
\end{lemma}

\begin{proof}
$r_1 > r_0$ ならば、規約により
$\wedge^{r_0 - r_1}\mathcal{F} = 0$ であり、一意な写像は同型にはなれない。
従って $r = r_0 - r_1 \geq 0$ と仮定してよい。写像を公式
$$
s_1 \wedge \ldots \wedge s_r \mapsto
t_1^\vee \wedge \ldots \wedge t_{r_1}^\vee \otimes
\alpha(t_1) \wedge \ldots \wedge \alpha(t_{r_1}) \wedge
\tilde s_1 \wedge \ldots \wedge \tilde s_r
$$
で定める。ここで $t_1, \ldots, t_{r_1}$ は $\mathcal{E}_1$ の局所基底、
$t_1^\vee, \ldots, t_{r_1}^\vee$ は対応する $\mathcal{E}_1^\vee$ の双対基底であり、
$s'_i$ は $s_i$ の $\mathcal{E}_0$ の切断への局所持ち上げである。
これが良定義であることの証明は省略する。

\medskip\noindent
$\mathcal{F}$ が階数 $r$ の局所自由加群ならば、写像が同型であることは直接確認できる。
逆に、写像が $x$ における茎上の同型であると仮定する。このとき
$\wedge^r\mathcal{F}_x$ は可逆である。従って $\mathcal{F}_x$ は高々
$r$ 個の元で生成される。これは $\alpha$ が $\mathfrak m_x$ を法として階数
$r$ をもつ場合に限り起こり得る。すなわち、$\alpha$ は $\mathfrak m_x$ を法として最大階数をもつ。
従って $\alpha$ は $x$ のある近傍で最大階数をもち、
$\mathcal{F}$ は所望のとおりその近傍で階数 $r$ の局所自由加群である。
\end{proof}

\begin{lemma}
\label{duality-lemma-fundamental-class-lci}
$Y$ を Noether スキームとする。$f : X \to Y$ を、
埋め込み $X \to P$ と固有滑らかな射 $P \to Y$ の合成として分解される局所完全交叉射とする。
$r$ を $X$ 上の局所定数関数で、
$\omega_{X/Y} = H^{-r}(f^!\mathcal{O}_Y)$ が
$f^!\mathcal{O}_Y$ の一意な非零コホモロジー層となるものとする。
補題 \ref{duality-lemma-lci-shriek} を参照する。このとき写像
$$
\wedge^r\Omega_{X/Y} \longrightarrow \omega_{X/Y}
$$
が存在し、これが点 $x$ における茎上の同型であることと、
$f$ が $x$ において滑らかであることは同値である。
\end{lemma}

\begin{proof}
仮定により $X$ は $Y$ 上コンパクト化可能なので $f^!$ が定義される。
第 \ref{duality-section-upper-shriek} 節を参照する。
$j : W \to P$ を、$X \to P$ が閉埋め込み $i : X \to W$ を経由するような開部分スキームとする。
また、$g : P \to Y$ を与えられた射とすると
$f^! = i^! \circ j^! \circ g^!$ である。
補題 \ref{duality-lemma-smooth-proper} により
$g^!\mathcal{O}_Y = \wedge^d\Omega_{P/Y}[d]$ である。
ここで $d$ は $Y$ 上の $P$ の相対次元を与える局所定数関数である。
$j^! = j^*$ である。また
$i^!\mathcal{O}_W = \wedge^c\mathcal{N}[-c]$ である。
ここで $c$ は $W$ における $X$ の余次元
（$X$ 上の局所定数関数）であり、$\mathcal{N}$ は Koszul 正則埋め込み
$i$ の法層である。補題 \ref{duality-lemma-regular-immersion} を参照する。
以上を組み合わせると
$$
f^!\mathcal{O}_Y =
\left(\wedge^c\mathcal{N} \otimes_{\mathcal{O}_X}
\wedge^d\Omega_{P/Y}|_X\right)[d - c]
$$
を得る。ここでは補題 \ref{duality-lemma-perfect-comparison-shriek} も用いた。
従って $X$ 上の局所定数関数として $r = d|_X - c$ である。
$X \to P$ の余法層は加群 $\mathcal{I}/\mathcal{I}^2$ である。
ここで $\mathcal{I} \subset \mathcal{O}_W$ は $i$ のイデアル層である。
『スキームの射』の第 \ref{morphisms-section-conormal-sheaf} 節を参照する。
『スキームの射』の補題 \ref{morphisms-lemma-differentials-relative-immersion}
の標準完全列
$$
\mathcal{I}/\mathcal{I}^2 \to
\Omega_{P/Y}|_X \to \Omega_{X/Y} \to 0
$$
を考える。補題 \ref{duality-lemma-determinant} を適用して所望の写像を得る。

\medskip\noindent
$f$ が $x$ において滑らかならば、$\Omega_{X/Y}$ は $x$ において階数
$r$ の局所自由加群なので、補題 \ref{duality-lemma-determinant} により写像は同型である。
逆に、写像が $x$ における茎上の同型であると仮定する。
補題により、$X$ を $x$ の開近傍で置き換えた後
$\Omega_{X/Y}$ は階数 $r$ の自由加群である。
一方、$X = \Spec(A)$, $Y = \Spec(R)$ であり、
$A = R[x_1, \ldots, x_n]/(f_1, \ldots, f_m)$、
$f_1, \ldots, f_m$ は Koszul 正則列であると仮定してもよい
（局所完全交叉射の定義から従う）。
明らかに $r = n - m$ である。従って偏微分行列
$\partial f_j/\partial x_i$ は $x$ における剰余体上で階数 $m$ をもつ。
変数を並べ替えることにより、
$(\partial f_j/\partial x_i)_{1 \leq i, j \leq m}$
の行列式は $x$ のある開近傍で可逆であると仮定してよい。
従って $R \to A$ はこの点で滑らかである。例えば『可換代数』の例 \ref{algebra-example-make-standard-smooth} を参照する。
\end{proof}

\begin{lemma}
\label{duality-lemma-fundamental-class-almost-lci}
$f : X \to Y$ をスキームの射とし、$r \geq 0$ とする。次を仮定する。
\begin{enumerate}
\item $Y$ は Cohen-Macaulay である（『スキームの性質』の定義 \ref{properties-definition-Cohen-Macaulay}）。
\item $f$ は $X \to P \to Y$ と分解され、第一の射は埋め込み、第二の射は滑らかかつ固有である。
\item $x \in X$ かつ $\dim(\mathcal{O}_{X, x}) \leq 1$ ならば、
$f$ は $x$ において Koszul である（『射の続論』の定義 \ref{more-morphisms-definition-lci}）。
\item $\xi$ が $X$ の既約成分の生成点ならば
$\text{trdeg}_{\kappa(f(\xi))} \kappa(\xi) = r$ である。
\end{enumerate}
このとき $\omega_{X/Y} = H^{-r}(f^!\mathcal{O}_Y)$ とおけば写像
$$
\wedge^r\Omega_{X/Y} \longrightarrow \omega_{X/Y}
$$
が存在し、$f$ が滑らかな軌跡上で同型である。
\end{lemma}

\begin{proof}
$U \subset X$ を、$f$ が局所完全交叉射となる開部分スキームとする。
仮定 (4) により $f$ はすべての生成点で相対次元 $r$ をもつので、
補題 \ref{duality-lemma-fundamental-class-lci} の局所定数関数は値 $r$ の定数関数であり、写像
$$
\wedge^r\Omega_{X/Y}|_U = \wedge^r \Omega_{U/Y}
\longrightarrow
\omega_{U/Y} = \omega_{X/Y}|_U
$$
を得る。これは $f$ の滑らかな点で同型である
（滑らかな射は局所完全交叉射なので、この軌跡は $U$ に含まれる）。
補題 \ref{duality-lemma-shriek-over-CM} と $Y$ が Cohen-Macaulay であるという仮定により、
加群 $\omega_{X/Y}$ は $(S_2)$ である。
条件 (3) により $U$ は余次元 $1$ のすべての点を含むので、
『因子』の補題 \ref{divisors-lemma-depth-2-hartog} を用いると
$j_*\omega_{U/Y} = \omega_{X/Y}$ である。
従って $U$ 上の写像は $X$ に延長し、証明が完了する。
\end{proof}






\section{連接加群の零延長}
\label{duality-section-extension-by-zero}

\noindent
本節およびこれに続くいくつかの節の内容は、\cite{RD} の Deligne による付録に見いだせる。

\medskip\noindent
本節では $j : U \to X$ を Noether スキームの開埋め込みとする。
以下のように構成される
$D^b_{\textit{Coh}}(\mathcal{O}_X)$ 内の逆系 $(K_n)$ を考える。
$\mathcal{F}^\bullet$ を連接
$\mathcal{O}_X$-加群の有界複体とする。$\mathcal{I} \subset \mathcal{O}_X$
を $V(\mathcal{I}) = X \setminus U$ を満たす準連接イデアル層とする。
このとき
$$
K_n = \mathcal{I}^n\mathcal{F}^\bullet
$$
と置ける。より正確には、$K_n$ は $D^b_{\textit{Coh}}(\mathcal{O}_X)$
の対象であって、次数 $q$ の項が $\mathcal{F}^q$ の連接部分加群
$\mathcal{I}^n\mathcal{F}^q$ である複体によって表される。
写像 $\ldots \to K_3 \to K_2 \to K_1$ は $U$ への制限上で同型を誘導する。
このような系を {\it Deligne 系} と呼ぶことにする。

\begin{lemma}
\label{duality-lemma-lift-map}
$j : U \to X$ を Noether スキームの開埋め込みとする。
$(K_n)$ を Deligne 系とし、定値系 $(K_n|_U)$ の値を
$K \in D^b_{\textit{Coh}}(\mathcal{O}_U)$ と書く。
$L$ を $D^b_{\textit{Coh}}(\mathcal{O}_X)$ の対象とする。
このとき $\colim \Hom_X(K_n, L) = \Hom_U(K, L|_U)$ である。
\end{lemma}

\begin{proof}
$L \to M \to N \to L[1]$ を
$D^b_{\textit{Coh}}(\mathcal{O}_X)$ の完全三角とする。このとき可換図式
$$
\xymatrix{
\ldots \ar[r] &
\colim \Hom_X(K_n, L) \ar[r] \ar[d] &
\colim \Hom_X(K_n, M) \ar[r] \ar[d] &
\colim \Hom_X(K_n, N) \ar[r] \ar[d] &
\ldots \\
\ldots \ar[r] &
\Hom_U(K, L|_U) \ar[r] &
\Hom_U(K, M|_U) \ar[r] &
\Hom_U(K, N|_U) \ar[r] &
\ldots
}
$$
を得る。その各行は『導来圏』の補題 \ref{derived-lemma-representable-homological} および
『可換代数』の補題 \ref{algebra-lemma-directed-colimit-exact}
により完全である。従って、補題の主張が
$N[-1]$, $L$, $N$, $L[1]$ に対して成り立てば、5 項補題により $M$ に対しても成り立つ。
そこで $L$ の標準切断に対する完全三角
（『導来圏』の注意 \ref{derived-remark-truncation-distinguished-triangle} を参照）を用いることにより、
$L$ がただ一つの非零コホモロジー層しかもたない場合に帰着する。

\medskip\noindent
連接 $\mathcal{O}_X$-加群の有界複体 $\mathcal{F}^\bullet$ と
$X \setminus U$ を切り出す準連接イデアル
$\mathcal{I} \subset \mathcal{O}_X$ を、
$K_n$ が $\mathcal{I}^n\mathcal{F}^\bullet$ により表されるように選ぶ。
「愚直」切断を用いると、各項で分裂する両立した複体の短完全列
$$
0 \to \sigma_{\geq a + 1} \mathcal{I}^n\mathcal{F}^\bullet \to
\mathcal{I}^n\mathcal{F}^\bullet \to
\sigma_{\leq a} \mathcal{I}^n\mathcal{F}^\bullet \to 0
$$
を得る。これらはさらに
$D^b_{\textit{Coh}}(\mathcal{O}_X)$ における完全三角の両立する系に対応する。
上と同様に論じることにより、$\mathcal{F}^\bullet$ がただ一つの非零項しかもたない場合に帰着する。
従って次の段落で扱う場合に帰着する。

\medskip\noindent
連接 $\mathcal{O}_X$-加群 $\mathcal{F}$ と連接
$\mathcal{O}_X$-加群 $\mathcal{G}$ が与えられたとき、標準写像
$$
\colim \Ext^i_X(\mathcal{I}^n\mathcal{F}, \mathcal{G})
\longrightarrow
\Ext^i_U(\mathcal{F}|_U, \mathcal{G}|_U)
$$
がすべての $i \geq 0$ に対して同型であることを示せばよい。
$i = 0$ の場合、これは『スキームのコホモロジー』の補題 \ref{coherent-lemma-homs-over-open} である。$i > 0$ と仮定する。

\medskip\noindent
単射性。$\xi \in \Ext^i_X(\mathcal{I}^n\mathcal{F}, \mathcal{G})$
を、その $U$ への制限が零である元とする。
ある $m \geq n$ が存在して、$\xi$ の
$\mathcal{I}^m\mathcal{F} = \mathcal{I}^{m - n}\mathcal{I}^n\mathcal{F}$
への制限が零になることを示す必要がある。
$\mathcal{F}$ を $\mathcal{I}^n\mathcal{F}$ で置き換えれば、
$n = 0$、すなわち
$\xi \in \Ext^i_X(\mathcal{F}, \mathcal{G})$
であってその $U$ への制限が零であると仮定してよい。
『スキームの導来圏』の命題 \ref{perfect-proposition-DCoh}
により
$D^b_{\textit{Coh}}(\mathcal{O}_X) = D^b(\textit{Coh}(\mathcal{O}_X))$
である。従って $\Ext$ 群は連接
$\mathcal{O}_X$-加群の Abel 圏で計算できる。
このことから、$\xi \circ \alpha = 0$ を満たす全射
$\alpha : \mathcal{F}'' \to \mathcal{F}$ が存在する
（ここで $i > 0$ を用いる）。
$\mathcal{F}' = \Ker(\alpha)$ と置けば、短完全列
$$
0 \to \mathcal{F}' \to \mathcal{F}'' \to \mathcal{F} \to 0
$$
を得る。従って $\xi$ はある元
$\xi' \in \Ext^{i - 1}_X(\mathcal{F}', \mathcal{G})$
の像であり、その $U$ への制限は
$\Ext^{i - 1}_U(\mathcal{F}''|_U, \mathcal{G}|_U) \to
\Ext^{i - 1}_U(\mathcal{F}'|_U, \mathcal{G}|_U)$
の像に属する。
Artin-Rees により逆系 $(\mathcal{I}^n\mathcal{F}')$ と
$(\mathcal{I}^n \mathcal{F}'' \cap \mathcal{F}')$ は pro-同型である。
『スキームのコホモロジー』の補題 \ref{coherent-lemma-Artin-Rees} を参照する。
両立する短完全列の系
$$
0 \to
\mathcal{F}' \cap \mathcal{I}^n\mathcal{F}'' \to
\mathcal{I}^n\mathcal{F}'' \to
\mathcal{I}^n\mathcal{F} \to 0
$$
があるので、可換図式
$$
\xymatrix{
\colim \Ext^{i - 1}_X(\mathcal{I}^n\mathcal{F}'', \mathcal{G})
\ar[r] \ar[d] &
\colim \Ext^{i - 1}_X(\mathcal{F}' \cap \mathcal{I}^n\mathcal{F}'', \mathcal{G})
\ar[r] \ar[d] &
\colim \Ext^i_X(\mathcal{I}^n\mathcal{F}, \mathcal{G})
\ar[d] \\
\Ext^{i - 1}_U(\mathcal{F}''|_U, \mathcal{G}|_U) \ar[r] &
\Ext^{i - 1}_U(\mathcal{F}'|_U, \mathcal{G}|_U) \ar[r] &
\Ext^{i - 1}_U(\mathcal{F}|_U, \mathcal{G}|_U)
}
$$
を得る。その各行は完全である。$i$ に関する帰納法と上記の逆系についての注意により、
左の二つの垂直矢印は同型である。
いま $\xi$ は右上の群の元を与え、それは中央上の群における $\xi'$ の像である。
この元はさらに、左下の群のある元から来る右下の群の元へ写る。
従って、ある $n$ に対して $\xi$ が
$\Ext^i_X(\mathcal{I}^n\mathcal{F}, \mathcal{G})$
において零へ写ることが分かり、所望の結論を得る。

\medskip\noindent
全射性。$\xi \in \Ext^i_U(\mathcal{F}|_U, \mathcal{G}|_U)$ とする。
上と同様に、$i > 0$ を用いて連接 $\mathcal{O}_U$-加群の全射
$\mathcal{H} \to \mathcal{F}|_U$ を、
$\xi$ が $\Ext^i_U(\mathcal{H}, \mathcal{G}|_U)$ において零へ写るように選べる。
すると、$U$ への制限が
$\mathcal{H} \to \mathcal{F}|_U$ である連接
$\mathcal{O}_X$-加群の写像 $\varphi : \mathcal{F}'' \to \mathcal{F}$
を見いだせる。『スキームの性質』の補題 \ref{properties-lemma-extend-finite-presentation} を参照する。
この補題は $\varphi$ が全射であることを保証しないが、これは問題にならない
（少し余分に作業すれば全射な $\varphi$ を選ぶこともできる）。
$\mathcal{F}' = \Ker(\varphi)$ と書く。短完全列
$$
0 \to \mathcal{F}'|_U \to \mathcal{F}''|_U \to \mathcal{F}|_U \to 0
$$
より、$\xi$ は
$\Ext^{i - 1}_U(\mathcal{F}'|_U, \mathcal{G}|_U)$
における $\xi'$ の像である。
$i$ に関する帰納法により、ある $n$ が存在して、$\xi'$ はある $\xi'_n$ の
$\Ext^{i - 1}_X(\mathcal{I}^n\mathcal{F}', \mathcal{G})$
における像である。
Artin-Rees により、ある $m \geq n$ を
$\mathcal{F}' \cap \mathcal{I}^m\mathcal{F}'' \subset
\mathcal{I}^n\mathcal{F}'$ となるように選べる。短完全列
$$
0 \to \mathcal{F}' \cap \mathcal{I}^m\mathcal{F}'' \to
\mathcal{I}^m\mathcal{F}'' \to \mathcal{I}^m\Im(\varphi) \to 0
$$
を用いると、$\xi'_n$ の
$\Ext^{i - 1}_X(\mathcal{F}' \cap \mathcal{I}^m\mathcal{F}'', \mathcal{G})$
における像は、境界写像により
$\Ext^i_X(\mathcal{I}^m\Im(\varphi), \mathcal{G})$ の元 $\xi_m$
へ写り、この元は $\xi$ へ写る。
$\Im(\varphi)$ と $\mathcal{F}$ は $U$ 上で一致するので、
$\mathcal{F}/\mathcal{I}^m\Im(\varphi)$ は $X \setminus U$ 上に台をもつ。
従って、ある $l \geq m$ が存在して
$\mathcal{I}^l\mathcal{F} \subset \mathcal{I}^m\Im(\varphi)$ となる。
『スキームのコホモロジー』の補題 \ref{coherent-lemma-power-ideal-kills-sheaf} を参照する。
$\xi_m$ の
$\Ext^i_X(\mathcal{I}^l\mathcal{F}, \mathcal{G})$ における像を取ればよい。
\end{proof}

\begin{lemma}
\label{duality-lemma-lift-map-plus}
補題 \ref{duality-lemma-lift-map} の結果は
$L \in D^+_{\textit{Coh}}(\mathcal{O}_X)$ に対しても成り立つ。
\end{lemma}

\begin{proof}
実際、$(K_n)$ が Deligne 系ならば、ある
$b \in \mathbf{Z}$ が存在して、$i > b$ に対して $H^i(K_n) = 0$ となる。
すると $\Hom(K_n, L) = \Hom(K_n, \tau_{\leq b}L)$ かつ
$\Hom(K, L) = \Hom(K, \tau_{\leq b}L)$ である。
従って $\tau_{\leq b}L$ に補題の結果を用いればよい。
\end{proof}

\begin{lemma}
\label{duality-lemma-extension-by-zero}
$j : U \to X$ を Noether スキームの開埋め込みとする。
\begin{enumerate}
\item $(K_n)$ と $(L_n)$ を Deligne 系とする。
定値系 $(K_n|_U)$ と $(L_n|_U)$ の値を $K$ と $L$ とする。
$D(\mathcal{O}_U)$ の射 $\alpha : K \to L$ が与えられたとき、
$U$ への制限が $\alpha$ である
$D^b_{\textit{Coh}}(\mathcal{O}_X)$ の pro-系の射
$(K_n) \to (L_n)$ が一意に存在する。
\item $K \in D^b_{\textit{Coh}}(\mathcal{O}_U)$ が与えられたとき、
$(K_n|_U)$ が値 $K$ の定値系となる Deligne 系 $(K_n)$ が存在する。
\item (2) の $D^b_{\textit{Coh}}(\mathcal{O}_X)$ の pro-対象 $(K_n)$ は、
（pro-対象として）一意な同型を除いて一意である。
\end{enumerate}
\end{lemma}

\begin{proof}
(1) は補題 \ref{duality-lemma-lift-map} と、pro-系の間の射が、それらの余表現する関手の間の射と同じであるという事実から直ちに従う。
『圏論』の注意 \ref{categories-remark-pro-category-copresheaves} を参照する。

\medskip\noindent
$K$ を (2) のものとする。$U$ への制限が $K$ と同型である
$K' \in D^b_{\textit{Coh}}(\mathcal{O}_X)$ を選べる。
『スキームの導来圏』の補題 \ref{perfect-lemma-lift-coherent} を参照する。
『スキームの導来圏』の命題 \ref{perfect-proposition-DCoh}
により、$K'$ を連接 $\mathcal{O}_X$-加群の有界複体
$\mathcal{F}^\bullet$ で表せる。
消滅軌跡が $X \setminus U$ である準連接イデアル層
$\mathcal{I} \subset \mathcal{O}_X$ を選ぶ
（例えば $X \setminus U$ 上の被約誘導部分スキーム構造に対応する $\mathcal{I}$ を選ぶ）。
このとき、本節の導入部と同様に、複体
$\mathcal{I}^n\mathcal{F}^\bullet$ によって表される対象を $K_n$ と置ける。

\medskip\noindent
(3) は (1) と (2) から直ちに従う。
\end{proof}

\begin{lemma}
\label{duality-lemma-extension-by-zero-triangle}
$j : U \to X$ を Noether スキームの開埋め込みとする。
$$
K \to L \to M \to K[1]
$$
を $D^b_{\textit{Coh}}(\mathcal{O}_U)$ の完全三角とする。
このとき、$D^b_{\textit{Coh}}(\mathcal{O}_X)$ における完全三角の逆系
$$
K_n \to L_n \to M_n \to K_n[1]
$$
であって、$(K_n)$, $(L_n)$, $(M_n)$ が Deligne 系であり、
これらの完全三角の $U$ への制限が出発点の完全三角と同型になるものが存在する。
\end{lemma}

\begin{proof}
$(K_n)$ を補題 \ref{duality-lemma-extension-by-zero} (2) のものとする。
$U$ への制限が $L$ である
$D^b_{\textit{Coh}}(\mathcal{O}_X)$ の対象 $L'$ を選ぶ
（補題が示すように、これは可能である）。
補題 \ref{duality-lemma-lift-map} により、ある $n$ と $X$ 上の射
$K_n \to L'$ を、その $U$ への制限が与えられた矢印
$K \to L$ となるように選べる。
従って $D^b_{\textit{Coh}}(\mathcal{O}_X)$ の射 $K' \to L'$ で、
その $U$ への制限が与えられた矢印 $K \to L$ となるものが存在する。

\medskip\noindent
『スキームの導来圏』の命題 \ref{perfect-proposition-DCoh}
により、$K' \to L'$ を表す連接 $\mathcal{O}_X$-加群の有界複体の射
$\alpha^\bullet : \mathcal{F}^\bullet \to
\mathcal{G}^\bullet$ を選べる。
消滅軌跡が $X \setminus U$ である準連接イデアル層
$\mathcal{I} \subset \mathcal{O}_X$ を選ぶ。
そこで
$K_n = \mathcal{I}^n\mathcal{F}^\bullet$
および
$L_n = \mathcal{I}^n\mathcal{G}^\bullet$
と置く。$\alpha^\bullet$ は複体の射
$\alpha_n^\bullet : \mathcal{I}^n\mathcal{F}^\bullet \to
\mathcal{I}^n\mathcal{G}^\bullet$ を誘導する。
『導来圏』の第 \ref{derived-section-cones} 節
における錐の構成から、
$$
C(\alpha_n)^\bullet = \mathcal{I}^nC(\alpha^\bullet)
$$
は明らかであり、従って $M_n = C(\alpha_n)^\bullet$ と置ける。
すなわち、両立する完全三角の系
（『導来圏』の第 \ref{derived-section-canonical-delta-functor} 節の議論を参照）
$$
K_n \to L_n \to M_n \to K_n[1]
$$
を得る。その $U$ への制限は、公理 TR3 および『導来圏』の補題 \ref{derived-lemma-third-isomorphism-triangle}
により、出発点の完全三角と同型である。
\end{proof}

\begin{remark}
\label{duality-remark-extension-by-zero}
$j : U \to X$ を Noether スキームの開埋め込みとする。
$K \in D^b_{\textit{Coh}}(\mathcal{O}_U)$ に、その $U$ への制限が $K$ である
Deligne 系を対応させることにより、関手
$$
Rj_! :
D^b_{\textit{Coh}}(\mathcal{O}_U)
\longrightarrow
\text{Pro-}D^b_{\textit{Coh}}(\mathcal{O}_X)
$$
が定まる。これは補題 \ref{duality-lemma-extension-by-zero-triangle} により「完全」であり、
補題 \ref{duality-lemma-lift-map} により関手
$j^* : D^b_{\textit{Coh}}(\mathcal{O}_X) \to
D^b_{\textit{Coh}}(\mathcal{O}_U)$ の「左随伴」である。
\end{remark}

\begin{remark}
\label{duality-remark-extension-by-zero-linear-pro-system}
$(A_n)$ と $(B_n)$ を圏 $\mathcal{C}$ の逆系とする。
$(A_n)$ から $(B_n)$ への線形 pro-射とは、ある固定された整数
$c, d \geq 1$ に対し、すべての $n \geq 1$ についての射の両立する族
$\varphi_n : A_{cn + d} \to B_n$ によって与えられるものとする。
$(\varphi_n : A_{cn + d} \to B_n)$ と
$(\psi_n : A_{c'n + d'} \to B_n)$ が同じ射を定めるとは、
ある $c'' \geq \max(c, c')$ と $d'' \geq \max(d, d')$ が存在して、
誘導される二つの射 $A_{c'' n + d''} \to B_n$ がすべての $n$ に対して一致することとする。
$U$ 上で所与の値をもつ Deligne 系 $(K_n)$ は、線形 pro-同型を除いて良定義であると思われる。
これが必要になれば、その時点でここに注意深い定式化と証明を与える。
\end{remark}

\begin{lemma}
\label{duality-lemma-deligne-system-2-out-of-3}
$j : U \to X$ を Noether スキームの開埋め込みとする。
$$
K_n \to L_n \to M_n \to K_n[1]
$$
を $D^b_{\textit{Coh}}(\mathcal{O}_X)$ における完全三角の逆系とする。
$(K_n)$ と $(M_n)$ が Deligne 系と pro-同型ならば、$(L_n)$ もそうである。
\end{lemma}

\begin{proof}
系 $(K_n|_U)$ と $(M_n|_U)$ は定値系と pro-同型なので、本質的に定値である。
それらの値を $K$ と $M$ と書く。『導来圏』の補題 \ref{derived-lemma-essentially-constant-2-out-of-3}
により、逆系 $L_n|_U$ も本質的に定値である。その値を $L$ と書く。
$N \in D^b_{\textit{Coh}}(\mathcal{O}_X)$ とし、可換図式
$$
\xymatrix{
\ldots \ar[r] &
\colim \Hom_X(M_n, N) \ar[r] \ar[d] &
\colim \Hom_X(L_n, N) \ar[r] \ar[d] &
\colim \Hom_X(K_n, N) \ar[r] \ar[d] &
\ldots \\
\ldots \ar[r] &
\Hom_U(M, N|_U) \ar[r] &
\Hom_U(L, N|_U) \ar[r] &
\Hom_U(K, N|_U) \ar[r] &
\ldots
}
$$
を考える。補題 \ref{duality-lemma-lift-map} と、同型な ind-系は同じ帰納極限をもつという事実により、
中央の垂直矢印の左右それぞれ二つの垂直矢印は同型である。
5 項補題により、中央の垂直矢印も同型である。
ここで $(L'_n)$ を、その $U$ への制限が値 $L$ の定値系となる Deligne 系とすれば
（これは補題 \ref{duality-lemma-extension-by-zero} により存在する）、
$\colim \Hom_X(L'_n, N) = \Hom_U(L, N|_U)$ も成り立つ。
従って pro-系 $(L_n)$ と $(L'_n)$ は、『圏論』の注意 \ref{categories-remark-pro-category-copresheaves} により pro-同型である。
\end{proof}

\begin{lemma}
\label{duality-lemma-consequence-Artin-Rees-bis}
$X$ を Noether スキームとする。$\mathcal{I} \subset \mathcal{O}_X$
を準連接イデアル層とし、$\mathcal{F}^\bullet$ を連接
$\mathcal{O}_X$-加群の複体とする。$p \in \mathbf{Z}$ とする。
$\mathcal{H} = H^p(\mathcal{F}^\bullet)$ および
$\mathcal{H}_n = H^p(\mathcal{I}^n\mathcal{F}^\bullet)$ と置く。
このとき標準的な $\mathcal{O}_X$-加群の写像
$$
\ldots \to \mathcal{H}_3 \to \mathcal{H}_2 \to \mathcal{H}_1 \to \mathcal{H}
$$
がある。ある $c > 0$ が存在して、$n \geq c$ に対し
$\mathcal{H}_n \to \mathcal{H}$ の像は
$\mathcal{I}^{n - c}\mathcal{H}$ に含まれ、さらに標準的な
$\mathcal{O}_X$-加群の写像
$\mathcal{I}^n\mathcal{H} \to \mathcal{H}_{n - c}$ があって、合成
$$
\mathcal{I}^n \mathcal{H} \to \mathcal{H}_{n - c} \to
\mathcal{I}^{n - 2c}\mathcal{H}
\quad\text{and}\quad
\mathcal{H}_n \to \mathcal{I}^{n - c}\mathcal{H} \to \mathcal{H}_{n - 2c}
$$
は標準的な写像である。特に、逆系
$(\mathcal{H}_n)$ と $(\mathcal{I}^n\mathcal{H})$ は
$\textit{Mod}(\mathcal{O}_X)$ の pro-対象として同型である。
\end{lemma}

\begin{proof}
$X$ がアフィンならば、代数に翻訳すると、これは『代数の続論』の補題 \ref{more-algebra-lemma-consequence-Artin-Rees-bis} である。
一般の場合には、その補題の証明とまったく同様に論じ、代数における Artin-Rees への参照を
『スキームのコホモロジー』の補題 \ref{coherent-lemma-Artin-Rees}
への参照で置き換える。詳細は省略する。
\end{proof}

\begin{lemma}
\label{duality-lemma-characterize-extension-by-zero-algebra}
$j : U \to X$ を Noether スキームの開埋め込みとする。
$a \leq b$ を整数とする。$(K_n)$ を
$D^b_{\textit{Coh}}(\mathcal{O}_X)$ の逆系で、
$i \not \in [a, b]$ に対して $H^i(K_n) = 0$ となるものとする。
次は同値である。
\begin{enumerate}
\item $(K_n)$ は Deligne 系と pro-同型である。
\item すべての $p \in \mathbf{Z}$ に対し、連接
$\mathcal{O}_X$-加群 $\mathcal{F}$ が存在して、pro-系
$(H^p(K_n))$ と $(\mathcal{I}^n\mathcal{F})$ が pro-同型となる。
\end{enumerate}
\end{lemma}

\begin{proof}
(1) を仮定する。(2) を示すには、$(K_n)$ が Deligne 系であると仮定してよい。
定義により、連接 $\mathcal{O}_X$-加群の有界複体
$\mathcal{F}^\bullet$ と、$X \setminus U$ を切り出す準連接イデアル層を、
$K_n$ が $\mathcal{I}^n\mathcal{F}^\bullet$ によって表されるように選べる。
従って結果は補題 \ref{duality-lemma-consequence-Artin-Rees-bis} から従う。

\medskip\noindent
(2) を仮定する。$(K_n)$ が (1) のとおりであることを $b - a$ に関する帰納法で示す。
$a = b$ ならば、(1) は本質的に仮定そのものから従う。
$a < b$ ならば、両立する完全三角の系
$$
\tau_{\leq a}K_n \to K_n \to \tau_{\geq a + 1}K_n \to (\tau_{\leq a}K_n)[1]
$$
を考える。『導来圏』の注意 \ref{derived-remark-truncation-distinguished-triangle} を参照する。
$b - a$ に関する帰納法により、$\tau_{\leq a}K_n$ と
$\tau_{\geq a + 1}K_n$ は Deligne 系と pro-同型である。
補題 \ref{duality-lemma-deligne-system-2-out-of-3} により結論を得る。
\end{proof}

\begin{lemma}
\label{duality-lemma-characterize-extension-by-zero}
$j : U \to X$ を Noether スキームの開埋め込みとする。
$(K_n)$ を $D^b_{\textit{Coh}}(\mathcal{O}_X)$ の逆系とする。
$X = W_1 \cup \ldots \cup W_r$ を開被覆とする。次は同値である。
\begin{enumerate}
\item $(K_n)$ は Deligne 系と pro-同型である。
\item 各 $i$ に対し、制限 $(K_n|_{W_i})$ は開埋め込み
$U \cap W_i \to W_i$ に関する Deligne 系と pro-同型である。
\end{enumerate}
\end{lemma}

\begin{proof}
$r$ に関する帰納法による。$r = 1$ ならば結果は明らかである。
$r > 1$ と仮定し、$V = W_1 \cup \ldots \cup W_{r - 1}$ と置く。
帰納法により $(K_n|_V)$ は Deligne 系である。
従って次の段落での議論に帰着する。

\medskip\noindent
$X = V \cup W$ が開被覆であり、$(K_n|_W)$ と $(K_n|_V)$ が
Deligne 系と pro-同型であると仮定する。
$(K_n)$ が Deligne 系と pro-同型であることを示す必要がある。
$(K_n|_{V \cap W})$ は Deligne 系と pro-同型である
（Deligne 系の構成から、開集合への制限がそれを保つことは直ちに従う）。
特に pro-系
$(K_n|_{U \cap V})$,
$(K_n|_{U \cap W})$, および
$(K_n|_{U \cap V \cap W})$
は本質的に定値である。
『層のコホモロジー』の補題 \ref{cohomology-lemma-exact-sequence-j-star} の完全三角と
『導来圏』の補題 \ref{derived-lemma-essentially-constant-2-out-of-3}
により、$(K_n|_U)$ も本質的に定値である。
この系の値を $K \in D^b_{\textit{Coh}}(\mathcal{O}_U)$ と書く。
$L$ を $D^b_{\textit{Coh}}(\mathcal{O}_X)$ の対象とし、図式
$$
\xymatrix{
\colim \Ext^{-1}(K_n|_V, L|_V) \oplus
\colim \Ext^{-1}(K_n|_W, L|_W) \ar[r] \ar[d] &
\Ext^{-1}(K|_{U \cap V}, L|_{U \cap V}) \oplus
\Ext^{-1}(K|_{U \cap W}, L|_{U \cap W}) \ar[d] \\
\colim \Ext^{-1}(K_n|_{V \cap W}, L|_{V \cap W}) \ar[r] \ar[d] &
\Ext^{-1}(K|_{U \cap V \cap W}, L|_{U \cap V \cap W}) \ar[d] \\
\colim \Hom(K_n, L) \ar[d] \ar[r] &
\Hom(K|_U, L|_U) \ar[d] \\
\colim \Hom(K_n|_V, L|_V) \oplus \colim \Hom(K_n|_W, L|_W) \ar[r] \ar[d] &
\Hom(K|_{U \cap V}, L|_{U \cap V}) \oplus
\Hom(K|_{U \cap W}, L|_{U \cap W}) \ar[d] \\
\colim \Hom(K_n|_{V \cap W}, L|_{V \cap W}) \ar[r] &
\Hom(K|_{U \cap V \cap W}, L|_{U \cap V \cap W})
}
$$
を考える。垂直列は『層のコホモロジー』の補題 \ref{cohomology-lemma-mayer-vietoris-hom} と、フィルター付き帰納極限が完全であるという事実により完全である。
中央以外のすべての水平矢印は補題 \ref{duality-lemma-lift-map} と、
pro-同型な系が同じ帰納極限をもつという事実により同型である。
従って 5 項補題により中央の矢印も同型である。
これより $(K_n)$ は $K$ に対する Deligne 系と pro-同型である。
すなわち、$(K'_n)$ を、その $U$ への制限が値 $K$ の定値系となる Deligne 系とすれば
（これは補題 \ref{duality-lemma-extension-by-zero} により存在する）、
$\colim \Hom_X(K'_n, L) = \Hom_U(K, L|_U)$ も成り立つ。
従って pro-系 $(K_n)$ と $(K'_n)$ は『圏論』の注意 \ref{categories-remark-pro-category-copresheaves} により pro-同型である。
\end{proof}

\begin{lemma}
\label{duality-lemma-tensoring-Deligne-system}
$j : U \to X$ を Noether スキームの開埋め込みとする。
$\mathcal{I} \subset \mathcal{O}_X$ を
$V(\mathcal{I}) = X \setminus U$ を満たす準連接イデアル層とする。
$K$ を $D^b_{\textit{Coh}}(\mathcal{O}_X)$ の対象とする。このとき
$$
K \otimes_{\mathcal{O}_X}^\mathbf{L} \mathcal{I}^n
$$
は、$U$ 上で定値 $K|_U$ をもつ Deligne 系と pro-同型である。
\end{lemma}

\begin{proof}
補題 \ref{duality-lemma-characterize-extension-by-zero} により問題は $X$ 上局所的である。
従って $X$ は Noether 環のスペクトルであると仮定してよい。
この場合、主張は代数版である『代数の続論』の補題 \ref{more-algebra-lemma-tensoring-Deligne-system} から従う。
\end{proof}







\section{コンパクト台コホモロジーの準備}
\label{duality-section-preliminaries-compactly-supported}

\noindent
状況 \ref{duality-situation-shriek} において、$f : X \to Y$ を圏
$\textit{FTS}_S$ の射とする。前節の構成を用いて関手
$$
Rf_! :
D^b_{\textit{Coh}}(\mathcal{O}_X)
\longrightarrow
\text{Pro-}D^b_{\textit{Coh}}(\mathcal{O}_Y)
$$
を構成する。$f$ が開埋め込みならば、これは注意 \ref{duality-remark-extension-by-zero} の関手に帰着し、一般には $f$ のコンパクト化を用いて構成される。
これを行う前に、構成が良定義であることを示すため次の補題が必要である。

\begin{lemma}
\label{duality-lemma-well-defined-pre}
$f : X \to Y$ を Noether スキームの固有射とする。
$V \subset Y$ を開部分スキームとし、$U = f^{-1}(V)$ と置く。図式
$$
\xymatrix{
U \ar[r]_j \ar[d]_g & X \ar[d]^f \\
V \ar[r]^{j'} & Y
}
$$
を考える。このとき、注意 \ref{duality-remark-extension-by-zero} の
$Rj_!$ と $Rj'_!$ に対して、関手
$D^b_{\textit{Coh}}(\mathcal{O}_U) \to
\text{Pro-}D^b_{\textit{Coh}}(\mathcal{O}_Y)$
の標準同型 $Rj'_! \circ Rg_* \to Rf_* \circ Rj_!$ が存在する。
\end{lemma}

\begin{proof}[第一証明]
$K$ を $D^b_{\textit{Coh}}(\mathcal{O}_U)$ の対象とする。
$(K_n)$ を $U \to X$ に対する Deligne 系で、その $U$ への制限が値 $K$ の定値系であるものとする。
これはもちろん、$(K_n)$ が
$\text{Pro-}D^b_{\textit{Coh}}(\mathcal{O}_X)$ において $Rj_!K$ を表すことを意味する。
$Rj'_!Rg_*K$ と $Rf_*Rj_!K$ のいずれも、$V$ への制限は値 $Rg_*K$ の定値 pro-対象である。
前者については直ちに分かり、後者については
$(Rf_*K_n)|_V = Rg_*(K_n|_U) = Rg_*K$ から従う。
補題 \ref{duality-lemma-extension-by-zero} における Deligne 系の一意性により、
$(Rf_*K_n)$ が Deligne 系と pro-同型であることを示せば十分である。
同じ補題により、得られる同型が関手的であることも分かる。

\medskip\noindent
$(Rf_*K_n)$ が Deligne 系と pro-同型であることの証明。
まず、この問題は $K$ に対応する Deligne 系 $(K_n)$ の選択に依存しない
（これは前述の一意性による）。
補題 \ref{duality-lemma-extension-by-zero-triangle} および
\ref{duality-lemma-deligne-system-2-out-of-3} により、
$D^b_{\textit{Coh}}(\mathcal{O}_U)$ に完全三角
$$
K \to L \to M \to K[1]
$$
があり、結果が $K$ と $M$ に対して成り立つならば、$L$ に対しても成り立つ。
標準切断の完全三角（『導来圏』の注意 \ref{derived-remark-truncation-distinguished-triangle}）を用いて、
次の段落で扱う問題に帰着する。

\medskip\noindent
$\mathcal{F}$ を連接 $\mathcal{O}_X$-加群とする。
$\mathcal{J} \subset \mathcal{O}_Y$ を $Y \setminus V$ を切り出す準連接イデアル層とする。
$f^*\mathcal{J}^n \otimes \mathcal{F} \to \mathcal{F}$ の像を
$\mathcal{J}^n\mathcal{F}$ と書く。
$(Rf_*(\mathcal{J}^n\mathcal{F}))$ が Deligne 系であることを示す必要がある。
補題 \ref{duality-lemma-characterize-extension-by-zero} により問題は $Y$ 上局所的である。
従って $Y = \Spec(A)$ はアフィンであり、$\mathcal{J}$ はイデアル
$I \subset A$ に対応すると仮定してよい。
補題 \ref{duality-lemma-characterize-extension-by-zero-algebra}
により、コホモロジー加群の逆系
$(H^p(X, I^n\mathcal{F}))$ が、ある有限 $A$-加群 $M$ に対する逆系
$(I^n M)$ と pro-同型であることを示せば十分である。
これは『スキームのコホモロジー』の補題 \ref{coherent-lemma-cohomology-powers-ideal-application} で示されている。
\end{proof}

\begin{proof}[第二証明]
$K$ を $D^b_{\textit{Coh}}(\mathcal{O}_U)$ の対象とする。
$L$ を $D^b_{\textit{Coh}}(\mathcal{O}_Y)$ の対象とする。
$K$ と $L$ に関して関手的な全単射
$$
\Hom_{\text{Pro-}D^b_{\textit{Coh}}(\mathcal{O}_Y)}(Rj'_!Rg_*K, L)
\longrightarrow
\Hom_{\text{Pro-}D^b_{\textit{Coh}}(\mathcal{O}_Y)}(Rf_*Rj_!K, L)
$$
を構成する。$K$ を固定すると、これは『圏論』の注意 \ref{categories-remark-pro-category-copresheaves} により pro-対象の同型
$Rf_*Rj_!K \to Rj'_!Rg_*K$
を定め、さらに $K$ を動かすと関手の同型を定める。
実際にこの同型を得るため、次の関手的な等式列を用いる。
\begin{align*}
\Hom_{\text{Pro-}D^b_{\textit{Coh}}(\mathcal{O}_Y)}(Rj'_!Rg_*K, L)
& =
\Hom_V(Rg_*K, L|_V) \\
& =
\Hom_U(K, g^!(L|_V)) \\
& =
\Hom_U(K, f^!L|_U)) \\
& =
\Hom_{\text{Pro-}D^b_{\textit{Coh}}(\mathcal{O}_X)}(Rj_!K, f^!L) \\
& =
\Hom_{\text{Pro-}D^b_{\textit{Coh}}(\mathcal{O}_Y)}(Rf_*Rj_!K, L)
\end{align*}
第一の等式は補題 \ref{duality-lemma-lift-map} による。
第二の等式は、$g$ が $f$ の $V$ への基底変換として固有であり、従って構成により
$g^!$ が押し出しの右随伴だからである。第 \ref{duality-section-upper-shriek} 節を参照せよ。
第三の等式は補題 \ref{duality-lemma-restrict-before-or-after} による
$g^!(L|_V) = f^!L|_U$ から従う。
補題 \ref{duality-lemma-shriek-coherent} により
$f^!L$ は $D^+_{\textit{Coh}}(\mathcal{O}_X)$ に属するので、第四の等式は 
補題 \ref{duality-lemma-lift-map-plus} から従う。第五の等式も、$f$ が固有であるため
$f^!$ が $Rf_*$ の右随伴であることから従う。
\end{proof}

\begin{lemma}
\label{duality-lemma-well-defined}
$j : U \to X$ を Noether スキームの開埋め込みとする。
$j' : U \to X'$ を $X$ 上の $U$ のコンパクト化（証明を参照）とし、
$f : X' \to X$ を構造射とする。このとき、注意 \ref{duality-remark-extension-by-zero} の $Rj_!$ と $Rj'_!$ に対して、関手
$D^b_{\textit{Coh}}(\mathcal{O}_U) \to
\text{Pro-}D^b_{\textit{Coh}}(\mathcal{O}_X)$
の標準同型 $Rj_! \to Rf_* \circ R(j')_!$ が存在する。
\end{lemma}

\begin{proof}
$X'$ が $X$ 上の $U$ のコンパクト化であるとは、正確には
$f : X' \to X$ が固有であり、$j'$ が開埋め込みであり、かつ
$j = f \circ j'$ であることを意味する。『平坦性の続論』の第 \ref{flat-section-compactify} 節を参照せよ。
$j'(U) = f^{-1}(j(U))$ ならば、補題は補題 \ref{duality-lemma-well-defined-pre} から直ちに従う。
$j'(U) \not = f^{-1}(j(U))$ ならば、$X'' \subset X'$ を
$j' : U \to X'$ のスキーム論的閉包とし、
$j'' : U \to X''$ を対応する開埋め込みとする。図式
$$
\xymatrix{
& & X'' \ar[d]^{f'} \\
& & X' \ar[d]^f \\
U \ar[rr]^j \ar[rru]^{j'} \ar[rruu]^{j''} & & X
}
$$
を考える。『平坦性の続論』の補題 \ref{flat-lemma-compactifications-cofiltered} part (c) と上の議論により、同型
$Rf'_* \circ Rj''_! = Rj'_!$ および $R(f \circ f')_* \circ Rj''_! = Rj_!$
を得る。$R(f \circ f')_* = Rf_* \circ Rf'_*$ であるから結論が従う。
\end{proof}

\begin{remark}
\label{duality-remark-covariance-open-j-lower-shriek}
$X$ を Noether スキームとし、$X \supset U \supset U'$ を開部分スキームとする。
$j : U \to X$ および $j' : U' \to X$ を包含射と書く。
$D^b_{\textit{Coh}}(\mathcal{O}_U)$ の $K$ に関して関手的な標準射
$$
Rj'_!(K|_{U'}) \longrightarrow Rj_!K
$$
が存在することを主張する。実際、補題 \ref{duality-lemma-lift-map} により、
$D^b_{\textit{Coh}}(\mathcal{O}_X)$ の任意の $L$ に対して射
\begin{align*}
\Hom_{\text{Pro-}D^b_{\textit{Coh}}(\mathcal{O}_X)}(Rj_!K, L)
& =
\Hom_U(K, L|_U) \\
& \to
\Hom_{U'}(K|_{U'}, L|_{U'}) \\
& =
\Hom_{\text{Pro-}D^b_{\textit{Coh}}(\mathcal{O}_X)}(Rj'_!(K|_{U'}), L)
\end{align*}
を得る。これは $L$ と $K'$ に関して関手的である。
$L$ に関する関手性と『圏論』の注意 \ref{categories-remark-pro-category-copresheaves} により、標準射
$Rj'_!(K|_{U'}) \to Rj_!K$ を得る。これは上の矢印の $K$ に関する関手性により
$K$ に関して関手的である。

\medskip\noindent
この矢印を明示的に構成しよう。$\mathcal{F}^\bullet$ を連接
$\mathcal{O}_X$-加群の有界複体で、その $U$ への制限が導来圏において
$K$ を表すものとする。このような複体が常に存在することは補題 \ref{duality-lemma-extension-by-zero} の証明で見た。
$\mathcal{I}$, resp.\ $\mathcal{I}'$ をそれぞれ $X$ 上の準連接イデアル層で、
$V(\mathcal{I}) = X \setminus U$, resp.\ $V(\mathcal{I}') = X \setminus U'$
を満たすものとする。$\mathcal{I}$ を $\mathcal{I} + \mathcal{I}'$ で置き換えることにより、
$\mathcal{I}' \subset \mathcal{I}$ と仮定してよい。
構成により $Rj_!K$, resp.\ $Rj'_!(K|_{U'})$ は、
$D^b_{\textit{Coh}}(\mathcal{O}_X)$ の逆系 $(K_n)$, resp.\ $(K'_n)$ で表される。ただし
$$
K_n = \mathcal{I}^n\mathcal{F}^\bullet
\quad\text{resp.}\quad
K'_n = (\mathcal{I}')^n\mathcal{F}^\bullet
$$
である。上で構成した射が、包含
$(\mathcal{I}')^n \subset \mathcal{I}^n$ から得られる射
$K'_n \to K_n$ により与えられることは明らかである。
\end{remark}






\section{連接加群のコンパクト台コホモロジー}
\label{duality-section-compactly-supported}

\noindent
状況 \ref{duality-situation-shriek} において、$f : X \to Y$ を
$\textit{FTS}_S$ の射とする。関手
$$
Rf_! : D^b_{\textit{Coh}}(\mathcal{O}_X)
\longrightarrow
\text{Pro-}D^b_{\textit{Coh}}(\mathcal{O}_Y)
$$
を次のように定義する。『平坦性の続論』の定理 \ref{flat-theorem-nagata} および
補題 \ref{flat-lemma-compactifyable} により、$Y$ 上のコンパクト化
$j : X \to \overline{X}$ を選ぶことができる。
$\overline{f} : \overline{X} \to Y$ を構造射と書く。注意 \ref{duality-remark-extension-by-zero} の $Rj_!$ を用い、
$K \in D^b_{\textit{Coh}}(\mathcal{O}_X)$ に対して
$$
Rf_!K  = R\overline{f}_* Rj_! K
$$
と置く。

\begin{lemma}
\label{duality-lemma-lower-shriek-well-defined}
関手 $Rf_!$ は、同型を除いてコンパクト化の選択に依存しない。
\end{lemma}

\noindent
実際、関手 $Rf_!$ は $f^!$ の「左随伴」として特徴づけられ、
これにより一意な同型を除いて決まる。

\begin{proof}
$Y$ 上の $X$ のコンパクト化の圏を考える。『平坦性の続論』の定理 \ref{flat-theorem-nagata} および補題 \ref{flat-lemma-compactifications-cofiltered} および \ref{flat-lemma-compactifyable} により、この圏は余フィルター的である。
コンパクト化の各選択
$$
j : X \to \overline{X},\quad \overline{f} : \overline{X} \to Y
$$
に対し、上の構成は関手 $R\overline{f}_* \circ Rj_!$ を対応させる。
$Y$ 上のコンパクト化 $j_i : X \to \overline{X}_i$ の間の射
$g : \overline{X}_1 \to \overline{X}_2$ が与えられたとする。このとき同型
$$
R\overline{f}_{2, *} \circ Rj_{2, !} =
R\overline{f}_{2, *} \circ Rg_* \circ j_{1, !} =
R\overline{f}_{1, *} \circ Rj_{1, !}
$$
を得る。第一の等式では補題 \ref{duality-lemma-well-defined} を用いた。
従ってこの関手は、同型を除いてコンパクト化の選択に依存しない。
\end{proof}

\begin{proposition}
\label{duality-proposition-duality-compactly-supported}
状況 \ref{duality-situation-shriek} において、$f : X \to Y$ を
$\textit{FTS}_S$ の射とする。このとき関手 $Rf_!$ と $f^!$ は次の意味で随伴する。
任意の $K \in D^b_{\textit{Coh}}(\mathcal{O}_X)$ と
$L \in D^+_{\textit{Coh}}(\mathcal{O}_Y)$ に対して、$K$ と $L$ に関して双関手的に
$$
\Hom_X(K, f^!L) =
\Hom_{\text{Pro-}D^+_{\textit{Coh}}(\mathcal{O}_Y)}(Rf_!K, L)
$$
が成り立つ。
\end{proposition}

\begin{proof}
$Y$ 上のコンパクト化 $j : X \to \overline{X}$ を選び、
$\overline{f} : \overline{X} \to Y$ を構造射と書く。このとき
\begin{align*}
\Hom_X(K, f^!L)
& =
\Hom_X(K, j^*\overline{f}{}^!L) \\
& =
\Hom_{\text{Pro-}D^+_{\textit{Coh}}(\mathcal{O}_{\overline{X}})}
(Rj_!K, \overline{f}{}^!L) \\
& =
\Hom_{\text{Pro-}D^+_{\textit{Coh}}(\mathcal{O}_Y)}(Rf_*Rj_!K, L) \\
& =
\Hom_{\text{Pro-}D^+_{\textit{Coh}}(\mathcal{O}_Y)}(Rf_!K, L)
\end{align*}
を得る。第一の等式は第 \ref{duality-section-upper-shriek} 節における
$f^!$ の構成から直ちに従う。補題 \ref{duality-lemma-shriek-coherent} により
$\overline{f}{}^!L$ は $D^+_{\textit{Coh}}(\mathcal{O}_{\overline{X}})$ に属するので、
第二の等式は補題 \ref{duality-lemma-lift-map-plus} から従う。
$\overline{f}$ は固有であるから、構成により関手 $\overline{f}{}^!$ は押し出しの右随伴である。
これが第三の等式を与える。第四の等式は $Rf_! = Rf_* Rj_!$ による。
\end{proof}

\begin{lemma}
\label{duality-lemma-compactly-supported-triangle}
状況 \ref{duality-situation-shriek} において、$f : X \to Y$ を
$\textit{FTS}_S$ の射とする。
$$
K \to L \to M \to K[1]
$$
を $D^b_{\textit{Coh}}(\mathcal{O}_X)$ の完全三角とする。このとき
$D^b_{\textit{Coh}}(\mathcal{O}_Y)$ の完全三角の逆系
$$
K_n \to L_n \to M_n \to K_n[1]
$$
で、pro-系 $(K_n)$, $(L_n)$, $(M_n)$ がそれぞれ
$Rf_!K$, $Rf_!L$, $Rf_!M$ を与えるものが存在する。
\end{lemma}

\begin{proof}
$Y$ 上のコンパクト化 $j : X \to \overline{X}$ を選び、
$\overline{f} : \overline{X} \to Y$ を構造射と書く。
開埋め込み $j$ と与えられた完全三角に対応する
$D^b_{\textit{Coh}}(\mathcal{O}_{\overline{X}})$ の完全三角の逆系
$$
\overline{K}_n \to \overline{L}_n \to \overline{M}_n \to \overline{K}_n[1]
$$
を補題 \ref{duality-lemma-extension-by-zero-triangle} のように選ぶ。
$K_n = R\overline{f}_*\overline{K}_n$ と置き、$L_n$, $M_n$ についても同様に置く。
これは $Rf_!$ の定義そのものにより所望の性質をもつ。
\end{proof}

\begin{remark}
\label{duality-remark-compose-inverse-systems}
$\mathcal{C}$ を圏とする。$\mathcal{C}$ の pro-対象の圏における逆系の逆系
$$
\ldots \xrightarrow{\alpha_4} (M_{3, n}) \xrightarrow{\alpha_3} (M_{2, n})
\xrightarrow{\alpha_2} (M_{1, n})
$$
が与えられたとする。すなわち、矢印 $\alpha_i$ は pro-対象の射である。
『圏論』の例 \ref{categories-example-pro-morphism-inverse-systems} により、
各 $\alpha_i$ は対 $(m_i, a_i)$ で表せる。ここで
$m_i : \mathbf{N} \to \mathbf{N}$ は増加関数であり、
$a_{i, n} : M_{i, m_i(n)} \to M_{i - 1, n}$ は $\mathcal{C}$ の射であって、図式
$$
\xymatrix{
\ldots \ar[r] &
M_{i, m_i(3)} \ar[d]^{a_{i, 3}} \ar[r] &
M_{i, m_i(2)} \ar[d]^{a_{i, 2}} \ar[r] &
M_{i, m_i(1)} \ar[d]^{a_{i, 1}} \\
\ldots \ar[r] &
M_{i - 1, 3} \ar[r] &
M_{i - 1, 2} \ar[r] &
M_{i - 1, 1}
}
$$
を可換にする。$m_i(n)$ を $\max(n, m_i(n))$ で置き換え、それに応じて射
$a_i(n)$ を調整することにより（参照した例と同様に）、
$m_i(n) \geq n$ と仮定してよい。このとき、一般項を
$$
M_k = M_{k, m_k(m_{k - 1}(\ldots (m_2(k))\ldots))}
$$
とする逆系
$$
\ldots \to
M_{4, m_4(m_3(m_2(4)))} \to
M_{3, m_3(m_2(3))} \to
M_{2, m_2(2)} \to
M_{1, 1}
$$
を考える。$\mathcal{C}$ の任意の対象 $N$ に対して
$$
\colim_i \colim_n \Mor_\mathcal{C}(M_{i, n}, N) =
\colim_k \Mor_\mathcal{C}(M_k, N) 
$$
が成り立つ。詳細は省略する。換言すれば、逆系 $(M_k)$ は
$$
\colim_i \Mor_{\text{Pro-}\mathcal{C}}((M_{i, n}), N) =
\Mor_{\text{Pro-}\mathcal{C}}((M_k), N)
$$
という性質をもつ。『圏論』の注意 \ref{categories-remark-pro-category-copresheaves} の議論により、
この性質は逆系 $(M_k)$ を pro-同型を除いて決定する。
この方法により、$\text{Pro-}\mathcal{C}$ のある種の逆系を、
可算な添字圏をもつ pro-対象に変換できる。
\end{remark}

\begin{remark}
\label{duality-remark-composition-lower-shriek}
状況 \ref{duality-situation-shriek} において、$f : X \to Y$ と $g : Y \to Z$ を
$\textit{FTS}_S$ の合成可能な射とする。合成
$$
Rg_! \circ Rf_! :
D^b_{\textit{Coh}}(\mathcal{O}_X)
\longrightarrow
\text{Pro-}D^b_{\textit{Coh}}(\mathcal{O}_Z)
$$
を定義しよう。$Rf_!$ の構成そのものにより、
$D^b_{\textit{Coh}}(\mathcal{O}_X)$ の $K$ に対する出力 $Rf_!K$ は、
$D^b_{\textit{Coh}}(\mathcal{O}_Y)$ の逆系 $(M_n)$ の pro-同型類である。
$Rg_!$ も同様に構成されるので、
$$
\ldots \to Rg_!M_3 \to Rg_!M_2 \to Rg_!M_1
$$
は $\text{Pro-}D^b_{\textit{Coh}}(\mathcal{O}_Y)$ の逆系である。
注意 \ref{duality-remark-compose-inverse-systems} の議論により、
$D^b_{\textit{Coh}}(\mathcal{O}_Z)$ の逆系の一意な pro-同型類で、
$Rg_! Rf_! K$ と書かれ、
$$
\Hom_{\text{Pro-}D^b_{\textit{Coh}}(\mathcal{O}_Z)}(Rg_!Rf_!K, L) =
\colim_n \Hom_{\text{Pro-}D^b_{\textit{Coh}}(\mathcal{O}_Z)}(Rg_!M_n, L)
$$
を満たすものが存在する。この構成の $K$ に関する関手性は次の補題から直ちに従うので、
その確認に必要な議論は省略する。
\end{remark}

\begin{lemma}
\label{duality-lemma-composition-lower-shriek}
状況 \ref{duality-situation-shriek} において、$f : X \to Y$ と $g : Y \to Z$ を
$\textit{FTS}_S$ の合成可能な射とする。注意 \ref{duality-remark-composition-lower-shriek} の記法のもとで
$Rg_! \circ Rf_! = R(g \circ f)_!$ が成り立つ。
\end{lemma}

\begin{proof}
『圏論』の注意 \ref{categories-remark-pro-category-copresheaves} の議論により、
$D^b_{\textit{Coh}}(\mathcal{O}_Z)$ の $L$ への $\Hom$ を計算した結果が同じであることを示せば十分である。
注意 \ref{duality-remark-composition-lower-shriek} の記法を用いて計算すると
\begin{align*}
\Hom_Z(Rg_!Rf_!K, L)
& =
\colim_n \Hom_Z(Rg_!M_n, L) \\
& =
\colim_n \Hom_Y(M_n, g^!L) \\
& =
\Hom_Y(Rf_!K, g^!L) \\
& =
\Hom_X(K, f^!g^!L) \\
& =
\Hom_X(K, (g \circ f)^!L) \\
& =
\Hom_Z(R(g \circ f)_!K, L)
\end{align*}
を得る。第一の等式は $Rg_!Rf_!K$ の定義である。
第二の等式は $g$ に対する命題 \ref{duality-proposition-duality-compactly-supported} である。
第三の等式は $Rf_!K$ が $(M_n)$ で与えられることによる。
第四の等式は $f$ に対する命題 \ref{duality-proposition-duality-compactly-supported} である。
第五の等式は補題 \ref{duality-lemma-upper-shriek-composition} である。
第六の等式は $g \circ f$ に対する命題 \ref{duality-proposition-duality-compactly-supported} である。
\end{proof}

\begin{remark}
\label{duality-remark-covariance-open-lower-shriek}
状況 \ref{duality-situation-shriek} において、$f : X \to Y$ を
$\textit{FTS}_S$ の射とし、$U \subset X$ を開集合とする。
$g = f|_U : U \to Y$ と置く。このとき
$D^b_{\textit{Coh}}(\mathcal{O}_X)$ の $K$ に関して関手的な標準射
$$
Rg_!(K|_U) \longrightarrow Rf_!K
$$
が存在し、少なくとも次の三通りに定義できる。
\begin{enumerate}
\item $i : U \to X$ を包含射と書く。補題 \ref{duality-lemma-composition-lower-shriek} により
$Rg_! = Rf_! \circ Ri_!$ である。注意 \ref{duality-remark-covariance-open-j-lower-shriek} の特殊な場合である射
$Ri_!(K|_U) \to K$ に $Rf_!$ を適用すればよい。
\item $X$ の $Y$ 上のコンパクト化 $j : X \to \overline{X}$ を選び、
$\overline{f} : \overline{X} \to Y$ を構造射とする。
$j' = j \circ i : U \to \overline{X}$ と置く。
$Rf_! = R\overline{f}_* \circ Rj_!$ および
$Rg_! = R\overline{f}_* \circ Rj'_!$ を用い、注意 \ref{duality-remark-covariance-open-j-lower-shriek} の射
$Rj'_!(K|_U) \to Rj_!K$ に $R\overline{f}_*$ を適用すればよい。
\item $L$ と $K$ に関して関手的な写像列
\begin{align*}
\Hom_{\text{Pro-}D^b_{\textit{Coh}}(\mathcal{O}_Y)}(Rf_!K, L)
& =
\Hom_X(K, f^!L) \\
& \to
\Hom_U(K|_U, f^!L|_U) \\
& =
\Hom_U(K|_U, g^!L) \\
& =
\Hom_{\text{Pro-}D^b_{\textit{Coh}}(\mathcal{O}_Y)}(Rg_!(K|_U), L)
\end{align*}
を用いる。ここでは命題 \ref{duality-proposition-duality-compactly-supported} を二度用い、また上付き shriek 関手の構成から
$g^! = i^* \circ f^!$ が成り立つことを用いた。
$L$ に関する関手性と『圏論』の注意 \ref{categories-remark-pro-category-copresheaves} により、
$\text{Pro-}D^b_{\textit{Coh}}(\mathcal{O}_Y)$ における標準射
$Rg_!(K|_U) \to Rf_!K$ を得る。これは上の矢印の $K$ に関する関手性により
$K$ に関して関手的である。
\end{enumerate}
これら三つの構成はいずれも同じ矢印を与える。詳細は省略する。
\end{remark}

\begin{remark}
\label{duality-remark-covariance-etale-lower-shriek}
注意 \ref{duality-remark-covariance-open-lower-shriek} で与えた
コンパクト台コホモロジーの共変性を \'etale 射へ一般化しよう。
すなわち状況 \ref{duality-situation-shriek} において、可換図式
$$
\xymatrix{
U \ar[rr]_h \ar[rd]_g & & X \ar[ld]^f \\
& Y
}
$$
が $\textit{FTS}_S$ において与えられ、$h$ が \'etale であるとする。このとき
$D^b_{\textit{Coh}}(\mathcal{O}_X)$ の $K$ に関して関手的な標準射
$$
Rg_!(h^*K) \longrightarrow Rf_!K
$$
が存在する。この変換を、$L$ と $K$ に関して関手的な写像列
\begin{align*}
\Hom_{\text{Pro-}D^b_{\textit{Coh}}(\mathcal{O}_Y)}(Rf_!K, L)
& =
\Hom_X(K, f^!L) \\
& \to
\Hom_U(h^*K, h^*(f^!L)) \\
& =
\Hom_U(h^*K, h^!f^!L) \\
& =
\Hom_U(h^*K, g^!L) \\
& =
\Hom_{\text{Pro-}D^b_{\textit{Coh}}(\mathcal{O}_Y)}(Rg_!(h^*K), L)
\end{align*}
を用いて定義する。ここでは命題 \ref{duality-proposition-duality-compactly-supported} を二度用い、補題 \ref{duality-lemma-shriek-etale} の等式 $h^* = h^!$ と、補題 \ref{duality-lemma-upper-shriek-composition} の等式 $h^! \circ f^! = g^!$ を用いた。
$L$ に関する関手性と『圏論』の注意 \ref{categories-remark-pro-category-copresheaves} により、
$\text{Pro-}D^b_{\textit{Coh}}(\mathcal{O}_Y)$ における標準射
$Rg_!(h^*K) \to Rf_!K$ を得る。これは上の矢印の $K$ に関する関手性により
$K$ に関して関手的である。
\end{remark}

\begin{remark}
\label{duality-remark-covariance-lower-shriek}
注意 \ref{duality-remark-covariance-open-lower-shriek} および
\ref{duality-remark-covariance-etale-lower-shriek} では、コンパクト台コホモロジーの構成が
開埋め込みと \'etale 射に関して共変であることを見た。
実際、適切な一般性は次の通りである。$\textit{FTS}_S$ の可換図式
$$
\xymatrix{
U \ar[rr]_h \ar[rd]_g & & X \ar[ld]^f \\
& Y
}
$$
が与えられ、$h$ が平坦かつ準有限であるとき、標準変換
$$
Rg_! \circ h^* \longrightarrow Rf_!
$$
が存在する。注意 \ref{duality-remark-covariance-etale-lower-shriek} と同様に、
この射は $D^+_{\textit{Coh}}(\mathcal{O}_X)$ 上の関手の変換
$h^* \to h^!$ を用いて構成できる。平坦準有限射 $h$ の相対双対化層を
$\omega_{U/X} = h^!\mathcal{O}_X$ とすると
$h^!K = h^*K \otimes \omega_{U/X}$ であることを想起しよう
（補題 \ref{duality-lemma-perfect-comparison-shriek} および
\ref{duality-lemma-flat-quasi-finite-shriek} を参照）。
$\omega_{U/X}$ は、『判別式と差イデアル』の注意 \ref{discriminant-remark-relative-dualizing-for-quasi-finite}
で構成される相対双対化加群と同じであることが、『判別式と差イデアル』の補題 \ref{discriminant-lemma-compare-dualizing} により分かる。
従って『判別式と差イデアル』の注意 \ref{discriminant-remark-relative-dualizing-for-flat-quasi-finite}
で構成されるトレース元
$\tau_{U/X} : \mathcal{O}_U \to \omega_{U/X}$
を用いてこの変換を定義できる。
必要になれば、この結果をここで厳密に定式化し証明する。
\end{remark}







\section{コンパクト台コホモロジーに対する双対性}
\label{duality-section-duality-compactly-supported}

\noindent
$k$ を体とし、$U$ を $k$ 上有限型の分離スキームとする。
$K$ を $D^b_{\textit{Coh}}(\mathcal{O}_U)$ の対象とする。
$K$ のコンパクト台コホモロジー $H^i_c(U, K)$ を次のように定義する。
$k$ 上固有なスキームへの開埋め込み $j : U \to X$ と、
$j : U \to X$ に対する Deligne 系 $(K_n)$ で、その $U$ への制限が値 $K$ の定値系であるものを選ぶ。このとき
$$
H^i_c(U, K) = \lim H^i(X, K_n)
$$
と置く。これを極限位相による位相 $k$-ベクトル空間とみなす
（『代数の続論』の第 \ref{more-algebra-section-topological-ring} 節を参照）。
ここでいくつか注意が必要である。

\medskip\noindent
第一に、この定義は $X$ と $(K_n)$ の選択に依存しない。
実際、$p : U \to \Spec(k)$ を構造射とすると、
$Rp_!K = (R\Gamma(X, K_n))$ が $D(k)$ において pro-同型を除いて良定義であることは既に分かっている。
従って $H^i_c(U, K)$ を定める極限も同様である。

\medskip\noindent
第二に、式
$$
H^i(R\lim R\Gamma(X, K_n)) = R\Gamma(X, R\lim K_n)
$$
を用いる方が自然に見えるかもしれないが、これも同じ答えを与える。
$k$-ベクトル空間 $H^j(X, K_n)$ は有限次元なので、これらの逆系は Mittag--Leffler 条件を満たし、
『層のコホモロジー』の補題 \ref{cohomology-lemma-RGamma-commutes-with-Rlim} の
$R^1\lim$ 項は消えるからである。

\medskip\noindent
$U' \subset U$ を開部分スキームとすると、
$D^b_{\textit{Coh}}(\mathcal{O}_U)$ の $K$ に関して関手的な標準写像
$$
H^i_c(U', K|_{U'}) \longrightarrow H^i_c(U, K)
$$
が存在する。例えば注意 \ref{duality-remark-covariance-open-lower-shriek} を参照せよ。
実際、注意 \ref{duality-remark-covariance-etale-lower-shriek} により、より一般に
$k$ 上有限型の分離スキームの \'etale 射 $U' \to U$ に対してこのような写像が存在する。

\medskip\noindent
$V$ を $k$-ベクトル空間とする。$\Hom_k(V, k)$ には次のように位相を入れる。
有限次元 $k$-部分ベクトル空間のフィルター付き合併として $V = \bigcup V_i$ と書き、
$\Hom_k(V, k) = \lim \Hom_k(V_i, k)$ に極限位相を入れる。
$\dim_k V < \infty$ ならば $\Hom_k(V, k)$ の位相は離散的である。
より一般に、$V = \colim_n V_n$ が有限次元ベクトル空間の帰納極限として書かれるならば、
位相ベクトル空間として $\Hom_k(V, k) = \lim \Hom_k(V_n, k)$ である。

\begin{lemma}
\label{duality-lemma-duality-compact-support}
$k$ を体とし、$p : U \to \Spec(k)$ を有限型分離射とする。
$\omega_{U/k}^\bullet = p^!\mathcal{O}_{\Spec(k)}$ と置く。このとき
$D^b_{\textit{Coh}}(\mathcal{O}_U)$ の $K$ に関して関手的な位相 $k$-ベクトル空間の標準同型
$$
\Hom_k(H^i(U, K), k) =
H^{-i}_c(U, R\SheafHom_{\mathcal{O}_U}(K, \omega_{U/k}^\bullet))
$$
が存在する。
\end{lemma}

\begin{proof}
$k$ 上のコンパクト化 $j : U \to X$ を選ぶ。
$\mathcal{I} \subset \mathcal{O}_X$ を
$V(\mathcal{I}) = X \setminus U$ を満たす準連接イデアル層とする。
『スキームの導来圏』の命題 \ref{perfect-proposition-DCoh} により、
$K = M|_U$ を満たす $M \in D^b_{\textit{Coh}}(\mathcal{O}_X)$ を選べる。
補題 \ref{duality-lemma-lift-map} により
$$
H^i(U, K) =
\Ext^i_U(\mathcal{O}_U, M|_U) =
\colim \Ext^i_X(\mathcal{I}^n, M) =
\colim H^i(X, R\SheafHom_{\mathcal{O}_X}(\mathcal{I}^n, M))
$$
である。$\mathcal{I}^n$ は連接 $\mathcal{O}_X$-加群なので
$\mathcal{I}^n$ は $D^-_{\textit{Coh}}(\mathcal{O}_X)$ に属する。従って
$R\SheafHom_{\mathcal{O}_X}(\mathcal{I}^n, M)$ は、『スキームの導来圏』の補題 \ref{perfect-lemma-coherent-internal-hom} により
$D^+_{\textit{Coh}}(\mathcal{O}_X)$ に属する。

\medskip\noindent
$q : X \to \Spec(k)$ を構造射とし、
$\omega_{X/k}^\bullet = q^!\mathcal{O}_{\Spec(k)}$ と置く。第 \ref{duality-section-duality-proper-over-field} 節を参照せよ。すると
\begin{align*}
\Hom_k(
&
H^i(X, R\SheafHom_{\mathcal{O}_X}(\mathcal{I}^n, M)), k) \\
& =
\Ext^{-i}_X(R\SheafHom_{\mathcal{O}_X}(\mathcal{I}^n, M),
\omega_{X/k}^\bullet) \\
& =
H^{-i}(X, R\SheafHom_{\mathcal{O}_X}(R\SheafHom_{\mathcal{O}_X}(
\mathcal{I}^n, M), \omega_{X/k}^\bullet))
\end{align*}
を得る。これは補題 \ref{duality-lemma-duality-proper-over-field} による。
補題 \ref{duality-lemma-internal-hom-evaluate-isom} の (1) により、標準写像
$$
R\SheafHom_{\mathcal{O}_X}(M, \omega_{X/k}^\bullet)
\otimes_{\mathcal{O}_X}^\mathbf{L} \mathcal{I}^n
\longrightarrow
R\SheafHom_{\mathcal{O}_X}(R\SheafHom_{\mathcal{O}_X}(
\mathcal{I}^n, M), \omega_{X/k}^\bullet)
$$
は同型である。$p^!$ は $q^!$ と $U$ への制限との合成として構成されるので、
$\omega^\bullet_{U/k} = \omega^\bullet_{X/k}|_U$ である。
従って $R\SheafHom_{\mathcal{O}_X}(M, \omega_{X/k}^\bullet)$ は
$D^b_{\textit{Coh}}(\mathcal{O}_X)$ の対象で、その $U$ への制限は
$R\SheafHom_{\mathcal{O}_U}(K, \omega_{U/k}^\bullet)$ である。
補題 \ref{duality-lemma-tensoring-Deligne-system} により
$$
\lim H^{-i}(X, R\SheafHom_{\mathcal{O}_X}(M, \omega_{X/k}^\bullet)
\otimes_{\mathcal{O}_X}^\mathbf{L} \mathcal{I}^n)
$$
は補題の等式の右辺を表す。以上で得た同型を組み合わせると補題の同型を得る。
\end{proof}

\begin{lemma}
\label{duality-lemma-duality-compact-support-restrict-open}
補題 \ref{duality-lemma-duality-compact-support} の記法のもとで、
$U' \subset U$ を開部分スキームとする。このとき図式
$$
\xymatrix{
\Hom_k(H^i(U, K), k) \ar[rr] & &
H^{-i}_c(U, R\SheafHom_{\mathcal{O}_U}(K, \omega_{U/k}^\bullet)) \\
\Hom_k(H^i(U', K|_{U'}), k) \ar[rr] \ar[u] & &
H^{-i}_c(U', R\SheafHom_{\mathcal{O}_{U'}}(K, \omega_{U'/k}^\bullet)) \ar[u]
}
$$
は可換である。ここで水平矢印は補題 \ref{duality-lemma-duality-compact-support} の同型であり、左の垂直矢印は制限写像
$H^i(U, K) \to H^i(U', K|_{U'})$ の双対写像である。
右の垂直矢印は注意 \ref{duality-remark-covariance-open-lower-shriek} のものである
（補題の前の議論を参照）。
\end{lemma}

\begin{proof}
読者にはこの証明を読み飛ばすことを強く勧める。補題 \ref{duality-lemma-duality-compact-support} の証明と同様に $X$ と $M$ を選ぶ。
$R\SheafHom$ と $\otimes^\mathbf{L}$ の添字 $\mathcal{O}_X$ は省略する。
補題 \ref{duality-lemma-duality-compact-support} の証明と同様に
$$
H^i(U, K) = \colim H^i(X, R\SheafHom(\mathcal{I}^n, M))
$$
および
$$
H^i(U', K|_{U'}) = \colim H^i(X, R\SheafHom((\mathcal{I}')^n, M))
$$
と書く。ここで注意 \ref{duality-remark-covariance-open-j-lower-shriek} の議論のように
$\mathcal{I}' \subset \mathcal{I}$ を選び、写像
$H^i(U, K) \to H^i(U', K|_{U'})$ が写像
$(\mathcal{I}')^n \to \mathcal{I}^n$ によって誘導されるようにする。
同様に
$$
H^i_c(U, R\SheafHom(K, \omega_{U/k}^\bullet)) =
\lim H^i(X,  R\SheafHom(M, \omega_{X/k}^\bullet)
\otimes^\mathbf{L} \mathcal{I}^n)
$$
および
$$
H^i_c(U', R\SheafHom(K|_{U'}, \omega_{U'/k}^\bullet)) =
\lim H^i(X,  R\SheafHom(M, \omega_{X/k}^\bullet)
\otimes^\mathbf{L} (\mathcal{I}')^n)
$$
と書く。このとき矢印
$H^i_c(U', R\SheafHom(K|_{U'}, \omega_{U'/k}^\bullet))
\to H^i_c(U, R\SheafHom(K, \omega_{U/k}^\bullet))$
も同じく写像 $(\mathcal{I}')^n \to \mathcal{I}^n$ から得られる。
図式
$$
\xymatrix{
R\SheafHom(M, \omega_{X/k}^\bullet)
\otimes^\mathbf{L} \mathcal{I}^n
\ar[rr] & &
R\SheafHom(R\SheafHom(\mathcal{I}^n, M), \omega_{X/k}^\bullet) \\
R\SheafHom(M, \omega_{X/k}^\bullet) \otimes^\mathbf{L} (\mathcal{I}')^n
\ar[rr] \ar[u] & &
R\SheafHom(R\SheafHom((\mathcal{I}')^n, M), \omega_{X/k}^\bullet) \ar[u]
}
$$
は可換である。『層のコホモロジー』の補題 \ref{cohomology-lemma-internal-hom-evaluate} における水平矢印の構成が三つの各変数に関して関手的だからである。
従って最後に、図式
$$
\xymatrix{
\Hom_k(H^i(X, R\SheafHom(\mathcal{I}^n, M)), k) \ar[r] &
H^{-i}(X, R\SheafHom(R\SheafHom(
\mathcal{I}^n, M), \omega_{X/k}^\bullet)) \\
\Hom_k(H^i(X, R\SheafHom((\mathcal{I}')^n, M)), k) \ar[r] \ar[u] &
H^{-i}(X, R\SheafHom(R\SheafHom(
(\mathcal{I}')^n, M), \omega_{X/k}^\bullet)) \ar[u]
}
$$
が可換であることに帰着する。これは双対性の同型
$$
\Hom_k(H^i(X, L), k) = \Ext^{-i}_X(L, \omega_{X/k}^\bullet) =
H^{-i}(X, R\SheafHom(L, \omega_{X/k}^\bullet))
$$
が $D_\QCoh(\mathcal{O}_X)$ の $L$ に関して関手的であることから従う。
\end{proof}

\begin{lemma}
\label{duality-lemma-h0-compactly-supported}
$X$ を体 $k$ 上の固有スキームとする。
$K \in D^b_{\textit{Coh}}(\mathcal{O}_X)$ が $i < 0$ に対して
$H^i(K) = 0$ を満たすとし、$\mathcal{F} = H^0(K)$ と置く。
$Z \subset X$ を閉部分集合で、補集合を $U = X \setminus U$ とする。このとき
$$
H^0_c(U, K|_U) \subset H^0(X, \mathcal{F})
$$
は、$Z$ のある開近傍上で消える $\mathcal{F}$ の大域切断全体からなる。
\end{lemma}

\begin{proof}
注意 \ref{duality-remark-covariance-open-lower-shriek} の写像
$H^0_c(U, K|_U) \to H^0_X(X, K) = H^0(X, K) = H^0(X, \mathcal{F})$
を考える。これを調べるため、$K$ を
$i < 0$ に対して $\mathcal{F}^i = 0$ を満たす有界複体
$\mathcal{F}^\bullet$ で表す。定義により
$$
H^0_c(U, K|_U) = \lim H^0(X, \mathcal{I}^n\mathcal{F}^\bullet)
= \lim \Ker(
H^0(X, \mathcal{I}^n\mathcal{F}^0) \to
H^0(X, \mathcal{I}^n\mathcal{F}^1))
$$
である。Artin--Rees（『スキームのコホモロジー』の補題 \ref{coherent-lemma-Artin-Rees}）により、これは
$\lim H^0(X, \mathcal{I}^n\mathcal{F})$ と同じである。
従って矢印 $H^0_c(U, K|_U) \to H^0(X, \mathcal{F})$ は単射であり、その像は任意の
$n$ に対して部分層 $\mathcal{I}^n\mathcal{F}$ に含まれる
$\mathcal{F}$ の大域切断全体からなる。
これらが $Z$ のある近傍上で消える切断として特徴づけられることは、
$\mathcal{F}$ の茎で Krull の共通部分定理（『可換代数』の補題 \ref{algebra-lemma-intersect-powers-ideal-module-zero}）を用いることから従う。
関数の場合については『可換代数』の注意 \ref{algebra-remark-intersection-powers-ideal} の議論を参照せよ。
\end{proof}




\section{Lichtenbaum の定理}
\label{duality-section-lichtenbaum}

\noindent
以下の定理は Lichtenbaum が予想し、Grothendieck が証明した
（\cite{Hartshorne-local-cohomology} を参照）。Kleiman による非常に見事な証明が
\cite{Kleiman-Lichtenbaum} にある。台付きコホモロジーの場合への一般化は
\cite{Lyubeznik-Lichtenbaum} に見られる。議論の最も興味深い部分は次の補題の証明に含まれている。

\begin{lemma}
\label{duality-lemma-lichtenbaum}
$U$ を多様体とし、$\mathcal{F}$ を連接 $\mathcal{O}_U$-加群とする。
$H^d(U, \mathcal{F})$ が零でないならば $\dim(U) \geq d$ であり、
等号が成り立つならば $U$ は固有である。
\end{lemma}

\begin{proof}
『層のコホモロジー』の命題 \ref{cohomology-proposition-vanishing-Noetherian} の
Grothendieck の消滅結果により $\dim(U) \geq d$ である。
$\dim(U) = d$ と仮定する。$U$ が $X$ で稠密となるコンパクト化
$U \to X$ を選ぶ。これは『平坦性の続論』の定理 \ref{flat-theorem-nagata} および補題 \ref{flat-lemma-compactifyable} により可能である。
$X$ をその簡約化で置き換えると、$X$ は次元 $d$ の固有多様体となり、
$U$ が固有であることと $U = X$ であることは同値である。
$Z = X \setminus U$ と置く。$Z$ が空でなければ
$H^d(U, \mathcal{F})$ が零であることを示す。

\medskip\noindent
$U$ への制限が $\mathcal{F}$ である連接 $\mathcal{O}_X$-加群
$\mathcal{G}$ を選ぶ。『スキームの性質』の補題 \ref{properties-lemma-lift-finite-presentation} を参照せよ。
第 \ref{duality-section-duality-proper-over-field} 節のように、
$\omega_X^\bullet$ を $X$ の双対化複体とし、
$\omega_U^\bullet = \omega_X^\bullet|_U$ と置く。
補題 \ref{duality-lemma-duality-compact-support} により、
$H^d(U, \mathcal{F})$ は
$$
H^{-d}_c(U, R\SheafHom_{\mathcal{O}_U}(\mathcal{F}, \omega_U^\bullet))
$$
と双対である。補題 \ref{duality-lemma-duality-proper-over-field} により、
$\omega_X^\bullet$ のコホモロジー層は次数 $< -d$ で消え、
$H^{-d}(\omega_X^\bullet) = \omega_X$ は $(S_2)$ を満たし台が $X$ である連接
$\mathcal{O}_X$-加群である。特に $\omega_X$ は捩れなしである。
『因子』の補題 \ref{divisors-lemma-torsion-free-finite-noetherian-domain} を参照せよ。
従ってコホモロジー層
$$
H^{-d}(R\SheafHom_{\mathcal{O}_X}(\mathcal{G}, \omega_X^\bullet)) =
\SheafHom(\mathcal{G}, \omega_X)
$$
は捩れなしである。『因子』の補題 \ref{divisors-lemma-hom-into-torsion-free} を参照せよ。
従ってこの層には、$X$ の空でない開集合上で消える零でない切断は存在しない
（そのような切断は捩れ切断である）。補題 \ref{duality-lemma-h0-compactly-supported} により
$H^{-d}_c(U, R\SheafHom_{\mathcal{O}_U}(\mathcal{F}, \omega_U^\bullet))$
は零であり、従って所望の通り $H^d(U, \mathcal{F})$ は零である。
\end{proof}

\begin{theorem}
\label{duality-theorem-lichtenbaum}
$X$ を体 $k$ 上有限型の空でない分離スキームとし、
$d = \dim(X)$ と置く。次は同値である。
\begin{enumerate}
\item $X$ 上の任意の連接 $\mathcal{O}_X$-加群 $\mathcal{F}$ に対して
$H^d(X, \mathcal{F}) = 0$ である。
\item $X$ 上の任意の準連接 $\mathcal{O}_X$-加群 $\mathcal{F}$ に対して
$H^d(X, \mathcal{F}) = 0$ である。
\item 次元 $d$ の既約成分 $X' \subset X$ で $k$ 上固有なものは存在しない。
\end{enumerate}
\end{theorem}

\begin{proof}
次元 $d$ で固有な既約成分 $X' \subset X$ が存在すると仮定し、これを整閉部分スキームとみなす。
$\omega_{X'}$ を $X'$ の $k$ 上の双対化加群とする。補題 \ref{duality-lemma-duality-proper-over-field} を参照せよ。同補題により
$H^d(X', \omega_{X'})$ は $H^0(X', \mathcal{O}_{X'})$ と双対であるから零でない。
従って $H^d(X, \omega_{X'}) = H^d(X', \omega_{X'})$ は零でなく、(1) は成り立たない。
これにより (1) は (3) を含意する。

\medskip\noindent
(3) が (1) を含意することを示そう。
$\mathcal{F}$ を $H^d(X, \mathcal{F})$ が零でない連接 $\mathcal{O}_X$-加群とする。
『スキームのコホモロジー』の補題 \ref{coherent-lemma-coherent-filter} のようなフィルトレーション
$$
0 = \mathcal{F}_0 \subset \mathcal{F}_1 \subset
\ldots \subset \mathcal{F}_m = \mathcal{F}
$$
を選ぶ。完全列
$$
H^d(X, \mathcal{F}_i) \to H^d(X, \mathcal{F}_{i + 1}) \to
H^d(X, \mathcal{F}_{i + 1}/\mathcal{F}_i)
$$
を得る。従って、ある $i \in \{1, \ldots, m\}$ に対して
$H^d(X, \mathcal{F}_{i + 1}/\mathcal{F}_i)$ は零でない。
フィルトレーションの選び方により、整閉部分スキーム $Z \subset X$ と、
零でない連接イデアル層 $\mathcal{I} \subset \mathcal{O}_Z$ で
$H^d(Z, \mathcal{I})$ が零でないものが存在する。
補題 \ref{duality-lemma-lichtenbaum} により $\dim(Z) = d$ かつ $Z$ は $k$ 上固有であり、
これは (3) に矛盾する。従って (3) は (1) を含意する。

\medskip\noindent
最後に、任意の Noether スキーム $X$ に対して (1) と (2) が同値であることを示す。
(2) が (1) を含意することは明らかである。逆に (1) を仮定し、
$\mathcal{F}$ を準連接 $\mathcal{O}_X$-加群とする。
『スキームの性質』の補題 \ref{properties-lemma-quasi-coherent-colimit-finite-type} により、
$\mathcal{F} = \colim \mathcal{F}_i$ を連接部分加群のフィルター付き帰納極限として書ける。
『層のコホモロジー』の補題 \ref{cohomology-lemma-quasi-separated-cohomology-colimit} により
$H^d(X, \mathcal{F}) = \colim H^d(X, \mathcal{F}_i) = 0$ である。
従って (2) が成り立つ。
\end{proof}
